
%%FILE blueprint/src/chapter/ch01_overview.tex
\chapter{Overview}
\label{chap:overview}

This blueprint records the formalization plan for the low individual degree
test paper of Ji, Natarajan, Vidick, Wright, and Yuen
\cite{Ji2020LowIndividualDegree}.  The paper is available as
arXiv:2009.12982.

\begin{theorem}[Raz--Safra]\label{thm:raz-safra}
	\uses{def:polyfunc}
	Suppose Provers~$\mathrm{A}$ and~$\mathrm{B}$ pass the $k=2$ surface-versus-point low-degree test with probability $1-\eps$. Then there exists a degree-$d$ polynomial $g \colon \F_q^m \to \F_q$ such that
	\[
		\Prb_{u \sim \F_q^m}[g(u)=a] \ge 1-\eps-\poly(m)\cdot\poly(d/q).
	\]
	Here $a$ denotes the prover's point answer.
\end{theorem}
\begin{remark}
	The paper theorem is an external classical soundness input and is not used by
	the current low individual degree test formalization route.
\end{remark}
\begin{proof}
	\uses{def:polyfunc}
	See Raz--Safra [RS97].
\end{proof}

% External classical source theorem recorded for orientation; the Lean-facing
% conditional interface is stated immediately below.
\begin{theorem}[Polishchuk--Spielman]\label{thm:classical-test-soundness}
	\uses{def:low-individual-degree-test}
	Suppose Provers~$\mathrm{A}$ and~$\mathrm{B}$ pass the low individual degree test with probability $1-\eps$. Then there exists a polynomial $g \colon \F_q^m \to \F_q$ with individual degree $d$ such that
	\[
		\Prb_{u \sim \F_q^m}[g(u)=a] \ge 1-\poly(m)\cdot\bigl(\poly(\eps)+\poly(d/q)\bigr).
	\]
	Here $a$ denotes the prover's point answer.
\end{theorem}
\begin{remark}
	The paper theorem is an external classical soundness input.  The nearby Lean
	corollary below records the specialized conditional interface currently present
	in the repository; it is not used to mark this source theorem as formalized.
\end{remark}
\begin{proof}
	See Polishchuk--Spielman [PS94].
\end{proof}

\begin{proposition}[Conditional Polishchuk--Spielman corollary]\label{prop:lean-classical-test-soundness-interface}
	\lean{MIPStarRE.LDT.Test.TwoProverClassicalLIDPassCondition, MIPStarRE.LDT.Test.PolishchukSpielmanClassicalSoundnessStatement, MIPStarRE.LDT.Test.classicalTestSoundness}
	\leanok
	\uses{thm:classical-test-soundness}
	The Lean theorem \texttt{MIPStarRE.LDT.Test.classicalTestSoundness} is a
	conditional corollary of the quoted Polishchuk--Spielman theorem.  It assumes
	the modeled deterministic low individual degree test pass condition
	\texttt{MIPStarRE.LDT.Test.TwoProverClassicalLIDPassCondition} and a
	specialized external hypothesis
	\texttt{MIPStarRE.LDT.Test.PolishchukSpielmanClassicalSoundnessStatement}
	at the caller's chosen slack bound.  Under these hypotheses it returns a
	repository polynomial whose evaluations agree with the point answers up to
	that slack.
\end{proposition}
\begin{proof}
	\leanok
	The Lean proof applies the specialized external Polishchuk--Spielman
	hypothesis to the modeled low individual degree test pass condition.
\end{proof}

% Paper source: references/ldt-paper/introduction.tex, lines 236--264.
\begin{remark}[A degree-$(d+1)$ strategy that passes with probability $1-\frac{1}{m}$]\label{ex:bad-individual-degree-example}
		Fix the polynomial $h(x_1, \ldots, x_m) = x_1^{d+1}$. Consider the following classical strategy for the degree-$d$ low individual degree test. On a point query $u \in \F_q^m$, Prover~$\mathrm{A}$ answers $a = h(u)$. On an axis-parallel line query $\ell = \{u + x e_i : x \in \F_q\}$, Prover~$\mathrm{B}$ returns the restriction $h\vert_{\ell}$ when $i \neq 1$, since $h$ is then constant along $\ell$, and returns the zero polynomial when $i = 1$.

		The verifier certainly accepts whenever $i \neq 1$, which happens with probability $1 - \frac{1}{m}$. Thus the strategy passes the degree-$d$ test with probability at least $1-\eps$ for $\eps = \frac{1}{m}$. However, Prover~$\mathrm{A}$ is answering according to a polynomial whose individual degree in the first variable is $d+1$. By Schwartz--Zippel, any polynomial $g$ of individual degree at most $d$ agrees with $h$ on at most a $\frac{d+1}{q}$ fraction of all inputs, so
		\[
			\Prb_{u \sim \F_q^m}[g(u)=a] \le 1 - m\eps + \frac{d+1}{q}.
		\]
		This shows that the $m$-dependent slack in Theorems~\ref{thm:classical-test-soundness} and~\ref{thm:main-informal} is unavoidable up to polynomial factors.
\end{remark}

% Informal source theorem; its proof points to the formal theorem below.
\begin{theorem}[Main theorem, informal]\label{thm:main-informal}
	\uses{def:low-individual-degree-test, def:projective-strategy, def:polymeasurements}
	If $d > 0$ and a two-prover projective strategy passes the $(m,q,d)$ low individual degree test with probability $1-\eps$, then there are global polynomial measurements whose evaluations agree with the point answers except with error $\poly(m)(\poly(\eps)+\poly(d/q))$.
\end{theorem}
\begin{proof}
	\uses{thm:main-formal}
	This is the informal form of Theorem~\ref{thm:main-formal}.
\end{proof}


%%FILE blueprint/src/chapter/ch02_test.tex
\chapter{The low individual degree test}
\label{chap:test}
\label{sec:the-test}

\section{The game and its strategies}

\begin{definition}[Roles]\label{def:roles}
	\lean{MIPStarRE.LDT.Role, MIPStarRE.LDT.Role.other}
	\leanok
	A role is an element of $\{\mathrm A,\mathrm B\}$. For a role $r \in \{\mathrm A,\mathrm B\}$, write $\overline r$ for the other element of $\{\mathrm A,\mathrm B\}$.
\end{definition}

\begin{definition}[Low individual degree test]\label{def:low-individual-degree-test}
	\lean{MIPStarRE.LDT.ProjStrat.lowIndividualDegreeFailureProbability, MIPStarRE.LDT.ProjStrat.PassesLowIndividualDegreeTest}
	\leanok
	\uses{def:roles}
	Fix integers $m,d \ge 0$ and a prime power $q$. The $(m,q,d)$ low individual degree test has three equally likely subtests.
	In the axis-parallel lines test, choose a uniformly random role $r \in \{\mathrm A,\mathrm B\}$, a uniformly random point $u \in \F_q^m$, and a uniformly random coordinate $i \in \{1,\dots,m\}$. Let $\ell=\{u+t e_i : t \in \F_q\}$. Player $r$ receives $\ell$ and returns a degree-$d$ univariate polynomial $f\colon \ell \to \F_q$; Player $\overline r$ receives $u$ and returns a field element $a \in \F_q$. The verifier accepts when $f(u)=a$.
	In the self-consistency test, both provers receive the same uniformly random point $u \in \F_q^m$ and must return the same field element.
	In the diagonal lines test, choose a uniformly random role $r \in \{\mathrm A,\mathrm B\}$, a uniformly random point $u \in \F_q^m$, a uniformly random index $i \in \{1,\dots,m\}$, and a uniformly random direction vector $v \in \F_q^m$ whose last $m-i$ coordinates are $0$. Let $\ell=\{u+t v : t \in \F_q\}$. Player $r$ receives $\ell$ and returns a degree-$md$ univariate polynomial $f\colon \ell \to \F_q$; Player $\overline r$ receives $u$ and returns a field element $a \in \F_q$. The verifier accepts when $f(u)=a$.
\end{definition}

\begin{lemma}[Branch-average form of the classical test]\label{fig:test}
	\lean{MIPStarRE.LDT.Test.TwoProverClassicalLIDStrategy.lowIndividualDegreeAcceptanceProbability_eq_branchAverage}
	\leanok
	\uses{def:low-individual-degree-test}
	For every classical strategy, the acceptance probability of the low individual degree test is
	\[
		\frac13\left(\frac{P^{\mathrm{line,point}}_{\mathrm{axis}}+P^{\mathrm{point,line}}_{\mathrm{axis}}}{2}
		+P_{\mathrm{self}}
		+\frac{P^{\mathrm{line,point}}_{\mathrm{diag}}+P^{\mathrm{point,line}}_{\mathrm{diag}}}{2}\right),
	\]
	where the five terms are the acceptance probabilities of the two axis-parallel role orderings, the self-consistency branch, and the two diagonal role orderings.
\end{lemma}
\begin{proof}
	\leanok
	\proves{fig:test}
	\uses{def:low-individual-degree-test}
	Expand the definition of the total acceptance probability as the average of
	the three subtests.  The axis-parallel and diagonal branches each first
	choose one of the two role orderings uniformly, so their branch acceptance
	probabilities are the corresponding averages displayed above.
\end{proof}

\begin{definition}[Symmetric projective strategy]\label{def:symmetric-projective-strategy}
	\lean{MIPStarRE.LDT.SymStrat}
	\leanok
	A symmetric projective strategy for the $(m,q,d)$ low individual degree test consists of a permutation-invariant bipartite state $\ket{\psi} \in \mathcal H \ot \mathcal H$, a projective point measurement $A^u=\{A^u_a\}_{a \in \F_q}$ for each $u \in \F_q^m$, a projective axis-parallel line measurement $B^\ell=\{B^\ell_f\}$ for each axis-parallel line $\ell$, and a projective diagonal-line measurement $L^\ell=\{L^\ell_f\}$ for each line $\ell$.
\end{definition}

\begin{definition}[General projective strategy]\label{def:projective-strategy}
	\lean{MIPStarRE.LDT.ProjStrat}
	\leanok
	A general projective strategy for the $(m,q,d)$ low individual degree test consists of a bipartite state $\ket{\psi} \in \mathcal H_{\mathrm A} \ot \mathcal H_{\mathrm B}$ together with point, axis-parallel line, and diagonal-line projective measurements for each prover separately. We write these measurements as $A^{\mathrm A}, B^{\mathrm A}, L^{\mathrm A}$ on the first prover and $A^{\mathrm B}, B^{\mathrm B}, L^{\mathrm B}$ on the second.
\end{definition}

\begin{lemma}[Bipartite strategy block identities]\label{lem:bipartite-strategy-block-identities}
	\lean{MIPStarRE.LDT.ProjStrat.localDirectSumBlock_nonneg, MIPStarRE.LDT.ProjStrat.roleBlock_mul, MIPStarRE.LDT.ProjStrat.roleBlock_nonneg}
	\leanok
	\uses{def:projective-strategy}
	The block operators obtained from a bipartite projective strategy are
	positive semidefinite, and the role-indexed blocks multiply according to the
	Kronecker delta on roles.
\end{lemma}
\begin{proof}
	\leanok
	The assertions are the block-diagonal matrix calculations for the direct-sum
	and role-register operators. Positivity follows by writing each positive
	operator as $C^*C$ and placing the corresponding square root in the same
	block decomposition.
\end{proof}

\begin{lemma}[Heterogeneous role-register unsymmetrization]\label{lem:heterogeneous-role-register-unsymmetrization}
	\lean{MIPStarRE.LDT.ProjStrat.qBipartiteConsDefect_roleRegisterProjMeas_arbitrary_eq_average}
	\lean{MIPStarRE.LDT.ProjStrat.qBipartiteConsDefect_extractRoleRegisterBob_le_two_symm}
	\lean{MIPStarRE.LDT.ProjStrat.qBipartiteConsDefect_extractRoleRegisterAlice_le_two_symm}
	\leanok
	\uses{def:projective-strategy, lem:bipartite-strategy-block-identities}
	Let \(\psi\) be a normalized bipartite state on
	\(\mathcal H_{\mathrm A}\otimes\mathcal H_{\mathrm B}\), let
	\(M^{\mathrm A}\) and \(M^{\mathrm B}\) be projective measurements on the two
	prover spaces with a common outcome type, and let \(G\) be an arbitrary
	measurement on the heterogeneous role-register space
	\[
		\{\mathrm A,\mathrm B\}\times(\mathcal H_{\mathrm A}\oplus\mathcal H_{\mathrm B}).
	\]
	After applying the direct-sum role-register symmetrization, the consistency
	defect between the symmetrized point measurement and \(G\) is the average of
	the two defects obtained from the occupied principal blocks of \(G\).  In
	particular, each extracted defect is at most twice the symmetrized defect.
\end{lemma}
\begin{proof}
	\leanok
	The proof is a trace-compression calculation.  The \(AB\)-supported component
	of the symmetrized density pairs with an arbitrary role-register observable
	only through the principal block indexed by
	\((\mathrm A,\mathrm{inl})\) on the left and
	\((\mathrm B,\mathrm{inr})\) on the right, and similarly for the \(BA\)
	component.  Summing over outcomes gives the average-of-two-defects identity,
	and nonnegativity of consistency defects gives the two factor-two
	consequences.
\end{proof}

\begin{lemma}[Heterogeneous role-register symmetrization]\label{lem:heterogeneous-role-register-symmetrization}
	\lean{MIPStarRE.LDT.ProjStrat.roleRegisterSymmStrategy}
	\lean{MIPStarRE.LDT.ProjStrat.roleRegisterSymmStrategy_is_good_three_mul}
	\leanok
	\uses{def:projective-strategy, def:symmetric-projective-strategy, def:good-strategy}
	Let
	\[
		(\psi,A^{\mathrm A},B^{\mathrm A},L^{\mathrm A},
		A^{\mathrm B},B^{\mathrm B},L^{\mathrm B})
	\]
	be a two-space projective strategy passing the low individual degree test
	with error at most \(\eps\).  The heterogeneous role-register construction
	on
	\[
		\{\mathrm A,\mathrm B\}\times
		(\mathcal H_{\mathrm A}\oplus\mathcal H_{\mathrm B})
	\]
	produces a symmetric strategy that is
	\((3\eps,3\eps,3\eps)\)-good.
\end{lemma}
\begin{proof}
	\leanok
	The role-register state is the direct-sum version of the symmetrized state in
	the paper proof.  The axis-parallel and diagonal branches are exactly the
	role averages of the original strategy, and the self-consistency branch is
	the original point-agreement branch.  Since the test chooses among its three
	branches with equal probability, each branch is bounded by \(3\eps\).
\end{proof}

\begin{remark}\label{rem:strategy-convention}
	Throughout this work, the term strategy refers to a projective strategy, and the term symmetric strategy refers to a symmetric projective strategy.
\end{remark}

\begin{definition}[Good strategy]\label{def:good-strategy}
	\lean{MIPStarRE.LDT.SymStrat.IsGood}
	\leanok
	\uses{def:low-individual-degree-test}
	A strategy is $(\eps,\delta,\gamma)$-good if it passes the axis-parallel lines test with probability at least $1-\eps$, the self-consistency test with probability at least $1-\delta$, and the diagonal lines test with probability at least $1-\gamma$.
\end{definition}

\begin{remark}[Good-strategy characterization]\label{rem:good-strat-characterization}
	Using notation which will be introduced in Chapter~\ref{chap:preliminaries} below, a symmetric strategy is $(\eps,\delta,\gamma)$-good if and only if it satisfies the following three conditions. For $\ell$ and $u$ as in the axis-parallel lines test,
	\[
		A_a^u \ot I \simeq_\eps I \ot B^\ell_{[f(u)=a]},
		\quad
		\text{and}
		\quad
		A_a^u \ot I \simeq_\delta I \ot A_a^u.
	\]
	And for $\ell$ and $u$ as in the diagonal lines test,
	\[
		A_a^u \ot I \simeq_\gamma I \ot L^\ell_{[f(u)=a]}.
	\]
\end{remark}

\begin{definition}[Last-direction line notation]\label{not:conditioned-on-last-direction}
	\lean{MIPStarRE.LDT.lastDirectionLine, MIPStarRE.LDT.lastDirectionMeasurementFamily}
	\leanok
	\uses{def:symmetric-projective-strategy}
	For $u \in \F_q^m$, the last-direction axis-parallel line in $\F_q^{m+1}$ is
	\[
		\ell_u = \{(u,x) \mid x \in \F_q\}.
	\]
	If $(\psi,A,B,L)$ is a symmetric strategy for the $(m+1,q,d)$ low individual degree test, we write $B_f^u$ for the operator $B_f^{\ell_u}$. For a function $f \colon \ell_u \to \F_q$, we also write $f(x)$ for $f(u,x)$.
\end{definition}

\begin{definition}[Restricted diagonal lines test]\label{def:restricted-diagonal-lines-test}
	\lean{MIPStarRE.LDT.RestrictedDiagonalSample, MIPStarRE.LDT.MainInductionStep.RestrictedSymStrat.diagonalFailureProbability}
	\leanok
	Consider the $(m,q,d)$ low individual degree test. For $j \in \{1,\dots,m\}$, the $j$-restricted diagonal lines test is the diagonal lines test conditioned on $i=j$. In particular, the $m$-restricted diagonal lines test is the diagonal lines test in which the sampled line is a uniformly random line in $\F_q^m$.
\end{definition}

\section{Formal soundness target}

\begin{theorem}[Quantum soundness of the low individual degree test]\label{thm:main-formal}
	\lean{MIPStarRE.LDT.Test.mainFormal}
	\leanok
	\uses{def:low-individual-degree-test, def:projective-strategy, def:polymeasurements, def:simeq}
	Let $(\psi, A^{\mathrm A}, B^{\mathrm A}, L^{\mathrm A},
	A^{\mathrm B}, B^{\mathrm B}, L^{\mathrm B})$ be a projective strategy that
	passes the $(m,q,d)$ low individual degree test with probability at least
	$1-\eps$. Let \(k>0\) be an integer with \(k \ge 400md\), and set
	\[
		\nu = 100000 k^2 m^4\left(\eps^{1/40000} + (d/q)^{1/40000} + e^{-k/(2560000m^2)}\right).
	\]
	Then there exist projective measurements
	$G^{\mathrm A}, G^{\mathrm B} \in \polymeas{m}{q}{d}$ with the following
	properties:
	\begin{enumerate}
		\item on average over $u \sim \F_q^m$,
		      \[
			      A_a^{\mathrm A,u} \ot I \simeq_\nu I \ot G^{\mathrm B}_{[g(u)=a]},
			      \qquad
			      I \ot A_a^{\mathrm B,u} \simeq_\nu G^{\mathrm A}_{[g(u)=a]} \ot I;
		      \]
		\item
		      \[
			      G_g^{\mathrm A} \ot I \simeq_\nu I \ot G_g^{\mathrm B}.
		      \]
	\end{enumerate}
\end{theorem}
\begin{proof}
	\leanok
	\uses{prop:main-formal-source-reduction, prop:main-formal-source-small-error}
	The saturated-error branch is formalized by
	\texttt{mainFormal\_trivial\_witness}.  The remaining branch is the
	small-error source-boundary proposition
	Proposition~\ref{prop:main-formal-source-small-error}.  The
	factor \(400\) and the nonzero sampling condition are the corrected source
	boundaries recorded in \cite{gap:issue-906-main-formal-k-bound} and
	\cite{gap:issue-422-main-formal-zero-k-boundary}.
\end{proof}

\begin{proposition}[Source-boundary reduction for the final theorem]\label{prop:main-formal-source-reduction}
	\lean{MIPStarRE.LDT.Test.mainFormalConclusion}
	\leanok
	This proposition records the saturated-error reduction for the corrected
	two-space theorem
	Theorem~\ref{thm:main-formal}.  It proves the saturated-error branch by the
	trivial consistency bound and reduces the remaining non-vacuous branch to
	Proposition~\ref{prop:main-formal-source-small-error}.  It is not
	an additional hypothesis of the paper theorem.
\end{proposition}
\begin{proof}
	\leanok
	\proves{prop:main-formal-source-reduction}
	\uses{prop:main-formal-source-small-error}
	The saturated-error branch is already proved.  The small-error branch is
	Proposition~\ref{prop:main-formal-source-small-error}.
\end{proof}

\begin{proposition}[Small-error source-boundary conclusion for the final theorem]\label{prop:main-formal-source-small-error}
	\lean{MIPStarRE.LDT.Test.mainFormal_smallErrorConclusion}
	\leanok
	This is the non-vacuous branch of the corrected two-space theorem.
	It proves the conclusions of Theorem~\ref{thm:main-formal} under
	\(k\ge 400md\), \(k>0\), and \(\nu<1\).  It is not a source assumption.  Its
	proof uses the two-space role-register reduction and the scalar absorption
	available after excluding the zero-sampling boundary.
\end{proposition}
\begin{proof}
	\leanok
	\proves{prop:main-formal-source-small-error}
	\uses{prop:main-formal-source-two-space-role-register}
	The proof applies the checked two-space role-register construction and then
	absorbs its explicit errors into \(\nu\) at the nonzero sampling boundary.
\end{proof}

\begin{proposition}[Two-space role-register reduction for the source theorem]\label{prop:main-formal-source-two-space-role-register}
	\lean{MIPStarRE.LDT.Test.mainFormalConclusion_ofRoleRegisterScalarBoundary}
	\leanok
	\uses{lem:heterogeneous-role-register-unsymmetrization}
	The printed final theorem starts from a projective strategy on two possibly
	different local Hilbert spaces.  The source proof uses the role-register
	symmetrization and then unsymmetrizes the polynomial measurement.  Lean
	formalizes this reduction for the general two-space source strategy.  The
	trace-level factor-two unsymmetrization estimate for arbitrary heterogeneous
	role-register measurements is formalized in
	Lemma~\ref{lem:heterogeneous-role-register-unsymmetrization}.  The
	heterogeneous role-register symmetrization may be fed into the source-shaped
	main-induction theorem under the corrected large-\(k\) hypothesis
	\(k\ge 400md\), and the resulting role-register consistency estimate can be
	unsymmetrized into the two point-consistency estimates on the original
	two-space strategy with the expected factor \(2\).

	The remaining steps of the source-boundary passage are also checked.  The
	point-agreement branch is available for the paper-faithful two-space
	strategy, and the Schwartz--Zippel Step~5 theorem has been generalized from
	the same-space carrier to a bipartite state on \(H_A\otimes H_B\).  The
	heterogeneous triangle/SDD comparison theorem gives the complete-measurement
	full-polynomial consistency statement at the end of the paper's Step~5
	calculation.  The two heterogeneous orthonormalization applications produce
	projective submeasurements on the two local Hilbert spaces, and the completion
	step widens the resulting left- and right-factor state-dependent-distance
	estimates to the orthonormalize-and-complete error appearing in the paper.
	The repaired polynomial line-169 consistency relations use the
	Cauchy--Schwarz loss from the pre-completion orthonormalization estimates,
	not a new hypothesis.  The final point-evaluation triangle postprocesses
	these relations and the completed \(Q_A,Q_B\) projective consistency estimate
	by evaluation at a point, combines them by the heterogeneous triangle
	inequality, and absorbs the explicit pre-absorption errors into the final
	\(\nu=\texttt{mainFormalError}\) bound under the nonzero scalar-cascade
	boundary \(0<k\).  This is the small-error branch used in the corrected
	two-space source theorem, where \(k>0\) is part of the theorem statement.
\end{proposition}
\begin{proof}
	\proves{prop:main-formal-source-two-space-role-register}
	\lean{MIPStarRE.LDT.Test.mainFormalConclusion_ofRoleRegisterScalarBoundary}
	\leanok
	\uses{lem:heterogeneous-role-register-symmetrization, lem:heterogeneous-role-register-unsymmetrization}
	\uses{thm:main-induction}
	The proof is the linked heterogeneous role-register chain, ending in
	\texttt{mainFormalConclusion\_ofRoleRegisterScalarBoundary}.  It
	combines the source-shaped main-induction theorem, the heterogeneous
	unsymmetrization estimates, the step making measurements projective on both
	local Hilbert spaces, the completion and line-169 consistency steps, and the
	final scalar absorption under the nonzero sampling boundary.
\end{proof}

\begin{remark}[Final-theorem zero-sampling boundary]\label{rem:main-formal-source-k-range-boundary}
	The paper states Theorem~\ref{thm:main-formal} for \(k\ge md\), while the
	corrected formal route uses the confirmed large-\(k\) condition
	\(k\ge 400md\) and a nonzero \(k\) boundary for the scalar cascade.  The
	factor \(400\) is now treated as a statement correction, following
	Theorem~\ref{thm:main-induction}.  The other corrected final-theorem boundary is
	the zero-sampling case allowed when \(d=0\), since \(400md\le k\) then permits
	\(k=0\).  At this excluded boundary the displayed final error itself collapses to \(0\);
	the formal calculation of this equality is linked in Lean, and the
	obstruction is documented in \cite{gap:issue-422-main-formal-zero-k-boundary}.
\end{remark}

\begin{remark}[$k$-bound boundary for the public theorem]\label{rem:main-formal-k-boundary}
	The paper statement uses the weaker hypothesis $k \ge md$, but the proof
	invokes the Section~6 pasting theorem with side condition $k \ge 400\,md$.
	This side condition is not implied by $k \ge md$.  The Lean theorem
	therefore records the stronger large-$k$ assumption required by the pasting
	step, together with the scalar-cascade boundary \(k>0\), rather than hiding
	the side-condition gap.  The helper
	\texttt{mainFormal\_trivial\_witness} remains only for already-saturated
	branches where $\nu \ge 1$ makes the three $\simeq_\nu$ conclusions vacuous;
	it is not a replacement for the zero-sampling boundary \(k=0\) when \(md=0\).
	The source-boundary reduction proposition proves the saturated-error branch
	and delegates the non-vacuous branch to the checked small-error proposition;
	neither declaration is an extra hypothesis of Theorem~\ref{thm:main-formal}.
\end{remark}


%%FILE blueprint/src/chapter/ch03_preliminaries.tex
\chapter{Preliminaries}
\label{chap:preliminaries}
% Paper label sec:prelims kept for cross-reference fidelity; resolves to chapter number.
\label{sec:prelims}

\section{Polynomials and measurements}

For complex numbers $\alpha,\beta \in \C$, write $\alpha \approx_\eps \beta$ when $|\alpha-\beta|\le \eps$. In particular,
\begin{equation}\label{eq:triangle-inequality-for-numbers}
	\text{if $\alpha \approx_{\eps} \beta$ and $\beta \approx_{\delta} \gamma$, then $\alpha \approx_{\eps+\delta}\gamma$.}
\end{equation}

Fix a prime power $q=p^t$. Write $\omega=e^{2\pi i/p}$, and let $\tr:\F_q\to\F_p$ denote the finite-field trace.

\begin{definition}[Finite field trace]\label{def:ff-trace}
	\lean{MIPStarRE.LDT.Preliminaries.ffTrace}
	\leanok
	The finite-field trace is the map $\tr:\F_q\to\F_p$ defined by
	\[
		\tr[x] = \sum_{\ell=0}^{t-1} x^{p^\ell}.
	\]
\end{definition}

\begin{lemma}\label{prop:fourier-fact-scalar}
	\lean{MIPStarRE.LDT.Preliminaries.fourier_fact_scalar}
	\leanok
	\uses{def:ff-trace}
	Let $a \in \F_q$. Then
	\begin{equation*}
		\E_{\bx \sim \F_q} \omega^{\tr[\bx \cdot a]}
		= \left\{\begin{array}{rl}
			1 & \text{if $a = 0$}, \\
			0 & \text{otherwise}.
		\end{array}\right.
	\end{equation*}
\end{lemma}
\begin{proof}
	\leanok
	If $a=0$, then $\tr[\bx\cdot a]=0$ for every $\bx$, so the average is $1$. If $a\neq 0$, choose $y \in \F_q$ such that $\tr[a\cdot y]\neq 0$, and set
	\[
		C=\E_{\bx \sim \F_q} \omega^{\tr[\bx \cdot a]}.
	\]
	Translation invariance of the uniform distribution gives
	\[
		C=\E_{\bx \sim \F_q} \omega^{\tr[(\bx+y)\cdot a]}
		= C \cdot \omega^{\tr[y\cdot a]}.
	\]
	Since $\omega^{\tr[y\cdot a]}\neq 1$, it follows that $C=0$.
\end{proof}

\begin{lemma}\label{prop:fourier-fact-vector}
	\lean{MIPStarRE.LDT.Preliminaries.fourier_fact_vector}
	\leanok
	\uses{def:ff-trace, prop:fourier-fact-scalar}
	Let $v \in \F_q^m$. Then
	\begin{equation*}
		\E_{\bu \sim \F_q^m} \omega^{\tr[\bu \cdot v]}
		= \left\{\begin{array}{rl}
			1 & \text{if $v = 0$}, \\
			0 & \text{otherwise}.
		\end{array}\right.
	\end{equation*}
\end{lemma}
\begin{proof}
	\leanok
	By linearity of the trace and independence of the coordinates,
	\[
		\E_{\bu \sim \F_q^m} \omega^{\tr[\bu \cdot v]}
		= \prod_{i=1}^m \E_{\bu_i \sim \F_q} \omega^{\tr[\bu_i \cdot v_i]}.
	\]
	Proposition~\ref{prop:fourier-fact-scalar} shows that the $i$-th factor is $1$ when $v_i=0$ and $0$ otherwise, so the product is $1$ exactly when $v=0$.
\end{proof}

\begin{definition}[Low individual degree polynomials]\label{def:polyfunc}
	\lean{MIPStarRE.LDT.Preliminaries.polyFunc, MIPStarRE.LDT.Polynomial}
	\leanok
	Let $\polyfunc{m}{q}{d}$ be the set of polynomials $g \in \F_q[x_1,\dots,x_m]$ whose degree in each variable is at most $d$, viewed as functions $\F_q^m \to \F_q$ via evaluation. (Over finite fields, distinct low-degree polynomials can induce the same function; we identify elements of $\polyfunc{m}{q}{d}$ with their polynomial representatives, not their functional equivalence classes.)
\end{definition}

\begin{remark}[Individual-degree convention]\label{rem:individual-degree-convention}
	Here, ``individual degree $d$'' means degree at most $d$ in each coordinate. In
	particular,
	\[
		\polyfunc{m}{q}{d} \subseteq \polyfunc{m}{q}{d+1}.
	\]
\end{remark}

\begin{lemma}[Schwartz--Zippel lemma~\cite{Sch80,Zip79}]\label{lem:schwartz-zippel-total-degree}
	\lean{MIPStarRE.LDT.Preliminaries.schwartzZippel_totalDegree}
	\leanok
	Let $g, h:\F_q^m \rightarrow \F_q$ be two distinct polynomials of total degree~$d$.
	Then
	\begin{equation*}
		\Prb_{\bx \sim \F_q^m}[g(\bx) = h(\bx)] \leq \frac{d}{q}.
	\end{equation*}
\end{lemma}
\begin{proof}
	\leanok
	Apply the Schwartz--Zippel lemma to the nonzero polynomial $g-h$, which has total degree at most $d$.
\end{proof}

\begin{lemma}[Schwartz--Zippel for individual degree]\label{lem:schwartz-zippel-individual}
	\lean{MIPStarRE.LDT.Preliminaries.schwartzZippel_individualDegree}
	\leanok
	\uses{def:polyfunc, lem:schwartz-zippel-total-degree}
	If $g,h \in \polyfunc{m}{q}{d}$ are distinct, then
	\[
		\Prb_{u \sim \F_q^m}[g(u)=h(u)] \le \frac{md}{q}.
	\]
\end{lemma}
\begin{proof}
	\leanok
	The difference $g-h$ is a nonzero polynomial of total degree at most $md$. Apply Lemma~\ref{lem:schwartz-zippel-total-degree}.
\end{proof}

\begin{lemma}[Polynomial agreement bound on coded points]\label{lem:polynomial-agreement-mdq}
	\lean{MIPStarRE.LDT.Preliminaries.polynomialAgreement_avg_le_mdq}
	\leanok
	\uses{lem:schwartz-zippel-individual}
	Let $g, g' : \mathrm{Polynomial}\,\mathrm{params}$ be two distinct full polynomial outcomes with individual degrees at most $d$. Then
	\[
		\E_{u \sim \mathrm{Point}\,\mathrm{params}}\bigl[\mathbf{1}[g(u) = g'(u)]\bigr] \le \frac{md}{q}.
	\]
	This packages the $md/q$ loss term used in the \texttt{mainFormal} self-consistency cascade (\texttt{inductive\_step.tex}, lines 119--133) and in \texttt{comMain} (\texttt{commutativity-G.tex}).
\end{lemma}
\begin{proof}
	\leanok
	Transport along the equivalence between coded $\F_q$ points and the underlying scalar function space, then apply Lemma~\ref{lem:schwartz-zippel-individual}.
\end{proof}

\begin{lemma}[Axis-line polynomial agreement bound]\label{lem:axis-line-polynomial-agreement-mdq}
	\lean{MIPStarRE.LDT.Preliminaries.axisLinePolynomialAgreement_avg_eq_scalarDomain, MIPStarRE.LDT.Preliminaries.axisLinePolynomialAgreement_avg_le_mdq}
	\leanok
	Let $f,h : \mathrm{AxisLinePolynomial}\,\mathrm{params}$ be two line-polynomial outcomes with distinct underlying polynomials. Then
	\[
		\E_{t \sim \F_q}\bigl[\mathbf{1}[f(t)=h(t)]\bigr] \le \frac{md}{q}.
	\]
	The native univariate root-counting bound is $d/q$; the statement deliberately pads this to the paper's ambient $md/q$ loss for Lemma~\ref{lem:generalize-b}.
\end{lemma}
\begin{proof}
	\leanok
	Reindex coded line parameters to the scalar field, count the roots of the nonzero univariate polynomial $f-h$, and use $m\ge 1$ to pad $d/q$ to $md/q$.
\end{proof}

\begin{definition}[Measurements and submeasurements]\label{def:submeasurement}
	\lean{MIPStarRE.LDT.SubMeas, MIPStarRE.LDT.Measurement, MIPStarRE.LDT.ProjSubMeas}
	\leanok
	Let $\mathcal H$ be a Hilbert space and let $\mathcal A$ be a set of outcomes. A submeasurement on $\mathcal A$ is a family $A=\{A_a\}_{a \in \mathcal A}$ of Hermitian positive semidefinite operators on $\mathcal H$ such that $\sum_a A_a \le I$. It is a measurement when $\sum_a A_a = I$, and it is projective when each $A_a$ is an idempotent projection. (In the Lean formalization, outcome sets are assumed finite throughout.)
\end{definition}

\begin{definition}[Low-degree polynomial measurements]\label{def:polymeasurements}
	\lean{MIPStarRE.LDT.SubMeas, MIPStarRE.LDT.Measurement, MIPStarRE.LDT.Polynomial}
	\leanok
	\uses{def:polyfunc, def:submeasurement}
	We write $\polysub{m}{q}{d}$ for the submeasurements indexed by $\polyfunc{m}{q}{d}$, and $\polymeas{m}{q}{d}$ for the corresponding measurements.
\end{definition}

\begin{definition}[Post-processing]\label{def:post-processing}
	\lean{MIPStarRE.LDT.postprocess}
	\leanok
	\uses{def:submeasurement}
	Let $A=\{A_a\}_{a \in \mathcal A}$ be a family of operators and let $f\colon \mathcal A \to \mathcal B$. The post-processed family $A_{[f(a)=b]}$ is defined by
	\[
		A_{[f(a)=b]} = \sum_{a\, :\, f(a)=b} A_a.
	\]
\end{definition}

\begin{proposition}[Post-processing preserves measurements]\label{prop:post-processing-preserves}
	\lean{MIPStarRE.LDT.Preliminaries.postprocessPreservesMeasurements, MIPStarRE.Quantum.Submeasurement.postprocess_sum_eq, MIPStarRE.Quantum.Measurement.postprocess}
	\leanok
	\uses{def:post-processing, def:submeasurement}
	Let $A=\{A_a\}_{a \in \mathcal A}$ be a family of operators and let
	$f\colon \mathcal A \to \mathcal B$. Then
	\[
		\sum_a A_a = \sum_b A_{[f(a)=b]}.
	\]
	Consequently, if $\{A_a\}$ is a submeasurement, respectively a measurement, then
	$\{A_{[f(a)=b]}\}$ is again a submeasurement, respectively a measurement.
\end{proposition}
\begin{proof}
	\leanok
	Each $a \in \mathcal A$ contributes exactly once to the right-hand side, namely in
	the summand indexed by $b=f(a)$.
\end{proof}

\begin{remark}[Post-processing notation]\label{rem:post-processing-notation}
	In the notation $A_{[f(a)=b]}$, the measurement outcome is the variable to which
	the function is applied. For example, if $G=\{G_g\} \in \polysub{m}{q}{d}$ and
	$u \in \F_q^m$, then
	\[
		G_{[g(u)=a]} = \sum_{g\,:\,g(u)=a} G_g.
	\]
	Here $g$ is the measurement outcome, and $u$ is the evaluation point.
\end{remark}

\begin{definition}[Completion of a submeasurement]\label{def:measurement-completion}
	\lean{MIPStarRE.LDT.Preliminaries.completeAtOutcome}
	\leanok
	\uses{def:submeasurement}
	Let $A=\{A_a^x\}_{a \in \mathcal A}$ be a submeasurement. Its completion is the measurement $\widehat A$ with outcome set $\widehat{\mathcal A}=\mathcal A \cup \{\bot\}$ given by $\widehat A_a^x=A_a^x$ for $a \in \mathcal A$ and $\widehat A_\bot^x = I-\sum_a A_a^x$.
\end{definition}

\section{Consistency and state-dependent distance}
\label{sec:comparing-measurements}


\begin{definition}[Consistency]\label{def:simeq}
	\lean{MIPStarRE.LDT.ConsRel}
	\leanok
	Let $\ket{\psi} \in \mathcal H_{\mathrm A} \ot \mathcal H_{\mathrm B}$. Let $A=\{A_a^x\}$ and $B=\{B_a^x\}$ be submeasurements with the same answer set, and let $\mathcal D$ be a distribution on the question set. We write
	\[
		A_a^x \ot I \simeq_\delta I \ot B_a^x
	\]
	when
	\begin{equation}\label{eq:no-big-Oh}
		\E_{x \sim \mathcal D} \sum_{a \ne b} \bra{\psi} A_a^x \ot B_b^x \ket{\psi} \le \delta.
	\end{equation}
\end{definition}

\begin{lemma}[Characterization of a good symmetric strategy]\label{lem:good-strategy-characterization}
	\lean{MIPStarRE.LDT.Preliminaries.goodStrategyCharacterization}
	\leanok
	\uses{def:good-strategy, def:simeq, prop:simeq-for-measurements}
	A symmetric projective strategy $(\psi,A,B,L)$ is $(\eps,\delta,\gamma)$-good if and only if the following hold on the relevant test distributions:
	\[
		A_a^u \ot I \simeq_\eps I \ot B^{\ell}_{[f(u)=a]},
		\qquad
		A_a^u \ot I \simeq_\delta I \ot A_a^u,
		\qquad
		A_a^u \ot I \simeq_\gamma I \ot L^{\ell}_{[f(u)=a]}.
	\]
\end{lemma}
\begin{proof}
	\leanok
	Each subtest accepts exactly when the two measurement outcomes agree. Because the relevant families are measurements, Lemma~\ref{prop:simeq-for-measurements} rewrites those acceptance probabilities as consistency bounds.
\end{proof}

\begin{lemma}[Consistency for measurements]\label{prop:simeq-for-measurements}
	\lean{MIPStarRE.LDT.Preliminaries.simeqForMeasurements}
	\leanok
	\uses{def:simeq}
	If $A=\{A_a^x\}$ and $B=\{B_a^x\}$ are measurements, then
	\[
		A_a^x \ot I \simeq_\delta I \ot B_a^x
		\qquad\text{if and only if}\qquad
		\E_x \sum_a \bra{\psi} A_a^x \ot B_a^x \ket{\psi} \ge 1-\delta.
	\]
\end{lemma}
\begin{proof}
	\leanok
	Expand the off-diagonal sum by writing $\sum_{a \ne b} A_a^x \ot B_b^x = \sum_b (I-A_b^x) \ot B_b^x$.
\end{proof}

\begin{definition}[State-dependent distance]\label{def:approx_delta}
	\lean{MIPStarRE.LDT.SDDRel}
	\leanok
	Let $\ket{\psi} \in \mathcal H$, let $A=\{A_a^x\}$ and $B=\{B_a^x\}$ be families of operators on $\mathcal H$, and let $\mathcal D$ be a distribution on the question set. We write
	\[
		A_a^x \approx_\delta B_a^x
	\]
	when
	\[
		\E_{x \sim \mathcal D} \sum_a \lVert (A_a^x-B_a^x)\ket{\psi} \rVert^2 \le \delta.
	\]
\end{definition}

In particular, if $A^x_a \ot I \simeq_{\eps} I \ot C^x_a$, the question is what hypothesis on $A$ and $B$ allows one to conclude that
\begin{equation}\label{eq:can-we-use-approx-delta-to-derive-this}
	B^x_a \ot I \simeq_{\eps} I \ot C^x_a.
\end{equation}

\begin{lemma}[Transferring consistency through state-dependent distance]\label{prop:triangle-sub}
	\lean{MIPStarRE.LDT.Preliminaries.triangleSub_heterogeneous}
	\leanok
	\uses{def:simeq, def:approx_delta, prop:simeq-for-measurements}
	Let $A=\{A_a^x\}$ and $B=\{B_a^x\}$ be measurements and let $C=\{C_a^x\}$ be a submeasurement. If $A_a^x \ot I \simeq_\delta I \ot C_a^x$ and $A_a^x \ot I \approx_\eps B_a^x \ot I$, then
	\[
		B_a^x \ot I \simeq_{\delta+\sqrt\eps} I \ot C_a^x.
	\]
\end{lemma}
\begin{proof}
	\leanok
	Rewrite the inconsistency with $C$ as total mass minus diagonal overlap. The total mass is unchanged because $A$ and $B$ are measurements, and the diagonal overlap changes by at most $\sqrt\eps$ by Cauchy--Schwarz.
\end{proof}

\begin{lemma}[Consistency implies state-dependent distance for measurements]\label{prop:simeq-to-approx}
	\lean{MIPStarRE.LDT.Preliminaries.simeqToApprox_heterogeneous, MIPStarRE.LDT.Preliminaries.approxToSimeq_heterogeneous}
	\leanok
	\uses{def:simeq, def:approx_delta, prop:simeq-for-measurements}
	If $A$ and $B$ are measurements and $A_a^x \ot I \simeq_\delta I \ot B_a^x$, then
	\[
		A_a^x \ot I \approx_{2\delta} I \ot B_a^x.
	\]
	If both measurements are projective, the converse also holds: if $A_a^x \ot I \approx_{2\delta} I \ot B_a^x$ then $A_a^x \ot I \simeq_\delta I \ot B_a^x$.
\end{lemma}
\begin{proof}
	\leanok
	Expand the squared norm of $(A_a^x \ot I - I \ot B_a^x)\ket{\psi}$ and use the diagonal-overlap formula from Lemma~\ref{prop:simeq-for-measurements}; projectivity turns the diagonal inequality into an equality, giving the converse.
\end{proof}

\begin{remark}\label{rem:simeq-not-approx-submeas}
	The implication above does not extend to submeasurements. For example, if
	$A_a^x=0$ for all $a$, then $A_a^x \ot I \simeq_0 I \ot B_a^x$, but
	\[
		\E_{\bx} \sum_a \lVert (A^{\bx}_a \ot I - I \ot B^{\bx}_a)\ket{\psi} \rVert^2
		= \E_{\bx} \sum_a \lVert (I \ot B^{\bx}_a)\ket{\psi} \rVert^2,
	\]
	which is nonzero unless $(I \ot B_a^x)\ket{\psi}=0$ for all $x$ and $a$.
\end{remark}

\section{Consistency from state-dependent distance}
\label{sec:consistency-from-state-dependent-distance}

\begin{theorem}[Closeness of inner products]\label{prop:closeness-of-ip}
	\lean{MIPStarRE.LDT.Preliminaries.closenessOfIP, MIPStarRE.LDT.Preliminaries.closenessOfIPAdjoint}
	\leanok
	\uses{def:approx_delta}
	Let $\{A^x_a\}$, $\{B^x_a\}$, and $\{C^x_{a,b}\}$ be matrices. Suppose that $A^x_a \approx_\gamma B^x_a$ and that for all~$x$,
	\[
		\sum_a \Big(\sum_b C^x_{a,b}\Big) \Big(\sum_b C^x_{a,b}\Big)^\dagger \le I.
	\]
	Then
	\begin{equation}
		\E_{\bx} \sum_{a,b} \bra{\psi} C^{\bx}_{a,b} A^{\bx}_a  \ket{\psi} \approx_{\sqrt{\gamma}} \E_{\bx} \sum_{a,b} \bra{\psi} C^{\bx}_{a,b} B^{\bx}_a  \ket{\psi}. \label{eq:closeness3}
	\end{equation}
	Similarly, suppose that $(A^x_a)^\dagger \approx_\gamma (B^x_a)^\dagger$ and that for all~$x$,
	\[
		\sum_a \Big(\sum_b C^x_{a,b}\Big)^\dagger \Big(\sum_b C^x_{a,b}\Big) \le I.
	\]
	Then
	\begin{equation}
		\E_{\bx} \sum_{a,b} \bra{\psi} A^{\bx}_a  C^{\bx}_{a,b} \ket{\psi} \approx_{\sqrt{\gamma}} \E_{\bx} \sum_{a,b} \bra{\psi} B^{\bx}_a  C^{\bx}_{a,b}  \ket{\psi}. \label{eq:closeness4}
	\end{equation}
\end{theorem}
\begin{proof}
	\leanok
	To prove \eqref{eq:closeness3}, write the difference as
	\[
		\E_{\bx} \sum_a \bra{\psi} \Big(\sum_b C^\bx_{a,b}\Big) (A^\bx_a-B^\bx_a) \ket{\psi}.
	\]
	Cauchy--Schwarz bounds its magnitude by
	\[
		\Big( \E_{\bx} \sum_a \bra{\psi} \Big(\sum_b C^{\bx}_{a,b}\Big)\Big(\sum_b C^{\bx}_{a,b}\Big)^\dagger \ket{\psi} \Big)^{1/2}
		\Big( \E_{\bx} \sum_a \bra{\psi} (A^{\bx}_a - B^{\bx}_a)^\dagger (A^{\bx}_a - B^{\bx}_a) \ket{\psi} \Big)^{1/2},
	\]
	and the two hypotheses bound these factors by $1$ and $\sqrt{\gamma}$ respectively. For \eqref{eq:closeness4}, rewrite the difference as
	\[
		\E_{\bx} \sum_{a,b} \bra{\psi} (C^{\bx}_{a,b})^\dagger \big((A^{\bx}_a)^\dagger-(B^{\bx}_a)^\dagger\big) \ket{\psi},
	\]
	and apply \eqref{eq:closeness3} to the adjoint families.
\end{proof}

\begin{theorem}\label{prop:easy-approx-from-approx-delta}
	\lean{MIPStarRE.LDT.Preliminaries.easyApproxFromApproxDelta}
	\leanok
	\uses{def:submeasurement, def:approx_delta}
	Let $A = \{A^x_a\}$, $B = \{B^x_a\}$, and $C = \{C^x_a\}$ be submeasurements such that $A^x_a \approx_{\delta} B^x_a$. Then
	\begin{equation*}
		\E_{\bx} \sum_a \bra{\psi} A^{\bx}_a  C^{\bx}_a \ket{\psi}
		\approx_{\sqrt{\delta}} \E_{\bx} \sum_a \bra{\psi} B^{\bx}_a  C^{\bx}_a \ket{\psi}.
	\end{equation*}
\end{theorem}
\begin{proof}
	\leanok
	The difference has magnitude
	\[
		\Big| \E_{\bx} \sum_a \bra{\psi} (A^{\bx}_a - B^{\bx}_a) C^{\bx}_a \ket{\psi}\Big|.
	\]
	By Cauchy--Schwarz this is at most
	\[
		\Big(\E_{\bx} \sum_a \bra{\psi} (A^{\bx}_a - B^{\bx}_a)^2 \ket{\psi}\Big)^{1/2}
		\Big(\E_{\bx} \sum_a \bra{\psi} (C^{\bx}_a)^2 \ket{\psi}\Big)^{1/2}.
	\]
	The first factor is at most $\sqrt{\delta}$ by hypothesis, and the second is at most $1$ because $C$ is a submeasurement.
\end{proof}

\begin{theorem}\label{prop:cab-approx-delta}
	\lean{MIPStarRE.LDT.Preliminaries.cabApproxDelta}
	\leanok
	\uses{def:approx_delta}
	Let $\{A^x_a\}$, $\{B^x_a\}$, and $\{C^x_{a,b}\}$ be matrices. Suppose that $A^{x}_a \approx_\delta B^{x}_a$ and that for all $x$ and $a$,
	\[
		\sum_b (C^{x}_{a,b})^\dagger (C^{x}_{a,b}) \leq I.
	\]
	Then
	\[
		C^{x}_{a,b} A^x_{a} \approx_{\delta} C^{x}_{a,b} B^{x}_a.
	\]
\end{theorem}
\begin{proof}
	\leanok
	The error term is
	\[
		\E_{\bx} \sum_{a,b} \bra{\psi}
		(A^{\bx}_a - B^{\bx}_a)^\dagger
		(C^{\bx}_{a,b})^\dagger C^{\bx}_{a,b}
		(A^{\bx}_a - B^{\bx}_a) \ket{\psi}.
	\]
	Using $\sum_b (C^{x}_{a,b})^\dagger (C^{x}_{a,b}) \le I$, this is bounded by
	\[
		\E_{\bx} \sum_a \bra{\psi}
		(A^{\bx}_a - B^{\bx}_a)^\dagger
		(A^{\bx}_a - B^{\bx}_a) \ket{\psi}
		\le \delta.
	\]
\end{proof}

\begin{lemma}[Triangle inequality for vectors squared]\label{prop:triangle-inequality-for-vectors-squared}
	\lean{MIPStarRE.LDT.Preliminaries.triangleInequalityForVectorsSquared}
	\leanok
	For vectors $\ket{\psi_1},\dots,\ket{\psi_k}$,
	\[
		\bigl\lVert \ket{\psi_1}+\cdots+\ket{\psi_k} \bigr\rVert^2 \le k(\lVert\ket{\psi_1}\rVert^2+\cdots+\lVert\ket{\psi_k}\rVert^2).
	\]
\end{lemma}
\begin{proof}
	\leanok
	First,
	\begin{equation}\label{eq:prop-for-real-numbers}
		(x_1+\cdots+x_k)^2 \le k(x_1^2+\cdots+x_k^2)
	\end{equation}
	for all real numbers $x_1,\dots,x_k$. Apply \eqref{eq:prop-for-real-numbers} to $x_i=\lVert\ket{\psi_i}\rVert$ and combine it with the triangle inequality.
\end{proof}

\begin{lemma}[Triangle inequality for $\approx_\delta$]\label{prop:triangle-inequality-for-approx_delta}
	\lean{MIPStarRE.LDT.Preliminaries.sddOpRel_chain}
	\leanok
	\uses{def:approx_delta, prop:triangle-inequality-for-vectors-squared}
	Suppose $(A_i)_a^x \approx_{\delta_i} (A_{i+1})_a^x$ for $i=1,\dots,k$. Then
	\[
		(A_1)_a^x \approx_{k(\delta_1+\cdots+\delta_k)} (A_{k+1})_a^x.
	\]
\end{lemma}
\begin{proof}
	\leanok
	Expand $(A_1-A_{k+1})\ket{\psi}$ as a telescoping sum and apply Lemma~\ref{prop:triangle-inequality-for-vectors-squared}.
\end{proof}

\begin{remark}[Lean finite-step triangle helpers]\label{rem:lean-finite-step-triangle-helpers}
	The general operator-family theorem above is complemented in Lean by reusable
	binary and three-step helper lemmas for operator expectations, `qSDD`, and
	`SDDRel`.  The three-step `SDDRel` helper is the specialization used in the
	Step~6 projectivization chain in Theorem~\ref{thm:main-formal}; this remark is
	only a cross-reference and makes no separate proof-level completeness claim.
\end{remark}

\begin{lemma}[Triangle inequality for $\simeq_\delta$]\label{prop:simeq-triangle-inequality}
	\lean{MIPStarRE.LDT.Preliminaries.simeqTriangleInequality_heterogeneous}
	\leanok
	\uses{prop:simeq-to-approx, prop:triangle-inequality-for-approx_delta, prop:triangle-sub}
	Suppose $A_a^x \ot I \simeq_\eps I \ot B_a^x$, $C_a^x \ot I \simeq_\delta I \ot B_a^x$, and $C_a^x \ot I \simeq_\gamma I \ot D_a^x$, where all four families are measurements. Then
	\[
		A_a^x \ot I \simeq_{\eps + 2\sqrt{\delta+\gamma}} I \ot D_a^x.
	\]
\end{lemma}
\begin{proof}
	\leanok
	Convert the two consistencies involving $C$ to state-dependent distance, compose them by the triangle inequality, and transfer the result back to consistency by Lemma~\ref{prop:triangle-sub}.
\end{proof}

\begin{lemma}[Data processing for consistency]\label{prop:simeq-data-processing}
	\lean{MIPStarRE.LDT.Preliminaries.simeqDataProcessing_heterogeneous}
	\leanok
	\uses{def:post-processing, def:simeq}
	If $A=\{A_a^x\}$ and $B=\{B_a^x\}$ are measurements such that $A_a^x \ot I \simeq_\delta I \ot B_a^x$, and $f$ is a map on the answer set, then
	\[
		A_{[f(a)=b]}^x \ot I \simeq_\delta I \ot B_{[f(a)=b]}^x.
	\]
\end{lemma}
\begin{proof}
	\leanok
	Merging outcome classes can only decrease the off-diagonal mass.
\end{proof}

\begin{lemma}[Consistency controls a submeasurement]\label{prop:cons-sub-meas}
	\lean{MIPStarRE.LDT.Preliminaries.consSubMeas_heterogeneous}
	\leanok
	\uses{def:simeq, def:approx_delta, prop:triangle-inequality-for-approx_delta}
	Let $A=\{A_a^x\}$ be a submeasurement and $B=\{B_a^x\}$ a measurement such that $A_a^x \ot I \simeq_\gamma I \ot B_a^x$. Then
	\begin{equation}
		A_a^x \ot I \approx_\gamma A_a^x \ot B_a^x \approx_\gamma A^x \ot B_a^x,
		\label{eq:closeness5}
	\end{equation}
	where $A^x=\sum_a A_a^x$. In particular,
	\[
		A_a^x \ot I \approx_{4\gamma} A^x \ot B_a^x.
	\]
\end{lemma}
\begin{proof}
	\leanok
	For the first approximation in \eqref{eq:closeness5}, bound
	\[
		\E_{\bx} \sum_a \bra{\psi} \bigl(A^{\bx}_a \ot (I-B^{\bx}_a)\bigr)^2 \ket{\psi}
	\]
	by the inconsistency between $A$ and $B$. The second approximation is identical, applied to
	\[
		\E_{\bx} \sum_a \bra{\psi} \bigl((A^\bx-A^\bx_a)\ot B^\bx_a\bigr)^2 \ket{\psi}.
	\]
	Then compose the two bounds with Lemma~\ref{prop:triangle-inequality-for-approx_delta}.
\end{proof}

\begin{lemma}[Switch-sandwich]\label{prop:switch-sandwich}
	\lean{MIPStarRE.LDT.Preliminaries.switchSandwich}
	\leanok
	\uses{def:approx_delta}
	Suppose $A=\{A_a^x\}$ is a projective submeasurement satisfying
	\begin{equation}
		A_a^x \ot I \approx_\delta I \ot A_a^x.
		\label{eq:Aapproxd}
	\end{equation}
	Then for every operator $B$ with $0 \le B \le I$,
	\begin{equation}
		\E_x \sum_a \bra{\psi} A_a^x B A_a^x \ot I \ket{\psi}
		\approx_{2\sqrt\delta}
		\E_x \sum_a \bra{\psi} B \ot A_a^x \ket{\psi}
		\approx_{\sqrt\delta}
		\E_x \sum_a \bra{\psi} BA_a^x \ot I \ket{\psi}.
		\label{eq:switch-sandwich}
	\end{equation}
\end{lemma}
\begin{proof}
	\leanok
	We prove the first approximation in \eqref{eq:switch-sandwich} in two steps. First,
	\begin{equation}
		\E_{\bx} \sum_a \bra{\psi} A^{\bx}_a B A^{\bx}_a \ot I \ket{\psi}
		\approx_{\sqrt{\delta}}
		\E_{\bx} \sum_a \bra{\psi} A^{\bx}_a B \ot A^{\bx}_a \ket{\psi}.
		\label{eq:shift-right-A}
	\end{equation}
	This is obtained by applying Cauchy--Schwarz to the difference and using $B^2 \le B \le I$ together with \eqref{eq:Aapproxd}. For the second step, move the remaining left copy of $A^\bx_a$ across the bipartition in the same way; projectivity then collapses the sandwich and yields the first approximation in \eqref{eq:switch-sandwich}. The second approximation in \eqref{eq:switch-sandwich} is the same argument with the rightmost factor $A_a^\bx$ omitted.
\end{proof}

\section{Strong self-consistency}

\begin{lemma}[Completeness transfer to a projective approximation]\label{prop:completeness-transfer-projective-P}
	\lean{MIPStarRE.LDT.Preliminaries.completenessTransferProjectiveP}
	\leanok
	\uses{def:approx_delta, prop:easy-approx-from-approx-delta}
	Let $A=\{A_a^x\}$ be a submeasurement and $P=\{P_a^x\}$ a projective submeasurement such that $A_a^x \ot I \approx_\eps P_a^x \ot I$. Then
	\[
		\bra{\psi} A \ot I \ket{\psi} \ge \bra{\psi} P \ot I \ket{\psi} - 2\sqrt\eps.
	\]
\end{lemma}
\begin{proof}
	\leanok
	Since $P$ is projective,
	\[
		\bra{\psi} P \ot I \ket{\psi}
		= \E_{\bx} \sum_a \bra{\psi} (P^{\bx}_a)^2 \ot I \ket{\psi}.
	\]
	Apply Proposition~\ref{prop:easy-approx-from-approx-delta} twice to replace $(P^\bx_a)^2$ first by $A^\bx_aP^\bx_a$ and then by $(A^\bx_a)^2$, and conclude with $(A_a^\bx)^2 \le A_a^\bx$.
\end{proof}

\begin{definition}[Strong self-consistency]\label{def:strong-self-consistency}
	\lean{MIPStarRE.LDT.SSCRel, MIPStarRE.LDT.BipartiteSSCRel}
	\leanok
	Let $\ket{\psi}$ be permutation-invariant and let $A=\{A_a^x\}$ be a submeasurement. We say that $A$ is $\delta$-strongly self-consistent when
	\[
		\E_x \sum_a \bra{\psi} A_a^x \ot A_a^x \ket{\psi}
		\ge
		\bra{\psi} A \ot I \ket{\psi} - \delta.
	\]
\end{definition}

\begin{lemma}[Strong self-consistency implies ordinary self-consistency]\label{prop:other-two-notions-of-self-consistency}
	\lean{MIPStarRE.LDT.Preliminaries.otherTwoNotionsOfSelfConsistency, MIPStarRE.LDT.Preliminaries.bipartiteSSCRel_iff_consRel_self_measurement}
	\leanok
	\uses{def:strong-self-consistency, def:simeq}
	If $A$ is $\delta$-strongly self-consistent, then $A_a^x \ot I \simeq_\delta I \ot A_a^x$. If $A$ is a measurement, the converse also holds.
\end{lemma}
\begin{proof}
	\leanok
	Rewrite the off-diagonal mass as the difference between the total mass of $A$ and the diagonal overlap in the strong self-consistency inequality.
	When $A$ is a full measurement, the total mass is the identity, so this inequality becomes an equality; hence the ordinary self-consistency defect and the strong self-consistency defect coincide, giving the converse.
\end{proof}

\begin{lemma}[Strong self-consistency and state-dependent distance]\label{prop:two-notions-of-self-consistency}
	\lean{MIPStarRE.LDT.Preliminaries.twoNotionsOfSelfConsistency}
	\leanok
	\uses{def:strong-self-consistency, def:approx_delta}
	If $A$ is $\delta$-strongly self-consistent, then
	\[
		A_a^x \ot I \approx_{2\delta} I \ot A_a^x.
	\]
	If $A$ is projective, the converse also holds.
\end{lemma}
\begin{proof}
	\leanok
	Expanding the squared norm gives
	\begin{align}
		      & \E_{\bx} \sum_a \lVert(A^{\bx}_a \otimes I - I \otimes A^{\bx}_a)\ket{\psi}\rVert^2 \nonumber                                                                  \\
		=~    & 2\cdot \Big( \E_{\bx} \sum_a \bra{\psi} (A^{\bx}_a)^2 \otimes I \ket{\psi} -  \E_{\bx} \sum_a \bra{\psi} A^{\bx}_a \otimes A^{\bx}_a \ket{\psi}\Big) \nonumber \\
		\leq~ & 2\cdot \Big( \E_{\bx} \sum_a \bra{\psi} A^{\bx}_a \otimes I \ket{\psi} -  \E_{\bx} \sum_a \bra{\psi} A^{\bx}_a \otimes A^{\bx}_a \ket{\psi}\Big).
		\label{eq:here's-where-projectivity-would-help}
	\end{align}
	The strong self-consistency bound controls the last line by $2\delta$. If $A$ is projective, then \eqref{eq:here's-where-projectivity-would-help} is an equality, which gives the converse.
\end{proof}

\begin{lemma}[Post-processing preserves strong self-consistency]\label{prop:two-notions-of-self-consistency-after-evaluation}
	\lean{MIPStarRE.LDT.Preliminaries.twoNotionsOfSelfConsistencyAfterEvaluation}
	\leanok
	\uses{def:post-processing, def:strong-self-consistency, prop:two-notions-of-self-consistency}
	If $A$ is $\delta$-strongly self-consistent and $f$ is a function on its answer set, then
	\[
		A_{[f(a)=b]}^x \ot I \approx_{2\delta} I \ot A_{[f(a)=b]}^x.
	\]
\end{lemma}
\begin{proof}
	\leanok
	The same computation as in Lemma~\ref{prop:two-notions-of-self-consistency} gives
	\begin{align}
		      & \E_{\bx} \sum_b \lVert(A^{\bx}_{[f(a)=b]} \otimes I - I \otimes A^{\bx}_{[f(a)=b]}) \ket{\psi}\rVert^2 \nonumber                            \\
		\leq~ & 2\cdot \Big( \bra{\psi} A \otimes I \ket{\psi} -  \E_{\bx} \sum_b \bra{\psi} A^{\bx}_{[f(a)=b]} \otimes A^{\bx}_{[f(a)=b]} \ket{\psi}\Big).
		\label{eq:finishing-this-up}
	\end{align}
	The post-processed diagonal term dominates $\E_{\bx} \sum_a \bra{\psi} A^\bx_a \otimes A^\bx_a \ket{\psi}$, so strong self-consistency bounds \eqref{eq:finishing-this-up} by $2\delta$.
\end{proof}

\begin{lemma}[Completeness transfer for a strongly self-consistent approximation]\label{prop:completeness-transfer-self-consistent-A}
	\lean{MIPStarRE.LDT.Preliminaries.completenessTransferSelfConsistentA}
	\leanok
	\uses{def:strong-self-consistency, def:approx_delta, prop:easy-approx-from-approx-delta}
	Let $A$ be a $\delta$-strongly self-consistent submeasurement and let $B$ be a submeasurement such that $A_a^x \ot I \approx_\eps B_a^x \ot I$. Then
	\[
		\bra{\psi} B \ot I \ket{\psi} \ge \bra{\psi} A \ot I \ket{\psi} - \delta - 2\sqrt\eps.
	\]
\end{lemma}
\begin{proof}
	\leanok
	Bound $\bra{\psi} B \ot I \ket{\psi}$ from below by $\E_{\bx}\sum_a \bra{\psi} B^\bx_a \ot B^\bx_a \ket{\psi}$. Then apply Proposition~\ref{prop:easy-approx-from-approx-delta} twice to replace this by $\E_{\bx}\sum_a \bra{\psi} A^\bx_a \ot A^\bx_a \ket{\psi}$, and finish with strong self-consistency of $A$.
\end{proof}

\begin{lemma}[Self-consistency implies data-processing closeness]\label{prop:self-consistency-implies-data-processing}
	\lean{MIPStarRE.LDT.Preliminaries.selfConsistencyImpliesDataProcessing}
	\leanok
	\uses{def:post-processing, def:strong-self-consistency, def:approx_delta, prop:two-notions-of-self-consistency-after-evaluation, prop:completeness-transfer-projective-P, prop:easy-approx-from-approx-delta, prop:triangle-inequality-for-approx_delta}
	Let $A$ be a $\delta$-strongly self-consistent submeasurement and let $P$ be a projective submeasurement such that $P_a^x \ot I \approx_\eps A_a^x \ot I$. Then for every function $f$,
	\[
		P_{[f(a)=b]}^x \ot I \approx_{8\delta + 8\sqrt\eps} A_{[f(a)=b]}^x \ot I.
	\]
\end{lemma}
\begin{proof}
	\leanok
	First prove
	\begin{equation}
		P^x_{[f(a)=b]} \ot I \approx_{2\delta + 4\sqrt{\eps}} I \ot A_{[f(a)=b]}^x.
		\label{eq:what-we-want-to-prove-but-on-wrong-side}
	\end{equation}
	Expanding the squared norm and using projectivity gives
	\begin{align}
		      & \E_{\bx} \sum_b \lVert(P^{\bx}_{[f(a)=b]} \ot I - I \ot A_{[f(a)=b]}^{\bx}) \ket{\psi} \rVert^2 \nonumber \\
		\leq~ & \bra{\psi} P \ot I \ket{\psi} + \bra{\psi} I \ot A \ket{\psi}
		- 2\cdot \E_{\bx} \sum_b \bra{\psi}  P^{\bx}_{[f(a)=b]} \ot A^{\bx}_{[f(a)=b]} \ket{\psi}.
		\label{eq:gonna-handle-third-term}
	\end{align}
	Proposition~\ref{prop:completeness-transfer-projective-P} bounds the first term by $\bra{\psi} A \ot I \ket{\psi}+2\sqrt{\eps}$, and Proposition~\ref{prop:easy-approx-from-approx-delta} together with strong self-consistency bounds the third term below by $\bra{\psi} A \ot I \ket{\psi}-\delta-\sqrt{\eps}$. Substituting into \eqref{eq:gonna-handle-third-term} yields \eqref{eq:what-we-want-to-prove-but-on-wrong-side}. Now combine \eqref{eq:what-we-want-to-prove-but-on-wrong-side} with
	\[
		A^x_{[f(a)=b]} \ot I \approx_{2\delta} I \ot A^x_{[f(a)=b]}
	\]
	from Lemma~\ref{prop:two-notions-of-self-consistency-after-evaluation}, and conclude by Lemma~\ref{prop:triangle-inequality-for-approx_delta}.
\end{proof}

\begin{lemma}[Self-consistency moves one copy to the same side]\label{lem:self-consistency-same-side-square}
	\lean{MIPStarRE.LDT.Preliminaries.bipartiteSSCSquaredMass}
	\leanok
	\uses{def:strong-self-consistency}
	Let $A=\{A_a\}$ be a $\zeta$-strongly self-consistent submeasurement. Then
	\[
		\sum_a \bra{\psi} A_a^2 \ot I \ket{\psi}
		\ge
		\sum_a \bra{\psi} A_a \ot I \ket{\psi} - \zeta.
	\]
\end{lemma}
\begin{proof}
	\leanok
	Apply Cauchy--Schwarz to the two-sided overlap $\sum_a \bra{\psi} A_a \ot A_a \ket{\psi}$ and compare it with the strong self-consistency lower bound.
\end{proof}

\begin{lemma}[Bounding the missing mass of a close submeasurement]\label{lem:completion-missing-mass-bound}
	\lean{MIPStarRE.LDT.Preliminaries.completionMissingMassBound}
	\leanok
	\uses{def:strong-self-consistency, def:approx_delta, lem:self-consistency-same-side-square}
	Let $A=\{A_a\}$ be a measurement that is $\zeta$-strongly self-consistent, and let $B=\{B_a\}$ be a submeasurement such that $A_a \ot I \approx_\delta B_a \ot I$. Writing $B=\sum_a B_a$, one has
	\[
		\lVert (I-B) \ot I\ket{\psi} \rVert^2 \le 2\sqrt\delta + \zeta.
	\]
\end{lemma}
\begin{proof}
	\leanok
	Since $0 \le I-B \le I$, the left-hand side is at most $1-\bra{\psi} B \ot I \ket{\psi}$. Compare $\sum_a \bra{\psi} B_a^2 \ot I \ket{\psi}$ with $\sum_a \bra{\psi} A_a B_a \ot I \ket{\psi}$ and then with $\sum_a \bra{\psi} A_a^2 \ot I \ket{\psi}$ using the $\approx_\delta$ hypothesis, and finish with Lemma~\ref{lem:self-consistency-same-side-square}.
\end{proof}

\begin{lemma}[Completing a close submeasurement to a full measurement]\label{prop:completing-to-measurement}
	\lean{MIPStarRE.LDT.Preliminaries.completingToMeasurement}
	\leanok
	\uses{def:measurement-completion, def:strong-self-consistency, def:approx_delta, prop:triangle-inequality-for-vectors-squared, prop:easy-approx-from-approx-delta, prop:cool-prop}
	Let $A=\{A_a\}$ be a measurement that is $\zeta$-strongly self-consistent, let $B=\{B_a\}$ be a submeasurement such that $A_a \ot I \approx_\delta B_a \ot I$, and let $C$ be the measurement obtained by adding the missing mass $I-B$ to one distinguished answer $a^*$. Then
	\[
		A_a \ot I \approx_{2\delta + 4\sqrt\delta + 2\zeta} C_a \ot I.
	\]
\end{lemma}
\begin{proof}[Proof of \ref{prop:completing-to-measurement}]
	\proves{prop:completing-to-measurement}
	\leanok
	\uses{prop:triangle-inequality-for-vectors-squared, prop:easy-approx-from-approx-delta, prop:cool-prop}
	By Lemma~\ref{prop:triangle-inequality-for-vectors-squared},
	\[
		\sum_a \lVert(A_a-C_a)\ot I \ket{\psi} \rVert^2
		\leq 2\delta + 2 \cdot \lVert(I-B)\ot I \ket{\psi}\rVert^2.
	\]
	It remains to bound the residual term. Since $0 \le I-B \le I$,
	\begin{align}
		\lVert(I-B) \ot I \ket{\psi}\rVert^2
		 & \leq 1 -\sum_a \bra{\psi} B_a^2 \ot I \ket{\psi}.
		\label{eq:to-return-later-whatevs}
	\end{align}
	By Proposition~\ref{prop:easy-approx-from-approx-delta},
	\[
		\sum_a \bra{\psi} B_a^2 \ot I \ket{\psi}
		\approx_{\sqrt{\delta}} \sum_a \bra{\psi} (A_a B_a)\ot I \ket{\psi}
		\approx_{\sqrt{\delta}} \sum_a \bra{\psi} A_a^2\ot I \ket{\psi}.
	\]
	Proposition~\ref{prop:cool-prop} gives
	\[
		\sum_a \bra{\psi} A_a^2\ot I \ket{\psi} \geq \bra{\psi} A \ot I \ket{\psi} - \zeta = 1-\zeta.
	\]
	Hence \eqref{eq:to-return-later-whatevs} implies
	\[
		\lVert(I-B) \ot I \ket{\psi}\rVert^2 \le 2\sqrt\delta + \zeta.
	\]
	Substituting this into the first bound proves the claim.
\end{proof}

% prop:cool-prop from the paper is the same result as lem:self-consistency-same-side-square above.
% We keep the label for cross-reference fidelity.
\begin{theorem}\label{prop:cool-prop}
	\lean{MIPStarRE.LDT.Preliminaries.bipartiteSSCSquaredMass}
	\leanok
	\uses{lem:self-consistency-same-side-square}
	This is the same statement as Lemma~\ref{lem:self-consistency-same-side-square}: for a $\zeta$-strongly self-consistent submeasurement $A$,
	\[
		\sum_a \bra{\psi} (A_a)^2 \ot I \ket{\psi} \geq \sum_a \bra{\psi} A_a \ot I \ket{\psi} -\zeta.
	\]
\end{theorem}
\begin{proof}
	\leanok
	\uses{lem:self-consistency-same-side-square}
	Immediate from Lemma~\ref{lem:self-consistency-same-side-square}.
\end{proof}


%%FILE blueprint/src/chapter/ch04_projective.tex
\chapter{Making measurements projective}
\label{chap:projective}

% Source background theorem. The orthonormalization route below is the one used
% downstream for state-dependent consistency estimates.
\begin{theorem}[Naimark dilation]\label{thm:naimark}
	\lean{MIPStarRE.LDT.MakingMeasurementsProjective.NaimarkTensorProductCorrelationStatement}
	\lean{MIPStarRE.LDT.MakingMeasurementsProjective.naimarkTensorProductCorrelation}
	\leanok
	% The Lean declaration constructs the Naimark data from the checked
	% one-measurement theorem and the two-sided trace identity
	% OneMeasNaimarkData.twoSidedCorrelationPreservation.
	\uses{def:submeasurement}
	Let $\ket{\psi}$ be a state in $\mathcal H_{\mathrm A} \ot \mathcal H_{\mathrm B}$.
	Let $A = \{A^x_a\}$ be a sub-measurement acting on $\mathcal H_{\mathrm A}$
	and $B=\{B^y_b\}$ be a sub-measurement acting on $\mathcal H_{\mathrm B}$.
	Then there exist Hilbert spaces $\mathcal H_{\mathrm A_{\mathsf{aux}}}$ and $\mathcal H_{\mathrm B_{\mathsf{aux}}}$,
	a state $\ket{\mathsf{aux}} \in \mathcal H_{\mathrm A_{\mathsf{aux}}} \ot \mathcal H_{\mathrm B_{\mathsf{aux}}}$,
	and projective sub-measurements $\widehat{A} = \{\widehat{A}^x_a\}$ and $\widehat{B}=\{\widehat{B}^y_b\}$
	acting on $\mathcal H_{\mathrm A} \ot \mathcal H_{\mathrm A_{\mathsf{aux}}}$
	and $\mathcal H_{\mathrm B} \ot \mathcal H_{\mathrm B_{\mathsf{aux}}}$, respectively,
	such that, if $\ket{\widehat{\psi}} = \ket{\psi} \ot \ket{\mathsf{aux}}$, then for all $x,y,a,b$,
	\[
		\bra{\psi} A^x_a \ot B^y_b \ket{\psi}
		=
		\bra{\widehat{\psi}} \widehat{A}^x_a \ot \widehat{B}^y_b \ket{\widehat{\psi}}.
	\]
	In addition, $\ket{\mathsf{aux}}$ is a product state.
	The word ``measurement'' in the source statement is read here in the
	projective-submeasurement sense used by Lemma~\ref{lem:naimark-helper}: for a
	general sub-measurement, the residual mass is carried by the additional
	\(\bot\) outcome in the auxiliary construction, so the original outcomes need
	not sum to the identity.
\end{theorem}
\begin{proof}
	\uses{lem:naimark-helper}
	\leanok

	For each question $x$, choose a local auxiliary state
	$\ket{\mathsf{aux}_{\mathrm A,x}}$ of dimension one larger than the number of
	outcomes of $A^x$, and apply Lemma~\ref{lem:naimark-helper} to obtain a
	projective sub-measurement $\widetilde A^x$. Do the same for each question $y$
	on Bob's side, obtaining auxiliary states $\ket{\mathsf{aux}_{\mathrm B,y}}$
	and projective sub-measurements $\widetilde B^y$.

	Let
	\[
		\ket{\mathsf{aux}}
		=
		\Big(\bigotimes_x \ket{\mathsf{aux}_{\mathrm A,x}}\Big)
		\ot
		\Big(\bigotimes_y \ket{\mathsf{aux}_{\mathrm B,y}}\Big),
		\qquad
		\ket{\widehat{\psi}} = \ket{\psi} \ot \ket{\mathsf{aux}},
	\]
	and let $\widehat A^x_a$ and $\widehat B^y_b$ be obtained by letting
	$\widetilde A^x_a$ or $\widetilde B^y_b$ act on the relevant auxiliary factor
	and the identity act on every other auxiliary factor. Because the auxiliary
	state is a product state and the two dilations act on disjoint tensor factors,
	the correlation for fixed $(x,y)$ reduces to the compression identity from
	Lemma~\ref{lem:naimark-helper} on the $x$- and $y$-registers. Hence
	\[
		\bra{\widehat{\psi}} \widehat A^x_a \ot \widehat B^y_b \ket{\widehat{\psi}}
		=
		\bra{\psi} A^x_a \ot B^y_b \ket{\psi}.
	\]
\end{proof}

\begin{remark}[Lean questionwise Naimark interface]\label{rem:lean-questionwise-naimark}
	In addition to Theorem~\ref{thm:naimark}, Lean records the questionwise
	one-measurement interface obtained from Lemma~\ref{lem:naimark-helper}.  For
	every indexed family \(A^x\) and \(B^y\), it produces local dilation data for
	each question and proves the corresponding single-outcome
	marginal-preservation identities.  This is a useful Lean-only auxiliary
	interface for downstream arguments, not a substitute for the full
	tensor-product correlation theorem, which is proved separately by
	\texttt{naimarkTensorProductCorrelation}.
\end{remark}

\begin{remark}[Lean auxiliary declarations for Naimark]\label{rem:lean-naimark-auxiliary-declarations}
	The checked Lean proof of Theorem~\ref{thm:naimark} uses the standard
	one-measurement auxiliary register
	\[
		\mathrm{Option}(\mathcal A),
	\]
	with the all-\(\bot\) vector as the initial auxiliary state.  A separate
	Naimark unitary is chosen for each question, so the proof does not need a
	distinct auxiliary tensor factor for every question.  This is equivalent for
	the theorem's conclusion, since no commutation relation between different
	questions' dilated measurements is required.  The declaration
	\texttt{OneMeasNaimarkData.twoSidedCorrelationPreservation} proves the
	four-register trace identity which turns the two local compression identities
	into the bipartite correlation-preservation identity used by
	\texttt{naimarkTensorProductCorrelation}.
\end{remark}

\begin{theorem}[Orthogonalization lemma]\label{thm:orthonormalization}
	\lean{MIPStarRE.LDT.MakingMeasurementsProjective.orthonormalization}
	\leanok
	% The theorem now proves the paper constant and is axiom-clean.  Issue #1032
	% still tracks the heterogeneous proof of
	% Lemma~\ref{lem:orthonormalization-main-lemma}, which is a separate
	% auxiliary lemma and is not used as a proof dependency here.
	% In Lean, normalization is an explicit hypothesis because the ambient state
	% record contains positive semidefinite density matrices of arbitrary
	% normalized trace; this matches the paper's convention that a state is
	% normalized.
	% The paper states the constant $100\zeta^{1/4}$, and Lean now proves
	% this source theorem directly.  The completion-route variant with the
	% weaker $120\zeta^{1/4}$ bound is recorded separately as
	% `orthonormalizationCompletionRoute`.
	\uses{def:strong-self-consistency, def:approx_delta}
	Let $\ket{\psi}$ be a permutation-invariant state,
	and let $A = \{A_a\}$ be a sub-measurement with strong self-consistency
	\[
		\sum_a \bra{\psi} A_a \ot A_a \ket{\psi}
		\geq
		\sum_a \bra{\psi} A_a \ot I \ket{\psi} - \zeta.
	\]
	Then there exists a projective sub-measurement $P = \{P_a\}$ such that
	\[
		A_a \ot I \approx_{100 \zeta^{1/4}} P_a \ot I.
	\]
\end{theorem}
\begin{proof}
	\lean{MIPStarRE.LDT.MakingMeasurementsProjective.orthonormalization}
	\leanok
	\uses{lem:orthonormalization-main-lemma, lem:locality-preserving-projectivization,
		prop:other-two-notions-of-self-consistency,
		prop:triangle-inequality-for-approx_delta}
	Complete the sub-measurement by adding the missing outcome and use
	permutation invariance to convert strong self-consistency into the
	cross-consistency hypothesis of
	Lemma~\ref{lem:orthonormalization-main-lemma}.  Applying the
	locality-preserving projectivization repair gives the required projective
	sub-measurement with the paper constant \(100\zeta^{1/4}\).
\end{proof}

\begin{remark}[Lean right-register completion helpers]\label{rem:lean-right-register-completion-helpers}
	The inductive-step application of the orthogonalization lemma uses
	completion on both tensor factors.  The formal development records the
	analytic completion statement on the left register; the shared preliminary
	helper compares right- and left-tensor placements at the level of the
	underlying state-dependent-distance form.  Its submeasurement and indexed
	lemmas transport state-dependent distance to the right register under the
	permutation-invariance hypothesis on the state.  The last linked declaration
	is a Lean-only construction from two already constructed
	orthonormalize-and-complete statements and a pre-projective consistency
	proof.  It is not a source theorem and it is not an additional hypothesis of
	Theorem~\ref{thm:main-formal}; it records only the tensor-factor bookkeeping
	needed to form the later line-156 projectivization transition.
\end{remark}

\begin{remark}[Lean line-169 projectivization match-mass declarations]\label{rem:lean-line169-projectivization-match-mass}
	Lean exposes the quantitative obstruction at the projectivization transition used
	in the inductive step.  The existing line-156 transition gives only the generic
	consequence of Proposition~\ref{prop:triangle-sub}: from
	$G^{\mathrm A}\simeq_{\zeta_1}G^{\mathrm B}$ and
	$G^{\mathrm A}\approx_{\zeta_2}Q^{\mathrm A}$ one obtains
	$Q^{\mathrm A}\simeq_{\zeta_1+\sqrt{\zeta_2}}G^{\mathrm B}$, and similarly
	on the other tensor factor.  This direct square-root-loss route is an
	immediate specialization of the general Lean declarations
	\texttt{MIPStarRE.LDT.Preliminaries.triangleSub} and
	\texttt{MIPStarRE.LDT.Preliminaries.triangleSub\_right} to the handoff
	fields, and is no longer maintained as separate named projectivization API.

	Lean also proves a more local repaired transport that keeps the comparison at
	the pre-completion stage.  If $P^{\mathrm A}$ is the orthonormalized
	projective sub-measurement with $G^{\mathrm A}\approx_{\epsilon}P^{\mathrm A}$,
	then Proposition~\ref{prop:easy-approx-from-approx-delta} bounds the diagonal
	match-mass loss by $\sqrt{\epsilon}$, and canonical completion can only add
	nonnegative diagonal match mass.  Thus the declarations
	\texttt{leftConsistency\_of\_completion\_and\_sdd} and
	\texttt{rightConsistency\_of\_completion\_and\_sdd} prove
	$Q^{\mathrm A}\simeq_{\zeta_1+\sqrt{\epsilon}}G^{\mathrm B}$ and its Bob-side
	mirror from the pre-projective $\zeta_1$ consistency plus the
	pre-completion $\approx_{\epsilon}$ comparison.  In particular, substituting
	the checked orthonormalization error
	$\epsilon=100\zeta_1^{1/4}$ yields the repaired line-169 bound
	$\zeta_1+10\zeta_1^{1/8}$ recorded by
	\texttt{leftConsistency\_with\_orthonormalization\_loss} and its mirror.
	This is much sharper than the direct completion-level
	$\zeta_1+\sqrt{\zeta_2}$ route, but it is still not the paper's exact
	$\zeta_1$.

	Lean also isolates the completion fact used inside this sharper repair: the
	canonical completion of a projective submeasurement can only increase its
	diagonal match mass against a fixed partner submeasurement.  This is the formal
	declaration
	\texttt{ProjectivizationMatchMassMonotonicity.completeAtOutcomeProj\_left\_matchMass\_ge}.
	The sharper route therefore pays its square-root loss before completion and then
	uses completion only in this monotone direction, with no additional match-mass
	penalty.

	The older exact-$\zeta_1$ construction-level monotonicity statement has been
	retired.  The current Section~5 API proves the checked repaired transport above,
	and the active orthonormalization route now uses those repaired completion-transport estimates rather
	than a separate exact match-mass witness.
\end{remark}

\subsection{Naimark dilation}\label{sec:naimark}

\begin{lemma}\label{lem:naimark-helper}
	\lean{MIPStarRE.LDT.MakingMeasurementsProjective.oneMeasNaimark}
	\leanok
	\uses{def:submeasurement}
	Let $A = \{A_a\}$ be a sub-measurement with $k$ distinct outcomes $a \in \mathcal A$, and let $\ket{\mathsf{aux}} \in \C^{k+1}$ be any state. Then there exists a projective sub-measurement $\widehat{A} = \{\widehat{A}_a\}$ such that for each outcome $a$,
	\[
		(I \ot \bra{\mathsf{aux}}) \cdot \widehat{A}_a \cdot (I \ot \ket{\mathsf{aux}})
		=
		A_a.
	\]
\end{lemma}
\begin{proof}
	\leanok
	Consider an orthonormal basis for $\C^{k+1}$ consisting of vectors $\ket{a}$ for $a \in \mathcal A$ and the vector $\ket{\bot}$. Let $U$ be any unitary such that for each vector $\ket{\phi}$,
	\[
		U \cdot (\ket{\phi} \ot \ket{\mathsf{aux}})
		=
		\sum_{a \in \mathcal A} ((A_a)^{1/2}\ket{\phi}) \ot \ket{a}
		+
		((I-A)^{1/2}\ket{\phi}) \ot \ket{\bot},
	\]
	where $A=\sum_a A_a$. Then
	\[
		U \cdot (I \ot \ket{\mathsf{aux}})
		=
		\sum_{a \in \mathcal A} (A_a)^{1/2} \ot \ket{a}
		+
		(I-A)^{1/2} \ot \ket{\bot},
	\]
	so for any $a \in \mathcal A$,
	\[
		(I \ot \bra{a}) \cdot U \cdot (I \ot \ket{\mathsf{aux}})
		=
		(A_a)^{1/2}.
	\]
	Define
	\[
		\widehat{A}_a = U^\dagger \cdot (I \ot \ket{a}\bra{a}) \cdot U.
	\]
	Each $\widehat{A}_a$ is a projector, and
	\begin{align*}
		(I \ot \bra{\mathsf{aux}}) \cdot \widehat{A}_a \cdot (I \ot \ket{\mathsf{aux}})
		 & =
		(I \ot \bra{\mathsf{aux}}) \cdot U^\dagger \cdot (I \ot \ket{a}) \cdot (I \ot \bra{a}) \cdot U \cdot (I \ot \ket{\mathsf{aux}}) \\
		 & =
		(A_a)^{1/2} \cdot (A_a)^{1/2}
		=
		A_a.
	\end{align*}
\end{proof}

% Paper source: references/ldt-paper/orthonormalization.tex, lines 189--271.
\begin{remark}[Naimark dilation does not preserve $\approx_\delta$]\label{ex:easy-but-long}
		Let $\ket{\psi}$ be any state in $\C^d \ot \C^d$, and consider the two-outcome measurements $A = \{A_0, A_1\}$ and $B = \{B_0, B_1\}$ with
		\[
			A_0 = A_1 = B_0 = B_1 = \tfrac{1}{2} I.
		\]
		For each outcome $a \in \{0,1\}$, the operators $A_a \ot I$ and $I \ot B_a$ are both equal to $\tfrac{1}{2} I \ot I$, so $A_a \ot I \approx_0 I \ot B_a$ on the state $\ket{\psi}$.

		Now apply the Naimark construction with a qubit auxiliary space on each side and auxiliary state $\ket{+} = \tfrac{1}{\sqrt 2}(\ket{0}+\ket{1})$. Since $A$ and $B$ are genuine measurements rather than sub-measurements, the extra $\ket{\bot}$ direction from \ref{lem:naimark-helper} is unnecessary in this example. The unitary from \ref{lem:naimark-helper} can then be taken to be the identity, so the dilated measurements are
		\[
			\widehat A_a = I \ot \ket{a}\bra{a}, \qquad \widehat B_a = I \ot \ket{a}\bra{a}.
		\]
		If $\ket{\widehat{\psi}} = \ket{\psi} \ot \ket{+} \ot \ket{+}$, then a direct calculation gives
		\[
			\sum_{a \in \{0,1\}} \bigl\lVert (\widehat A_a \ot I - I \ot \widehat B_a) \ket{\widehat{\psi}} \bigr\rVert^2 = 1.
		\]
		Thus on the state $\ket{\widehat{\psi}}$ the dilated family satisfies $\widehat A_a \ot I \approx_\delta I \ot \widehat B_a$ only for $\delta \ge 1$, even though the original family satisfied the same relation with $\delta = 0$. This is why we use Theorem~\ref{thm:orthonormalization} rather than simply applying Theorem~\ref{thm:naimark}: Naimark preserves the outcome probabilities, but it need not preserve the state-dependent $\approx_\delta$ estimates needed later.
\end{remark}

\subsection{Orthogonalization lemma}\label{sec:orthogonalization}

\begin{lemma}[Orthogonalization lemma for measurements]\label{lem:orthonormalization-main-lemma}
	\lean{MIPStarRE.LDT.MakingMeasurementsProjective.orthonormalizationMainLemma}
	\leanok
	\uses{def:simeq, def:approx_delta}
	Let $\ket{\psi}$ be a state, and let $A = \{A_a\}$ and $B = \{B_a\}$ be measurements such that
	\[
		A_a \ot I \simeq_\zeta I \ot B_a
	\]
	on state $\ket{\psi}$. Here $\ket{\psi}$ is not assumed to be permutation-invariant.
	Then there exists a projective sub-measurement $P = \{P_a\}$ such that
	\[
		A_a \ot I \approx_{84 \zeta^{1/4}} P_a \ot I.
	\]
\end{lemma}

\begin{lemma}[Formalized orthogonalization envelope for measurements]\label{lem:orthonormalization-main-lemma-formalized-envelope}
	\lean{MIPStarRE.LDT.MakingMeasurementsProjective.orthonormalizationMeasurement_of_consistency_from_projectivizationRepair}
	\leanok
	\uses{lem:orthonormalization-main-lemma, lem:locality-preserving-projectivization, def:simeq, def:approx_delta}
	This same-space corollary weakens Lemma~\ref{lem:orthonormalization-main-lemma}
	to the public orthonormalization envelope \(100\zeta^{1/4}\).  Thus, if \(A=\{A_a\}\) and
	\(B=\{B_a\}\) are measurements on the same finite-dimensional space satisfying
	\[
		A_a \ot I \simeq_\zeta I \ot B_a
	\]
	on a normalized bipartite state \(\ket{\psi}\), then there exists a projective
	sub-measurement \(P=\{P_a\}\) such that
	\[
		A_a \ot I \approx_{100\zeta^{1/4}} P_a \ot I .
	\]
	This is the corrected Lean theorem proved from the Section~5
	projectivization repair construction.  The sharper heterogeneous source
	statement \(84\zeta^{1/4}\) is now formalized as
	Lemma~\ref{lem:orthonormalization-main-lemma}.
\end{lemma}
\begin{proof}
	\lean{MIPStarRE.LDT.MakingMeasurementsProjective.orthonormalizationMeasurement_of_consistency_from_projectivizationRepair}
	\leanok
	\uses{lem:orthonormalization-main-lemma, lem:locality-preserving-projectivization,
		prop:simeq-for-measurements, prop:triangle-inequality-for-approx_delta}
	The consistency hypothesis first gives the source almost-projective estimate
	for \(A_a\ot I\).  The locality-preserving projectivization repair then
	produces a projective sub-measurement on the original space, and the scalar
	bounds convert the \(84\)-constant repair estimate into the public
	\(100\zeta^{1/4}\) envelope.
\end{proof}

\begin{remark}[Lean auxiliary declarations for the orthogonalization lemma]
	\label{rem:lean-orthonormalization-main-declarations}
	The Lean declaration
	\texttt{MIPStarRE.LDT.MakingMeasurementsProjective.orthonormalizationMainLemma}
	proves Lemma~\ref{lem:orthonormalization-main-lemma} without adding
	spectral-truncation or repair hypotheses.  The proof first converts the
	cross-consistency hypothesis into the source almost-projective estimate for
	the measurement \(A_a \otimes I\), and then applies the locality-preserving
	\(Q/X/\widehat X/P\) construction in its heterogeneous left-register form.
	Lemma~\ref{lem:orthonormalization-main-lemma-formalized-envelope} remains as
	the same-space corollary with the weaker public \(100\zeta^{1/4}\) envelope.
	The measurement-level declarations
	\texttt{orthonormalizationMeasurement\_of\_consistency} and
	\texttt{orthonormalizationMeasurement} are proved corollaries, not
	conditional wrappers: they derive the projective submeasurement from the
	displayed consistency or self-consistency hypotheses.
\end{remark}

\begin{lemma}[Locality-preserving projectivization]\label{lem:locality-preserving-projectivization}
	\lean{MIPStarRE.LDT.MakingMeasurementsProjective.leftLiftedProjectivizationRepair}
	\uses{lem:projective-non-measurement, def:approx_delta}
	\leanok
	Paper origin: \texttt{references/ldt-paper/orthonormalization.tex} lines 534--860 (rank reduction and the $Q$-side setup, including \ref{lem:xa-t} and \ref{lem:qa-restated}) and 862--1194 (the $X/\widehat X/P$ algebra: \ref{lem:X-squared}, \ref{lem:X-hat-squared}, \ref{lem:X-times-X-hat}, \ref{lem:squared-difference}, \ref{lem:P-projectivity}, \ref{lem:P-Q-approx}, plus the final triangle-inequality calculation producing the $84\zeta^{1/4}$ bound). Together these are the late repair stage of the orthogonalization-lemma proof.

	Given a normalized bipartite state $\ket{\psi}$, a measurement
	$A = \{A_a\}$, and the source almost-projective estimate for the
	left-lifted family $\{A_a \ot I\}$, there exists a locality-preserving
	projective sub-measurement with error $84\zeta^{1/4}$.

	More precisely, if $\ket{\psi}$ is a normalized bipartite state, if
	$A = \{A_a\}$ is a measurement, and if
	\[
		\sum_a \bra{\psi} \bigl((A_a \ot I) - (A_a \ot I)^2\bigr) \ket{\psi} \le \zeta,
	\]
	then there exists a projective sub-measurement $P = \{P_a\}$ on the underlying
	space such that
	\[
		A_a \ot I \approx_{84\zeta^{1/4}} P_a \ot I.
	\]
	The locality-preserving form -- output $P_a \ot I$ rather than an arbitrary
	lifted family -- is recorded in
	Remark~\ref{rem:lean-orthonormalization-main-declarations} and is the form
	required by the unconditional orthonormalization theorem
	(\ref{thm:orthonormalization}).
\end{lemma}
\begin{proof}
	\leanok
	\uses{lem:projective-non-measurement, lem:projective-low-rank-sum,
		lem:xa-t, lem:qa-restated, lem:X-squared, lem:X-hat-squared,
		lem:X-times-X-hat, lem:squared-difference, lem:P-projectivity, lem:P-Q-approx,
		prop:triangle-inequality-for-approx_delta}
	Pass to the left marginal state, apply the rank-reduction and
	$Q/X/\widehat X/P$ construction there, and transport the resulting
	$\approx_{84\zeta^{1/4}}$ estimate back to left lifts via the left-marginal
	expectation formula.
\end{proof}

\begin{proposition}[Measurement-level orthogonalization from consistency]\label{prop:orthonormalization-measurement-of-consistency}
	\lean{MIPStarRE.LDT.MakingMeasurementsProjective.orthonormalizationMeasurement_of_consistency}
	\leanok
	\uses{def:strong-self-consistency, def:approx_delta, lem:orthonormalization-main-lemma}
	This is the same-space measurement corollary recorded in
	Remark~\ref{rem:lean-orthonormalization-main-declarations}.  Let
	\(\ket{\psi}\) be a normalized bipartite state, let \(A=\{A_a\}\) and
	\(B=\{B_a\}\) be measurements on the same finite-dimensional space, and
	let \(0 \leq \zeta\).  Assume the bipartite consistency estimate
	\[
		A_a \ot I \simeq_{\zeta} I \ot B_a .
	\]
	Then there exists a projective sub-measurement \(P=\{P_a\}\) such that
	\[
		A_a \ot I \approx_{100\zeta^{1/4}} P_a \ot I .
	\]
\end{proposition}
\begin{proof}
	\leanok
	\uses{lem:orthonormalization-main-lemma}
	The proof applies Lemma~\ref{lem:orthonormalization-main-lemma} and weakens
	the \(84\zeta^{1/4}\) estimate to the public orthonormalization envelope.
	The spectral-truncation and locality-preserving repair steps are internal to
	that proof; they are not hypotheses of this declaration.
\end{proof}

\begin{proposition}[Measurement-level orthogonalization from self-consistency]\label{prop:orthonormalization-measurement}
	\lean{MIPStarRE.LDT.MakingMeasurementsProjective.orthonormalizationMeasurement}
	\leanok
	\uses{def:strong-self-consistency, def:approx_delta,
		prop:orthonormalization-measurement-of-consistency}
	This is the self-consistency specialization of
	Proposition~\ref{prop:orthonormalization-measurement-of-consistency}.
	Let \(\ket{\psi}\) be a normalized bipartite state, let \(A=\{A_a\}\) be a
	measurement, and let \(0\leq\zeta\).  If \(A\) is bipartite strongly
	self-consistent with error \(\zeta\), then there exists a projective
	sub-measurement \(P=\{P_a\}\) such that
	\[
		A_a \ot I \approx_{100 \zeta^{1/4}} P_a \ot I.
	\]
\end{proposition}
\begin{proof}
	\leanok
	\uses{prop:orthonormalization-measurement-of-consistency}
	Since \(A\) is complete, the self-consistency defect is the consistency
	defect of \(A\) with itself.  The proof therefore applies the consistency
	corollary with \(B=A\).  No spectral-truncation witness or repair witness is
	assumed at this interface.
\end{proof}

\begin{proposition}[Completion-route orthonormalization]\label{prop:orthonormalization-completion-route}
	\lean{MIPStarRE.LDT.MakingMeasurementsProjective.orthonormalizationCompletionRoute}
	\leanok
	\uses{def:strong-self-consistency, def:approx_delta, prop:other-two-notions-of-self-consistency,
		lem:locality-preserving-projectivization}

	The paper theorem \ref{thm:orthonormalization} is proved in Lean with the
	\(100\zeta^{1/4}\) constant.  The declaration checked here is the
	completion-route theorem with the weaker bound
	\(120\zeta^{1/4}\), recorded in
	\texttt{orthonormalizationCompletionRouteError}.  See
	\cite{gap:issue-1032-orthonormalization-constant}.

	Let \(\ket{\psi}\) be a normalized permutation-invariant state, and let
	\(A=\{A_a\}_{a\in\mathcal A}\) be a sub-measurement which is bipartite
	strongly self-consistent with error \(\zeta\).  Then there exists a
	projective sub-measurement \(P=\{P_a\}_{a\in\mathcal A}\) such that
	\[
		A_a \ot I \approx_{120\zeta^{1/4}} P_a \ot I.
	\]
\end{proposition}
\begin{proof}
	\leanok
	\uses{prop:other-two-notions-of-self-consistency, lem:locality-preserving-projectivization}
	Write \(A=\sum_a A_a\).  The strong self-consistency assumption implies
	\[
		\bra{\psi} A \ot A \ket{\psi}
		\geq
		\bra{\psi} A \ot I \ket{\psi} - \zeta,
	\]
	so
	\[
		\bra{\psi} A \ot (I-A) \ket{\psi} \leq \zeta.
	\]
	Rearranging once more gives
	\[
		\bra{\psi} (I-A) \ot (I-A) \ket{\psi}
		\geq
		\bra{\psi} (I-A) \ot I \ket{\psi} - \zeta.
	\]

	Complete \(A\) to the measurement
	\(\widehat A=\{\widehat A_a\}_{a \in \widehat{\mathcal A}}\) on
	\(\widehat{\mathcal A}=\mathcal A \cup \{\bot\}\) by setting
	\(\widehat A_a=A_a\) for \(a \in \mathcal A\) and
	\(\widehat A_{\bot}=I-A\).  Adding the two preceding inequalities shows
	that \(\widehat A\) is \(2\zeta\)-strongly self-consistent.  By
	\ref{prop:other-two-notions-of-self-consistency}, this gives the source
	almost-projective estimate
	\[
		\sum_a \bra{\psi}
		\bigl((\widehat A_a\ot I)-(\widehat A_a\ot I)^2\bigr)
		\ket{\psi}
		\le 4\zeta.
	\]
	Apply \ref{lem:locality-preserving-projectivization} to obtain a projective
	sub-measurement \(\widehat P\) such that
	\[
		\widehat A_a \ot I \approx_{84(4\zeta)^{1/4}} \widehat P_a \ot I.
	\]
	Discard the extra outcome and set \(P_a=\widehat P_a\) for
	\(a \in \mathcal A\).  Then
	\[
		A_a \ot I \approx_{84(4\zeta)^{1/4}} P_a \ot I,
	\]
	and \(84 \cdot 4^{1/4} \leq 120\), so the completion-route bound follows.
\end{proof}

\begin{proof}[Proof of \ref{lem:orthonormalization-main-lemma}]
	\proves{lem:orthonormalization-main-lemma}
	\lean{MIPStarRE.LDT.MakingMeasurementsProjective.orthonormalizationMainLemma}
	\leanok
	\uses{prop:simeq-for-measurements, lem:projective-non-measurement, lem:projective-low-rank-sum, lem:Q-completeness, lem:sqrt-Q-completeness, lem:q-almost-projective, lem:xa-t, lem:qa-restated, lem:X-squared, lem:X-expression-to-Q-expression, lem:pa-restated, lem:X-hat-squared, lem:X-times-X-hat, lem:squared-difference, lem:P-projectivity, lem:P-Q-approx, prop:triangle-inequality-for-approx_delta}
	The statement is trivial when $\zeta > 1/4$, so assume
	\begin{equation}\label{eq:assumption-on-zeta}
		\zeta \leq 1/4.
	\end{equation}
	By \ref{prop:simeq-for-measurements},
	\[
		\sum_a \bra{\psi} A_a \ot B_a \ket{\psi} \geq 1-\zeta.
	\]
	Apply Cauchy--Schwarz to
	\[
		\sum_a \bra{\psi} (A_a \ot I)\cdot(I \ot B_a)\ket{\psi}
	\]
	to obtain
	\[
		\sum_a \bra{\psi} (A_a)^2 \ot I \ket{\psi} \geq (1-\zeta)^2 \geq 1-2\zeta.
	\]
	Because $A$ is a measurement, this rewrites as
	\begin{equation}\label{eq:A-looks-projective}
		\sum_a \bra{\psi} (A_a-(A_a)^2) \ot I \ket{\psi} \leq 2\zeta.
	\end{equation}

	\begin{lemma}\label{lem:trunc-inequality}
		\lean{MIPStarRE.LDT.MakingMeasurementsProjective.truncationInequality}
		\leanok
		For any $x \in [0, 1]$,
		\begin{equation*}
			(x - \mathsf{trunc}_\delta(x))^2 \leq \frac{1}{\delta} \cdot (x - x^2).
		\end{equation*}
	\end{lemma}
	\begin{proof}
		\leanok
		If $x \geq 1-\delta$, then $\mathsf{trunc}_\delta(x)=1$. Since $\delta \leq 1/2$ below, one also has $x \geq \delta$, so
		\[
			(x-\mathsf{trunc}_\delta(x))^2
			=
			(1-x)^2
			\leq
			(1-x)\frac{x}{\delta}
			=
			\frac{1}{\delta}(x-x^2).
		\]
		If $x < 1-\delta$, then $\mathsf{trunc}_\delta(x)=0$, and similarly
		\[
			(x-\mathsf{trunc}_\delta(x))^2
			=
			x^2
			\leq
			x\frac{1-x}{\delta}
			=
			\frac{1}{\delta}(x-x^2).
		\]
	\end{proof}

	\begin{lemma}[Rounding to projectors]\label{lem:projective-non-measurement}
		\lean{MIPStarRE.LDT.MakingMeasurementsProjective.projectiveNonMeasurement}
		\lean{MIPStarRE.LDT.MakingMeasurementsProjective.projectiveNonMeasurement_of_sourceAlmostProjective_full}
		\leanok
		There exists a set of projective matrices $\{R_a\}$ such that
		\begin{equation*}
			A_a \ot I \approx_{2\sqrt{\zeta}} R_a \ot I.
		\end{equation*}
		and
		\begin{equation*}
			R:= \sum_a R_a \leq (1+2\sqrt{\zeta}) \cdot I.
		\end{equation*}
	\end{lemma}
	\begin{proof}
		\leanok
		\uses{lem:trunc-inequality}
		Write the eigendecomposition of $A_a$ as
		\[
			A_a = \sum_i \lambda_{a,i}\ket{u_{a,i}}\bra{u_{a,i}}.
		\]
		Choose
		\begin{equation}\label{eq:bound-on-delta}
			0 < \delta \leq 1/2,
		\end{equation}
		define $\mathsf{trunc}_\delta:[0,1] \to \{0,1\}$ by $\mathsf{trunc}_\delta(x)=1$ when $x \geq 1-\delta$ and $0$ otherwise, and set
		\[
			R_a = \mathsf{trunc}_\delta(A_a)
			=
			\sum_i \mathsf{trunc}_\delta(\lambda_{a,i}) \ket{u_{a,i}}\bra{u_{a,i}}.
		\]
		By \ref{lem:trunc-inequality},
		\[
			(A_a-R_a)^2
			\leq
			\frac{1}{\delta}(A_a-(A_a)^2).
		\]
		Summing expectations and using \ref{eq:A-looks-projective} gives
		\[
			\sum_a \Vert (A_a-R_a)\ot I \ket{\psi}\Vert^2
			\leq
			\frac{2\zeta}{\delta},
		\]
		so $A_a \ot I \approx_{2\zeta/\delta} R_a \ot I$.

		For every $x \in [0,1]$, one has $\mathsf{trunc}_\delta(x) \leq x/(1-\delta)$, hence
		\[
			R_a \leq \frac{1}{1-\delta} A_a.
		\]
		Summing over $a$ and using that $A$ is a measurement,
		\[
			R = \sum_a R_a \leq \frac{1}{1-\delta} I \leq (1+2\delta)I,
		\]
		where the last estimate uses $\delta \leq 1/2$. Setting $\delta=\sqrt{\zeta}$ gives the claim, and \ref{eq:assumption-on-zeta} ensures $\delta \leq 1/2$.
	\end{proof}

	\begin{remark}[Lean auxiliary declarations for rounding to projectors]\label{rem:lean-projective-non-measurement-auxiliary}
		The source-facing Lean declarations for
		Lemma~\ref{lem:projective-non-measurement} are the proposition
		\texttt{projectiveNonMeasurement} and the construction theorem
		\texttt{projectiveNonMeasurement\_of\_sourceAlmostProjective\_full}.
		The declarations listed here are the endpoint and regime splits, the
		\texttt{AlmostProjMeasStatement} witness forms, and the spectral-truncation
		statement conversions used by later Section~5 arguments.  They are not additional
		hypotheses in the paper statement of
		Lemma~\ref{lem:projective-non-measurement}.
	\end{remark}

	\begin{lemma}[Combinatorial averaging for the $r > d$ branch]\label{lem:small-overlap-averaging}
		\lean{MIPStarRE.LDT.MakingMeasurementsProjective.Truncation.exists_large_subset_ordered}
		\lean{MIPStarRE.LDT.MakingMeasurementsProjective.Truncation.card_mul_sum_small_le}
		\lean{MIPStarRE.LDT.MakingMeasurementsProjective.Truncation.card_univ_mul_sum_compl_le}
		\lean{MIPStarRE.LDT.MakingMeasurementsProjective.Truncation.sum_small_le_four_sqrt}
		\leanok
		Let $f : \alpha \to \mathbb{R}$ be a function on a finite index set $\alpha$ with
		$|\alpha| = r$, and let $L \subseteq \alpha$ with $|L| = d$ and $d < r$. Write
		$S = L^\complement$ and assume $f(s) \leq f(l)$ for every $s \in S$ and $l \in L$.
		Then the pairwise double-counting inequality
		\[
			|L|\, \sum_{s \in S} f(s) \leq |S|\, \sum_{l \in L} f(l)
		\]
		holds, and hence also
		\[
			r\, \sum_{s \in S} f(s) \leq |S|\, \sum_{x \in \alpha} f(x).
		\]
		If in addition $r \leq (1+2\sqrt{\zeta})\,d$, $\sum_{x \in \alpha} f(x) \leq
		1 + 2\sqrt{\zeta}$, and $0 \leq \zeta \leq 1/4$, then
		\[
			\sum_{s \in S} f(s) \leq 4\sqrt{\zeta}.
		\]
	\end{lemma}
	\begin{proof}
		\leanok
		Choose $L$ to be a set of $d$ indices on which the sum of $f$ is maximal. If some $s \in S$ and $l \in L$ satisfied $f(l) < f(s)$, replacing $l$ by $s$ would increase this sum, so every small index is bounded by every large index. Summing the inequalities $f(s) \leq f(l)$ over all pairs $(s,l) \in S \times L$ gives
		\[
			|L|\sum_{s \in S} f(s) \leq |S|\sum_{l \in L} f(l).
		\]
		Since $\alpha$ is the disjoint union of $L$ and $S$, adding $|S|\sum_{s \in S} f(s)$ to both sides yields
		\[
			r\sum_{s \in S} f(s) \leq |S|\sum_{x \in \alpha} f(x).
		\]
		The hypotheses $|L|=d$ and $r \leq (1+2\sqrt{\zeta})d$ imply $|S|/r \leq 2\sqrt{\zeta}$. Dividing by $r$ and using $\sum_{x \in \alpha} f(x) \leq 1+2\sqrt{\zeta}$ gives
		\[
			\sum_{s \in S} f(s) \leq 2\sqrt{\zeta}(1+2\sqrt{\zeta}).
		\]
		Finally, $\zeta \leq 1/4$ implies $\sqrt{\zeta} \leq 1/2$, and hence $2\sqrt{\zeta}(1+2\sqrt{\zeta}) \leq 4\sqrt{\zeta}$.
	\end{proof}


	\begin{lemma}[Rank reduction]\label{lem:projective-low-rank-sum}
		\lean{MIPStarRE.LDT.MakingMeasurementsProjective.projectiveLowRankSum}
		\leanok
		There exists a set of projection matrices $\{Q_a\}$ such that
		\begin{equation*}
			A_a \ot I \approx_{12\sqrt{\zeta}} Q_a \ot I.
		\end{equation*}
		and
		\begin{equation*}
			Q:= \sum_a Q_a \leq (1+2\sqrt{\zeta}) \cdot I.
		\end{equation*}
		Furthermore,
		\begin{equation*}
			\sum_a \mathrm{rank}(Q_a) \leq d.
		\end{equation*}
	\end{lemma}
	\begin{proof}
		\leanok
		\uses{lem:projective-non-measurement, lem:small-overlap-averaging, prop:triangle-inequality-for-approx_delta}
		\lean{MIPStarRE.LDT.MakingMeasurementsProjective.projectiveLowRankSum_of_projectors}
		\lean{MIPStarRE.LDT.MakingMeasurementsProjective.projectiveLowRankSum_of_rank_bound}
		\lean{MIPStarRE.LDT.MakingMeasurementsProjective.projectiveLowRankSum_of_spectralTruncationStatement}
		\lean{MIPStarRE.LDT.MakingMeasurementsProjective.projectiveLowRankSum_truncate}
		Let $\{R_a\}$ be given by \ref{lem:projective-non-measurement}. Write $r_a=\mathrm{rank}(R_a)$ and $r=\sum_a r_a$. If $r \leq d$, take $Q_a=R_a$. Otherwise,
		\begin{equation}\label{eq:bound-on-r}
			r
			=
			\sum_a r_a
			=
			\sum_a \Tr(R_a)
			=
			\Tr(R)
			\leq
			(1+2\sqrt{\zeta})d.
		\end{equation}
		Let $\ket{v_{a,1}},\ldots,\ket{v_{a,r_a}}$ be an orthonormal basis for the range of $R_a$, so
		\[
			R_a = \sum_{i=1}^{r_a} \ket{v_{a,i}}\bra{v_{a,i}}.
		\]
		Define the overlaps
		\[
			o_{a,i} = \bra{\psi} (\ket{v_{a,i}}\bra{v_{a,i}} \ot I)\ket{\psi}.
		\]
		Their total satisfies
		\begin{equation}\label{eq:total-overlap}
			\sum_a \sum_{i=1}^{r_a} o_{a,i}
			=
			\bra{\psi} R \ot I \ket{\psi}
			\leq
			1+2\sqrt{\zeta}.
		\end{equation}
		Let $\mathsf{Large}$ be the set of the $d$ largest overlaps $o_{a,i}$, breaking ties arbitrarily, and let $\mathsf{Small}$ be the complement. Then
		\begin{equation}\label{eq:size-of-small-is-small}
			|\mathsf{Small}|
			=
			r-d
			\leq
			2\sqrt{\zeta}\,d
			\leq
			2\sqrt{\zeta}\,r,
		\end{equation}
		by \ref{eq:bound-on-r}. Averaging the overlaps over the $r$ indices gives
		\begin{equation}\label{eq:small-overlaps}
			\sum_{(a,i)\in \mathsf{Small}} o_{a,i}
			\leq
			\frac{|\mathsf{Small}|}{r}\sum_{a,i} o_{a,i}
			\leq
			4\sqrt{\zeta},
		\end{equation}
		using \ref{eq:size-of-small-is-small}, \ref{eq:total-overlap}, and \ref{eq:assumption-on-zeta}.

		For each $a$, let $\mathsf{Large}_a=\{i : (a,i)\in \mathsf{Large}\}$ and define
		\[
			Q_a = \sum_{i \in \mathsf{Large}_a} \ket{v_{a,i}}\bra{v_{a,i}}.
		\]
		Then $\sum_a \mathrm{rank}(Q_a)=|\mathsf{Large}| \leq d$, and $Q_a \leq R_a$ for every $a$, so
		\[
			Q = \sum_a Q_a \leq \sum_a R_a = R \leq (1+2\sqrt{\zeta})I.
		\]
		Also,
		\[
			R_a-Q_a = \sum_{i \in \mathsf{Small}_a} \ket{v_{a,i}}\bra{v_{a,i}}
		\]
		is a projector, hence
		\[
			\sum_a \Vert (R_a-Q_a)\ot I \ket{\psi}\Vert^2
			=
			\sum_{(a,i)\in\mathsf{Small}} o_{a,i}
			\leq
			4\sqrt{\zeta}
		\]
		by \ref{eq:small-overlaps}. Therefore $R_a \ot I \approx_{4\sqrt{\zeta}} Q_a \ot I$. Combining this with \ref{lem:projective-non-measurement} and the triangle inequality for $\approx_\delta$ gives
		\[
			A_a \ot I \approx_{12\sqrt{\zeta}} Q_a \ot I.
		\]
	\end{proof}

	Henceforth let $Q=\{Q_a\}$ be the family given by \ref{lem:projective-low-rank-sum}.

	\begin{lemma}[Completeness of~$Q$]\label{lem:Q-completeness}
		\lean{MIPStarRE.LDT.MakingMeasurementsProjective.qCompleteness}
		\leanok
		\begin{equation*}
			\bra{\psi} Q \ot I \ket{\psi}
			\geq
			1 - 11 \zeta^{1/4}.
		\end{equation*}
	\end{lemma}
	\begin{proof}
		\leanok
		\uses{lem:projective-low-rank-sum}
		Since each $Q_a$ is projective,
		\begin{equation}\label{eq:Q-for-an-A}
			\bra{\psi} Q \ot I \ket{\psi}
			=
			\sum_a \bra{\psi} (Q_a)^2 \ot I \ket{\psi}
			\approx_{5\zeta^{1/4}}
			\sum_a \bra{\psi} (Q_a \cdot A_a)\ot I \ket{\psi}.
		\end{equation}
		By Cauchy--Schwarz,
		\[
			\sum_a \bra{\psi} (Q_a)^2 \ot I \ket{\psi} \leq 1+2\sqrt{\zeta}
			\quad\text{and}\quad
			\sum_a \bra{\psi} (Q_a-A_a)^2 \ot I \ket{\psi} \leq 12\sqrt{\zeta}
		\]
		from \ref{lem:projective-low-rank-sum}.
		A second Cauchy--Schwarz estimate yields
		\begin{equation}\label{eq:another-Q-for-A}
			\sum_a \bra{\psi} (Q_a \cdot A_a)\ot I \ket{\psi}
			\approx_{4\zeta^{1/4}}
			\sum_a \bra{\psi} (A_a)^2 \ot I \ket{\psi}.
		\end{equation}
		Combining \ref{eq:Q-for-an-A} and \ref{eq:another-Q-for-A} with \ref{eq:A-looks-projective} gives
		\[
			\bra{\psi} Q \ot I \ket{\psi}
			\geq
			\sum_a \bra{\psi} A_a \ot I \ket{\psi} - 2\zeta - 9\zeta^{1/4}
			=
			1 - 11\zeta^{1/4},
		\]
		since $\zeta \leq \zeta^{1/4}$ under \ref{eq:assumption-on-zeta}.
	\end{proof}

	\begin{lemma}[Completeness of~$\sqrt{Q}$]\label{lem:sqrt-Q-completeness}
		\lean{MIPStarRE.LDT.MakingMeasurementsProjective.sqrtQCompleteness}
		\leanok
		\begin{equation*}
			\bra{\psi} \sqrt{Q} \ot I \ket{\psi}
			\geq
			1 - 12 \zeta^{1/4}.
		\end{equation*}
	\end{lemma}
	\begin{proof}
		\leanok
		\uses{lem:Q-completeness, lem:projective-low-rank-sum}
		Diagonalize $Q=\sum_i \nu_i \ket{u_i}\bra{u_i}$. Because $Q \leq (1+2\sqrt{\zeta})I$ by \ref{lem:projective-low-rank-sum}, each $\nu_i$ is at most $1+2\sqrt{\zeta}$. Therefore
		\[
			\sqrt{Q}
			=
			\sum_i \sqrt{\nu_i}\ket{u_i}\bra{u_i}
			\geq
			\frac{1}{\sqrt{1+2\sqrt{\zeta}}} Q.
		\]
		The scalar estimate
		\[
			\frac{1}{\sqrt{1+2\sqrt{\zeta}}} \geq 1-\sqrt{\zeta}
		\]
		gives $\sqrt{Q} \geq (1-\sqrt{\zeta})Q$. Hence \ref{lem:Q-completeness} implies
		\[
			\bra{\psi} \sqrt{Q} \ot I \ket{\psi}
			\geq
			(1-\sqrt{\zeta})(1-11\zeta^{1/4})
			\geq
			1-12\zeta^{1/4}.
		\]
	\end{proof}

	\begin{lemma}[$Q$ is almost projective]\label{lem:q-almost-projective}
		\lean{MIPStarRE.LDT.MakingMeasurementsProjective.qAlmostProjective}
		\leanok
		\begin{equation*}
			\sum_a (Q_a \cdot Q \cdot Q_a - Q_a) \leq 4\sqrt{\zeta} \cdot I.
		\end{equation*}
	\end{lemma}
	\begin{proof}
		\leanok
		\uses{lem:projective-low-rank-sum}
		Since $Q \leq (1+2\sqrt{\zeta})I$ by \ref{lem:projective-low-rank-sum},
		\begin{align*}
			\sum_a Q_a \cdot Q \cdot Q_a - \sum_a Q_a
			 & \leq
			(1+2\sqrt{\zeta}) \sum_a Q_a \cdot I \cdot Q_a - \sum_a Q_a \\
			 & =
			(1+2\sqrt{\zeta})Q - Q                                      \\
			 & =
			2\sqrt{\zeta}\,Q                                            \\
			 & \leq
			2\sqrt{\zeta}(1+2\sqrt{\zeta})I                             \\
			 & \leq
			4\sqrt{\zeta}\,I,
		\end{align*}
		using \ref{eq:assumption-on-zeta} in the last step.
	\end{proof}

	\begin{definition}[Matrix decomposition of~$Q_a$]\label{def:matrix-decomposition-Q}
		\lean{MIPStarRE.LDT.MakingMeasurementsProjective.matrixDecompositionQ, MIPStarRE.LDT.MakingMeasurementsProjective.sigmaProjMeas, MIPStarRE.LDT.MakingMeasurementsProjective.sigmaFinProjMeas, MIPStarRE.LDT.MakingMeasurementsProjective.sigmaRangeEmbedding, MIPStarRE.LDT.MakingMeasurementsProjective.sigmaFinRangeEmbedding}
		\lean{MIPStarRE.Quantum.ProjectorRangeONB.rankOne_mulVec_vec}
		\lean{MIPStarRE.Quantum.ProjectorRangeONB.mulVec_vec}
		\leanok
		\uses{lem:projective-low-rank-sum}
		For each $a$, let
		\[
			m_a = \mathrm{rank}(Q_a),
		\]
		and let $\ket{v_{a,1}},\ldots,\ket{v_{a,m_a}}$ be an orthonormal basis for the range of $Q_a$, so that
		\[
			Q_a = \sum_{i=1}^{m_a} \ket{v_{a,i}}\bra{v_{a,i}}.
		\]
		Let
		\[
			m = \sum_a m_a,
		\]
		and let $\{\ket{a,i}\}$ be an orthonormal basis of $\C^m$, indexed by all $a$ and all
		$1 \leq i \leq m_a$. For each $a$, define
		\[
			X_a = \sum_{i=1}^{m_a} \ket{a,i}\bra{v_{a,i}},
			\qquad
			T_a = \sum_{i=1}^{m_a} \ket{a,i}\bra{a,i},
		\]
		and define
		\begin{equation}\label{eq:looks-like-singular-value-decomposition}
			X = \sum_a X_a = \sum_a \sum_{i=1}^{m_a} \ket{a,i}\bra{v_{a,i}}.
		\end{equation}
		Let $T = \{T_a\}$. Then $T$ is a projective measurement on $\C^m$.
	\end{definition}
	\begin{proof}
		\lean{MIPStarRE.LDT.MakingMeasurementsProjective.matrixDecompositionQ,
			MIPStarRE.LDT.MakingMeasurementsProjective.sigmaProjMeas,
			MIPStarRE.LDT.MakingMeasurementsProjective.sigmaFinProjMeas,
			MIPStarRE.LDT.MakingMeasurementsProjective.sigmaRangeEmbedding,
			MIPStarRE.LDT.MakingMeasurementsProjective.sigmaFinRangeEmbedding}
		\leanok
		\uses{lem:projective-low-rank-sum}
		Choose an orthonormal basis of the range of each projector \(Q_a\).
		The rank-one expansion gives the displayed decomposition, and the
		orthogonality of the labelled basis vectors makes the family \(T\)
		a projective measurement.
	\end{proof}

	\begin{definition}[SVD of~$X$]\label{def:svd-of-X}
		\lean{MIPStarRE.LDT.MakingMeasurementsProjective.svdOfX}
		\leanok
		\lean{MIPStarRE.LDT.MakingMeasurementsProjective.rectangularSvdX}
		\leanok
		\lean{MIPStarRE.LDT.MakingMeasurementsProjective.rectangularSvdXHat}
		\leanok
		\lean{MIPStarRE.LDT.MakingMeasurementsProjective.rectangularSvdDomainMiddle}
		\leanok
		\lean{MIPStarRE.LDT.MakingMeasurementsProjective.rectangularSvdRangeMiddle}
		\leanok
		\lean{MIPStarRE.LDT.MakingMeasurementsProjective.sigmaRangeCarrier}
		\leanok
		\lean{MIPStarRE.LDT.MakingMeasurementsProjective.QXPLayerData.ofRankReductionAndSvdIdentities}
		\leanok
		\lean{MIPStarRE.LDT.MakingMeasurementsProjective.QXPLayerData.ofSigmaRangeAndSvdIdentities}
		\leanok
		\lean{MIPStarRE.LDT.MakingMeasurementsProjective.QXPLayerData.ofRankReductionAndRectangularSvdUnitaryGroup}
		\leanok
		\lean{MIPStarRE.LDT.MakingMeasurementsProjective.QXPLayerData.ofSigmaRangeAndRectangularSvdUnitaryGroup}
		\leanok
		\lean{MIPStarRE.LDT.MakingMeasurementsProjective.exists_qxpLayerData_ofRankReductionSigmaRangeAndSvdIdentities}
		\leanok
		\lean{MIPStarRE.LDT.MakingMeasurementsProjective.exists_qxpLayerData_ofRankReductionSigmaRangeAndSvdIdentities_with_x_coisometry}
		\leanok
		\lean{MIPStarRE.LDT.MakingMeasurementsProjective.exists_qxpLayerData_ofRankReductionSigmaRangeAndRectangularSvdUnitaryGroup}
		\leanok
		\uses{def:matrix-decomposition-Q}
		The paper presents this step by choosing a singular value decomposition
		\[
			X = U \cdot \Sigma_{m \times d} \cdot V^\dagger.
		\]
		The formalization has two closely related auxiliary declarations.  The reduced
		datum records the downstream consequences of the choice: the matrices $X$
		and $\widehat X$, the identities $X^\dagger X=Q$,
		$\widehat X\widehat X^\dagger=I$, and $X^\dagger\widehat X=\sqrt Q$,
		together with the $Q_a=X^\dagger T_aX$ restatement.  The rectangular-SVD
		lemmas take explicit matrices $U,V,\Sigma,I_{m\times d}$ and derive the
		two stored $\widehat X$ identities from the matrix calculation below.
	\end{definition}
	\begin{proof}
		\lean{MIPStarRE.LDT.MakingMeasurementsProjective.svdOfX,
			MIPStarRE.LDT.MakingMeasurementsProjective.QXPLayerData.ofRankReductionAndSvdIdentities,
			MIPStarRE.LDT.MakingMeasurementsProjective.QXPLayerData.ofSigmaRangeAndSvdIdentities,
			MIPStarRE.LDT.MakingMeasurementsProjective.QXPLayerData.ofRankReductionAndRectangularSvdUnitaryGroup,
			MIPStarRE.LDT.MakingMeasurementsProjective.QXPLayerData.ofSigmaRangeAndRectangularSvdUnitaryGroup,
			MIPStarRE.LDT.MakingMeasurementsProjective.exists_qxpLayerData_ofRankReductionSigmaRangeAndSvdIdentities,
			MIPStarRE.LDT.MakingMeasurementsProjective.exists_qxpLayerData_ofRankReductionSigmaRangeAndRectangularSvdUnitaryGroup}
		\leanok
		\uses{def:matrix-decomposition-Q}
		Apply the finite-dimensional singular value decomposition to \(X\) and
		record exactly the consequences used later: the identities for
		\(X^\dagger X\), \(\widehat X\widehat X^\dagger\),
		\(X^\dagger\widehat X\), and the restatement of the \(Q_a\)'s.  The
		unitary-group constructors are the same rectangular-SVD calculation, with
		the square factors taken as elements of the relevant unitary groups
		rather than as matrices accompanied by separate unitarity equations.  In
		the sigma-range specialization, the auxiliary row space is the lifted
		finite-enumeration carrier associated to the ranks of the \(Q_a\)'s, which
		is definitionally the carrier of the canonical sigma-range \(Q\)-layer.
	\end{proof}

	\begin{remark}[Formalization support for orthogonal projector families]\label{rem:motivation-for-P}
		If the family $\{Q_a\}$ were a projective measurement, then the vectors
		$\ket{v_{a,i}}$, over all $a$ and all $1 \leq i \leq m_a$, would be orthonormal.
		Then \eqref{eq:looks-like-singular-value-decomposition} would already be a singular value
		decomposition of $X$, with $\Sigma = I_{m \times d}$. Equivalently,
		\[
			X = U \cdot I_{m \times d} \cdot V^\dagger.
		\]
		This motivates replacing $\Sigma$ by $I$ in the definition of $\widehat X$ and $P$.
	\end{remark}

	\begin{definition}[Definition of~$P$]\label{def:projective-P}
		\lean{MIPStarRE.LDT.MakingMeasurementsProjective.projectiveP}
		\leanok
		\uses{def:matrix-decomposition-Q, def:svd-of-X}
		The paper defines
		\[
			\widehat{X} = U \cdot I_{m \times d} \cdot V^\dagger.
		\]
		We take $\widehat X$ to be the reduced matrix supplied by the stored datum.
		For each $a$, define
		\[
			\widehat{X}_a = T_a \cdot \widehat{X},
			\qquad
			P_a = \widehat{X}_a^\dagger \cdot \widehat{X}_a.
		\]
	\end{definition}
	\begin{proof}
		\lean{MIPStarRE.LDT.MakingMeasurementsProjective.projectiveP}
		\leanok
		\uses{def:matrix-decomposition-Q, def:svd-of-X}
		The definition is the formal \(Q/X/\widehat X/P\) construction: once
		\(\widehat X\) is chosen from the stored singular-value data, each
		\(P_a\) is defined as the positive operator
		\(\widehat X_a^\dagger\widehat X_a\).
	\end{proof}

	\begin{lemma}\label{lem:xa-t}
		\lean{MIPStarRE.LDT.MakingMeasurementsProjective.xa_t}
		\leanok
		\uses{def:matrix-decomposition-Q}
		For each $a$,
		\[
			X_a = T_a \cdot X.
		\]
	\end{lemma}
	\begin{proof}
		\leanok
		Multiply the definition of $X$ by $T_a$ on the left and use that $T_a$ selects exactly the basis vectors $\ket{a,i}$:
		\[
			T_a \cdot X
			=
			\Big(\sum_{i=1}^{m_a} \ket{a,i}\bra{a,i}\Big)
			\Big(\sum_b \sum_{j=1}^{m_b} \ket{b,j}\bra{v_{b,j}}\Big)
			=
			\sum_{i=1}^{m_a} \ket{a,i}\bra{v_{a,i}}
			=
			X_a.
		\]
	\end{proof}

	\begin{lemma}[$Q_a$ restated]\label{lem:qa-restated}
		\lean{MIPStarRE.LDT.MakingMeasurementsProjective.sigmaRangeEmbedding_qa_eq, MIPStarRE.LDT.MakingMeasurementsProjective.sigmaFinRangeEmbedding_qa_eq, MIPStarRE.LDT.MakingMeasurementsProjective.qaRestated}
		\leanok
		For each $a$,
		\begin{equation*}
			Q_a = X_a^\dagger \cdot X_a = X^\dagger \cdot T_a \cdot X = X_a^\dagger \cdot X.
		\end{equation*}
	\end{lemma}
	\begin{proof}
		\leanok
		\uses{lem:xa-t}
		The identity $X_a^\dagger X_a = Q_a$ follows by expanding the basis sums. The remaining equalities come from \ref{lem:xa-t} and the projectivity of $T$:
		\[
			X_a^\dagger \cdot X_a
			=
			(X^\dagger \cdot T_a)(T_a \cdot X)
			=
			X^\dagger \cdot T_a \cdot X
			=
			X_a^\dagger \cdot X.
		\]
	\end{proof}

	\begin{lemma}[$X$ squared]\label{lem:X-squared}
		\lean{MIPStarRE.LDT.MakingMeasurementsProjective.xSquared_of_qa_eq, MIPStarRE.LDT.MakingMeasurementsProjective.sigmaFinRangeEmbedding_gram_right}
		\lean{MIPStarRE.LDT.MakingMeasurementsProjective.qtotal_posSemidef_of_x_squared}
		\lean{MIPStarRE.LDT.MakingMeasurementsProjective.qtotal_eigenvalues_nonneg_of_x_squared}
		\lean{MIPStarRE.LDT.MakingMeasurementsProjective.sqrt_eq_sum_sqrt_eigenvalues_vecMulVec}
		\lean{MIPStarRE.LDT.MakingMeasurementsProjective.positiveGramSpectrumRightRows}
		\lean{MIPStarRE.LDT.MakingMeasurementsProjective.positive_gram_spectrum_right_rows_mul_conjTranspose}
		\lean{MIPStarRE.LDT.MakingMeasurementsProjective.positive_gram_spectrum_image_rows_mixed_eq_sqrt}
		\lean{MIPStarRE.LDT.MakingMeasurementsProjective.matrix_image_eq_zero_of_nonpositive_gram_eigenvalue}
		\lean{MIPStarRE.LDT.MakingMeasurementsProjective.adjoint_image_eq_zero_of_orthogonal_positive_gram_images}
		\leanok
		\begin{equation*}
			X^\dagger \cdot X = Q.
		\end{equation*}
	\end{lemma}
	\begin{proof}
		\leanok
		\uses{lem:qa-restated}
		Use that $T$ is a measurement and then \ref{lem:qa-restated}:
		\[
			X^\dagger \cdot X
			=
			\sum_a X^\dagger \cdot T_a \cdot X
			=
			\sum_a Q_a
			=
			Q.
		\]
		The additional paper identities involving $U,V,\Sigma$ are the usual SVD
		computations.  The Lean development also records the spectral calculation
		which will identify the mixed positive-Gram-image product with
		$\sqrt Q$: the positive eigenvalues give the nonzero columns, while the
		zero eigenspace has zero coefficient in the square-root expansion and is
		killed by $X$.  The conjugate right-eigenvector rows over the positive
		spectrum are also orthonormal.  Equivalently, any row-space vector
		orthogonal to the positive Gram images is killed by $X^\dagger$.
	\end{proof}

	\begin{lemma}\label{lem:X-expression-to-Q-expression}
		\lean{MIPStarRE.LDT.MakingMeasurementsProjective.xExpressionToQExpression}
		\leanok
		For each $a$,
		\begin{equation*}
			X_a^\dagger \cdot (X \cdot X^\dagger - I_{m \times m})^2 \cdot X_a
			=
			Q_a \cdot Q \cdot Q_a - Q_a.
		\end{equation*}
	\end{lemma}
	\begin{proof}
		\leanok
		\uses{lem:X-squared, lem:qa-restated}
		Expand the square and apply \ref{lem:X-squared} and \ref{lem:qa-restated} term by term:
		\[
			X_a^\dagger (X \cdot X^\dagger \cdot X \cdot X^\dagger) X_a
			=
			Q_a \cdot Q \cdot Q_a,
		\]
		while
		\[
			X_a^\dagger (X \cdot X^\dagger) X_a
			=
			(X_a^\dagger \cdot X)(X^\dagger \cdot X_a)
			=
			Q_a \cdot Q_a
			=
			Q_a,
		\]
		and $X_a^\dagger I_{m \times m} X_a = X_a^\dagger X_a = Q_a$. Substituting these identities into
		\[
			(X \cdot X^\dagger - I_{m \times m})^2
			=
			X \cdot X^\dagger \cdot X \cdot X^\dagger - 2X \cdot X^\dagger + I_{m \times m}
		\]
		gives the claim.
	\end{proof}

	\begin{lemma}[$P_a$ restated]\label{lem:pa-restated}
		\lean{MIPStarRE.LDT.MakingMeasurementsProjective.paRestated}
		\leanok
		For each $a$,
		\begin{equation*}
			P_a = \widehat{X}^\dagger \cdot T_a \cdot \widehat{X} = \widehat{X}_a^\dagger \cdot \widehat{X}.
		\end{equation*}
	\end{lemma}
	\begin{proof}
		\leanok
		\uses{def:projective-P}
		Since $T$ is projective,
		\[
			P_a
			=
			\widehat{X}_a^\dagger \cdot \widehat{X}_a
			=
			(\widehat{X}^\dagger \cdot T_a)(T_a \cdot \widehat{X})
			=
			\widehat{X}^\dagger \cdot T_a \cdot \widehat{X}
			=
			\widehat{X}_a^\dagger \cdot \widehat{X}.
		\]
	\end{proof}

	\begin{lemma}[$\widehat{X}$ squared]\label{lem:X-hat-squared}
		\lean{Matrix.mul_conjTranspose_eq_one_of_orthonormal_rows}
		\lean{MIPStarRE.LDT.MakingMeasurementsProjective.orthonormal_normalized_image_of_adjoint_comp_eigenvectors}
		\lean{MIPStarRE.LDT.MakingMeasurementsProjective.orthonormal_normalized_matrix_image_of_gram_eigenvectors}
		\lean{MIPStarRE.LDT.MakingMeasurementsProjective.orthonormal_normalized_matrix_image_of_positive_gram_spectrum}
		\lean{MIPStarRE.LDT.MakingMeasurementsProjective.normalizedMatrixImageRows}
		\lean{MIPStarRE.LDT.MakingMeasurementsProjective.positiveGramSpectrumImageRows}
		\lean{MIPStarRE.LDT.MakingMeasurementsProjective.normalized_matrix_image_rows_mul_conjTranspose}
		\lean{MIPStarRE.LDT.MakingMeasurementsProjective.normalized_matrix_image_rows_transpose_mixed}
		\lean{MIPStarRE.LDT.MakingMeasurementsProjective.normalized_matrix_image_rows_mul_conjTranspose_of_positive_gram_spectrum}
		\lean{MIPStarRE.LDT.MakingMeasurementsProjective.positive_gram_spectrum_image_rows_transpose_mixed}
		\lean{MIPStarRE.LDT.MakingMeasurementsProjective.positive_gram_spectrum_card_le_rows}
		\lean{MIPStarRE.LDT.MakingMeasurementsProjective.exists_unitaryGroup_rows_extending_orthonormal}
		\lean{MIPStarRE.LDT.MakingMeasurementsProjective.exists_unitaryGroup_with_positive_gram_spectrum_rows}
		\lean{MIPStarRE.LDT.MakingMeasurementsProjective.exists_unitaryGroup_with_positive_gram_spectrum_rows_of_card}
		\lean{MIPStarRE.LDT.MakingMeasurementsProjective.exists_rectangular_coisometry_extending_orthonormal_rows}
		\lean{MIPStarRE.LDT.MakingMeasurementsProjective.exists_rectangular_coisometry_with_positive_gram_spectrum_right_rows}
		\lean{MIPStarRE.LDT.MakingMeasurementsProjective.transpose_unitaryGroup_mul_rectangular_coisometry}
		\lean{MIPStarRE.LDT.MakingMeasurementsProjective.unitaryGroup_mul_rectangular_coisometry}
		\lean{MIPStarRE.LDT.MakingMeasurementsProjective.rectangular_coisometry_mul_conjTranspose_unitaryGroup}
		\lean{MIPStarRE.LDT.MakingMeasurementsProjective.exists_unitaryGroup_with_qxp_positive_gram_spectrum_rows}
		\lean{MIPStarRE.LDT.MakingMeasurementsProjective.xHat_of_positive_gram_spectrum_unitaryGroup_choices}
		\lean{MIPStarRE.LDT.MakingMeasurementsProjective.qtotal_isHermitian_of_x_squared}
		\lean{MIPStarRE.LDT.MakingMeasurementsProjective.rectangularSvd_xHat_coisometry_unitaryGroup}
		\lean{MIPStarRE.LDT.MakingMeasurementsProjective.exists_xHat_of_positive_gram_spectrum}
		\lean{MIPStarRE.LDT.MakingMeasurementsProjective.exists_xHat_of_sigmaFinRangeEmbedding_positiveGram}
		\lean{MIPStarRE.LDT.MakingMeasurementsProjective.exists_qxpLayerData_ofRankReductionSigmaRangePositiveGram}
		\lean{MIPStarRE.LDT.MakingMeasurementsProjective.exists_qxpLayerData_ofRankReductionSigmaRangePositiveGram_with_x_coisometry}
		\lean{MIPStarRE.LDT.MakingMeasurementsProjective.exists_sigmaFin_xHat_coisometry_of_sum_le}
		\lean{MIPStarRE.LDT.MakingMeasurementsProjective.sigmaFinXHatCoisometry_spec}
		\leanok
		\uses{def:svd-of-X}
		\begin{equation*}
			\widehat{X} \cdot \widehat{X}^\dagger = I_{m \times m}.
		\end{equation*}
	\end{lemma}
	\begin{proof}
		\leanok
		The paper proves this by multiplying out
		$\widehat{X}=U\cdot I_{m\times d}\cdot V^\dagger$ and using the unitarity of
		$U$ and $V$.  In Lean, the dimension hypothesis $m\leq d$ already gives a
		rectangular matrix with orthonormal rows on the sigma auxiliary space; the
		chosen matrix \texttt{sigmaFinXHatCoisometry} and its specification record this
		existence result as a hypothesis for the future construction.  The formal lemma
		\texttt{orthonormal\_normalized\_image\_of\_adjoint\_comp\_eigenvectors}
		records the singular-vector normalization calculation, while
		\texttt{normalized\_matrix\_image\_rows\_mul\_conjTranspose} combines the
		corresponding row-coisometry identity, and
		\texttt{positive\_gram\_spectrum\_image\_rows\_transpose\_mixed} records the
		adjoint action $X^\dagger(\lambda^{-1/2}Xv)=\lambda^{1/2}v$ on the positive
		spectral part of the Gram operator.  The formal lemma
		\texttt{exists\_unitaryGroup\_with\_positive\_gram\_spectrum\_rows} records the
		next finite-dimensional step: after choosing distinct row positions for the
		positive spectrum, these normalized image rows extend to a square unitary
		group element on the row space.  The companion lemma
		\texttt{exists\_rectangular\_coisometry\_extending\_orthonormal\_rows}
		records the analogous rectangular completion for prescribed right rows
		when the dimension bound permits it.  The formal lemma
		\texttt{unitaryGroup\_mul\_rectangular\_coisometry} and its right-multiplication
		companion
		\texttt{rectangular\_coisometry\_mul\_conjTranspose\_unitaryGroup} record the
		two unitary-preservation steps in the paper's multiplication.  The formal
		lemma \texttt{rectangularSvd\_xHat\_coisometry\_unitaryGroup} applies these
		preservation steps to the rectangular-SVD expression.  The
		unconditional positive-Gram route is
		\texttt{exists\_xHat\_of\_sigmaFinRangeEmbedding\_positiveGram}, which
		constructs $\widehat X$ from the rank-reduction witness and proves
		$\widehat X\widehat X^\dagger=I$ together with the mixed square-root
		identity below.
	\end{proof}

	\begin{lemma}[$X$ times $\widehat{X}$]\label{lem:X-times-X-hat}
		\lean{MIPStarRE.LDT.MakingMeasurementsProjective.adjoint_image_eq_zero_of_unitary_positive_gram_completion_row}
		\lean{MIPStarRE.LDT.MakingMeasurementsProjective.positive_gram_selected_left_unitary_mixed_column}
		\lean{MIPStarRE.LDT.MakingMeasurementsProjective.positive_gram_completion_left_unitary_mixed_column_eq_zero}
		\lean{MIPStarRE.LDT.MakingMeasurementsProjective.positive_gram_polar_extension_mixed_eq_positive_rows}
		\lean{MIPStarRE.LDT.MakingMeasurementsProjective.positive_gram_polar_extension_mixed_eq_sqrt_unitaryGroup}
		\lean{MIPStarRE.LDT.MakingMeasurementsProjective.rectangularSvd_xHat_mixed_raw_unitaryGroup}
		\lean{MIPStarRE.LDT.MakingMeasurementsProjective.rectangularSvd_xHat_mixed_raw}
		\lean{MIPStarRE.LDT.MakingMeasurementsProjective.rectangularSvd_x_mul_xHat_conjTranspose_raw_unitaryGroup}
		\lean{MIPStarRE.LDT.MakingMeasurementsProjective.rectangularSvd_x_mul_xHat_conjTranspose_raw}
		\lean{MIPStarRE.LDT.MakingMeasurementsProjective.exists_xHat_of_positive_gram_spectrum}
		\lean{MIPStarRE.LDT.MakingMeasurementsProjective.exists_xHat_of_sigmaFinRangeEmbedding_positiveGram}
		\lean{MIPStarRE.LDT.MakingMeasurementsProjective.rectangularSvd_xHat_mixed_of_sqrtQ_unitaryGroup}
		\lean{MIPStarRE.LDT.MakingMeasurementsProjective.exists_xHat_of_rectangularSvd_unitaryGroup}
		\lean{MIPStarRE.LDT.MakingMeasurementsProjective.rectangularSvd_middle_eq_sqrt_of_square}
		\lean{MIPStarRE.LDT.MakingMeasurementsProjective.x_mul_xHat_adjoint_spectral_theorem}
		\lean{MIPStarRE.LDT.MakingMeasurementsProjective.QXPLayerData.ofRankReductionAndRectangularSvdSquareRootUnitaryGroup}
		\lean{MIPStarRE.LDT.MakingMeasurementsProjective.QXPLayerData.ofSigmaRangeAndRectangularSvdSquareRootUnitaryGroup}
		\lean{MIPStarRE.LDT.MakingMeasurementsProjective.exists_qxpLayerData_ofRankReductionSigmaRangeAndRectangularSvdSquareRootUnitaryGroup}
		\leanok
		\begin{equation*}
			X \cdot \widehat{X}^\dagger =
			U \cdot \Sigma_{m \times m} \cdot U^\dagger,
			\quad
			X^\dagger \cdot \widehat{X} = \sqrt{Q}.
		\end{equation*}
	\end{lemma}
	\begin{proof}
		\leanok
		\uses{lem:X-squared}
		The paper obtains this from the SVD formula for $X$ and the definition of
		$\widehat X$.  In Lean, the unconditional QXP data first proves that
		$X\widehat X^\dagger$ is Hermitian and positive.  The formal lemma
		\texttt{x\_mul\_xHat\_adjoint\_spectral\_theorem} then applies Mathlib's
		spectral theorem for Hermitian matrices, giving the displayed
		$U\Sigma U^\dagger$ form with the Mathlib eigenvector unitary and the
		diagonal matrix of real eigenvalues.  The auxiliary lemma
		\texttt{rectangularSvd\_x\_mul\_xHat\_conjTranspose\_raw} records the same
		calculation under explicit rectangular-SVD data, where the middle factor
		\(S I_{m\times d}^\dagger\) is the formal square matrix
		\(\Sigma_{m\times m}\).  The formal lemma
		\texttt{rectangularSvd\_xHat\_mixed\_of\_sqrtQ\_unitaryGroup} records the
		second matrix calculation in the form whose square-root target is the
		operator $Q$ supplied by the QXP construction.  The formal lemma
		\texttt{rectangularSvd\_middle\_eq\_sqrt\_of\_square} records the spectral
		step: a positive middle factor whose square is $Q$ is the positive square
		root of $Q$.  The reduced datum then records this mixed-product
		identity as a primitive consequence, which is exactly the part needed in
		\ref{lem:squared-difference} and \ref{lem:P-Q-approx}.
	\end{proof}

	\begin{remark}[Lean status for the QXP algebra]\label{rem:qxp-layer-lean-status}
		The Lean declarations linked in Lemmas~\ref{lem:X-squared},
		\ref{lem:X-hat-squared}, and~\ref{lem:X-times-X-hat} expose the
		corresponding matrix identities as construction theorems.  These
		declarations prove the identities from the \(Q\)-layer, the sigma-range
		realization, and the positive-Gram spectral construction, with the
		rectangular-SVD calculation retained as an auxiliary comparison.  Thus the
		local \(X/\widehat X/P\) algebra is formalized; the remaining chapter-level
		work is to compose this construction with the preceding orthonormalization
		step from the paper hypotheses.
	\end{remark}

	\begin{lemma}[Squared difference]\label{lem:squared-difference}
		\lean{MIPStarRE.LDT.MakingMeasurementsProjective.xHat_adjoint_mul_x_eq_sqrt}
		\lean{MIPStarRE.LDT.MakingMeasurementsProjective.x_mul_xHat_adjoint_isHermitian}
		\lean{MIPStarRE.LDT.MakingMeasurementsProjective.x_mul_xHat_adjoint_sq}
		\lean{MIPStarRE.LDT.MakingMeasurementsProjective.x_mul_xHat_adjoint_nonneg}
		\lean{MIPStarRE.LDT.MakingMeasurementsProjective.squaredDifference}
		\leanok
		\begin{equation*}
			(X - \widehat{X}) \cdot (X - \widehat{X})^\dagger \leq (X \cdot X^\dagger - I_{m \times m})^2.
		\end{equation*}
	\end{lemma}
	\begin{proof}
		\leanok
		\uses{lem:X-squared, lem:X-hat-squared, lem:X-times-X-hat}
		The paper diagonalizes the positive operator $X\widehat X^\dagger$ using the
		SVD notation.  The formal proof avoids storing that SVD explicitly: set
		$Y=X\widehat X^\dagger$.  The primitive identities imply that $Y$ is
		hermitian, nonnegative, and satisfies $Y^2=XX^\dagger$.  Hence
		$(X-\widehat X)(X-\widehat X)^\dagger=(Y-I)^2$, and positivity of $Y$ gives
		$(Y-I)^2\leq(Y^2-I)^2=(XX^\dagger-I)^2$.
	\end{proof}

	\begin{lemma}[Projectivity of~$P$]\label{lem:P-projectivity}
		\lean{MIPStarRE.LDT.MakingMeasurementsProjective.pProjectivity}
		\lean{MIPStarRE.LDT.MakingMeasurementsProjective.sum_pa_eq_xHat_adjoint_mul_xHat}
		\lean{MIPStarRE.LDT.MakingMeasurementsProjective.qxpProjSubMeas_total_eq_xHat_adjoint_mul_xHat}
		\lean{MIPStarRE.LDT.MakingMeasurementsProjective.qxpProjSubMeas_total_ev_eq_sum_pa_ev}
		\lean{MIPStarRE.LDT.MakingMeasurementsProjective.qxpProjSubMeas_total_le_of_outcome_le}
		\lean{MIPStarRE.LDT.MakingMeasurementsProjective.QXPLayerTotalDomination}
		\lean{MIPStarRE.LDT.MakingMeasurementsProjective.QXPLayerTotalDomination.of_outcome_le}
		\leanok
		$P = \{P_a\}$ forms a projective sub-measurement.
	\end{lemma}
	\begin{proof}
		\leanok
		\uses{lem:pa-restated, lem:X-hat-squared}
		For outcomes $a,b$,
		\begin{align*}
			P_a \cdot P_b
			 & =
			(\widehat{X}^\dagger \cdot T_a \cdot \widehat{X})
			(\widehat{X}^\dagger \cdot T_b \cdot \widehat{X})                              \\
			 & =
			\widehat{X}^\dagger \cdot T_a \cdot I_{m \times m} \cdot T_b \cdot \widehat{X} \\
			 & =
			\widehat{X}^\dagger \cdot T_a \cdot \widehat{X} \cdot \bone[a=b]               \\
			 & =
			P_a \cdot \bone[a=b],
		\end{align*}
		where the second line uses \ref{lem:X-hat-squared}, the third uses that $T$ is projective, and the last uses \ref{lem:pa-restated}. Summing over $a$ shows $\sum_a P_a \leq I$, so $P$ is a projective sub-measurement.  The formalization also records the sharper identity
		\[
			\sum_a P_a=\widehat X^\dagger\widehat X,
		\]
		and the resulting total-domination estimates used later in the monotone-total self-improvement route.
	\end{proof}

	\begin{lemma}[$P$ is close to~$Q$]\label{lem:P-Q-approx}
		\lean{MIPStarRE.LDT.MakingMeasurementsProjective.qxpMixedCrossSum_eq_sqrt}
		\lean{MIPStarRE.LDT.MakingMeasurementsProjective.pQApprox}
		\lean{MIPStarRE.LDT.MakingMeasurementsProjective.pQApprox_ofRankReductionSigmaRangeAndSvdIdentities}
		\lean{MIPStarRE.LDT.MakingMeasurementsProjective.pQApprox_ofRankReductionSigmaRangePositiveGram}
		\lean{MIPStarRE.LDT.MakingMeasurementsProjective.pQApprox_ofRankReductionSigmaRangePositiveGram_with_x_coisometry}
		\lean{MIPStarRE.LDT.MakingMeasurementsProjective.pQApprox_ofRankReductionSigmaRangeAndRectangularSvdSquareRootUnitaryGroup}
		\leanok
		\begin{equation*}
			Q_a \ot I \approx_{30\zeta^{1/4}} P_a \ot I.
		\end{equation*}
	\end{lemma}
	\begin{proof}
		\leanok
		\uses{lem:projective-low-rank-sum, lem:q-almost-projective, lem:sqrt-Q-completeness, lem:xa-t, lem:qa-restated, lem:X-expression-to-Q-expression, lem:pa-restated, lem:X-hat-squared, lem:X-times-X-hat, lem:squared-difference, lem:P-projectivity}
		The formal proof follows this argument directly from the primitive
		$X$, $\widehat X$, and $P$ identities; the $30\zeta^{1/4}$ estimate is
		derived here rather than stored as part of the reduced datum.  The
		unitary-group rectangular-SVD lemma combines the same estimate with the
		positive-square identification of the middle factor from
		\ref{lem:X-times-X-hat}; the square unitary factors are again carried as
		elements of the relevant unitary groups.  The
		positive-Gram sigma-space lemma also records the coisometry
		\(X X^\dagger=I\) for the original sigma embedding when the projective
		\(Q\)-family is subnormalized; this is the construction-level hypothesis
		later used to obtain the QXP-internal fresh-outcome comparison
		\(Q_{\mathrm{none}} \leq P_{\mathrm{none}}\).  An additional
		source-to-\(Q\) comparison is still needed for source-facing residual
		domination.
		The quantity to bound is
		\begin{equation}\label{eq:P-Q-thing-to-bound}
			\sum_a \bra{\psi} (Q_a-P_a)^2 \ot I \ket{\psi}
			=
			\bra{\psi} Q \ot I \ket{\psi}
			+
			\bra{\psi} P \ot I \ket{\psi}
			-
			\sum_a \bra{\psi} Q_aP_a \ot I \ket{\psi}
			-
			\sum_a \bra{\psi} P_aQ_a \ot I \ket{\psi},
		\end{equation}
		using the projectivity of $Q$ and $P$. The first term is at most $1+2\sqrt{\zeta}$ by \ref{lem:projective-low-rank-sum}, and the second is at most $1$ because $P$ is a sub-measurement.

		For the last two terms,
		\begin{equation}\label{eq:complex-conjugates}
			\sum_a \bra{\psi} Q_aP_a \ot I \ket{\psi}
			+
			\sum_a \bra{\psi} P_aQ_a \ot I \ket{\psi}
			=
			2 \cdot \mathfrak{R}\Big(\sum_a \bra{\psi} Q_aP_a \ot I \ket{\psi}\Big).
		\end{equation}
		By \ref{lem:qa-restated},
		\[
			\sum_a \bra{\psi} Q_aP_a \ot I \ket{\psi}
			=
			\sum_a \bra{\psi} (X_a^\dagger \cdot X \cdot P_a) \ot I \ket{\psi}.
		\]
		The main estimate is
		\begin{equation}\label{eq:add-a-hat}
			\sum_a \bra{\psi} (X_a^\dagger \cdot X \cdot P_a)\ot I \ket{\psi}
			\approx_{2\zeta^{1/4}}
			\sum_a \bra{\psi} (X_a^\dagger \cdot \widehat{X} \cdot P_a)\ot I \ket{\psi}.
		\end{equation}
		This follows by Cauchy--Schwarz:
		\[
			\Big|\sum_a \bra{\psi} ((X_a^\dagger \cdot (X-\widehat X)) \ot I)\cdot(P_a \ot I)\ket{\psi}\Big|
		\]
		is bounded by the square root of
		\[
			\sum_a \bra{\psi} \big(X_a^\dagger \cdot (X-\widehat X)(X-\widehat X)^\dagger \cdot X_a\big)\ot I \ket{\psi},
		\]
		times the square root of $\sum_a \bra{\psi}(P_a)^2 \ot I \ket{\psi} \leq 1$. The first factor is at most $4\sqrt{\zeta}$ by \ref{lem:squared-difference}, \ref{lem:X-expression-to-Q-expression}, and \ref{lem:q-almost-projective}, so the total error is at most $2\zeta^{1/4}$.

		The right-hand side of \ref{eq:add-a-hat} is real, because
		\[
			\sum_a \bra{\psi} (X_a^\dagger \cdot \widehat{X} \cdot P_a)\ot I \ket{\psi}
			=
			\bra{\psi} (X^\dagger \cdot \widehat{X}) \ot I \ket{\psi}
			=
			\bra{\psi} \sqrt{Q} \ot I \ket{\psi},
		\]
		where the first equality uses \ref{lem:xa-t}, \ref{lem:pa-restated}, and \ref{lem:X-hat-squared} and that $T$ is a measurement, and the second uses \ref{lem:X-times-X-hat}. Hence \ref{lem:sqrt-Q-completeness} and \ref{eq:add-a-hat} imply
		\[
			\mathfrak{R}\Big(\sum_a \bra{\psi} Q_aP_a \ot I \ket{\psi}\Big)
			\geq
			1-14\zeta^{1/4}.
		\]
		Substituting this into \ref{eq:complex-conjugates} and then into \ref{eq:P-Q-thing-to-bound} gives
		\[
			\sum_a \bra{\psi} (Q_a-P_a)^2 \ot I \ket{\psi}
			\leq
			(1+2\sqrt{\zeta}) + 1 - 2(1-14\zeta^{1/4})
			=
			2\sqrt{\zeta} + 28\zeta^{1/4}
			\leq
			30\zeta^{1/4}.
		\]
	\end{proof}

	Finally, \ref{lem:projective-low-rank-sum} gives
	\[
		A_a \ot I \approx_{12\sqrt{\zeta}} Q_a \ot I,
	\]
	and \ref{lem:P-Q-approx} gives
	\[
		Q_a \ot I \approx_{30\zeta^{1/4}} P_a \ot I.
	\]
	The triangle inequality for $\approx_\delta$ yields
	\[
		A_a \ot I \approx_{84\zeta^{1/4}} P_a \ot I.
	\]
\end{proof}


%%FILE blueprint/src/chapter/ch05_expansion.tex
\chapter{Expansion in the hypercube graph}
\label{chap:expansion}
\label{sec:expansion}

\section{The graph and its spectrum}


\begin{definition}[Rerandomizing one coordinate]\label{def:rerandomize-coord}
	\lean{MIPStarRE.LDT.ExpansionHypercubeGraph.rerandomizeCoord}
	\lean{MIPStarRE.LDT.ExpansionHypercubeGraph.rerandomizeCoordWeight, MIPStarRE.LDT.ExpansionHypercubeGraph.rerandomizeCoord_mass_eq_one}
	\lean{MIPStarRE.LDT.ExpansionHypercubeGraph.RerandomizeCoordSample, MIPStarRE.LDT.ExpansionHypercubeGraph.rerandomizeCoordSampleToPair}
	\lean{MIPStarRE.LDT.ExpansionHypercubeGraph.rerandomizeCoord_isProbability, MIPStarRE.LDT.ExpansionHypercubeGraph.avgOver_rerandomizeCoord_eq_uniform_sample, MIPStarRE.LDT.ExpansionHypercubeGraph.avgOver_rerandomizeCoord_eq_weight_sum}
	\leanok
	For $u \in \F_q^m$, an index $i \in \{1,\dots,m\}$, and $x \in \F_q$, let
	\[
		\operatorname{rerand}_i(u,x) = u + x e_i.
	\]
	This is the move obtained by rerandomizing the $i$-th coordinate of $u$ by a uniform additive shift.
	In Lean, the same edge distribution is represented as the push-forward of the uniform distribution
	on triples $(u,i,x)$ by the map that replaces the $i$-th coordinate by $x$; since $x$ is uniform,
	this is the same distribution as the additive-shift presentation.  The corresponding
	weight-sum identity identifies this push-forward distribution with the explicit edge
	coefficient used in the matrix proof.
\end{definition}

\begin{definition}[Hypercube graph]\label{def:hypercube-graph}
	\lean{MIPStarRE.LDT.ExpansionHypercubeGraph.IsHypercubeEdge, MIPStarRE.LDT.ExpansionHypercubeGraph.rerandomizeCoord}
	\leanok
	\uses{def:rerandomize-coord}
	The hypercube graph $C=(V,E)$ has vertex set $V=\F_q^m$, and an edge between $u$ and $v$ whenever they differ in at most one coordinate. A random edge $(u,v) \sim C$ is sampled by drawing $u \sim \F_q^m$, $i \in \{1,\dots,m\}$, and $x \in \F_q$ uniformly and then setting $v=\operatorname{rerand}_i(u,x)$.
\end{definition}

\begin{definition}[Adjacency matrix and Laplacian]\label{def:adjacency-laplacian}
	\lean{MIPStarRE.LDT.ExpansionHypercubeGraph.adjacency, MIPStarRE.LDT.ExpansionHypercubeGraph.laplacian}
	\leanok
	\uses{def:hypercube-graph}
	Let $M=q^m$. The normalized adjacency matrix of $C$ is
	\[
		K = \E_{(u,v)\sim C} \ket{u}\bra{v},
	\]
	and the Laplacian is
	\[
		L = \frac{1}{M}I-K.
	\]
\end{definition}

\begin{lemma}[Laplacian rewrite]\label{prop:laplacian-rewrite}
	\lean{MIPStarRE.LDT.ExpansionHypercubeGraph.laplacianRewrite}
	\uses{def:adjacency-laplacian}
	\leanok
	The Laplacian can be written as
	\[
		L = \frac12 \E_{(u,v)\sim C} (\ket{u}-\ket{v})(\bra{u}-\bra{v}).
	\]
\end{lemma}
\begin{proof}
	\lean{MIPStarRE.LDT.GlobalVariance.avgOver_rerandomizeCoord_fst, MIPStarRE.LDT.GlobalVariance.avgOver_rerandomizeCoord_snd}
	\leanok
	The two vertex marginals of a random edge are uniform on $\F_q^m$. Expanding the right-hand side yields $(1/M)I-K$.
\end{proof}

\begin{lemma}[Average of a nontrivial additive character]\label{lem:character-average-scalar}
	\lean{MIPStarRE.LDT.Preliminaries.fourier_fact_scalar}
	\leanok
	\uses{def:ff-trace}
	For $\beta \in \F_q$,
	\[
		\E_{x \sim \F_q} \omega^{\operatorname{tr}(\beta x)} =
		\begin{cases}
			1 & \text{if } \beta=0,    \\
			0 & \text{if } \beta\ne 0.
		\end{cases}
	\]
\end{lemma}
\begin{proof}
	\leanok
	If $\beta=0$ the character is constant. If $\beta\ne 0$, the map $x \mapsto \beta x$ permutes $\F_q$, so the average is the average of a nontrivial additive character, which vanishes.
\end{proof}

\begin{lemma}[Orthogonality of finite-field characters]\label{lem:character-average-vector}
	\lean{MIPStarRE.LDT.Preliminaries.fourier_fact_vector}
	\leanok
	For $\alpha,\beta \in \F_q^m$,
	\[
		\frac{1}{q^m}\sum_{u \in \F_q^m} \omega^{\operatorname{tr}(u\cdot(\beta-\alpha))} =
		\begin{cases}
			1 & \text{if } \alpha=\beta,    \\
			0 & \text{if } \alpha\ne \beta.
		\end{cases}
	\]
\end{lemma}
\begin{proof}
	\leanok
	\uses{prop:fourier-fact-vector}
	This is Lemma~\ref{prop:fourier-fact-vector} applied to $v=\beta-\alpha$.
\end{proof}

\begin{lemma}[Eigenvectors of the hypercube graph]\label{prop:eigenvectors}
	\lean{MIPStarRE.LDT.ExpansionHypercubeGraph.eigenvectors_orthonormality}
	\lean{MIPStarRE.LDT.ExpansionHypercubeGraph.eigenvectors_card}
	\lean{MIPStarRE.LDT.ExpansionHypercubeGraph.eigenvectors}
	\leanok
	\uses{def:adjacency-laplacian, lem:character-average-vector, lem:character-average-scalar}
	For each $\alpha \in \F_q^m$, define
	\[
		\ket{\varphi_\alpha}
		=
		\frac{1}{M^{1/2}} \cdot \sum_{u \in \F_q^m} \omega^{\operatorname{tr}(u\cdot \alpha)} \ket{u}.
	\]
	Then the following two statements hold.
	\begin{enumerate}
		\item The $\ket{\varphi_\alpha}$'s form an orthonormal basis of $\C^V$.\label{item:orthonormal}
		\item For each $\alpha \in \F_q^m$, $\ket{\varphi_\alpha}$ is an eigenvector for $K$ with eigenvalue $\frac{1}{M} \cdot \frac{m-|\alpha|}{m}$, where $|\alpha|$ is the number of nonzero coordinates of $\alpha$.\label{item:eigenvector}
	\end{enumerate}
\end{lemma}
\begin{proof}
	\leanok
	\uses{lem:character-average-vector, lem:character-average-scalar, def:hypercube-graph}
	First we prove \ref{item:orthonormal}. For $\alpha,\beta \in \F_q^m$,
	\[
		\braket{\varphi_\alpha}{\varphi_\beta}
		=
		\frac{1}{M} \sum_{u \in \F_q^m} \omega^{\operatorname{tr}(u\cdot(\beta-\alpha))}
		=
		\begin{cases}
			1 & \text{if } \alpha=\beta,    \\
			0 & \text{if } \alpha\ne \beta,
		\end{cases}
	\]
	by Lemma~\ref{lem:character-average-vector}. Thus the $\ket{\varphi_\alpha}$ form an orthonormal basis of $\C^V$.

	Next we prove \ref{item:eigenvector}. For $\alpha \in \F_q^m$,
	\begin{equation}
		\begin{aligned}
			K \cdot \ket{\varphi_\alpha}
			 & =
			\left(\E_{(u,v)\sim C} \ket{u}\bra{v}\right) \cdot \left(\frac{1}{M^{1/2}} \cdot \sum_{u \in \F_q^m} \omega^{\operatorname{tr}(u\cdot \alpha)} \ket{u}\right) \\
			 & =
			\frac{1}{M^{1/2}} \cdot \E_{(u,v)\sim C} \omega^{\operatorname{tr}(v\cdot \alpha)} \ket{u}.
		\end{aligned}
		\label{eq:eigenvector-calculation}
	\end{equation}
	By the edge-sampling rule from Definition~\ref{def:hypercube-graph}, we may write $v=u+x e_i$, where $i$ is uniformly random in $\{1,\dots,m\}$ and $x$ is uniformly random in $\F_q$. Therefore
	\[
		\eqref{eq:eigenvector-calculation}
		=
		\frac{1}{M^{1/2}} \cdot \E_{u,i,x} \omega^{\operatorname{tr}((u+x e_i)\cdot \alpha)} \ket{u}
		=
		\left(\E_{i,x}\omega^{\operatorname{tr}(x\alpha_i)}\right) \cdot \frac{1}{M^{1/2}} \cdot \E_u \omega^{\operatorname{tr}(u\cdot \alpha)} \ket{u}
		=
		\frac{1}{M}\left(\E_{i,x}\omega^{\operatorname{tr}(x\alpha_i)}\right) \cdot \ket{\varphi_\alpha}.
	\]
	Hence $\ket{\varphi_\alpha}$ is an eigenvector of $K$ with eigenvalue
	\[
		\frac{1}{M} \cdot \E_i \left[\E_x \omega^{\operatorname{tr}(x\alpha_i)}\right]
		=
		\frac{1}{M} \cdot \E_i \left[\bone[\alpha_i=0]\right]
		=
		\frac{1}{M} \cdot \frac{m-|\alpha|}{m},
	\]
	by Lemma~\ref{lem:character-average-scalar}.
\end{proof}

\begin{lemma}[Spectral gap of the Laplacian]\label{cor:laplacian-spectral-gap}
	\lean{MIPStarRE.LDT.ExpansionHypercubeGraph.laplacianSpectralGapOrdered}
	\leanok
	\uses{prop:eigenvectors}
	If $\lambda_1 \le \lambda_2 \le \cdots \le \lambda_{q^m}$ are the eigenvalues of $L$, then
	\[
		\lambda_1=0,
		\qquad
		\lambda_2 = \frac{1}{m q^m}.
	\]
\end{lemma}
\begin{proof}
	\leanok
	\lean{MIPStarRE.LDT.ExpansionHypercubeGraph.laplacianSpectralGap}
	\lean{MIPStarRE.LDT.ExpansionHypercubeGraph.laplacianEigenvalue_eq}
	\lean{MIPStarRE.LDT.ExpansionHypercubeGraph.hypercubeSpectralGap_le_laplacianEigenvalue}
	\lean{MIPStarRE.LDT.ExpansionHypercubeGraph.laplacianEigenvalue_of_weight_one}
	\lean{MIPStarRE.LDT.ExpansionHypercubeGraph.laplacianSpectralGapOrdered_of_list_bounds}
	\lean{MIPStarRE.LDT.ExpansionHypercubeGraph.laplacianSpectralGapOrdered_of_fourier_eigenvalue_order}
	\uses{prop:eigenvectors}
	Lemma~\ref{prop:eigenvectors} identifies the two largest eigenvalues of $K$, hence the two smallest eigenvalues of $L=(1/q^m)I-K$.
\end{proof}

\section{Local and global variance}

\begin{definition}[Local and global variance]\label{def:local-and-variance}
	\lean{MIPStarRE.LDT.ExpansionHypercubeGraph.localAndVariance}
	\leanok
	Let $\ket{\psi} \in \mathcal H_{\mathrm A} \ot \mathcal H_{\mathrm B}$ and let $0 \le A^u \le I$ be an operator for each $u \in \F_q^m$. The local variance is
	\[
		\mathbf{Var}_{\mathrm{local}}(A,\psi)
		= \frac12 \E_{(u,v)\sim C} \bra{\psi} (A^u-A^v)^2 \ot I \ket{\psi},
	\]
	and the global variance is
	\[
		\mathbf{Var}_{\mathrm{global}}(A,\psi)
		= \frac12 \E_{u,v\sim \F_q^m} \bra{\psi} (A^u-A^v)^2 \ot I \ket{\psi}.
	\]
\end{definition}

\begin{lemma}[Local variance as a Laplacian form]\label{lem:local-rewrite}
	\lean{MIPStarRE.LDT.ExpansionHypercubeGraph.localRewrite}
	\leanok
	\uses{def:adjacency-laplacian, def:local-and-variance}
	Define
	\[
		A_{\mathrm{comb}} = \sum_{u \in \F_q^m} \ket{u} \ot A^u \ot I.
	\]
	Then
	\[
		\Tr(A_{\mathrm{comb}}^\dagger (L \ot \ket{\psi}\bra{\psi}) A_{\mathrm{comb}})
		= \mathbf{Var}_{\mathrm{local}}(A,\psi).
	\]
\end{lemma}
\begin{proof}
	\leanok
	\uses{prop:laplacian-rewrite}
	For $u,v \in \F_q^m$,
	\begin{equation}
		\begin{aligned}
			((\bra{u}-\bra{v}) \ot \bra{\psi}) \cdot A_{\mathrm{comb}}
			 & =
			((\bra{u}-\bra{v}) \ot \bra{\psi}) \cdot \sum_{w \in \F_q^m} \ket{w} \ot A^w \ot I \\
			 & =
			\bra{\psi} \cdot ((A^u-A^v)\ot I).
		\end{aligned}
		\label{eq:edge-difference-action}
	\end{equation}
	Therefore
	\[
		A_{\mathrm{comb}}^\dagger (L \ot \ket{\psi}\bra{\psi}) A_{\mathrm{comb}}
		=
		\frac12 \E_{(u,v)\sim C} A_{\mathrm{comb}}^\dagger \bigl((\ket{u}-\ket{v})(\bra{u}-\bra{v}) \ot \ket{\psi}\bra{\psi}\bigr) A_{\mathrm{comb}}
	\]
	by Lemma~\ref{prop:laplacian-rewrite}, and hence
	\[
		A_{\mathrm{comb}}^\dagger (L \ot \ket{\psi}\bra{\psi}) A_{\mathrm{comb}}
		=
		\frac12 \E_{(u,v)\sim C} ((A^u-A^v)\ot I)\ket{\psi}\bra{\psi}((A^u-A^v)\ot I)
	\]
	by \eqref{eq:edge-difference-action}. Taking the trace gives
	\[
		\Tr(A_{\mathrm{comb}}^\dagger (L \ot \ket{\psi}\bra{\psi}) A_{\mathrm{comb}})
		=
		\frac12 \E_{(u,v)\sim C} \bra{\psi} (A^u-A^v)^2 \ot I \ket{\psi}
		=
		\mathbf{Var}_{\mathrm{local}}(A,\psi).
	\]
\end{proof}

\begin{lemma}[Global variance as the orthogonal Fourier mass]\label{lem:global-rewrite}
	\lean{MIPStarRE.LDT.ExpansionHypercubeGraph.globalRewrite}
	\leanok
	\uses{prop:eigenvectors, def:local-and-variance}
	Write
	\[
		A_{\mathrm{comb}} = \ket{\varphi_0} \ot A_0 + \ket{\varphi_\perp} \ot A_\perp,
	\]
	where $\ket{\varphi_\perp}$ is orthogonal to $\ket{\varphi_0}$. Then
	\[
		\frac{1}{q^m}\Tr((\bra{\varphi_\perp}\ot A_\perp)(I\ot \ket{\psi}\bra{\psi})(\ket{\varphi_\perp}\ot A_\perp))
		= \mathbf{Var}_{\mathrm{global}}(A,\psi).
	\]
\end{lemma}
\begin{proof}
	\leanok
	\uses{prop:eigenvectors}
	We first compute
	\[
		A_0
		=
		(\bra{\varphi_0}\ot I)\cdot A_{\mathrm{comb}}
		=
		\left(\frac{1}{M^{1/2}} \cdot \sum_{u \in \F_q^m} \bra{u}\right) \ot I \cdot \left(\sum_{u \in \F_q^m} \ket{u} \ot A^u \ot I\right)
		=
		\frac{1}{M^{1/2}} \cdot \sum_{u \in \F_q^m} A^u \ot I.
	\]
	Writing $A_{\mathrm{avg}}=\E_u A^u$, it follows that
	\[
		\ket{\varphi_0}\ot A_0
		=
		\sum_{u \in \F_q^m} \ket{u} \ot A_{\mathrm{avg}} \ot I,
		\qquad
		\ket{\varphi_\perp}\ot A_\perp
		=
		\sum_{u \in \F_q^m} \ket{u} \ot (A^u-A_{\mathrm{avg}}) \ot I.
	\]
	Hence
	\begin{equation}
		\frac{1}{M} \cdot \Tr((\bra{\varphi_\perp}\ot A_\perp)(I\ot \ket{\psi}\bra{\psi})(\ket{\varphi_\perp}\ot A_\perp))
		=
		\E_{u \sim \F_q^m} \bra{\psi} (A^u-A_{\mathrm{avg}})^2 \ot I \ket{\psi}.
		\label{eq:just-took-trace}
	\end{equation}
	Moreover,
	\[
		\E_{u \sim \F_q^m}(A^u-A_{\mathrm{avg}})^2
		=
		\E_{u \sim \F_q^m}\bigl((A^u)^2-(A_{\mathrm{avg}})^2\bigr)
		=
		\frac12 \E_{u,v \sim \F_q^m} (A^u-A^v)^2.
	\]
	Substituting this identity into \eqref{eq:just-took-trace} yields
	\[
		\frac{1}{M} \cdot \Tr((\bra{\varphi_\perp}\ot A_\perp)(I\ot \ket{\psi}\bra{\psi})(\ket{\varphi_\perp}\ot A_\perp))
		=
		\frac12 \E_{u,v \sim \F_q^m} \bra{\psi} (A^u-A^v)^2 \ot I \ket{\psi}
		=
		\mathbf{Var}_{\mathrm{global}}(A,\psi).
	\]
\end{proof}

\begin{lemma}[Local-to-global inequality]\label{lem:local-to-global}
	\lean{MIPStarRE.LDT.ExpansionHypercubeGraph.localToGlobalBipartite}
	\leanok
	\uses{cor:laplacian-spectral-gap, lem:local-rewrite, lem:global-rewrite}
	For every family $\{A^u\}$,
	\[
		\mathbf{Var}_{\mathrm{global}}(A,\psi) \le m\, \mathbf{Var}_{\mathrm{local}}(A,\psi).
	\]
\end{lemma}
\begin{proof}
	\leanok
	\lean{MIPStarRE.LDT.ExpansionHypercubeGraph.localToGlobal}
	\uses{lem:local-rewrite, cor:laplacian-spectral-gap, lem:global-rewrite}
	\begin{equation}
		\begin{aligned}
			A_{\mathrm{comb}}^\dagger (L \ot \ket{\psi}\bra{\psi}) A_{\mathrm{comb}}
			 & =
			(\bra{\varphi_0}\ot A_0+\bra{\varphi_\perp}\ot A_\perp)(L \ot \ket{\psi}\bra{\psi})(\ket{\varphi_0}\ot A_0+\ket{\varphi_\perp}\ot A_\perp) \\
			 & =
			(\bra{\varphi_\perp}\ot A_\perp)(L \ot \ket{\psi}\bra{\psi})(\ket{\varphi_\perp}\ot A_\perp)                                               \\
			 & =
			\bra{\varphi_\perp}L\ket{\varphi_\perp} \cdot A_\perp \ket{\psi}\bra{\psi} A_\perp.
		\end{aligned}
		\label{eq:used-0-eigenvector}
	\end{equation}
	The second step uses that $\ket{\varphi_0}$ is a $0$-eigenvector of $L$. Since $\ket{\varphi_\perp}$ is orthogonal to $\ket{\varphi_0}$, Lemma~\ref{cor:laplacian-spectral-gap} gives
	\[
		\bra{\varphi_\perp}L\ket{\varphi_\perp}
		\ge
		\frac{1}{mM} \cdot \braket{\varphi_\perp}{\varphi_\perp}.
	\]
	Therefore
	\[
		\mathbf{Var}_{\mathrm{local}}(A,\psi)
		=
		\Tr(A_{\mathrm{comb}}^\dagger (L \ot \ket{\psi}\bra{\psi}) A_{\mathrm{comb}})
		=
		\bra{\varphi_\perp}L\ket{\varphi_\perp} \cdot \Tr(A_\perp \ket{\psi}\bra{\psi} A_\perp)
	\]
	by Lemma~\ref{lem:local-rewrite} and \eqref{eq:used-0-eigenvector}, so
	\[
		\mathbf{Var}_{\mathrm{local}}(A,\psi)
		\ge
		\frac{1}{mM} \cdot \Tr((\bra{\varphi_\perp}\ot A_\perp)(I\ot \ket{\psi}\bra{\psi})(\ket{\varphi_\perp}\ot A_\perp))
		=
		\frac{1}{m}\mathbf{Var}_{\mathrm{global}}(A,\psi),
	\]
	where the last step is Lemma~\ref{lem:global-rewrite}. Rearranging gives the claim.
\end{proof}


%%FILE blueprint/src/chapter/ch06_variance.tex
\chapter{Global variance of the points measurements}
\label{chap:variance}
\label{sec:variance}

Throughout this chapter, $(\psi,A,B,L)$ is an $(\eps,\delta,\gamma)$-good symmetric strategy for the $(m,q,d)$ low individual degree test.


\begin{lemma}[Replacing a line answer by a fixed polynomial]\label{lem:generalize-b}
	\lean{MIPStarRE.LDT.GlobalVariance.generalizeBFromSchwartzZippel}
	\leanok
	\uses{def:polymeasurements, lem:schwartz-zippel-individual, lem:axis-line-polynomial-agreement-mdq}
	Let $G \in \polysub{m}{q}{d}$. On the axis-parallel line-test distribution,
	\[
		B^{\ell}_{[f(u)=g(u)]} \ot (G_g)^{1/2}
		\approx_{md/q}
		B^{\ell}_{g|_{\ell}} \ot (G_g)^{1/2}.
	\]
\end{lemma}
\begin{proof}
	\leanok
	\uses{lem:schwartz-zippel-individual, lem:axis-line-polynomial-agreement-mdq}
	Expand the $\approx_{md/q}$ norm on the weighted vector $(I\ot (G_g)^{1/2})\ket{\psi}$. Because $B^\ell$ is a projective measurement, only those line polynomials $f$ with $f\ne g|_\ell$ and $f(u)=g(u)$ contribute. Lemma~\ref{lem:axis-line-polynomial-agreement-mdq} is the univariate Schwartz--Zippel wrapper used in Lean for this line-restriction event; it supplies the paper's $md/q$ loss, and the remaining operator sum is bounded by the fact that $G$ is a submeasurement.
\end{proof}

\begin{remark}[Conditional variants for the collision-residual estimate]\label{rem:generalize-b-collision-residual-conditional}
	The paper proof of Lemma~\ref{lem:generalize-b} bounds the collision event
	directly by the Schwartz--Zippel estimate.  The following conditional variants
	instead assume this collision estimate, or one of the reindexing identities
	that implies it, as an explicit hypothesis.  These conditional statements
	make explicit the intermediate estimates still needed for this route; they
	are not additional hypotheses in the paper statement of
	Lemma~\ref{lem:generalize-b}.
\end{remark}

\begin{lemma}[Local variance of the points measurements]\label{lem:local-variance-of-points}
	\lean{MIPStarRE.LDT.GlobalVariance.localVarianceTransportChainError, MIPStarRE.LDT.GlobalVariance.localVarianceTransportChainError_le_localVarianceOfPointsError, MIPStarRE.LDT.GlobalVariance.localVarianceTransportChainBound, MIPStarRE.LDT.GlobalVariance.localVarianceOfPointsFromTransportChainBound, MIPStarRE.LDT.GlobalVariance.globalVarianceOfPointsFromTransportChainBound}
	\leanok
	\uses{lem:good-strategy-characterization, prop:simeq-to-approx, lem:generalize-b, prop:triangle-inequality-for-approx_delta}
	Let $G \in \polysub{m}{q}{d}$. On the hypercube edge distribution $(u,v) \sim C$,
	\begin{equation}
		A^u_{g(u)} \ot (G_g)^{1/2}
		\approx_{24\cdot\left(\eps+\delta+\frac{md}{q}\right)}
		A^v_{g(v)} \ot (G_g)^{1/2}.
		\label{eq:local-variance-of-points-equation}
	\end{equation}
\end{lemma}
\begin{proof}
	\lean{MIPStarRE.LDT.GlobalVariance.avgOver_axisParallelTestSample_update_eq_rerandomizeCoord}
	\leanok
	\uses{prop:simeq-to-approx, prop:cab-approx-delta, lem:generalize-b, lem:good-strategy-characterization, prop:triangle-inequality-for-approx_delta}
	Fix $g$ and keep the weight $(G_g)^{1/2}$ on the second tensor factor throughout. On the vector $(I\ot (G_g)^{1/2})\ket{\psi}$ one has the six-step chain
	\[
		A^u_{g(u)} \ot (G_g)^{1/2}
		\approx_{2\delta} I\ot (G_g)^{1/2} A^u_{g(u)}
		\approx_{2\eps} B^\ell_{[f(u)=g(u)]} \ot (G_g)^{1/2}
		\approx_{md/q} B^\ell_{g|_\ell} \ot (G_g)^{1/2}
		\approx_{md/q} B^\ell_{[f(v)=g(v)]} \ot (G_g)^{1/2}
		\approx_{2\eps} I\ot (G_g)^{1/2} A^v_{g(v)}
		\approx_{2\delta} A^v_{g(v)} \ot (G_g)^{1/2}.
	\]
	The first, second, fifth, and sixth comparisons use Lemma~\ref{prop:simeq-to-approx} together with Theorem~\ref{prop:cab-approx-delta} and the good-strategy hypotheses, and the middle pair is Lemma~\ref{lem:generalize-b}. Lemma~\ref{prop:triangle-inequality-for-approx_delta} sums these six weighted transports.
\end{proof}

\begin{lemma}[Equivalent local-variance bound]\label{lem:equivalent-local-variance}
	\lean{MIPStarRE.LDT.GlobalVariance.localVarianceDeviation_sum_le_localVarianceOfPointsError}
	\leanok
	\uses{lem:local-variance-of-points}
	The local-variance bound for the point measurements is
	\begin{equation}
		\sum_{g \in \polyfunc{m}{q}{d}} \E_{(u,v)\sim C} \bra{\psi} (A^u_{g(u)}-A^v_{g(v)})^2 \ot G_g \ket{\psi}
		\le
		24\left(\eps+\delta+\frac{md}{q}\right).
		\label{eq:equivalent-local-variance}
	\end{equation}
\end{lemma}
\begin{proof}
	\lean{MIPStarRE.LDT.GlobalVariance.localVarianceDeviation_sum_le_localVarianceOfPointsError}
	\leanok
	\uses{lem:local-variance-of-points}
	This is the squared-norm expansion of Lemma~\ref{lem:local-variance-of-points}
	after summing over the polynomial outcome~$g$.  The Lean proof expands the
	weighted operator difference, telescopes the six transports in the local
	variance proof, sums the resulting bounds over~$g$, and absorbs the
	transport-chain error into
	$24\left(\eps+\delta+\frac{md}{q}\right)$.
\end{proof}

\begin{remark}[Polynomial-sum formalization for \ref{eq:equivalent-local-variance}]\label{rem:equivalent-local-variance-formalization-inputs}
	The three referenced declarations and Lemma~\ref{lem:equivalent-local-variance} together provide the full polynomial-sum (cardinality-free) chain proving the equation.  The chain assembly lemma \texttt{localVarianceDeviation\_sum\_le\_localVarianceOfPointsError} telescopes the six operator differences via \texttt{ev\_sum\_conjTranspose\_mul\_sum\_le}, sums over~$g$, and applies the six individual \texttt{\_sum} bounds (including the reverse~\texttt{generalizeB} sum) to avoid the per-polynomial cardinality blow-up.
\end{remark}

\begin{lemma}[Global variance of the points measurements]\label{lem:global-variance-of-points}
	\lean{MIPStarRE.LDT.GlobalVariance.globalVarianceOfPoints}
	\leanok
	\uses{def:local-and-variance, lem:local-to-global, lem:equivalent-local-variance}
	Let $G \in \polysub{m}{q}{d}$. On the uniform distribution over independent $u,v \in \F_q^m$,
	\begin{equation}
		A^u_{g(u)} \ot (G_g)^{1/2}
		\approx_{24m\left(\eps+\delta+\frac{md}{q}\right)}
		A^v_{g(v)} \ot (G_g)^{1/2}.
		\label{eq:global-variance-of-points-equation}
	\end{equation}
\end{lemma}
\begin{proof}
	\leanok
	\lean{MIPStarRE.LDT.GlobalVariance.globalVarianceDeviation_sum_le_m_mul_localVarianceDeviation_sum, MIPStarRE.LDT.GlobalVariance.globalVarianceDeviation_sum_le_of_localVarianceDeviation_sum_le}
	\uses{def:local-and-variance, lem:local-to-global, lem:local-variance-of-points, lem:equivalent-local-variance}
	We want to bound
	\begin{equation}
		\E_{u,v \sim \F_q^m} \sum_{g \in \polyfunc{m}{q}{d}} \lVert (A^u_{g(u)}-A^v_{g(v)}) \ot (G_g)^{1/2} \ket{\psi} \rVert^2
		=
		\E_{u,v \sim \F_q^m} \sum_{g \in \polyfunc{m}{q}{d}} \bra{\psi} (A^u_{g(u)}-A^v_{g(v)})^2 \ot G_g \ket{\psi}.
		\label{eq:global-variance-target}
	\end{equation}
	For each $g \in \polyfunc{m}{q}{d}$, define
	\[
		A(g)^u := A^u_{g(u)},
		\qquad
		\ket{\psi_g} := (I \ot (G_g)^{1/2})\ket{\psi}.
	\]
	Then
	\[
		\eqref{eq:global-variance-target}
		=
		\sum_{g \in \polyfunc{m}{q}{d}} \E_{u,v \sim \F_q^m} \bra{\psi_g} (A(g)^u-A(g)^v)^2 \ot I \ket{\psi_g}
		=
		\sum_{g \in \polyfunc{m}{q}{d}} 2\,\mathbf{Var}_{\mathrm{global}}(A(g),\psi_g).
	\]
	Applying Lemma~\ref{lem:local-to-global} gives
	\[
		\eqref{eq:global-variance-target}
		\le
		\sum_{g \in \polyfunc{m}{q}{d}} 2m\,\mathbf{Var}_{\mathrm{local}}(A(g),\psi_g)
		=
		m \cdot \sum_{g \in \polyfunc{m}{q}{d}} \E_{(u,v)\sim C} \bra{\psi} (A^u_{g(u)}-A^v_{g(v)})^2 \ot G_g \ket{\psi}.
	\]
	The last quantity is at most $m \cdot 24\left(\eps+\delta+\frac{md}{q}\right)$ by Lemma~\ref{lem:equivalent-local-variance}, which is exactly \ref{eq:global-variance-of-points-equation}.
\end{proof}


%%FILE blueprint/src/chapter/ch07_self_improvement.tex
\chapter{Self-improvement}
\label{chap:self-improvement}

Throughout this chapter, $(\psi,A,B,L)$ is an $(\eps,\delta,\gamma)$-good symmetric strategy for the $(m,q,d)$ low individual degree test.

\section{Non-projective Output}
\label{sec:self-improvement}

A large part of this chapter is devoted to proving
Lemma~\ref{lem:self-improvement-helper}, which is a slightly weaker form of the
projective self-improvement theorem proved in Section~\ref{sec:self-improvement-projective}.
The key difference is that the output family~$H$ is only required to be a
submeasurement, rather than a projective submeasurement. Once this non-projective
statement is available, Theorem~\ref{thm:orthonormalization} upgrades~$H$ to a
projective family without losing control of the four quantitative conclusions.

Compared with the projective theorem, the helper lemma records strong
self-consistency using the explicit inequality from
Definition~\ref{def:strong-self-consistency}, keeps the boundedness conclusion in
terms of an auxiliary operator~$Z$, and produces a much smaller error parameter.

\begin{lemma}[Self-improvement with non-projective output]\label{lem:self-improvement-helper}
  \lean{MIPStarRE.LDT.SelfImprovement.selfImprovementHelper}
  \lean{MIPStarRE.LDT.SelfImprovement.self_improvement_helper_with_slackness_of_sdp_statement_with_slackness}
  \leanok
  \uses{def:polymeasurements, def:simeq, def:strong-self-consistency}
  Let $G \in \polymeas{m}{q}{d}$ be a measurement with the following property:
  \begin{enumerate}
    \item (Consistency with~$A$): \label{item:self-improvement-G-consistency} On average over $u \sim \F_q^{m}$,
    \[
      A^{u}_a \ot I \simeq_{\nu} I \ot G_{[g(u)=a]}.
    \]
  \end{enumerate}
  Let
  \[
    \zeta = 100m\cdot \Big(\eps^{1/2} + \delta^{1/2} + (d/q)^{1/2}\Big).
  \]
  Then there exists $H \in \polysub{m}{q}{d}$ with the following properties:
  \begin{enumerate}
    \item (Completeness): \label{item:self-improvement-completeness} If $H = \sum_h H_h$, then
    \[
      \bra{\psi} H \ot I \ket{\psi} \geq (1-\nu)-\zeta.
    \]
    \item (Consistency with~$A$): \label{item:self-improvement-A-consistency} On average over $u \sim \F_q^m$,
    \[
      A^u_a \ot I \simeq_{\zeta} I \ot H_{[h(u) = a]}.
    \]
    \item (Strong self-consistency): \label{item:self-improvement-self}
    \[
      \sum_{h} \bra{\psi} H_h \ot H_h \ket{\psi} \geq \bra{\psi} H \ot I \ket{\psi} - \zeta.
    \]
    \item (Boundedness): \label{item:self-improvement-boundedness} There exists a positive-semidefinite matrix~$Z$ such that
    \[
      \bra{\psi} Z \ot I \ket{\psi} -\E_{u} \sum_a \bra{\psi}  A^{u}_{a} \ot H_{[h(u)=a]} \ket{\psi} \leq \zeta
    \]
    and for each $h \in \polyfunc{m}{q}{d}$,
    \[
      Z \geq \left(\E_{u} A^{u}_{h(u)}\right).
    \]
  \end{enumerate}
\end{lemma}
\begin{proof}
  \uses{lem:sdp, lem:add-in-u, lem:global-variance-of-points, prop:simeq-to-approx, prop:simeq-for-measurements}
  \leanok
  Let $T=\{T_g\}$ and $Z$ be the optimal solutions to the SDPs~\eqref{eq:primal-objective} and~\eqref{eq:dual-objective} given by Lemma~\ref{lem:sdp}. Then $T$ is a measurement, and
  \begin{align}
    \forall g \in \polyfunc{m}{q}{d}: \qquad
    Z &\geq \left(\E_{u} A^{u}_{g(u)}\right), \label{eq:Z-greater-than-A}\\
    T_g \cdot Z &= T_g \cdot \left(\E_{u} A^{u}_{g(u)}\right). \label{eq:swap-Z-for-A}
  \end{align}
  For each $u \in \F_q^m$, define $H^u=\{H_h^u\}_{h \in \polyfunc{m}{q}{d}}$ by
  \[
    H_h^u := A^u_{h(u)} \cdot T_h \cdot A^u_{h(u)},
  \]
  and let
  \[
    H_h := \E_{u} H_h^u.
  \]
  The pointwise family $H^u$ is a submeasurement because
  \[
    \sum_h H_h^u
    = \sum_a A_a^u \cdot \Big(\sum_{h:h(u)=a} T_h\Big)\cdot A_a^u
    \le \sum_a (A_a^u)^2
    = I,
  \]
  where the inequality uses that~$T$ is a measurement and the last identity uses
  that~$A$ is projective. Averaging over~$u$ gives
  \[
    \sum_h H_h
    = \E_{u} \sum_h H_h^u
    \le I,
  \]
  so $H$ is also a submeasurement and therefore lies in~$\polysub{m}{q}{d}$.

  Set
  \[
    \zeta_{\mathrm{variance}} = 24m\cdot\Big(\eps + \delta + \frac{md}{q}\Big).
  \]
  Lemma~\ref{lem:add-in-u} below is the technical transfer from the averaged family~$H$ back to the sandwiched operators $A^u_{h(u)}T_hA^u_{h(u)}$. We use it in the completeness, consistency, strong self-consistency, and boundedness estimates that follow.

  We now verify the four conclusions of Lemma~\ref{lem:self-improvement-helper}.

  \emph{Proof of \ref{item:self-improvement-completeness}.}
  The completeness of~$H$ is
  \begin{align}
    \sum_h \bra{\psi} H_h \ot I \ket{\psi}
    &= \E_{u} \sum_h \bra{\psi} H^{u}_h \ot I \ket{\psi} \notag\\
    &= \E_{u} \sum_h \bra{\psi} \bigl(A^{u}_{h(u)} \cdot T_h \cdot A^{u}_{h(u)}\bigr) \ot I \ket{\psi}.
  \end{align}
  \lean{MIPStarRE.LDT.SelfImprovement.helper_mass_eq_avg_pointwise_sandwich_sum}
  Grouping the polynomials according to their value at~$u$ gives
  \begin{equation}\label{eq:group-by-point-value}
    \E_{u} \sum_h \bra{\psi} \bigl(A^{u}_{h(u)} \cdot T_h \cdot A^{u}_{h(u)}\bigr) \ot I \ket{\psi}
    = \E_{u} \sum_a \bra{\psi} \bigl(A^{u}_{a} \cdot T_{[h(u) = a]} \cdot A^{u}_{a}\bigr) \ot I \ket{\psi}.
  \end{equation}
  \lean{MIPStarRE.LDT.SelfImprovement.helperFiberOperator}
  \lean{MIPStarRE.LDT.SelfImprovement.helperFiberOperator_nonneg}
  \lean{MIPStarRE.LDT.SelfImprovement.helperFiberOperator_le_one}
  \lean{MIPStarRE.LDT.SelfImprovement.helperBracketedCompletenessQuantity}
  \lean{MIPStarRE.LDT.SelfImprovement.helper_pointwise_sandwich_sum_eq_bracketed}
  \lean{MIPStarRE.LDT.SelfImprovement.helper_mass_eq_avg_pointwise_bracketed_sum}
  \lean{MIPStarRE.LDT.SelfImprovement.helperBracketedCompletenessQuantity_eq_mass}
  We first move the leftmost copy of~$A^u_a$ across the bipartition:
  \begin{equation}\label{eq:first-point-projector-transport}
    \eqref{eq:group-by-point-value}
    \approx_{2\sqrt{\delta}}
    \E_{u} \sum_a \bra{\psi} \bigl(T_{[h(u)=a]} \cdot A^{u}_{a}\bigr) \ot A^{u}_{a} \ket{\psi}.
  \end{equation}
  The difference is bounded by Cauchy--Schwarz as
  in~\eqref{eq:move-first-point-projector-cauchy-schwarz}; the first factor is
  at most $\sqrt{2\delta}$ by self-consistency of~$A$, and the second is at
  most~$1$ because $T_{[h(u)=a]} \le I$.
  \lean{MIPStarRE.LDT.SelfImprovement.helper_first_move_second_factor_operator_le_one}
  \lean{MIPStarRE.LDT.SelfImprovement.helper_first_move_second_factor_le_one}
  \lean{MIPStarRE.LDT.SelfImprovement.helper_first_move_abs_sub_bracketed_le_two_sqrt_delta}
  We next remove the remaining copy of~$A^u_a$ on Bob's side:
  \begin{equation}\label{eq:first-moved-completeness}
    \eqref{eq:first-point-projector-transport}
    \approx_{\sqrt{\delta}}
    \E_{u} \sum_a \bra{\psi} \bigl(T_{[h(u)=a]} \cdot A^{u}_{a}\bigr) \ot I \ket{\psi}.
  \end{equation}
  The first Cauchy--Schwarz factor is at most~$1$ because the fiber operators form a submeasurement after grouping~$T$ by the value of~$h(u)$.
  \lean{MIPStarRE.LDT.SelfImprovement.helperFiberOperator_sum_eq_total}
  \lean{MIPStarRE.LDT.SelfImprovement.helper_second_move_first_factor_operator_le_one}
  \lean{MIPStarRE.LDT.SelfImprovement.helper_second_move_first_factor_le_one}
  Here the second Cauchy--Schwarz factor is
  \[
    \E_{u} \sum_{a\neq b} \bra{\psi} A^u_a \ot A^u_b \ket{\psi},
  \]
  which is at most~$\delta$ because~$A$ is self-consistent and projective. Reindexing by~$h$ and using~\eqref{eq:swap-Z-for-A},
  \lean{MIPStarRE.LDT.SelfImprovement.helper_second_move_second_factor_pointwise_le_qBipartiteSSCDefect}
  \lean{MIPStarRE.LDT.SelfImprovement.helper_second_move_second_factor_le_delta}
  \lean{MIPStarRE.LDT.SelfImprovement.helper_linearized_completeness_quantity_eq_fiber_sum}
  \lean{MIPStarRE.LDT.SelfImprovement.helper_second_move_pointwise_abs_le_sqrt}
  \lean{MIPStarRE.LDT.SelfImprovement.helper_second_move_abs_sub_first_moved_le_sqrt_delta}
  \[
    \eqref{eq:first-moved-completeness}
    = \sum_h \bra{\psi} \bigl(T_h \cdot \E_{u} A^u_{h(u)}\bigr) \ot I \ket{\psi}
    = \sum_h \bra{\psi} (T_h \cdot Z) \ot I \ket{\psi}
    = \bra{\psi} Z \ot I \ket{\psi}.
  \]
  \lean{MIPStarRE.LDT.SelfImprovement.helperFirstMovedCompletenessQuantity}
  \lean{MIPStarRE.LDT.SelfImprovement.helperLinearizedCompletenessQuantity}
  \lean{MIPStarRE.LDT.SelfImprovement.helper_linearized_completeness_eq_dual_mass_of_complementary_slackness}
  \lean{MIPStarRE.LDT.SelfImprovement.helper_linearized_completeness_quantity_eq_dual_mass_of_complementary_slackness}
  \lean{MIPStarRE.LDT.SelfImprovement.helper_slackness_eq_of_helper_with_slackness}
  Thus
  \begin{equation}\label{eq:h-versus-z-lower-bound}
    \sum_h \bra{\psi} H_h \ot I \ket{\psi} \geq \bra{\psi} Z \ot I \ket{\psi} - 3\sqrt{\delta}.
  \end{equation}
  \lean{MIPStarRE.LDT.SelfImprovement.helper_hhat_vs_z_of_cauchy_schwarz_and_complementary_slackness}
  \lean{MIPStarRE.LDT.SelfImprovement.input_consistency_match_mass_lower_bound}
  \lean{MIPStarRE.LDT.SelfImprovement.input_match_mass_eq_sdp_overlap}
  \lean{MIPStarRE.LDT.SelfImprovement.sdp_overlap_le_dual_mass}
  \lean{MIPStarRE.LDT.SelfImprovement.input_consistency_dual_mass_lower_bound}
  \lean{MIPStarRE.LDT.SelfImprovement.sdp_complementary_slackness_sum_eq_dual_mass}
  \lean{MIPStarRE.LDT.SelfImprovement.helper_sdp_complementary_slackness_sum_eq_dual_mass}
  \lean{MIPStarRE.LDT.SelfImprovement.helper_hhat_vs_z_of_bracketed_cauchy_schwarz_and_complementary_slackness}
  \lean{MIPStarRE.LDT.SelfImprovement.helper_hhat_vs_z_of_self_consistency_and_complementary_slackness}
  \lean{MIPStarRE.LDT.SelfImprovement.helper_hhat_vs_z_of_self_consistency_and_helper_slackness}
  \lean{MIPStarRE.LDT.SelfImprovement.helper_completeness_of_dual_mass_lower_bound}
  \lean{MIPStarRE.LDT.SelfImprovement.helper_completeness_of_input_consistency}
  \lean{MIPStarRE.LDT.SelfImprovement.helper_completeness_of_bracketed_cauchy_schwarz_input_consistency}
  \lean{MIPStarRE.LDT.SelfImprovement.helper_completeness_of_cauchy_schwarz_input_consistency}
  \lean{MIPStarRE.LDT.SelfImprovement.helper_completeness_of_self_consistency_complementary_slackness_input_consistency}
  \lean{MIPStarRE.LDT.SelfImprovement.helper_completeness_of_self_consistency_helper_slackness_input_consistency}
  \lean{MIPStarRE.LDT.SelfImprovement.helper_completeness_error_le_selfImprovementHelperError}
  Since $G$ is a measurement,~\eqref{eq:Z-greater-than-A} and Lemma~\ref{prop:simeq-for-measurements} give
  \[
    \bra{\psi} Z \ot I \ket{\psi}
    \ge \E_{u} \sum_a \bra{\psi} A^u_a \ot G_{[g(u)=a]} \ket{\psi}
    \ge 1-\nu.
  \]
  Therefore the completeness is at least $1-\nu-3\sqrt{\delta}$, which is bounded below by $(1-\nu)-\zeta$.

  \emph{Proof of \ref{item:self-improvement-A-consistency}.}
  The inconsistency of~$H$ with~$A$ is
  \begin{equation}\label{eq:point-consistency-expansion}
    \E_{u} \sum_{a \neq b} \bra{\psi} A^{u}_a \ot H_{[h(u) = b]} \ket{\psi}
    = \E_{u} \sum_{a, h: h(u) \neq a} \bra{\psi} A^{u}_a \ot H_h \ket{\psi}.
  \end{equation}
  The theorem-side off-diagonal selection used for this `add-in-u` application is
  \lean{MIPStarRE.LDT.SelfImprovement.pointConsistencyAddInUSelection}, and the
  corresponding left/right scalar quantities are
  \lean{MIPStarRE.LDT.SelfImprovement.addInULeftQuantity_pointConsistencySelection_eq_off_diagonal_avg}
  and
  \lean{MIPStarRE.LDT.SelfImprovement.addInURightQuantity_pointConsistencySelection_eq_zero}.
  The selected scalar chain for this specialization is recorded by
  \lean{MIPStarRE.LDT.SelfImprovement.pointConsistencyAddInUCSChainQ0}
  \lean{MIPStarRE.LDT.SelfImprovement.pointConsistencyAddInUCSChainQ1}
  \lean{MIPStarRE.LDT.SelfImprovement.pointConsistencyAddInUCSChainQ2}
  \lean{MIPStarRE.LDT.SelfImprovement.pointConsistencyAddInUCSChainQ3}
  \lean{MIPStarRE.LDT.SelfImprovement.pointConsistencyAddInUCSChainQ4}
  \lean{MIPStarRE.LDT.SelfImprovement.pointConsistencyAddInUCSChainQ0_eq_leftQuantity_averagedSandwiched}
  \lean{MIPStarRE.LDT.SelfImprovement.pointConsistencyAddInUCSChainQ4_eq_rightQuantity}
  \lean{MIPStarRE.LDT.SelfImprovement.pointConsistencyAddInU_transfer_of_selected_chain_bounds}
  \lean{MIPStarRE.LDT.SelfImprovement.pointConsistencyAddInU_transfer_of_selected_chain_selfConsistency_globalVariance}
  \lean{MIPStarRE.LDT.SelfImprovement.helper_point_consistency_of_selected_chain_selfConsistency_globalVariance}.
  Apply Lemma~\ref{lem:add-in-u} with $\mathcal O=\F_q$, $M=A$, and $S_u=\{(a,h):h(u)\neq a\}$. Then
  \[
    \eqref{eq:point-consistency-expansion}
    \approx_{4 \sqrt{\zeta_{\mathrm{variance}}}}
    \E_{u} \sum_{a, h: h(u) \neq a} \bra{\psi} \bigl(A^{u}_{h(u)} \cdot A^{u}_a \cdot A^{u}_{h(u)}\bigr)\ot T_h \ket{\psi},
  \]
  and the right-hand side is~$0$ because~$A$ is projective. Hence
  \begin{equation}\label{eq:explicit-bound-for-A-consistency}
    \E_{u} \sum_{a \neq b} \bra{\psi} A^{u}_a \ot H_{[h(u) = b]} \ket{\psi} \leq 4 \sqrt{\zeta_{\mathrm{variance}}}.
  \end{equation}
  Since
  \lean{MIPStarRE.LDT.SelfImprovement.pointConsistencyAddInU_off_diagonal_avg_le_helper_error_of_transfer}
  \lean{MIPStarRE.LDT.SelfImprovement.helper_point_consistency_error_eq_off_diagonal_avg}
  \lean{MIPStarRE.LDT.SelfImprovement.helper_point_consistency_of_pointConsistencyAddInU_transfer}
  \lean{MIPStarRE.LDT.SelfImprovement.helper_point_consistency_of_selected_chain_selfConsistency_globalVariance}
  \lean{MIPStarRE.LDT.SelfImprovement.helper_point_consistency_error_le_selfImprovementHelperError}
  \[
    4\sqrt{\zeta_{\mathrm{variance}}}
    = 4 \sqrt{24 m \cdot\Big(\eps + \delta + \frac{md}{q}\Big)}
    \le 20 m \cdot \Big(\eps^{1/2} + \delta^{1/2} + (d/q)^{1/2}\Big)
    \le \zeta,
  \]
  this proves \ref{item:self-improvement-A-consistency}.

  \emph{Proof of \ref{item:self-improvement-self}.}
  From the definition of $H_h^u$ and the projectivity of~$A$,
  \begin{align}
    H^u_h &= A^u_{h(u)} \cdot T_h \cdot A^u_{h(u)} = A^u_{h(u)} \cdot H^u_h \cdot A^u_{h(u)}, \label{eq:h-sandwich}\\
    A^u_{h(u)} \cdot H^u_{h'} \cdot A^u_{h(u)} &= H^u_{h'} \cdot A^u_{h(u)} = \bigl(A^u_{h(u)} \cdot T_{h'} \cdot A^u_{h(u)}\bigr)\cdot \bone{h(u) = h'(u)}. \label{eq:h-blt}
  \end{align}
  \lean{MIPStarRE.LDT.SelfImprovement.pointConditioned_sandwichedPolynomialOutcome_outer_eq_ite}
  \lean{MIPStarRE.LDT.SelfImprovement.ev_opTensor_pointConditioned_sandwichedPolynomialOutcome_outer_eq_indicator}
  \lean{MIPStarRE.LDT.SelfImprovement.polynomial_off_diagonal_outer_sandwich_eq_indicator_avg}
  The strong self-consistency expression is
  \begin{equation}\label{eq:self-consistency-expansion}
    \sum_{h\in\polyfunc{m}{q}{d}} \bra{\psi} H_h \ot H_h \ket{\psi}
    = \E_{u} \sum_{h \in \polyfunc{m}{q}{d}} \bra{\psi} H_h^{u} \ot H_h \ket{\psi}.
  \end{equation}
  Apply Lemma~\ref{lem:add-in-u} with $\mathcal O=\polyfunc{m}{q}{d}$, $M=H$, and $S_u=\{(h,h):h\in\polyfunc{m}{q}{d}\}$. This gives
  \begin{equation}\label{eq:add-point-conditioning}
    \eqref{eq:self-consistency-expansion}
    \approx_{4 \sqrt{\zeta_{\mathrm{variance}}}}
    \E_{u} \sum_{h \in \polyfunc{m}{q}{d}} \bra{\psi} \bigl(A^{u}_{h(u)} \cdot H_h^{u} \cdot A^{u}_{h(u)}\bigr) \ot T_h \ket{\psi}.
  \end{equation}
  We next enlarge the sum from $h'=h$ to all pairs $h,h'$:
  \begin{equation}\label{eq:enlarge-polynomial-sum}
    \eqref{eq:add-point-conditioning}
    \approx_{2\sqrt{\zeta_{\mathrm{variance}}} + \frac{md}{q}}
    \E_{u} \sum_{h, h' \in \polyfunc{m}{q}{d}} \bra{\psi} \bigl(A^{u}_{h(u)} \cdot H^{u}_{h'} \cdot A^{u}_{h(u)}\bigr) \ot T_h \ket{\psi}.
  \end{equation}
  The added off-diagonal contribution is
  \begin{align}
    \eqref{eq:enlarge-polynomial-sum}-\eqref{eq:add-point-conditioning}
    &= \E_{u} \sum_{h \neq h'} \bra{\psi} \bigl(A^{u}_{h(u)} \cdot H^{u}_{h'} \cdot A^{u}_{h(u)}\bigr) \ot T_h \ket{\psi} \notag\\
    &= \E_{u} \sum_{h \neq h'} \bra{\psi} \bigl(A^{u}_{h(u)} \cdot T_{h'} \cdot A^{u}_{h(u)}\bigr) \ot T_h \ket{\psi} \cdot \bone{h(u) = h'(u)}, \label{eq:added-indicator}
  \end{align}
  where the second equality is~\eqref{eq:h-blt}. To estimate~\eqref{eq:added-indicator}, first replace the left copy of $A^u_{h(u)}$ by $A^v_{h(v)}$:
  \begin{equation}\label{eq:swapped-u-for-v}
    \eqref{eq:added-indicator}
    \approx_{\sqrt{\zeta_{\mathrm{variance}}}}
    \E_{u,v} \sum_{h \neq h'} \bra{\psi} \bigl(A^{v}_{h(v)} \cdot T_{h'} \cdot A^{u}_{h(u)}\bigr) \ot T_h \ket{\psi} \cdot \bone{h(u) = h'(u)}.
  \end{equation}
  The corresponding Cauchy--Schwarz bound is
  \begin{multline}
    \Big|\E_{u,v} \sum_{h \neq h'} \bra{\psi} \bigl((A^{u}_{h(u)} - A^{v}_{h(v)}) \cdot T_{h'} \cdot A^{u}_{h(u)}\bigr) \ot T_h \ket{\psi} \cdot \bone{h(u) = h'(u)}\Big|\\
    \leq
    \sqrt{\E_{u,v} \sum_{h \neq h'} \bra{\psi} \bigl((A^{u}_{h(u)} - A^{v}_{h(v)}) \cdot T_{h'} \cdot (A^{u}_{h(u)} - A^{v}_{h(v)})\bigr) \ot T_h\ket{\psi}}\\
    \cdot \sqrt{\E_{u,v} \sum_{h \neq h'} \bra{\psi} \bigl(A^{u}_{h(u)} \cdot T_{h'} \cdot A^{u}_{h(u)}\bigr) \ot T_h\ket{\psi} \cdot \bone{h(u) = h'(u)}}. \label{eq:swapped-u-for-cauchy-schwarz}
  \end{multline}
  The first factor is bounded by $\sqrt{\zeta_{\mathrm{variance}}}$ via Lemma~\ref{lem:global-variance-of-points}, and the second by~$1$ because $T$ and $H^u$ are submeasurements.

  Repeating the same argument for the remaining copy of~$A$ yields
  \begin{equation}\label{eq:replace-second-point-conditioner}
    \eqref{eq:swapped-u-for-v}
    \approx_{\sqrt{\zeta_{\mathrm{variance}}}}
    \E_{u,v} \sum_{h \neq h'} \bra{\psi} \bigl(A^{v}_{h(v)} \cdot T_{h'} \cdot A^{v}_{h(v)}\bigr) \ot T_h \ket{\psi} \cdot \bone{h(u) = h'(u)}.
  \end{equation}
  The same intermediate estimate gives
  \begin{equation}\label{eq:swapped-indicator-total-bound}
    \E_{u,v} \sum_{h \neq h'} \bra{\psi}  \bigl(A^{v}_{h(v)} \cdot T_{h'} \cdot A^{v}_{h(v)}\bigr)\ot T_h\ket{\psi}
    \leq 1,
  \end{equation}
  so Schwartz--Zippel bounds the indicator average by $md/q$ and proves~\eqref{eq:enlarge-polynomial-sum}.
  \lean{MIPStarRE.LDT.SelfImprovement.polynomial_collision_pointMeasurement_sandwichTensor_avg_le_mdq}
  \lean{MIPStarRE.LDT.SelfImprovement.polynomial_off_diagonal_swapped_indicator_sandwich_avg_le_mdq}
  \lean{MIPStarRE.LDT.SelfImprovement.helperOffDiagonalOuterSandwichQuantity}
  \lean{MIPStarRE.LDT.SelfImprovement.helperOffDiagonalIndicatorQuantity}
  \lean{MIPStarRE.LDT.SelfImprovement.helperOffDiagonalOneSidedSwappedIndicatorQuantity}
  \lean{MIPStarRE.LDT.SelfImprovement.helperOffDiagonalSwappedIndicatorQuantity}
  \lean{MIPStarRE.LDT.SelfImprovement.helperOffDiagonalOuterSandwichQuantity_eq_indicator}
  \lean{MIPStarRE.LDT.SelfImprovement.helperOffDiagonalSwappedIndicatorQuantity_le_mdq}
  \lean{MIPStarRE.LDT.SelfImprovement.helperOffDiagonalIndicatorQuantity_le_swapped_of_abs_transports}
  \lean{MIPStarRE.LDT.SelfImprovement.helperOffDiagonalOuterSandwichQuantity_le_two_sqrt_variance_add_mdq}
  \lean{MIPStarRE.LDT.SelfImprovement.helperOffDiagonalOuterSandwichQuantity_le_two_sqrt_variance_add_mdq_of_abs_transports}

  Using~\eqref{eq:h-blt}, the right-hand side of~\eqref{eq:enlarge-polynomial-sum} is
  \begin{equation}\label{eq:remove-point-projector}
    \E_{u} \sum_{h, h' \in \polyfunc{m}{q}{d}} \bra{\psi} \bigl(H^{u}_{h'} \cdot A^{u}_{h(u)}\bigr) \ot T_h \ket{\psi}.
  \end{equation}
  \lean{MIPStarRE.LDT.SelfImprovement.sandwichedPolynomialOutcome_mul_pointConditioned_eq_ite}
  \lean{MIPStarRE.LDT.SelfImprovement.pointConditioned_sandwichedPolynomialOutcome_outer_eq_right}
  \lean{MIPStarRE.LDT.SelfImprovement.helperFullOuterSandwichQuantity}
  \lean{MIPStarRE.LDT.SelfImprovement.helperDeleteAQuantity}
  \lean{MIPStarRE.LDT.SelfImprovement.helperFullOuterSandwichQuantity_eq_release_add_offDiagonalOuterSandwichQuantity}
  \lean{MIPStarRE.LDT.SelfImprovement.helperFullOuterSandwichQuantity_eq_deleteAQuantity}
  Replace the remaining copy of $A^u_{h(u)}$ by $A^v_{h(v)}$:
  \begin{equation}\label{eq:replace-final-point-conditioner}
    \eqref{eq:remove-point-projector}
    \approx_{\sqrt{\zeta_{\mathrm{variance}}}}
    \E_{u,v} \sum_{h, h' \in \polyfunc{m}{q}{d}} \bra{\psi} \bigl(H^{u}_{h'} \cdot A^{v}_{h(v)}\bigr) \ot T_h \ket{\psi}.
  \end{equation}
  Here the first Cauchy--Schwarz factor is at most~$1$ because $T$ and $H^u$ are submeasurements, and the second is again bounded by Lemma~\ref{lem:global-variance-of-points}. Finally move this last copy of~$A^v_{h(v)}$ to Bob's side:
  \begin{equation}\label{eq:move-over-v}
    \eqref{eq:replace-final-point-conditioner}
    \approx_{\sqrt{2\delta}}
    \E_{u,v} \sum_{h, h' \in \polyfunc{m}{q}{d}} \bra{\psi} H^{u}_{h'}  \ot \bigl(T_h \cdot A^{v}_{h(v)}\bigr) \ket{\psi}.
  \end{equation}
  \lean{MIPStarRE.LDT.SelfImprovement.helperDeleteAClonedQuantity}
  \lean{MIPStarRE.LDT.SelfImprovement.helperMoveOverVQuantity}
  \lean{MIPStarRE.LDT.SelfImprovement.helperDeleteAQuantity_le_moveOverV_of_abs_transports}
  \lean{MIPStarRE.LDT.SelfImprovement.helperMoveOverVQuantity_le_deleteA_of_abs_transports}
  Substituting~\eqref{eq:swap-Z-for-A}, then~\eqref{eq:Z-greater-than-A}, and finally the already proved bound~\eqref{eq:explicit-bound-for-A-consistency}, one obtains
  \[
    \eqref{eq:move-over-v}
    \ge \sum_h \bra{\psi} H_h \ot I \ket{\psi} - 4\sqrt{\zeta_{\mathrm{variance}}}.
  \]
  Combining~\eqref{eq:add-point-conditioning}, \eqref{eq:enlarge-polynomial-sum}, \eqref{eq:replace-final-point-conditioner}, and~\eqref{eq:move-over-v} gives
  \[
    \sum_h \bra{\psi} H_h \ot H_h \ket{\psi}
    \ge \sum_h \bra{\psi} H_h \ot I \ket{\psi}
      - 11 \sqrt{\zeta_{\mathrm{variance}}}
      - \sqrt{2\delta}
      - \frac{md}{q}.
  \]
  The Lean formalization records this step as named internal bounds, starting
  with the structure
  \lean{MIPStarRE.LDT.SelfImprovement.HelperStrongSelfConsistencyBounds}
  and the exact residual-side expansion
  \lean{MIPStarRE.LDT.SelfImprovement.helper_mass_sub_release_eq_polynomial_off_diagonal}
  \lean{MIPStarRE.LDT.SelfImprovement.helperOffDiagonalBareQuantity}
  \lean{MIPStarRE.LDT.SelfImprovement.helperOffDiagonalBareQuantity_le_one}
  \lean{MIPStarRE.LDT.SelfImprovement.helper_residualLowerBound_of_offDiagonal_bound}
  \lean{MIPStarRE.LDT.SelfImprovement.helper_residualLowerBound_of_paper_chain_bound}
  \lean{MIPStarRE.LDT.SelfImprovement.helperOffDiagonalBareQuantity_le_paper_chain_of_scalar_transports}
  \lean{MIPStarRE.LDT.SelfImprovement.helper_strong_self_consistency_bounds_of_selfConsistency_localVariance_offDiagonal}
  \lean{MIPStarRE.LDT.SelfImprovement.helper_strong_self_consistency_bounds_of_selfConsistency_localVariance_paperChain}
  \lean{MIPStarRE.LDT.SelfImprovement.helper_strong_self_consistency_bounds_of_selfConsistency_localVariance_scalarTransports}
  and then uses the proved lemma
  \lean{MIPStarRE.LDT.SelfImprovement.helper_strong_self_consistency_of_helper_conclusion}
  together with
  \lean{MIPStarRE.LDT.SelfImprovement.helper_strong_self_consistency_bounds_of_selfConsistency_localVariance}
  to recover the final helper-stage self-consistency conclusion.  The
  theorem-level derivation is therefore not moved into an additional hypothesis.
  The paper's arithmetic shows that the right-hand error is at most
  \lean{MIPStarRE.LDT.SelfImprovement.helper_strong_self_consistency_error_le_selfImprovementHelperError}
  \[
    57m\cdot \Big(\eps^{1/2} + \delta^{1/2} + (d/q)^{1/2}\Big)
    \le \zeta,
  \]
  which proves \ref{item:self-improvement-self}.

  \emph{Proof of \ref{item:self-improvement-boundedness}.}
  The helper-stage boundedness quantity is
  \[
    \bra{\psi} Z \ot I \ket{\psi} -\E_{u} \sum_a \bra{\psi}  A^{u}_{a} \ot H_{[h(u)=a]} \ket{\psi}.
  \]
  The reindexing of the sum by~$h$ is the algebraic identity
  \lean{MIPStarRE.LDT.SelfImprovement.helperAgreementOperatorAtPoint_eq_sum_polynomial}
  \[
    \sum_a A^u_a \ot H_{[h(u)=a]} = \sum_h A^u_{h(u)} \ot H_h,
  \]
  whose averaged scalar form reads
  \lean{MIPStarRE.LDT.SelfImprovement.helper_agreement_average_ev_eq_polynomial_sum}
  \[
    \E_u \sum_a \bra{\psi}\,A^u_a \ot H_{[h(u)=a]}\,\ket{\psi}
      = \E_u \sum_h \bra{\psi}\,A^u_{h(u)} \ot H_h\,\ket{\psi}.
  \]
  Using $\sum_a A_a^u=I$, the pointwise difference between the right-placed
  total $I \ot H.\mathrm{total} = \sum_h I \ot H_h$ and the helper-agreement
  operator at $u$ reindexes as the off-diagonal sum
  \lean{MIPStarRE.LDT.SelfImprovement.helperAgreementOperatorAtPoint_off_diagonal_decomposition}
  \[
    \sum_h I \ot H_h - \sum_a A^u_a \ot H_{[h(u)=a]}
      = \sum_h \sum_{a \neq h(u)} A^u_a \ot H_h,
  \]
  whose averaged scalar form reads
  \lean{MIPStarRE.LDT.SelfImprovement.helper_boundedness_slack_average_ev_eq_off_diagonal_avg}
  \[
    \sum_h \bra{\psi}\,I \ot H_h\,\ket{\psi}
        - \E_u \sum_a \bra{\psi}\,A^u_a \ot H_{[h(u)=a]}\,\ket{\psi}
      = \E_u \sum_h \sum_{a \neq h(u)} \bra{\psi}\,A^u_a \ot H_h\,\ket{\psi}.
  \]
  Combined with~\eqref{eq:explicit-bound-for-A-consistency} this gives
  \lean{MIPStarRE.LDT.SelfImprovement.helper_boundedness_gap_eq_upper_gap_add_off_diagonal_avg}
  \lean{MIPStarRE.LDT.SelfImprovement.helper_boundedness_gap_le_selfImprovementHelperError}
  \lean{MIPStarRE.LDT.SelfImprovement.helper_boundedness_gap_le_selfImprovementHelperError_of_pointConsistencyAddInU_transfer}
  \lean{MIPStarRE.LDT.SelfImprovement.helper_upper_gap_rightTensor_le_three_sqrt_delta_of_helper_outputs}
  \lean{MIPStarRE.LDT.SelfImprovement.helper_boundedness_gap_le_selfImprovementHelperError_of_helper_outputs}
  \[
    \bra{\psi} Z \ot I \ket{\psi}
    -\E_{u} \sum_a \bra{\psi}  A^{u}_{a} \ot H_{[h(u)=a]} \ket{\psi}
    \le \bra{\psi} Z \ot I \ket{\psi} - \sum_h \bra{\psi} I \ot H_h \ket{\psi} + 4 \sqrt{\zeta_{\mathrm{variance}}}.
  \]
  By~\eqref{eq:h-versus-z-lower-bound}, the right-hand side is at most $3\sqrt{\delta}+4\sqrt{\zeta_{\mathrm{variance}}}$. Substituting the definition of $\zeta_{\mathrm{variance}}$ and collecting terms gives
  \lean{MIPStarRE.LDT.SelfImprovement.helper_boundedness_error_le_selfImprovementHelperError}
  \[
    23m\cdot\Big(\eps^{1/2} + \delta^{1/2} + (d/q)^{1/2}\Big)
    \le \zeta.
  \]
  Together with~\eqref{eq:Z-greater-than-A}, this proves \ref{item:self-improvement-boundedness}.
\end{proof}

\begin{remark}[Slackness-carrying self-improvement helper]\label{rem:self-improvement-helper-with-slackness-theorem}
  The formalization records a strengthened helper theorem whose output carries
  the complementary-slackness equations used in the completeness chain.  The
  previous top-level residual-domination assembly theorems have been removed:
  the full self-improvement conclusion is now stated only by the paper theorem
  below.  Its proof follows the paper's expectation-level total-difference
  transport route; issue~\#1642 records why the sharper generic operator-total
  strengthening is not supplied by the present route.
\end{remark}

\subsection{A Semidefinite Program}

\begin{lemma}[Uniform SDP feasible witness]\label{lem:sdp-uniform-feasible-witness}
  \lean{MIPStarRE.LDT.SelfImprovement.sdpStrictPrimalWeight}
  \lean{MIPStarRE.LDT.SelfImprovement.sdpStrictPrimalSubMeas}
  \lean{MIPStarRE.LDT.SelfImprovement.sdpStrictPrimalSubMeas_total}
  \lean{MIPStarRE.LDT.SelfImprovement.sdpStrictDualWitness}
  \lean{MIPStarRE.LDT.SelfImprovement.sdpStrictDualWitness_nonneg}
  \lean{MIPStarRE.LDT.SelfImprovement.one_le_sdpStrictDualWitness}
  \leanok
  \uses{def:polyfunc}
  The uniform primal family
  \(T_g=(2|\polyfunc{m}{q}{d}|)^{-1}I\) has total mass \((1/2)I\), and the dual
  operator \(Z=2I\) is positive semidefinite and dominates the identity.
\end{lemma}
\begin{proof}
  \leanok
  This is the elementary Slater-type feasible witness recorded in the basic
  self-improvement definitions.
\end{proof}

\begin{lemma}[Explicit matrix feasible bounds]\label{lem:sdp-matrix-feasible-bounds}
  \uses{lem:sdp-uniform-feasible-witness}
  \lean{MIPStarRE.LDT.SelfImprovement.matrixSdpStrictPrimalSubmeasurement}
  \lean{MIPStarRE.LDT.SelfImprovement.matrixSdpStrictPrimalSubmeasurement_sum_effect}
  \lean{MIPStarRE.LDT.SelfImprovement.matrixSdpStrictDualWitness}
  \lean{MIPStarRE.LDT.SelfImprovement.matrixSdpStrictDualWitness_nonneg}
  \lean{MIPStarRE.LDT.SelfImprovement.one_le_matrixSdpStrictDualWitness}
  \lean{MIPStarRE.LDT.SelfImprovement.matrixAveragedPointOperator_nonneg}
  \lean{MIPStarRE.LDT.SelfImprovement.matrixAveragedPointOperator_le_one}
  \lean{MIPStarRE.LDT.SelfImprovement.matrixSdpStrictDualWitness_dualFeasible}
  \lean{MIPStarRE.LDT.SelfImprovement.one_le_matrixSdpStrictDualWitness_dualSlack}
  \lean{MIPStarRE.LDT.SelfImprovement.matrixSdpDualPositive_of_dualFeasible}
  \lean{MIPStarRE.LDT.SelfImprovement.MatrixSdpFeasibleBounds}
  \lean{MIPStarRE.LDT.SelfImprovement.matrixSdpFeasibleBounds_canonical}
  \lean{MIPStarRE.LDT.SelfImprovement.one_le_matrixSdpCanonicalStrictDualConstraint}
  \lean{MIPStarRE.LDT.SelfImprovement.one_le_matrixSdpCanonicalStrictDualSlackBlockFamily}
  \lean{MIPStarRE.LDT.SelfImprovement.matrixSdpCanonicalStrictPrimalBlockMatrix_feasible}
  \lean{MIPStarRE.LDT.SelfImprovement.matrixSdpCanonicalStrictPrimalBlockMatrix_slack_half}
  \lean{MIPStarRE.LDT.SelfImprovement.matrixSdpCanonicalStrictPrimalBlockMatrix_objective}
  \lean{MIPStarRE.LDT.SelfImprovement.MatrixSdpCanonicalFeasibleBounds}
  \lean{MIPStarRE.LDT.SelfImprovement.matrixSdpCanonicalFeasibleBounds_canonical}
  \leanok
  In every concrete matrix realization, the uniform family
  \(T_g=(2|\polyfunc{m}{q}{d}|)^{-1}I\) has total mass \((1/2)I\).  The dual
  operator \(Z=2I\) dominates the identity and satisfies \(Z-A_g\geq I\), hence
  \(Z\geq A_g\), for every polynomial \(g\).  In the canonical block SDP, the
  corresponding block matrix is feasible, its slack block is \((1/2)I\), and
  the canonical dual slack for the same dual witness also dominates the block
  identity.  More generally, dual feasibility implies positivity of \(Z\), since
  the averaged point operators \(A_g\) are positive semidefinite.
\end{lemma}
\begin{proof}
  \leanok
  The formal proof is the matrix specialization of the uniform SDP witness
  recorded in the basic definitions.  The inequality \(A_g\leq I\) follows by
  averaging the point-measurement effects and using the submeasurement bound at
  each point.
\end{proof}

\begin{lemma}[Matrix slackness output]\label{lem:sdp-matrix-slackness-output}
  \uses{lem:sdp}
  \lean{MIPStarRE.LDT.SelfImprovement.MatrixSdpStatementWithSlackness}
  \lean{MIPStarRE.LDT.MakingMeasurementsProjective.MatrixSubmeasurement.toSubMeas_outcome}
  \lean{MIPStarRE.LDT.MakingMeasurementsProjective.MatrixSubmeasurement.toSubMeas_total}
  \lean{MIPStarRE.LDT.MakingMeasurementsProjective.MatrixMeasurement.toMeasurement}
  \lean{MIPStarRE.LDT.MakingMeasurementsProjective.MatrixMeasurement.toMeasurement_toSubMeas}
  \lean{MIPStarRE.LDT.MakingMeasurementsProjective.MatrixMeasurement.toMeasurement_outcome}
  \lean{MIPStarRE.LDT.SelfImprovement.matrixSdpPrimalTotalEqOne_extracted_of_canonicalSlackBlock_eq_zero}
  \lean{MIPStarRE.LDT.SelfImprovement.matrixSdpOptimalWitness_of_canonicalSaturatedComplementarySlackness}
  \lean{MIPStarRE.LDT.SelfImprovement.matrixSdpOptimalWitness_of_canonicalFeasibleSaturatedComplementarySlackness}
  \lean{MIPStarRE.LDT.SelfImprovement.matrixSdpStatementWithSlackness_of_canonicalSaturatedComplementarySlackness}
  \lean{MIPStarRE.LDT.SelfImprovement.matrixSdpStatementWithSlackness_of_canonicalFeasibleSaturatedComplementarySlackness}
  \lean{MIPStarRE.LDT.SelfImprovement.MatrixSdpCanonicalOptimalPair}
  \lean{MIPStarRE.LDT.SelfImprovement.MatrixSdpCanonicalOptimalPair.toMatrixSdpStatementWithSlackness}
  \lean{MIPStarRE.LDT.SelfImprovement.matrixSdpComplementarySlacknessEquation}
  \lean{MIPStarRE.LDT.MakingMeasurementsProjective.MatrixSubmeasurement.toSubMeas}
  \lean{MIPStarRE.LDT.SelfImprovement.matrixSdpPointRealizationOfStrategy}
  \lean{MIPStarRE.LDT.SelfImprovement.matrixAveragedPointOperator_ofPointRealization}
  \lean{MIPStarRE.LDT.SelfImprovement.matrixSdpDualSlackOperator_ofPointRealization}
  \lean{MIPStarRE.LDT.SelfImprovement.MatrixSdpStatementWithSlackness.exists_measurement_witness}
  \lean{MIPStarRE.LDT.SelfImprovement.MatrixSdpStatementWithSlackness.toSdpStatementWithSlackness}
  \lean{MIPStarRE.LDT.SelfImprovement.matrixSdpPointRealization_statementWithSlackness}
  \lean{MIPStarRE.LDT.SelfImprovement.SdpStatementWithSlackness.exists_measurement_witness}
  \leanok
  From a matrix-level strong-duality witness with complementary slackness of the
  shape asserted by Lemma~\ref{lem:sdp}, one obtains a complete measurement
  \(\{T_g\}\), a feasible dual operator \(Z\), equality of the primal and dual
  objectives, and the equations \(T_gZ=T_gA_g\) for all \(g\).  Lean isolates
  the saturated canonical optimal-pair output required from strong duality: a
  feasible canonical primal matrix, a dual-feasible operator with equal
  objective value, canonical complementary slackness, and vanishing of the
  extra canonical slack block.  This datum gives the matrix-level
  slackness-carrying SDP statement, whose witness is already a complete matrix
  measurement, and is then transported to the abstract Section~9 statement
  without assuming \(I \le Z\).  The displayed measurement witness is read
  directly from that abstract statement.  The canonical saturation datum is now
  the only matrix-level input used to pass from strong duality to the abstract
  SDP statement.  The paper-facing self-improvement helper consumes the
  resulting abstract SDP statement.
\end{lemma}
\begin{proof}
  \leanok
  The identity \(\sum_gT_g=I\) turns the optimal primal submeasurement into a
  measurement; Lean performs this completion before storing the matrix and
  abstract slackness statements.  The zero-defect form \(T_g(Z-A_g)=0\) is
  equivalent, by distributivity, to the displayed equation \(T_gZ=T_gA_g\).
\end{proof}

\begin{remark}[Removed residual-dominating construction theorems]\label{rem:sdp-matrix-residual-domination-assembly}
  Earlier Lean versions exposed top-level conditional construction theorems that
  combined the matrix slackness output with residual-dominating
  orthonormalization and QXP-repair inputs.  These declarations have been
  removed.  The matrix slackness material is retained only up to the
  slackness-carrying helper output, and the full projective-output conclusion is
  represented by the paper theorem
  \texttt{MIPStarRE.LDT.SelfImprovement.selfImprovement}.
\end{remark}

\begin{remark}[Removed reduced SDP interface]\label{rem:lean-reduced-sdp-dominance-interfaces}
  Earlier Lean versions also contained a reduced theorem \texttt{sdp} proving
  only the measurement-total and dual-feasibility fragment used before strong
  duality was supplied, as well as dominance-carrying declarations with the
  auxiliary condition \(I\le Z\).  These declarations have been removed, and
  they are not advertised as formalizations of Lemma~\ref{lem:sdp}.  The active
  Lean route for Lemma~\ref{lem:sdp} is the slackness-carrying statement
  \texttt{sdp\_statement\_with\_slackness}.
\end{remark}

\begin{lemma}[Dual SDP and complementary slackness]\label{lem:sdp}
  \lean{MIPStarRE.LDT.SelfImprovement.averagedPointOperator_nonneg}
  \lean{MIPStarRE.LDT.SelfImprovement.sdpDualPositive_of_dualFeasible}
  \lean{MIPStarRE.LDT.SelfImprovement.SdpOptimalPair.dualPositive}
  \lean{MIPStarRE.LDT.SelfImprovement.SdpOptimalPairWithSlackness}
  \lean{MIPStarRE.LDT.SelfImprovement.SdpOptimalPairWithSlackness.primal_total_operator}
  \lean{MIPStarRE.LDT.SelfImprovement.SdpOptimalPairWithSlackness.dual_positive}
  \lean{MIPStarRE.LDT.SelfImprovement.SdpOptimalPairWithSlackness.dual_feasible}
  \lean{MIPStarRE.LDT.SelfImprovement.SdpOptimalPairWithSlackness.primalMeasurement}
  \lean{MIPStarRE.LDT.SelfImprovement.SdpStatementWithSlackness}
  \lean{MIPStarRE.LDT.SelfImprovement.sdp_statement_with_slackness}
  \lean{MIPStarRE.LDT.SelfImprovement.sdp_slackness_measurement}
  \lean{MIPStarRE.LDT.SelfImprovement.MatrixSdpCanonicalBlockIndex}
  \lean{MIPStarRE.LDT.SelfImprovement.matrixSdpCanonicalBlockHilbertSpace}
  \lean{MIPStarRE.LDT.SelfImprovement.matrixSdpCanonicalDiagonalBlock}
  \lean{MIPStarRE.LDT.SelfImprovement.matrixSdpCanonicalConstraintOperator}
  \lean{MIPStarRE.LDT.SelfImprovement.matrixSdpCanonicalBlockDiagonal}
  \lean{Matrix.blockDiagonal_eq_sum_kronecker_diagonal}
  \lean{Matrix.blockDiagonal_nonneg}
  \lean{Matrix.blockDiagonal_nonneg_iff}
  \lean{Matrix.trace_blockDiagonal_mul}
  \lean{MIPStarRE.Quantum.norm_le_trace_re_of_nonneg}
  \lean{MIPStarRE.Quantum.isClosed_op_nonnegative}
  \lean{MIPStarRE.Quantum.opNonnegativeProperCone}
  \lean{MIPStarRE.LDT.SelfImprovement.matrixSdpCanonicalBlockDiagonal_eq_reindex_blockDiagonal}
  \lean{MIPStarRE.LDT.SelfImprovement.matrixSdpCanonicalBlockDiagonal_apply}
  \lean{MIPStarRE.LDT.SelfImprovement.matrixSdpCanonicalBlockDiagonal_zero}
  \lean{MIPStarRE.LDT.SelfImprovement.matrixSdpCanonicalBlockDiagonal_one}
  \lean{MIPStarRE.LDT.SelfImprovement.matrixSdpCanonicalBlockDiagonal_add}
  \lean{MIPStarRE.LDT.SelfImprovement.matrixSdpCanonicalBlockDiagonal_neg}
  \lean{MIPStarRE.LDT.SelfImprovement.matrixSdpCanonicalBlockDiagonal_sub}
  \lean{MIPStarRE.LDT.SelfImprovement.matrixSdpCanonicalBlockDiagonal_smul}
  \lean{MIPStarRE.LDT.SelfImprovement.matrixSdpCanonicalBlockDiagonal_sub_one}
  \lean{MIPStarRE.LDT.SelfImprovement.matrixSdpCanonicalBlockDiagonal_nonneg}
  \lean{MIPStarRE.LDT.SelfImprovement.matrixSdpCanonicalBlockDiagonal_nonneg_iff}
  \lean{MIPStarRE.LDT.SelfImprovement.matrixSdpCanonicalConstraintOperator_blockDiagonal}
  \lean{MIPStarRE.LDT.SelfImprovement.matrixSdpCanonical_trace_eq_sum_diagonalBlock}
  \lean{MIPStarRE.LDT.SelfImprovement.matrixSdpCanonicalBlockDiagonal_trace_mul}
  \lean{MIPStarRE.LDT.SelfImprovement.matrixSdpCanonicalBlockDiagonal_trace_mul_left}
  \lean{MIPStarRE.LDT.SelfImprovement.matrixSdpCanonicalDiagonalBlock_blockDiagonal_mul_left}
  \lean{MIPStarRE.LDT.SelfImprovement.matrixSdpCanonicalDiagonalBlock_mul_blockDiagonal_right}
  \lean{MIPStarRE.LDT.SelfImprovement.matrixSdpCanonicalBlockDiagonal_mul}
  \lean{MIPStarRE.LDT.SelfImprovement.matrixSdpCanonicalSlackOperator}
  \lean{MIPStarRE.LDT.SelfImprovement.matrixSdpCanonicalSlackOperator_nonneg}
  \lean{MIPStarRE.LDT.SelfImprovement.matrixSdpCanonicalPrimalBlockFamily}
  \lean{MIPStarRE.LDT.SelfImprovement.matrixSdpCanonicalPrimalBlockMatrix}
  \lean{MIPStarRE.LDT.SelfImprovement.matrixSdpCanonicalDiagonalBlock_primalBlockMatrix_some}
  \lean{MIPStarRE.LDT.SelfImprovement.matrixSdpCanonicalDiagonalBlock_primalBlockMatrix_none}
  \lean{MIPStarRE.LDT.SelfImprovement.matrixSdpCanonicalPrimalBlockMatrix_nonneg}
  \lean{MIPStarRE.LDT.SelfImprovement.matrixSdpCanonicalConstraintOperator_primalBlockMatrix}
  \lean{MIPStarRE.LDT.SelfImprovement.MatrixSdpCanonicalPrimalFeasible}
  \lean{MIPStarRE.LDT.SelfImprovement.matrixSdpCanonicalPrimalFeasible_isClosed}
  \lean{MIPStarRE.LDT.SelfImprovement.matrixSdpCanonicalPrimalBlockMatrix_feasible}
  \lean{MIPStarRE.LDT.SelfImprovement.matrixSdpCanonicalDiagonalBlock_nonneg}
  \lean{MIPStarRE.LDT.SelfImprovement.matrixSdpCanonicalPrimalFeasible_sum_diagonalBlock_some_le_one}
  \lean{MIPStarRE.LDT.SelfImprovement.matrixSdpCanonicalExtractedPrimalSubmeasurement}
  \lean{MIPStarRE.LDT.SelfImprovement.matrixSdpCanonicalExtractedPrimalSubmeasurement_effect}
  \lean{MIPStarRE.LDT.SelfImprovement.matrixSdpCanonicalSlackOperator_extractedPrimalSubmeasurement}
  \lean{MIPStarRE.LDT.SelfImprovement.matrixSdpCanonicalDiagonalBlock_primalBlockMatrix_extracted_some}
  \lean{MIPStarRE.LDT.SelfImprovement.matrixSdpCanonicalDiagonalBlock_primalBlockMatrix_extracted_none}
  \lean{MIPStarRE.LDT.SelfImprovement.matrixSdpCanonicalDiagonalBlock_primalBlockMatrix_extracted}
  \lean{MIPStarRE.LDT.SelfImprovement.matrixSdpCanonicalPrimalFeasible_extracts_submeasurement}
  \lean{MIPStarRE.LDT.SelfImprovement.matrixSdpCanonicalObjectiveBlockFamily}
  \lean{MIPStarRE.LDT.SelfImprovement.matrixSdpCanonicalObjectiveOperator}
  \lean{MIPStarRE.LDT.SelfImprovement.matrixSdpCanonicalDiagonalBlockCLM}
  \lean{MIPStarRE.LDT.SelfImprovement.matrixSdpCanonicalConstraintOperatorCLM}
  \lean{MIPStarRE.LDT.SelfImprovement.matrixSdpCanonicalPrimalObjectiveCLM}
  \lean{MIPStarRE.LDT.SelfImprovement.matrixSdpCanonicalDiagonalBlock_objectiveOperator_none}
  \lean{MIPStarRE.LDT.SelfImprovement.matrixSdpCanonicalDiagonalBlock_objectiveOperator_some}
  \lean{MIPStarRE.LDT.SelfImprovement.matrixSdpCanonicalDualOperatorBlockFamily}
  \lean{MIPStarRE.LDT.SelfImprovement.matrixSdpCanonicalDualOperator}
  \lean{MIPStarRE.LDT.SelfImprovement.matrixSdpCanonicalDualOperatorBlockFamily_apply}
  \lean{MIPStarRE.LDT.SelfImprovement.matrixSdpCanonicalDiagonalBlock_dualOperator}
  \lean{MIPStarRE.LDT.SelfImprovement.matrixSdpCanonicalDualSlackBlockFamily}
  \lean{MIPStarRE.LDT.SelfImprovement.matrixSdpCanonicalDualSlackBlockFamily_none}
  \lean{MIPStarRE.LDT.SelfImprovement.matrixSdpCanonicalDualSlackBlockFamily_some}
  \lean{MIPStarRE.LDT.SelfImprovement.matrixSdpCanonicalDualOperator_sub_objectiveOperator}
  \lean{MIPStarRE.LDT.SelfImprovement.matrixSdpCanonicalDualSlackBlockDiagonal_nonneg_of_dualFeasible}
  \lean{MIPStarRE.LDT.SelfImprovement.matrixSdpCanonicalDualSlackBlockDiagonal_nonneg_iff_dualFeasible}
  \lean{MIPStarRE.LDT.SelfImprovement.matrixSdpCanonicalDualConstraint_nonneg_iff_dualFeasible}
  \lean{MIPStarRE.LDT.SelfImprovement.matrixSdpCanonicalDualConstraint_nonneg_of_dualFeasible}
  \lean{MIPStarRE.LDT.SelfImprovement.matrixSdpDualFeasible_of_canonicalDualConstraint_nonneg}
  \lean{MIPStarRE.LDT.SelfImprovement.matrixSdpCanonicalStrictDualConstraint_nonneg}
  \lean{MIPStarRE.LDT.SelfImprovement.matrixSdpCanonicalDualFeasible_isClosed}
  \lean{MIPStarRE.LDT.SelfImprovement.matrixSdpCanonicalObjective_trace_primalBlockMatrix}
  \lean{MIPStarRE.LDT.SelfImprovement.matrixSdpCanonicalObjective_trace_extractedPrimalSubmeasurement}
  \lean{MIPStarRE.LDT.SelfImprovement.matrixSdpCanonicalObjective_trace_primalBlockMatrix_extracted}
  \lean{MIPStarRE.LDT.SelfImprovement.matrixSdpCanonicalDualOperator_trace_feasible}
  \lean{MIPStarRE.LDT.SelfImprovement.matrixSdpCanonicalDualGap_trace}
  \lean{MIPStarRE.Quantum.trace_mul_nonneg_of_nonneg}
  \lean{MIPStarRE.LDT.SelfImprovement.matrixSdpCanonicalPrimalImageCone}
  \lean{MIPStarRE.LDT.SelfImprovement.matrixSdpCanonicalPrimalImageCone_mem_of_nonnegative}
  \lean{MIPStarRE.LDT.SelfImprovement.matrixSdpCanonicalPrimalImageCone_identity_mem_iff_exists_feasible_objective}
  \lean{MIPStarRE.LDT.SelfImprovement.matrixSdpCanonicalNormalizedSeparatorDualMatrix_dualObjective_lt_of_sep}
  \lean{MIPStarRE.LDT.SelfImprovement.matrixSdpCanonicalSeparator_exists_dualFeasible_lt_of_feasible_lt}
  \lean{MIPStarRE.LDT.SelfImprovement.matrixSdpCanonicalStrongDuality}
  \lean{MIPStarRE.LDT.SelfImprovement.matrixSdpCanonicalWeakDuality}
  \lean{MIPStarRE.LDT.SelfImprovement.matrixSdpCanonicalDiagonalBlock_mul_dualSlack}
  \lean{MIPStarRE.LDT.SelfImprovement.matrixSdpCanonicalPrimalBlockMatrix_mul_dualSlack}
  \lean{MIPStarRE.LDT.SelfImprovement.matrixSdpCanonicalPrimalBlockMatrix_extracted_mul_dualSlack_of_canonical}
  \lean{MIPStarRE.LDT.SelfImprovement.matrixSdpComplementarySlacknessDefect_of_canonical}
  \lean{MIPStarRE.LDT.SelfImprovement.matrixSdpComplementarySlacknessDefect_extracted_of_canonical}
  \lean{MIPStarRE.LDT.SelfImprovement.matrixSdpCanonicalSlack_mul_dual_of_complementarySlackness}
  \leanok
  Let
  \[
    A_g = \E_{u \sim \F_q^m} A^u_{g(u)}.
  \]
  Then the primal is
  \begin{align}
    \sup &\quad \sum_g \,\Tr(T_g \cdot A_g) \label{eq:primal-objective}\\
    \text{s.t.} &\quad   T_g \geq 0\qquad\forall g\in\polyfunc{m}{q}{d}, \notag\\
    &\quad \sum_g T_g \leq I, \notag
  \end{align}
  and the dual is
  \begin{align}
    \inf &\quad \Tr(Z) \label{eq:dual-objective}\\
    \text{ s.t.} &\quad Z \geq A_g. \label{eq:dual-constraint}
  \end{align}
  These semidefinite programs are dual to each other. Moreover there is an optimal pair of solutions $\{T_g\}$ to~\eqref{eq:primal-objective} and $Z$ to~\eqref{eq:dual-objective} such that $\sum_g T_g = I$ and
  \begin{equation}\label{eq:slater}
    T_g Z  \,=\, T_g {A_g},\qquad\forall g\in \polyfunc{m}{q}{d}.
  \end{equation}
\end{lemma}
\begin{proof}
  \uses{lem:sdp-matrix-feasible-bounds}
  \leanok
  Rewrite~\eqref{eq:primal-objective} in canonical block form. Let $r$ be the dimension of the space on which~$A$ acts, let $M = |\polyfunc{m}{q}{d}|$, and fix an ordering $g_1,\ldots,g_M$ of the polynomials. Consider
  \begin{align}
    \sup &\quad \Tr(C^\dagger X) \label{eq:primal-canonical}\\
    \text{s.t.} &\quad \Tr(D_{ij}^\dagger X) = b_{ij} \qquad\forall i,j\in\{1,\ldots,r\}, \notag\\
    &\quad X \geq 0, \notag
  \end{align}
  where
  \[
    C = \sum_{i = 1}^M \ket{i}\bra{i} \ot A_{g_i},
  \]
  and, for $i,j \in \{1,\ldots,r\}$,
  \[
    D_{ij} = \sum_{k=1}^{M+1} \ket{k}\bra{k} \ot \ket{i}\bra{j},
    \qquad
    b_{ij}  = \left\{\begin{array}{rl}
      1 & \text{if } i = j,\\
      0 & \text{otherwise}.
    \end{array}\right.
  \]
  Writing $X=\sum_{i,j=1}^{M+1}\ket{i}\bra{j}\ot X_{ij}$, the constraints say precisely that $\sum_{i=1}^{M+1} X_{ii}=I$, hence $\sum_{i=1}^{M}X_{ii}\le I$, and the objective is
  \[
    \Tr(C^\dagger X)=\sum_{i=1}^M \Tr(A_{g_i}\cdot X_{ii}).
  \]
  Thus~\eqref{eq:primal-canonical} is equivalent to~\eqref{eq:primal-objective} under the identification $T_{g_i}=X_{ii}$.

  The canonical dual is
  \begin{align}
    \inf &\quad \sum_i\, z_{ij} b_{ij} \label{eq:dual-canonical}\\
    \text{s.t.} &\quad  \sum_{i,j} \,z_{ij} D_{ij} \geq C. \label{eq:dual-canonical-constraint}
  \end{align}
  Since
  \[
    \sum_{i,j=1}^r z_{ij} D_{ij}
    = \sum_{k=1}^{M+1} \ket{k}\bra{k} \ot \left(\sum_{i,j=1}^r z_{ij}\ket{i}\bra{j}\right)
    =: \sum_{k=1}^{M+1}\ket{k}\bra{k} \ot Z,
  \]
  the constraint~\eqref{eq:dual-canonical-constraint} is exactly $Z\ge A_{g_i}$ for every $i$, and the objective is $\Tr(Z)$.  In Lean, the latter statement is recorded as the trace-pairing identity
  \[
    \Re\Tr(D(Z)X)=\Re\Tr(Z)
  \]
  for every feasible canonical primal matrix \(X\), where \(D(Z)\) denotes the
  block-diagonal canonical dual operator.  Consequently the canonical
  primal-dual gap is the real trace pairing of \(X\) with the canonical dual
  slack \(D(Z)-C\).  Since the real trace pairing of positive semidefinite
  operators is nonnegative, Lean also records the resulting weak-duality
  inequality for any feasible canonical primal and dual pair.
  Hence~\eqref{eq:dual-canonical} is equivalent to~\eqref{eq:dual-objective}.

  Slater's condition holds because
  \[
    T_g = \frac{1}{2M}\cdot I \qquad\forall g\in\polyfunc{m}{q}{d},
    \qquad
    Z = 2I
  \]
  are strictly feasible for the primal and dual respectively. Therefore strong duality holds. For an optimal pair $(X,(z_{ij}))$ in canonical form, complementary slackness gives
  \begin{equation}\label{eq:complementary-slackness}
    X\Big( \sum_{i,j} \,z_{ij} D_{ij} - C \Big) \,=\, 0.
  \end{equation}
  We may assume that $X$ is block diagonal.
  Translating~\eqref{eq:complementary-slackness} back to the variables
  $\{T_g\}$ and~$Z$ yields $T_g(Z-A_g)=0$ for every~$g$ and \(S Z=0\), where
  \(S=I-\sum_gT_g\) is the slack block.
  The paper then treats this as forcing \(S=0\).  The Lean development records
  this saturation as part of the canonical optimal-pair datum for the block
  SDP.  The canonical calculation still uses the auxiliary expression
  \texttt{matrixSdpComplementarySlacknessDefect} to read off the polynomial
  blocks of \eqref{eq:complementary-slackness}.  The matrix optimal-witness
  interface, however, stores the displayed equation \(T_gZ=T_gA_g\) directly;
  \texttt{MatrixSdpOptimalWitness.primalMeasurement} expresses
  \(\sum_g T_g=I\) as a complete measurement.  The statement
  \texttt{MatrixSdpStatementWithSlackness} records this matrix-level
  strong-duality output and the abstract theorem
  \texttt{sdp\_statement\_with\_slackness} specializes it to the paper's
  self-improvement setting.  The explicit Slater witnesses used above are
  isolated in Lemma~\ref{lem:sdp-matrix-feasible-bounds} and separately recorded
  by \texttt{matrixSdpStrictPrimalSubmeasurement},
  \texttt{matrixSdpStrictDualWitness}, and
  \texttt{matrixSdpFeasibleBounds\_canonical}; the dual feasibility of \(2I\)
  follows from \texttt{matrixAveragedPointOperator\_le\_one}.
  The block-diagonal reduction is also recorded for arbitrary feasible
  canonical matrices: extracting the diagonal blocks preserves the objective
  and transfers canonical complementary slackness to the extracted paper primal
  submeasurement.
\end{proof}

\begin{lemma}[Averaging over the evaluation point]\label{lem:add-in-u}
  \lean{MIPStarRE.LDT.SelfImprovement.AddInUFullStatement}
  \lean{MIPStarRE.LDT.SelfImprovement.addInUFullStatement_of_isGood}
  \lean{MIPStarRE.LDT.SelfImprovement.selfConsistencyAddInUSelection}
  \lean{MIPStarRE.LDT.SelfImprovement.addInULeftQuantity_selfConsistencySelection_eq_matchMass}
  \lean{MIPStarRE.LDT.SelfImprovement.addInURightQuantity_selfConsistencySelection_eq_release}
  \lean{MIPStarRE.LDT.SelfImprovement.addInURightQuantity_selfConsistencySelection_eq_simplified}
  \lean{MIPStarRE.LDT.SubMeas.postprocess_id}
  \lean{MIPStarRE.LDT.rightTensor_le_one}
  \lean{MIPStarRE.LDT.SubMeas.opTensor_sum_filter_le_one}
  \lean{MIPStarRE.LDT.SelfImprovement.addInU_pointMeasurement_snd_selfConsistency}
  \lean{MIPStarRE.LDT.SelfImprovement.addInU_filtered_sandwiched_tensor_sum_le_one}
  \lean{MIPStarRE.LDT.SelfImprovement.addInUCSChainQ0}
  \lean{MIPStarRE.LDT.SelfImprovement.addInUCSChainQ1}
  \lean{MIPStarRE.LDT.SelfImprovement.addInUCSChainQ2}
  \lean{MIPStarRE.LDT.SelfImprovement.addInUCSChainQ3}
  \lean{MIPStarRE.LDT.SelfImprovement.addInUCSChainQ4}
  \lean{MIPStarRE.LDT.SelfImprovement.addInUSelectedCSChainQ0}
  \lean{MIPStarRE.LDT.SelfImprovement.addInUSelectedCSChainQ1}
  \lean{MIPStarRE.LDT.SelfImprovement.addInUSelectedCSChainQ2}
  \lean{MIPStarRE.LDT.SelfImprovement.addInUSelectedCSChainQ3}
  \lean{MIPStarRE.LDT.SelfImprovement.addInUSelectedCSChainQ4}
  \lean{MIPStarRE.LDT.SelfImprovement.add_in_u_cs_chain_q0_eq_match_mass}
  \lean{MIPStarRE.LDT.SelfImprovement.add_in_u_cs_chain_q4_eq_simplified_rhs}
  \lean{MIPStarRE.LDT.SelfImprovement.addInUSelectedCSChainQ0_eq_leftQuantity_averagedSandwiched}
  \lean{MIPStarRE.LDT.SelfImprovement.addInUSelectedCSChainQ4_eq_rightQuantity}
  \lean{MIPStarRE.LDT.SelfImprovement.addInU_cs_chain_step1_diff_eq}
  \lean{MIPStarRE.LDT.SelfImprovement.addInU_cs_chain_step2_diff_eq}
  \lean{MIPStarRE.LDT.SelfImprovement.addInU_cs_chain_step1_reverse_diff_eq}
  \lean{MIPStarRE.LDT.SelfImprovement.addInU_cs_chain_step2_reverse_diff_eq}
  \lean{MIPStarRE.LDT.SelfImprovement.addInU_cs_chain_step3_diff_eq}
  \lean{MIPStarRE.LDT.SelfImprovement.addInU_cs_chain_step4_diff_eq}
  \lean{MIPStarRE.LDT.SelfImprovement.addInU_selected_cs_chain_step1_diff_eq}
  \lean{MIPStarRE.LDT.SelfImprovement.addInU_selected_cs_chain_step2_diff_eq}
  \lean{MIPStarRE.LDT.SelfImprovement.addInU_selected_cs_chain_step3_diff_eq}
  \lean{MIPStarRE.LDT.SelfImprovement.addInU_selected_cs_chain_step4_diff_eq}
  \lean{MIPStarRE.LDT.SelfImprovement.addInU_selected_cs_chain_step1_reverse_diff_eq}
  \lean{MIPStarRE.LDT.SelfImprovement.addInU_selected_cs_chain_step2_reverse_diff_eq}
  \lean{MIPStarRE.LDT.SelfImprovement.addInU_selected_cs_chain_step3_reverse_diff_eq}
  \lean{MIPStarRE.LDT.SelfImprovement.addInU_selected_cs_chain_step4_reverse_diff_eq}
  \lean{MIPStarRE.LDT.SelfImprovement.addInU_selected_cs_chain_step1_abs_le_sqrt_two_delta}
  \lean{MIPStarRE.LDT.SelfImprovement.addInU_selected_cs_chain_step2_abs_le_sqrt_two_delta}
  \lean{MIPStarRE.LDT.SelfImprovement.addInU_selected_cs_chain_step3_abs_le_sqrt_globalVarianceDeviation_sum}
  \lean{MIPStarRE.LDT.SelfImprovement.addInU_selected_cs_chain_step4_abs_le_sqrt_globalVarianceDeviation_sum}
  \lean{MIPStarRE.LDT.SelfImprovement.addInU_selected_cs_chain_step3_abs_le_sqrt_of_globalVarianceDeviation_sum_le}
  \lean{MIPStarRE.LDT.SelfImprovement.addInU_selected_cs_chain_step4_abs_le_sqrt_of_globalVarianceDeviation_sum_le}
  \lean{MIPStarRE.LDT.SelfImprovement.addInU_selected_cs_chain_step34_abs_le_sqrt_of_globalVarianceDeviation_sum_le}
  \lean{MIPStarRE.LDT.SelfImprovement.addInU_cs_chain_step1_abs_le_sqrt_two_delta}
  \lean{MIPStarRE.LDT.SelfImprovement.addInU_cs_chain_step2_abs_le_sqrt_two_delta}
  \lean{MIPStarRE.LDT.SelfImprovement.add_in_u_cs_chain_q2_q3_factored_cs}
  \lean{MIPStarRE.LDT.SelfImprovement.add_in_u_cs_chain_q3_q4_factored_cs}
  \lean{MIPStarRE.LDT.SelfImprovement.add_in_u_cs_chain_q2_q3_self_energy_factor_le_one}
  \lean{MIPStarRE.LDT.SelfImprovement.add_in_u_cs_chain_q3_q4_self_energy_factor_le_one}
  \lean{MIPStarRE.LDT.SelfImprovement.add_in_u_cs_chain_q2_q3_variance_factor_le_globalVarianceDeviation_sum}
  \lean{MIPStarRE.LDT.SelfImprovement.add_in_u_cs_chain_q3_q4_variance_factor_le_globalVarianceDeviation_sum}
  \lean{MIPStarRE.LDT.SelfImprovement.add_in_u_cs_chain_q2_q3_le_sqrt_globalVarianceDeviation_sum}
  \lean{MIPStarRE.LDT.SelfImprovement.add_in_u_cs_chain_q3_q4_le_sqrt_globalVarianceDeviation_sum}
  \lean{MIPStarRE.LDT.SelfImprovement.add_in_u_cs_chain_q2_q3_le_sqrt_of_globalVarianceDeviation_sum_le}
  \lean{MIPStarRE.LDT.SelfImprovement.add_in_u_cs_chain_q3_q4_le_sqrt_of_globalVarianceDeviation_sum_le}
  \lean{MIPStarRE.LDT.SelfImprovement.add_in_u_cs_chain_q2_q3_le_sqrt_of_globalVarianceDeviation_sum_le_from_factor_bounds}
  \lean{MIPStarRE.LDT.SelfImprovement.add_in_u_cs_chain_q3_q4_le_sqrt_of_globalVarianceDeviation_sum_le_from_factor_bounds}
  \lean{MIPStarRE.LDT.SelfImprovement.add_in_u_cs_chain_global_variance_steps_of_sum_bound}
  \lean{MIPStarRE.LDT.SelfImprovement.add_in_u_cs_chain_global_variance_steps_of_local_sum_bound}
  \lean{MIPStarRE.LDT.SelfImprovement.add_in_u_cs_chain_global_variance_steps_of_sum_bound_from_factor_bounds}
  \lean{MIPStarRE.LDT.SelfImprovement.add_in_u_cs_chain_q2_q3_le_sqrt_of_localVarianceDeviation_sum_le_from_factor_bounds}
  \lean{MIPStarRE.LDT.SelfImprovement.add_in_u_cs_chain_q3_q4_le_sqrt_of_localVarianceDeviation_sum_le_from_factor_bounds}
  \lean{MIPStarRE.LDT.SelfImprovement.add_in_u_cs_chain_global_variance_steps_of_local_sum_bound_from_factor_bounds}
  \lean{MIPStarRE.LDT.SelfImprovement.add_in_u_simplified_transfer_of_cs_chain}
  \lean{MIPStarRE.LDT.SelfImprovement.add_in_u_selected_transfer_of_cs_chain}
  \lean{MIPStarRE.LDT.SelfImprovement.add_in_u_simplified_transfer_of_cs_chain_sqrt_form}
  \lean{MIPStarRE.LDT.SelfImprovement.add_in_u_simplified_transfer_of_cs_chain_local_variance_form}
  \lean{MIPStarRE.LDT.SelfImprovement.add_in_u_simplified_transfer_of_cs_chain_selfConsistency_local_variance_form}
  \lean{MIPStarRE.LDT.SelfImprovement.two_sqrt_two_mul_add_two_sqrt_le_four_sqrt}
  \lean{MIPStarRE.LDT.SelfImprovement.two_mul_delta_le_selfImprovementVarianceError}
  \lean{MIPStarRE.LDT.SelfImprovement.two_sqrt_two_delta_add_two_sqrt_selfImprovementVarianceError_le_addInUError}
  \lean{MIPStarRE.LDT.SelfImprovement.selfConsistencyDiagonalAddInU_of_transfer}
  \lean{MIPStarRE.LDT.SelfImprovement.selfConsistencyDiagonalAddInU_of_simplifiedTransfer}
	\leanok
  \uses{lem:global-variance-of-points, prop:simeq-to-approx}
  Let $T=\{T_h\}$ be an optimal solution of the primal SDP, let
  \[
    H_h = \E_{u \sim \F_q^m} A^u_{h(u)} T_h A^u_{h(u)},
  \]
  and set
  \[
    \zeta_{\mathrm{variance}} = 24m\cdot\Big(\eps + \delta + \frac{md}{q}\Big).
  \]
  Suppose $M = \{M^u_o\}$ is a submeasurement with outcomes in some set $\mathcal O$. For each $u \in \F_q^m$, let $S_u$ be a subset of $\mathcal O \times \polyfunc{m}{q}{d}$. Then
  \[
    \E_{u} \sum_{(o, h) \in S_{u}} \bra{\psi} M^{u}_o \ot H_h \ket{\psi}
    \approx_{4\sqrt{\zeta_{\mathrm{variance}}}}
    \E_{u} \sum_{(o, h) \in S_{u}} \bra{\psi} \bigl(A^{u}_{h(u)} \cdot M^{u}_o \cdot A^{u}_{h(u)}\bigr) \ot T_h \ket{\psi}.
  \]
\end{lemma}
The reduced variance-bound specialization
\lean{MIPStarRE.LDT.SelfImprovement.addInU}
keeps only the variance-bound consequence used later in the helper theorem; it
is not itself the full formal counterpart of Lemma~\ref{lem:add-in-u}.
\begin{proof}
  \leanok
  We begin by expanding
  \begin{align}
    \E_{u} \sum_{(o, h) \in S_{u}} \bra{\psi} M^{u}_o \ot H_h \ket{\psi}
    &= \E_{u,v} \sum_{(o, h) \in S_{u}} \bra{\psi} M^{u}_o \ot H^{v}_h \ket{\psi} \notag\\
    &= \E_{u,v} \sum_{(o, h) \in S_{u}} \bra{\psi} M^{u}_o \ot \bigl(A^{v}_{h(v)} \cdot T_h \cdot A^{v}_{h(v)}\bigr) \ket{\psi}. \label{eq:expand-that-H}
  \end{align}
  We claim that
  \begin{equation}\label{eq:move-first-point-projector}
    \eqref{eq:expand-that-H}
    \approx_{\sqrt{2\delta}}
    \E_{u,v} \sum_{(o, h) \in S_{u}} \bra{\psi} \bigl(A^{v}_{h(v)} \cdot M^{u}_o\bigr) \ot \bigl(T_h \cdot A^{v}_{h(v)}\bigr) \ket{\psi}.
  \end{equation}
  To show this, we bound the magnitude of the difference:
  \begin{multline}
    \Big|\E_{u,v} \sum_{(o, h) \in S_{u}} \bra{\psi} \bigl(A^{v}_{h(v)} \ot I - I \ot A^{v}_{h(v)}\bigr) \cdot \bigl(M^{u}_o \ot (T_h \cdot A^{v}_{h(v)})\bigr)\ket{\psi}\Big|\\
    \leq
    \sqrt{\E_{u,v} \sum_{(o, h) \in S_{u}} \bra{\psi} \bigl(A^{v}_{h(v)} \ot I - I \ot A^{v}_{h(v)}\bigr) \cdot (M^{u}_o \ot T_h) \cdot \bigl(A^{v}_{h(v)} \ot I - I \ot A^{v}_{h(v)}\bigr) \ket{\psi}}\\
    \cdot \sqrt{\E_{u,v} \sum_{(o, h) \in S_{u}} \bra{\psi}  M^{u}_o \ot \bigl(A^{v}_{h(v)} \cdot T_h\cdot  A^{v}_{h(v)}\bigr) \ket{\psi}}. \label{eq:move-first-point-projector-cauchy-schwarz}
  \end{multline}
  The first factor is controlled by summing over the value $a = h(v)$:
  \begin{align*}
    &\E_{u,v} \sum_{a \in \F_q} \bra{\psi} \bigl(A^v_a \ot I - I \ot A^v_a\bigr)
      \cdot \Bigl(\sum_{(o,h)\in S_u,\, h(v)=a} M^u_o \ot T_h\Bigr)
      \cdot \bigl(A^v_a \ot I - I \ot A^v_a\bigr)\ket{\psi} \\
    &\le \E_{u,v} \sum_{a \in \F_q} \bra{\psi} \bigl(A^v_a \ot I - I \ot A^v_a\bigr)^2
      \ket{\psi},
  \end{align*}
  where the inequality uses that $M^u$ and $T$ are submeasurements. By
  Lemma~\ref{prop:simeq-to-approx} and the self-consistency of~$A$, this is at most
  $2\delta$. The second factor is
  \[
    \E_{u,v} \sum_{(o,h)\in S_u} \bra{\psi} M^u_o \ot H^v_h \ket{\psi},
  \]
  which is at most~$1$ because $M^u$ and $H^v$ are submeasurements.

  Next, we claim that
  \begin{equation}\label{eq:move-second-point-projector}
    \eqref{eq:move-first-point-projector}
    \approx_{\sqrt{2\delta}}
    \E_{u,v} \sum_{(o, h) \in S_{u}} \bra{\psi} \bigl(A^{v}_{h(v)} \cdot M^{u}_o \cdot A^{v}_{h(v)}\bigr) \ot T_h  \ket{\psi}.
  \end{equation}
  Again, we bound the magnitude of the difference:
  \begin{align}\label{eq:move-second-point-projector-cauchy-schwarz}
    &\Big|\E_{u,v} \sum_{(o, h) \in S_{u}} \bra{\psi} \bigl((A^{v}_{h(v)} \cdot M^{u}_o) \ot T_h\bigr) \cdot \bigl(A^{v}_{h(v)} \ot I - I \ot A^{v}_{h(v)}\bigr)  \ket{\psi}\Big| \notag\\
    &\leq \sqrt{\E_{u,v} \sum_{(o, h) \in S_{u}} \bra{\psi} \bigl(A^{v}_{h(v)} \cdot M^{u}_o \cdot A^{v}_{h(v)}\bigr) \ot T_h  \ket{\psi}} \\
    &\qquad\cdot \sqrt{\E_{u,v} \sum_{(o, h) \in S_{u}} \bra{\psi} \bigl(A^{v}_{h(v)} \ot I - I \ot A^{v}_{h(v)}\bigr) \cdot (M^{u}_o \ot T_h) \cdot \bigl(A^{v}_{h(v)} \ot I - I \ot A^{v}_{h(v)}\bigr) \ket{\psi}}. \notag
  \end{align}
  The first factor is bounded by
  \begin{align*}
    &\E_{u,v} \sum_h \bra{\psi}
      \Bigl(A^v_{h(v)} \cdot \sum_{o:(o,h)\in S_u} M^u_o \cdot A^v_{h(v)}\Bigr)
      \ot T_h \ket{\psi} \\
    &\le \E_{u,v} \sum_h \bra{\psi} (A^v_{h(v)})^2 \ot T_h \ket{\psi}
      \le \E_{u,v} \sum_h \bra{\psi} I \ot T_h \ket{\psi}
      \le 1,
  \end{align*}
  using first that $M$ is a submeasurement, then that $A^v_{h(v)} \le I$, and
  finally that $T$ is a submeasurement ($\sum_g T_g \le I$). The second factor is
  exactly the same expression
  as in~\eqref{eq:move-first-point-projector-cauchy-schwarz}, hence at most
  $\sqrt{2\delta}$.

  Having moved both copies of~$A$ to the left-hand side, we next show that
  \begin{equation}\label{eq:replace-left-point-conditioner}
    \eqref{eq:move-second-point-projector}
    \approx_{\sqrt{\zeta_{\mathrm{variance}}}}
    \E_{u,v} \sum_{(o, h) \in S_{u}} \bra{\psi} \bigl(A^{u}_{h(u)} \cdot M^{u}_o \cdot A^{v}_{h(v)}\bigr) \ot T_h  \ket{\psi}.
  \end{equation}
  The difference is bounded by
  \begin{multline}\label{eq:replace-left-point-conditioner-cauchy-schwarz}
    \Big|\E_{u,v} \sum_{(o, h) \in S_{u}} \bra{\psi} \bigl((A^{v}_{h(v)} - A^{u}_{h(u)}) \cdot M^{u}_o \cdot A^{v}_{h(v)}\bigr) \ot T_h  \ket{\psi}\Big|\\
    \leq
    \sqrt{\E_{u,v} \sum_{(o, h) \in S_{u}} \bra{\psi} \bigl((A^{v}_{h(v)} - A^{u}_{h(u)}) \cdot M^{u}_o \cdot (A^{v}_{h(v)} - A^{u}_{h(u)})\bigr) \ot T_h  \ket{\psi}}\\
    \cdot \sqrt{\E_{u,v} \sum_{(o, h) \in S_{u}} \bra{\psi} \bigl(A^{v}_{h(v)} \cdot M^{u}_o \cdot A^{v}_{h(v)}\bigr) \ot T_h  \ket{\psi}}.
  \end{multline}
  The first factor can be rewritten as
  \begin{align*}
    &\E_{u,v} \sum_h \bra{\psi}
      \Bigl((A^v_{h(v)} - A^u_{h(u)}) \cdot
      \sum_{o:(o,h)\in S_u} M^u_o \cdot
      (A^v_{h(v)} - A^u_{h(u)})\Bigr) \ot T_h \ket{\psi} \\
    &\le \E_{u,v} \sum_h \bra{\psi}
      (A^v_{h(v)} - A^u_{h(u)})^2 \ot T_h \ket{\psi},
  \end{align*}
  because $M^u$ is a submeasurement. Lemma~\ref{lem:global-variance-of-points}
  bounds this by~$\zeta_{\mathrm{variance}}$ since $T \in \polysub{m}{q}{d}$. The
  second factor is exactly the first factor
  from~\eqref{eq:move-second-point-projector-cauchy-schwarz},
  so it is at most~$1$.

  Finally, we show that
  \begin{equation}\label{eq:replace-right-point-conditioner}
    \eqref{eq:replace-left-point-conditioner}
    \approx_{\sqrt{\zeta_{\mathrm{variance}}}}
    \E_{u,v} \sum_{(o, h) \in S_{u}} \bra{\psi} \bigl(A^{u}_{h(u)} \cdot M^{u}_o \cdot A^{u}_{h(u)}\bigr) \ot T_h  \ket{\psi}.
  \end{equation}
  The same Cauchy--Schwarz estimate
  as in~\eqref{eq:replace-left-point-conditioner-cauchy-schwarz}, with
  the final copy of $A^v_{h(v)}$ replaced by $A^u_{h(u)}$, gives the same bound
  $\sqrt{\zeta_{\mathrm{variance}}}$. Indeed, the factor containing
  $A^u_{h(u)} \cdot M^u_o \cdot A^u_{h(u)}$ is bounded by~$1$ by the same
  submeasurement argument as above, while the factor containing
  $(A^v_{h(v)} - A^u_{h(u)})^2$ is unchanged. Summing the four transports gives an
  error of $2\sqrt{2\delta} + 2\sqrt{\zeta_{\mathrm{variance}}}$, and the lemma
  follows from $2\delta \leq \zeta_{\mathrm{variance}}$.
\end{proof}

\section{Projective Output}
\label{sec:self-improvement-projective}

This final step applies orthonormalization to the non-projective family produced by
Lemma~\ref{lem:self-improvement-helper}. The bookkeeping here is entirely about
showing that completeness, consistency, self-consistency, and boundedness survive
the passage from~$\widehat H$ to a projective submeasurement~$H$.

\begin{remark}[Auxiliary self-improvement lemmas]\label{rem:self-improvement-lean-interfaces}
  Several auxiliary results are used to prove Theorem~\ref{thm:self-improvement}
  from the helper output and the orthonormalization theorem.  They record
  local intermediate estimates and final-field assembly lemmas.  The
  former orthonormalization spectral-transport record has been retired; the
  spectral-truncation part of the construction is now represented directly by
  \texttt{spectralTruncationStatement\_of\_sourceAlmostProjective}.  This result is
  not a locality-preserving repair theorem.  The theorem below records the
  paper-facing conclusion.  Its proof follows the expectation-level completeness
  argument from the paper; the obstruction to a generic operator-total
  strengthening is recorded in issue~\#1642.
\end{remark}

\begin{theorem}[Self-improvement]\label{thm:self-improvement}
  \lean{MIPStarRE.LDT.SelfImprovement.selfImprovement}
  \leanok
  \uses{def:polymeasurements, def:simeq, def:approx_delta, def:strong-self-consistency}
  Let $G \in \polymeas{m}{q}{d}$ be a measurement with the following properties:
  \begin{enumerate}
    \item (Consistency with~$A$): \label{item:self-improvement-projective-G-consistency} On average over $u \sim \F_q^{m}$,
    \[
      A^{u}_a \ot I \simeq_{\nu} I \ot G_{[g(u)=a]}.
    \]
  \end{enumerate}
  Let
  \[
    \zeta = 3000m\cdot \Big(\eps^{1/32} + \delta^{1/32} + (d/q)^{1/32}\Big).
  \]
  Then there exists a projective submeasurement $H \in \polysub{m}{q}{d}$ with the following properties:
  \begin{enumerate}
    \item (Completeness): \label{item:self-improvement-projective-completeness} If $H = \sum_h H_h$, then
    \[
      \bra{\psi} H \ot I \ket{\psi} \geq (1-\nu)-\zeta.
    \]
    \item (Consistency with~$A$): \label{item:self-improvement-projective-A-consistency} On average over $u \sim \F_q^m$,
    \[
      A^u_a \ot I \simeq_{\zeta} I \ot H_{[h(u) = a]}.
    \]
    \item (Strong self-consistency): \label{item:self-improvement-projective-self}
    \[
      H_h \ot I \approx_{\zeta} I \ot H_h.
    \]
    \item (Boundedness): \label{item:self-improvement-projective-boundedness} There exists a positive-semidefinite matrix~$Z$ such that
    \[
      \bra{\psi} Z \ot (I - H) \ket{\psi} \leq \zeta
    \]
    and for each $h \in \polyfunc{m}{q}{d}$,
    \[
      Z \geq \left(\E_{u} A^{u}_{h(u)}\right).
    \]
  \end{enumerate}
\end{theorem}
\begin{proof}
  \leanok
  \uses{lem:self-improvement-helper, thm:orthonormalization,
    prop:self-consistency-implies-data-processing,
    prop:completeness-transfer-self-consistent-A, prop:triangle-sub,
    prop:two-notions-of-self-consistency,
    prop:triangle-inequality-for-approx_delta, prop:easy-approx-from-approx-delta}
  The bound is trivial if one of $\eps$, $\delta$, or $d/q$ is at least~$1$, so assume all three are at most~$1$.  In Lean, the third assumption is recorded as $d \le q$, which is equivalent to $d/q \le 1$ since $q>0$. Apply Lemma~\ref{lem:self-improvement-helper} to~$G$ and let $\widehat H \in \polysub{m}{q}{d}$ be the output, with
  \[
    \widehat{\zeta} = 100m \cdot \Big(\eps^{1/2} + \delta^{1/2} + (d/q)^{1/2}\Big).
  \]
  By \ref{item:self-improvement-self} of Lemma~\ref{lem:self-improvement-helper},
  \[
    \sum_{h} \bra{\psi} \widehat{H}_h \ot \widehat{H}_h \ket{\psi}
    \geq \bra{\psi} \widehat{H} \ot I \ket{\psi} - \widehat{\zeta}.
  \]
  Therefore Theorem~\ref{thm:orthonormalization} supplies a projective submeasurement $H \in \polysub{m}{q}{d}$ such that
  \begin{equation}\label{eq:approx-between-H-with-and-without-hat}
    \widehat{H}_h \ot I \approx_{\widehat{\zeta}_{\mathrm{ortho}}} H_{h} \ot I,
  \end{equation}
  where
  \[
    \widehat{\zeta}_{\mathrm{ortho}} = 100 \widehat{\zeta}^{1/4}.
  \]
  In addition, Lemma~\ref{prop:self-consistency-implies-data-processing} implies that
  \begin{equation}\label{eq:approx-data-processed}
    \widehat{H}_{[h(u)=a]} \ot I \approx_{\widehat{\zeta}_{\mathrm{dataprocess}}} H_{[h(u)=a]} \ot I,
  \end{equation}
  where
  \[
    \widehat{\zeta}_{\mathrm{dataprocess}} = 8 \widehat{\zeta} + 8 \sqrt{\widehat{\zeta}_{\mathrm{ortho}}}.
  \]

  The four conclusions are transferred one by one.

  \emph{Proof of \ref{item:self-improvement-projective-completeness}.}
  \lean{MIPStarRE.LDT.SelfImprovement.completeness_transport_through_orthonormalization}
  \lean{MIPStarRE.LDT.SelfImprovement.final_fields_completeness_of_helper_completeness}
  \lean{MIPStarRE.LDT.SelfImprovement.final_fields_completeness_of_helper_completeness_of_small_errors}
  Lemma~\ref{prop:completeness-transfer-self-consistent-A} implies
  \[
    \bra{\psi} H \ot I \ket{\psi}
    \geq \bra{\psi} \widehat{H} \ot I \ket{\psi} - \widehat{\zeta} - 2\sqrt{\widehat{\zeta}_{\mathrm{ortho}}}
    \geq (1-\nu) - 2\widehat{\zeta} - 2\sqrt{\widehat{\zeta}_{\mathrm{ortho}}},
  \]
  where the second inequality is \ref{item:self-improvement-completeness} of Lemma~\ref{lem:self-improvement-helper}.

  \emph{Proof of \ref{item:self-improvement-projective-A-consistency}.}
  Lemma~\ref{prop:triangle-sub} applied to \ref{item:self-improvement-A-consistency} of Lemma~\ref{lem:self-improvement-helper} and the post-processed comparison~\eqref{eq:approx-data-processed} gives
  \[
    A^u_a \ot I \simeq_{\widehat{\zeta} + \sqrt{\widehat{\zeta}_{\mathrm{dataprocess}}}}
    I \ot H_{[h(u) = a]}.
  \]
  In Lean the corresponding transport is stated for submeasurements with the
  total-overlap displacement recorded explicitly; this is the term which
  vanishes in the measurement-valued form of Lemma~\ref{prop:triangle-sub}.
  The monotone-total route records a sufficient structural replacement: if the
  orthonormalization repair preserves the completed residual outcome, then the
  projective total is bounded by the helper total as an operator, hence also
  has no larger right-register expectation in the strategy state.  Under this
  structural hypothesis the alphabet-size displacement term is not introduced.
  The QXP-layer interface used by the formalization records both the total
  domination of the canonical projective family and the sharper
  residual comparison for the fresh option-completion outcome.  At present the
  formal QXP algebra only isolates the construction-level statement
  \(Q_{\mathrm{none}} \leq P_{\mathrm{none}}\), obtained when the repair
  preserves the fresh row block \(\widehat X_{\mathrm{none}}=X_{\mathrm{none}}\).
  The former generic ``RestrictSome'' monotone-total route would have required
  one to turn this into the source-facing inequality
  \(\widehat A_{\bot} \leq \widehat P_{\bot}\), by adding a comparison between
  the completed source residual and the ``Q''-layer fresh outcome.  The present
  formal development no longer uses that generic route, and the comparison is
  false without additional hypotheses; see
  \texttt{docs/reports/issue-1642-restrictsome-residual-domination-obstruction.md}.
  \lean{MIPStarRE.LDT.Preliminaries.triangleSub_right_subMeas_totalGap}
  \lean{MIPStarRE.LDT.Preliminaries.triangleSub_right_subMeas_total_le}
  \lean{MIPStarRE.LDT.MakingMeasurementsProjective.pa_eq_xHatA_adjoint_mul_xHatA}
  \lean{MIPStarRE.LDT.MakingMeasurementsProjective.qa_eq_pa_of_xHatA_eq_xa}
  \lean{MIPStarRE.LDT.MakingMeasurementsProjective.fresh_outcome_le_of_xHatA_eq_xa}
  \lean{MIPStarRE.LDT.MakingMeasurementsProjective.xHat_eq_x_of_x_mul_conjTranspose_eq_one_of_mixed_eq_gram}
  \lean{MIPStarRE.LDT.MakingMeasurementsProjective.xHatA_eq_xa_of_x_mul_conjTranspose_eq_one_of_mixed_eq_gram}
  \lean{MIPStarRE.LDT.MakingMeasurementsProjective.xHat_eq_x_of_x_mul_conjTranspose_eq_one}
  \lean{MIPStarRE.LDT.MakingMeasurementsProjective.xHatA_eq_xa_of_x_mul_conjTranspose_eq_one}
  \lean{MIPStarRE.LDT.MakingMeasurementsProjective.qxpProjSubMeas_total_eq_xHat_adjoint_mul_xHat}
  \lean{MIPStarRE.LDT.MakingMeasurementsProjective.qxpProjSubMeas_total_le_of_outcome_le}
  \lean{MIPStarRE.LDT.MakingMeasurementsProjective.QXPLayerTotalDomination}
  \lean{MIPStarRE.LDT.MakingMeasurementsProjective.QXPLayerTotalDomination.of_outcome_le}
  \lean{MIPStarRE.LDT.MakingMeasurementsProjective.QXPLayerTotalDomination.rightTensor_total_ev_le}
  \lean{MIPStarRE.LDT.SelfImprovement.final_fields_point_consistency_totalGap_natural}
  \lean{MIPStarRE.LDT.SelfImprovement.final_fields_point_consistency_totalGap}
  \lean{MIPStarRE.LDT.SelfImprovement.final_fields_point_consistency_totalGap_natural_of_total_difference}
  \lean{MIPStarRE.LDT.SelfImprovement.final_fields_point_consistency_totalGap_of_total_difference}
  \lean{MIPStarRE.LDT.SelfImprovement.final_fields_point_consistency_natural_of_total_le}
  \lean{MIPStarRE.LDT.SelfImprovement.final_fields_point_consistency_natural_of_total_expectation_le}
  \lean{MIPStarRE.LDT.SelfImprovement.final_fields_point_consistency_of_total_expectation_le_of_small_errors}
  \lean{MIPStarRE.LDT.opTensor_mono_right}
  \lean{MIPStarRE.LDT.rightTensor_mono}
  \lean{MIPStarRE.LDT.SelfImprovement.final_fields_point_consistency_natural_of_total_operator_le}
  \lean{MIPStarRE.LDT.SelfImprovement.final_fields_point_consistency_of_total_operator_le_of_small_errors}
  \lean{MIPStarRE.LDT.Preliminaries.subMeas_total_ev_gap_abs_le_sqrt_card_qSDD}
  \lean{MIPStarRE.LDT.SelfImprovement.final_fields_total_difference_le_sqrt_card_data}
  \lean{MIPStarRE.LDT.SelfImprovement.final_fields_point_consistency_totalGap_natural_of_data_processing}
  \lean{MIPStarRE.LDT.SelfImprovement.final_fields_point_consistency_totalGap_of_data_processing}
  Since \(A^u_{\mathrm{tot}}=I\) and postprocessing does not change the total
  operator, Lean also records the equivalent form in which this displacement is
  the single scalar difference
  \[
    \left|
      \langle \psi, I \ot H_{\mathrm{tot}} \psi\rangle
      -
      \langle \psi, I \ot \widehat H_{\mathrm{tot}} \psi\rangle
    \right|
  \]
  between the right-register totals of \(H\) and \(\widehat H\).  The
  outcomewise data-processing bound controls this scalar by applying the
  vector triangle inequality to the sum over \(a\in\F_q\); the formal estimate
  gives the additional term
  \[
    \sqrt{q\,\widehat{\zeta}_{\mathrm{dataprocess}}}.
  \]
  Thus the total-overlap term is no longer an independent hypothesis, although
  the numerical absorption of this alphabet-size loss is recorded separately.
  There is also a sharper formal route: if the right-register total of the
  projective family has no larger expectation than that of \(\widehat H\), then
  the total-overlap term can only improve the consistency defect.  Under this
  monotonicity hypothesis the transported error is again
  \(\widehat{\zeta}+\sqrt{\widehat{\zeta}_{\mathrm{dataprocess}}}\), and the
  standard final-stage numerical estimate absorbs it into \(\zeta\).

  \emph{Proof of \ref{item:self-improvement-projective-self}.}
  \lean{MIPStarRE.LDT.SelfImprovement.self_closeness_transport_through_orthonormalization}
  \lean{MIPStarRE.LDT.SelfImprovement.final_fields_self_closeness}
  \lean{MIPStarRE.LDT.SelfImprovement.final_fields_self_closeness_of_small_errors}
  Lemma~\ref{prop:two-notions-of-self-consistency} implies
  \[
    \widehat{H}_h \ot I \approx_{2\widehat{\zeta}} I \ot \widehat{H}_h.
  \]
  Combining this with~\eqref{eq:approx-between-H-with-and-without-hat} on both sides and summing the three transports with Lemma~\ref{prop:triangle-inequality-for-approx_delta} yields
  \[
    H_h \ot I \approx_{6 \widehat{\zeta} + 6\widehat{\zeta}_{\mathrm{ortho}}} I \ot H_h.
  \]

  \emph{Proof of \ref{item:self-improvement-projective-boundedness}.}
  \lean{MIPStarRE.LDT.SelfImprovement.projective_boundedness_gap_le_helper_boundedness_gap}
  \lean{MIPStarRE.LDT.SelfImprovement.final_fields_projective_residual_bound_natural}
  \lean{MIPStarRE.LDT.SelfImprovement.final_fields_projective_residual_bound}
  \lean{MIPStarRE.LDT.SelfImprovement.final_fields_projective_residual_bound_of_small_errors}
  \lean{MIPStarRE.LDT.SelfImprovement.final_fields_projective_residual_bound_of_helper_outputs}
  The boundedness quantity is
  \begin{align}
    &\bra{\psi} Z \ot (I - H) \ket{\psi} \notag\\
    ={}& \bra{\psi} Z \ot I \ket{\psi} - \sum_h \bra{\psi} Z \ot H_h \ket{\psi} \notag\\
    \leq{}& \bra{\psi} Z \ot I \ket{\psi} - \sum_h \bra{\psi} \E_{u} A^{u}_{h(u)} \ot H_h \ket{\psi} \notag\\
    ={}& \bra{\psi} Z \ot I \ket{\psi} -\E_{u} \sum_h \bra{\psi}  A^{u}_{h(u)} \ot H_h \ket{\psi} \notag\\
    ={}& \bra{\psi} Z \ot I \ket{\psi} -\E_{u} \sum_a \bra{\psi}  A^{u}_{a} \ot H_{[h(u)=a]} \ket{\psi}. \label{eq:projective-boundedness-expanded}
  \end{align}
  Proposition~\ref{prop:easy-approx-from-approx-delta} upgrades
  \lean{MIPStarRE.LDT.SelfImprovement.helper_boundedness_gap_transport_through_data_processing}
  \eqref{eq:approx-data-processed} to the scalar estimate
  \[
    \eqref{eq:projective-boundedness-expanded}
    \le \bra{\psi} Z \ot I \ket{\psi} -\E_{u} \sum_a \bra{\psi}  A^{u}_{a} \ot \widehat H_{[h(u)=a]} \ket{\psi} + \sqrt{\widehat{\zeta}_{\mathrm{dataprocess}}},
  \]
  and now \ref{item:self-improvement-boundedness} of Lemma~\ref{lem:self-improvement-helper} bounds this by $\widehat{\zeta} + \sqrt{\widehat{\zeta}_{\mathrm{dataprocess}}}$.
  The formal boundedness field
  \lean{MIPStarRE.LDT.SelfImprovement.final_fields_bounded}
  is the corresponding operator-monotonicity step: once the SDP dual witness
  satisfies $I\le Z$, the total mass of any left-placed submeasurement, and in
  particular of the final projective submeasurement, is bounded by $Z\ot I$ with
  zero boundedness error.

  Thus the four conclusions hold with errors
  \[
    2\widehat{\zeta} + 2\sqrt{\widehat{\zeta}_{\mathrm{ortho}}},
    \qquad
    \widehat{\zeta} + \sqrt{\widehat{\zeta}_{\mathrm{dataprocess}}},
    \qquad
    6 \widehat{\zeta} + 6\widehat{\zeta}_{\mathrm{ortho}},
    \qquad
    \widehat{\zeta} + \sqrt{\widehat{\zeta}_{\mathrm{dataprocess}}},
  \]
  respectively. The paper's exponent bookkeeping then shows the following
  bounds. The formal bookkeeping uses the three-term power sum
  $\eps^p + \delta^p + (d/q)^p$, its monotonicity on the unit interval, and the
  corresponding square-root estimate. The displayed Lean theorems record these
  arithmetic identities together with the completeness, self-closeness, and
  projective-residual absorptions
  $2 \widehat{\zeta} + 2\sqrt{\widehat{\zeta}_{\mathrm{ortho}}} \le \zeta$,
  $6 \widehat{\zeta} + 6\widehat{\zeta}_{\mathrm{ortho}} \le \zeta$ and
  $\widehat{\zeta} + \sqrt{\widehat{\zeta}_{\mathrm{dataprocess}}} \le \zeta$.
  The two threshold lemmas involving \(|\mathbb F_q|\) record the scalar part of
  the alphabet-size obstruction in the current final-fields point-consistency
  transport: the data-processing threshold contains \(8\widehat{\zeta}\), but
  the transported estimate still carries the factor \(\sqrt{|\mathbb F_q|}\).
  \lean{MIPStarRE.LDT.SelfImprovement.finalStagePowerSum}
  \lean{MIPStarRE.LDT.SelfImprovement.finalStagePowerSum_nonneg}
  \lean{MIPStarRE.LDT.SelfImprovement.finalStagePowerSum_le_of_exponent_ge}
  \lean{MIPStarRE.LDT.SelfImprovement.sqrt_finalStagePowerSum_le}
  \lean{MIPStarRE.LDT.SelfImprovement.selfImprovementHelperError_eq_finalStagePowerSum}
  \lean{MIPStarRE.LDT.SelfImprovement.selfImprovementError_eq_finalStagePowerSum}
  \lean{MIPStarRE.LDT.SelfImprovement.selfImprovementOrthogonalizationError_eq}
  \lean{MIPStarRE.LDT.SelfImprovement.selfImprovementDataProcessingError_eq}
  \lean{MIPStarRE.LDT.SelfImprovement.eight_selfImprovementHelperError_le_selfImprovementDataProcessingError}
  \lean{MIPStarRE.LDT.SelfImprovement.sqrt_card_mul_eight_helperError_le_sqrt_card_mul_dataProcessingError}
  \lean{MIPStarRE.LDT.SelfImprovement.final_fields_completeness_error_le_selfImprovementError}
  \lean{MIPStarRE.LDT.SelfImprovement.final_fields_self_closeness_error_le_selfImprovementError}
  \lean{MIPStarRE.LDT.SelfImprovement.final_fields_projective_residual_error_le_selfImprovementError}
  \[
    \widehat{\zeta}_{\mathrm{ortho}}
    \leq 400 m \cdot \Big(\eps^{1/8} + \delta^{1/8} + (d/q)^{1/8}\Big),
  \]
  \[
    \widehat{\zeta}_{\mathrm{dataprocess}}
    \leq 960m \cdot \Big(\eps^{1/16} + \delta^{1/16} + (d/q)^{1/16}\Big),
  \]
  and hence
  \[
    6 \widehat{\zeta} + 6\widehat{\zeta}_{\mathrm{ortho}} \le \zeta,
    \qquad
    \widehat{\zeta} + \sqrt{\widehat{\zeta}_{\mathrm{dataprocess}}} \le \zeta.
  \]
  Since $2\widehat{\zeta} + 2\sqrt{\widehat{\zeta}_{\mathrm{ortho}}} \le \widehat{\zeta}_{\mathrm{dataprocess}} \le \zeta$, all four bounds are at most~$\zeta$.
\end{proof}


%%FILE blueprint/src/chapter/ch08_commutativity.tex
\chapter{Commutativity}
\label{chap:commutativity}
\label{sec:g-comm}

\section{Commutativity of the point measurements}
\label{sec:commutativity-points}

\begin{definition}[Shared diagonal-line sample]\label{def:point-pair-shared-diagonal-line-distribution}
	\lean{MIPStarRE.LDT.CommutativityPoints.pointPairSharedDiagonalLineDistribution}
	\leanok
	Sample a uniformly random ordered point pair $(u,v)$ and a uniformly random
	parameter $t \in \F_q$. Package this data as the diagonal line whose direction
	is $v-u$ and whose parametrization visits $u$ at $t$ and $v$ at $t+1$.
\end{definition}

\begin{theorem}[Commutativity of the point measurements]\label{thm:commutativity-points}
	\lean{MIPStarRE.LDT.CommutativityPoints.commutativityPoints}
	\leanok
	\uses{def:good-strategy, def:approx_delta,
	      def:point-pair-shared-diagonal-line-distribution}
	Let $(\psi,A,B,L)$ be an $(\eps,\delta,\gamma)$-good symmetric strategy for the $(m,q,d)$ low individual degree test.
	On average over independent and uniformly random $u,v \sim \F_q^m$,
	\[
		(A_a^u A_b^v) \ot I \approx_{32\gamma m} (A_b^v A_a^u) \ot I.
	\]
\end{theorem}
\begin{proof}
	\leanok
	\uses{def:restricted-diagonal-lines-test, lem:good-strategy-characterization, prop:simeq-to-approx, prop:triangle-inequality-for-approx_delta}
	By Definition~\ref{def:restricted-diagonal-lines-test}, the strategy passes the $m$-restricted diagonal lines test with probability $1-\gamma m$. Hence
	\[
		A_a^u \otimes I \simeq_{\gamma m} I \otimes L^\ell_{[f(u)=a]}
	\]
	on average over uniformly random $u \sim \F_q^m$ and a uniformly random line $\ell$ in $\F_q^m$ containing $u$. By \ref{prop:simeq-to-approx},
	\begin{equation}\label{eq:point-diagonal-line-approx}
		A_a^u \otimes I \approx_{2\gamma m} I \otimes L^\ell_{[f(u)=a]}.
	\end{equation}
	Let $u$ and $v$ be independent uniformly random points in $\F_q^m$, and let $\ell$ be a uniformly random line containing both points. The marginal distribution on $(u,\ell)$ and on $(v,\ell)$ is the same as above, so
	\begin{align*}
		(A_a^u A_b^v) \ot I
		 & \approx_{2\gamma m} A_a^u \otimes L^\ell_{[f(v)=b]} \tag{by \ref{eq:point-diagonal-line-approx}}                           \\
		 & \approx_{2\gamma m} I \otimes \bigl(L^\ell_{[f(v)=b]} L^\ell_{[f(u)=a]}\bigr) \tag{by \ref{eq:point-diagonal-line-approx}} \\
		 & = I \otimes \bigl(L^\ell_{[f(u)=a]} L^\ell_{[f(v)=b]}\bigr) \tag{because~$L$ is projective}                                \\
		 & \approx_{2\gamma m} A_b^v \otimes L^\ell_{[f(u)=a]} \tag{by \ref{eq:point-diagonal-line-approx}}                           \\
		 & \approx_{2\gamma m} (A_b^v A_a^u) \ot I. \tag{by \ref{eq:point-diagonal-line-approx}}
	\end{align*}
	The theorem follows from \ref{prop:triangle-inequality-for-approx_delta}.
\end{proof}

\section{Commutativity of~$G$ after evaluation}
\label{sec:G-commutes-after-evaluation}

\begin{lemma}[Commutativity of~$G$ after evaluation]\label{lem:comm-data-processed-g}
	\lean{MIPStarRE.LDT.Commutativity.commDataProcessedG}
	\leanok
	\uses{def:good-strategy, def:polymeasurements, def:approx_delta, def:strong-self-consistency, clm:g-comm-stability, clm:g-comm-stability2}
	Let $(\psi,A,B,L)$ be an $(\eps,\delta,\gamma)$-good symmetric strategy for the $(m+1,q,d)$ low individual degree test.
	Let $\{G^x\} \in \polysub{m}{q}{d}$ be a collection of projective sub-measurements indexed by $x \in \F_q$ with the following properties:
	\begin{enumerate}
		\item (Consistency with~$A$)\label{item:data-processed-consistency}
		      On average over $(u,x) \sim \F_q^{m+1}$,
		      \[
			      A_a^{u,x} \ot I \simeq_\zeta I \ot G^x_{[g(u)=a]}.
		      \]
		\item (Strong self-consistency)\label{item:data-processed-self-consistency}
		      On average over $x \sim \F_q$,
		      \[
			      G_g^x \ot I \approx_\zeta I \ot G_g^x.
		      \]
		\item (Boundedness)\label{item:data-processed-boundedness}
		      There exists a positive-semidefinite matrix $Z^x$ for each $x \in \F_q$ such that
		      \[
			      \E_x \bra{\psi} (I-G^x) \ot Z^x \ket{\psi} \le \zeta
		      \]
		      and for each $x \in \F_q$ and $g \in \polyfunc{m}{q}{d}$,
		      \[
			      Z^x \ge \left(\E_u A_{g(u)}^{u,x}\right).
		      \]
	\end{enumerate}
	Let
	\[
		\nu = 48m\bigl(\gamma^{1/2}+\zeta^{1/2}\bigr).
	\]
	Then on average over independent and uniformly random $(u,x),(v,y) \sim \F_q^{m+1}$,
	\[
		G^x_{[g(u)=a]} G^y_{[h(v)=b]} \ot I \approx_\nu G^y_{[h(v)=b]} G^x_{[g(u)=a]} \ot I.
	\]
\end{lemma}
\begin{proof}
	\leanok
	\uses{prop:cons-sub-meas, prop:closeness-of-ip, prop:post-processing-preserves, prop:two-notions-of-self-consistency, prop:two-notions-of-self-consistency-after-evaluation, thm:commutativity-points}
	Write $G_a^{u,x}=G^x_{[g(u)=a]}$. Expanding the commutator square gives
	\begin{align}
		 & \E_{u,v,x,y} \sum_{a,b} \bra{\psi}
		\bigl(G_a^{u,x} G_b^{v,y} - G_b^{v,y} G_a^{u,x}\bigr)^\dagger
		\bigl(G_a^{u,x} G_b^{v,y} - G_b^{v,y} G_a^{u,x}\bigr)
		\ot I \ket{\psi}\nonumber             \\
		 & = 2 \E_{u,v,x,y} \sum_{a,b}
		\Bigl(\bra{\psi} G_b^{v,y} G_a^{u,x} G_b^{v,y} \ot I \ket{\psi}
		- \bra{\psi} G_a^{u,x} G_b^{v,y} G_a^{u,x} G_b^{v,y} \ot I \ket{\psi}\Bigr),
		\label{eq:gcom8}
	\end{align}
	using the projectivity of~$G$. To compare the two quartic terms, first note that
	\begin{equation}\label{eq:sum-of-gux}
		G^{u,x} = \sum_a G_a^{u,x} = \sum_a G^x_{[g(u)=a]} = G^x.
	\end{equation}
	The identity \eqref{eq:sum-of-gux} is exactly Proposition~\ref{prop:post-processing-preserves} for the evaluation map $g \mapsto g(u)$.
	By \ref{item:data-processed-consistency} and \ref{prop:cons-sub-meas},
	\begin{align}
		G_a^{u,x} \ot I
		 & \approx_{4\zeta} G^{u,x} \ot A_a^{u,x}\nonumber \\
		 & = G^x \ot A_a^{u,x}, \label{eq:add-an-a}
	\end{align}
	where the last equality is \ref{eq:sum-of-gux}. Applying \ref{prop:closeness-of-ip} once to the second term in \ref{eq:gcom8} gives
	\begin{align}
		                        & \E_{u,v,x,y} \sum_{a,b}
		\bra{\psi} G_a^{u,x} G_b^{v,y} G_a^{u,x} G_b^{v,y} \ot I \ket{\psi}\nonumber \\
		\approx_{2\sqrt{\zeta}} &
		\E_{u,v,x,y} \sum_{a,b}
		\bra{\psi} G_a^{u,x} G_b^{v,y} G_a^{u,x} G^y \ot A_b^{v,y} \ket{\psi}.
		\label{eq:apply-add-an-a-once}
	\end{align}
	The first boundedness-driven stability step is
	\begin{equation}
		\eqref{eq:apply-add-an-a-once}
		\approx_{\sqrt{\zeta}}
		\E_{u,v,x,y} \sum_{a,b}
		\bra{\psi} G_a^{u,x} G_b^{v,y} G_a^{u,x} \ot A_b^{v,y} \ket{\psi},
		\label{eq:gcom9}
	\end{equation}
	which is \ref{clm:g-comm-stability}. Applying \ref{eq:add-an-a} once more and then \ref{thm:commutativity-points}, we obtain
	\begin{align}
		\eqref{eq:gcom9}
		 & \approx_{2\sqrt{\zeta}}
		\E_{u,v,x,y} \sum_{a,b}
		\bra{\psi} G_a^{u,x} G_b^{v,y} G^x \ot A_a^{u,x} A_b^{v,y} \ket{\psi}\nonumber \\
		 & \approx_{6\sqrt{\gamma(m+1)}}
		\E_{u,v,x,y} \sum_{a,b}
		\bra{\psi} G_a^{u,x} G_b^{v,y} G^x \ot A_b^{v,y} A_a^{u,x} \ket{\psi}.
		\label{eq:commuting-answer-swapped}
	\end{align}
	The second boundedness-driven stability step is
	\begin{equation}
		\eqref{eq:commuting-answer-swapped}
		\approx_{\sqrt{\zeta}+6\sqrt{\gamma(m+1)}}
		\E_{u,v,x,y} \sum_{a,b}
		\bra{\psi} G_a^{u,x} G_b^{v,y} \ot A_a^{u,x} A_b^{v,y} \ket{\psi},
		\label{eq:gcom10}
	\end{equation}
	which is \ref{clm:g-comm-stability2}. Using \ref{eq:add-an-a} twice more and \ref{prop:closeness-of-ip},
	\begin{align}
		\eqref{eq:gcom10}
		 & \approx_{2\sqrt{\zeta}}
		\E_{u,v,x,y} \sum_{a,b}
		\bra{\psi} G_a^{u,x} G_b^{v,y} \ot A_b^{v,y} \ket{\psi}\nonumber \\
		 & \approx_{2\sqrt{\zeta}}
		\E_{u,v,x,y} \sum_{a,b}
		\bra{\psi} G_a^{u,x} G_b^{v,y} G_b^{v,y} \ot I \ket{\psi}.
		\label{eq:gonna-cite-this-in-just-a-bit}
	\end{align}
	Next, \ref{item:data-processed-self-consistency} and the projectivity of~$G$ imply via \ref{prop:two-notions-of-self-consistency} and \ref{prop:two-notions-of-self-consistency-after-evaluation} that
	\begin{equation}\label{eq:new-fact-that-i-derived}
		G_a^{u,x} \ot I \approx_\zeta I \ot G_a^{u,x}.
	\end{equation}
	Applying \ref{eq:new-fact-that-i-derived} twice with \ref{prop:closeness-of-ip} turns \ref{eq:gonna-cite-this-in-just-a-bit} into
	\[
		\E_{u,v,x,y} \sum_{a,b}
		\bra{\psi} G_b^{v,y} G_a^{u,x} G_b^{v,y} \ot I \ket{\psi},
	\]
	which is the first quartic term from \ref{eq:gcom8}. Collecting the ten approximation steps appearing between \ref{eq:gcom8} and \ref{eq:new-fact-that-i-derived} yields the bound
	\[
		48m\bigl(\gamma^{1/2}+\zeta^{1/2}\bigr),
	\]
	so the commutator norm is at most~$\nu$.
\end{proof}

\begin{lemma}\label{clm:g-comm-stability}
	\lean{MIPStarRE.LDT.IdxPolyFamily.SliceBoundednessInput.storedBoundedResidualBound, MIPStarRE.LDT.IdxPolyFamily.SliceBoundednessInput.averagedPoint_le_witness, MIPStarRE.LDT.Commutativity.gCommStability_overlap, MIPStarRE.LDT.Commutativity.gCommStability_scalar}
	\leanok
	\uses{def:approx_delta, def:strong-self-consistency}
	This claim is made under the same standing hypotheses used in the
	commutativity-of-$G$ theorem above.  The linked
	\texttt{SliceBoundednessInput} declarations are precisely the two displayed
	parts of item~\ref{item:data-processed-boundedness}: the averaged
	\((I-G^y)\otimes Z^y\) residual bound and the domination
	\(Z^y\ge \E_u A^{u,y}_{g(u)}\).
	\begin{align}
		                       & \E_{u,v,x,y} \sum_{a,b}
		\bra{\psi} G_a^{u,x} G_b^{v,y} G_a^{u,x} G^y \ot A_b^{v,y} \ket{\psi}
		\label{eq:g-comm-stab1}                          \\
		\approx_{\sqrt{\zeta}} &
		\E_{u,v,x,y} \sum_{a,b}
		\bra{\psi} G_a^{u,x} G_b^{v,y} G_a^{u,x} \ot A_b^{v,y} \ket{\psi}.
		\nonumber
	\end{align}
\end{lemma}
\begin{proof}
	\leanok
	For each $y \in \F_q$ and $g \in \polyfunc{m}{q}{d}$, define
	\[
		R_g^y = \E_{u,x} \sum_a G_a^{u,x} G_g^y G_a^{u,x}.
	\]
	Since $G$ is a sub-measurement, $R^y=\{R_g^y\}$ is also a sub-measurement. The difference between the two sides of the lemma is
	\begin{align}
		 & \E_{u,v,x,y} \sum_{a,b}
		\bra{\psi} G_a^{u,x} G_b^{v,y} G_a^{u,x} (I-G^y) \ot A_b^{v,y} \ket{\psi}\nonumber \\
		 & = \E_{v,y} \sum_g
		\bra{\psi} R_g^y (I-G^y) \ot A_{g(v)}^{v,y} \ket{\psi}.
		\label{eq:bound-this-right-now!}
	\end{align}
	Apply Cauchy--Schwarz to \ref{eq:bound-this-right-now!}. The first factor is at most~$1$ because $R^y$ is a sub-measurement. For the second factor, average over~$v$ and use \ref{item:data-processed-boundedness} to replace $\E_v A_{g(v)}^{v,y}$ by $Z^y$; the resulting quantity is at most
	\[
		\E_y \bra{\psi} (I-G^y) \ot Z^y \ket{\psi} \le \zeta.
	\]
	Hence the difference has magnitude at most $\sqrt{\zeta}$.
\end{proof}

\begin{lemma}\label{clm:g-comm-stability2}
	\lean{MIPStarRE.LDT.IdxPolyFamily.SliceBoundednessInput.storedBoundedResidualBound, MIPStarRE.LDT.IdxPolyFamily.SliceBoundednessInput.averagedPoint_le_witness, MIPStarRE.LDT.Commutativity.gCommStabilityTwo_overlap, MIPStarRE.LDT.Commutativity.gCommStabilityTwo_scalar}
	\leanok
	\uses{def:approx_delta, def:strong-self-consistency}
	This claim is made under the same standing hypotheses used in the
	commutativity-of-$G$ theorem above, in particular the boundedness item
	\ref{item:data-processed-boundedness}.  The linked
	\texttt{SliceBoundednessInput} declarations expose those boundedness
	hypotheses in Lean; they are not auxiliary assumptions added to the paper
	statement.
	\begin{align*}
		                                           & \E_{u,v,x,y} \sum_{a,b}
		\bra{\psi} G_a^{u,x} G_b^{v,y} G^x \ot A_a^{u,x} A_b^{v,y} \ket{\psi} \\
		\approx_{\sqrt{\zeta}+6\sqrt{\gamma(m+1)}} &
		\E_{u,v,x,y} \sum_{a,b}
		\bra{\psi} G_a^{u,x} G_b^{v,y} \ot A_a^{u,x} A_b^{v,y} \ket{\psi}.
	\end{align*}
\end{lemma}
\begin{proof}
	\leanok
	\uses{thm:commutativity-points}
	The difference between the two sides is
	\begin{align}
		                              & \E_{u,v,x,y} \sum_{a,b}
		\bra{\psi} G_a^{u,x} G_b^{v,y} (I-G^x) \ot A_a^{u,x} A_b^{v,y} \ket{\psi}\nonumber \\
		\approx_{6\sqrt{\gamma(m+1)}} &
		\E_{u,v,x,y} \sum_{a,b}
		\bra{\psi} G_a^{u,x} G_b^{v,y} (I-G^x) \ot A_b^{v,y} A_a^{u,x} \ket{\psi},
		\label{eq:just-got-commuted}
	\end{align}
	where \ref{thm:commutativity-points} commutes the point measurements on the right register. Expanding $G_a^{u,x}=\sum_{g:g(u)=a} G_g^x$, we rewrite
	\begin{align}
		\eqref{eq:just-got-commuted}
		 & = \E_{u,v,x,y} \sum_{g,b}
		\bra{\psi} G_g^x G_b^{v,y} (I-G^x) \ot A_b^{v,y} A_{g(u)}^{u,x} \ket{\psi}.
		\label{eq:g-comm-stab7}
	\end{align}
	Apply Cauchy--Schwarz to \ref{eq:g-comm-stab7}. The first factor is at most~$1$ because both $G$ and~$A$ are sub-measurements. In the second factor, average over~$u$ and use \ref{item:data-processed-boundedness} to replace $\E_u A_{g(u)}^{u,x}$ by $Z^x$; then sum over~$g$ and~$b$ using the sub-measurement property. This leaves
	\[
		\E_x \bra{\psi} (I-G^x) \ot Z^x \ket{\psi} \le \zeta.
	\]
	Hence \ref{eq:g-comm-stab7} has magnitude at most $\sqrt{\zeta}$, and the total loss is $\sqrt{\zeta}+6\sqrt{\gamma(m+1)}$.
\end{proof}

\section{Commutativity of~$G$}
\label{sec:G-commutes}

\begin{lemma}[Normalization condition for a sandwiched family]\label{lem:normalization-condition}
	\lean{MIPStarRE.LDT.Commutativity.normalizationCondition}
	\leanok
	\uses{def:submeasurement}
	Let $P=\{P_a\}$ be a sub-measurement and $Q=\{Q_b\}$ be a projective sub-measurement.
	Define $C_{a,b}=Q_b P_a Q_b$.
	Then
	\[
		\sum_a \left(\sum_b C_{a,b}\right)^\dagger \left(\sum_b C_{a,b}\right)
		=
		\sum_a \left(\sum_b C_{a,b}\right)\left(\sum_b C_{a,b}\right)^\dagger
		\le I.
	\]
\end{lemma}
\begin{proof}
	\leanok
	Expand the square, use the orthogonality $Q_b Q_{b'}=0$ for $b \neq b'$, and then sum over the sub-measurement~$P$.
\end{proof}

\begin{theorem}[Commutativity of~$G$]\label{thm:com-main}
	\lean{MIPStarRE.LDT.Commutativity.comMain}
	\leanok
	\uses{def:good-strategy, def:polymeasurements, def:approx_delta, def:strong-self-consistency}
	Let $(\psi,A,B,L)$ be an $(\eps,\delta,\gamma)$-good symmetric strategy for the $(m+1,q,d)$ low individual degree test.
	Let $\{G^x\}_{x \in \F_q}$ denote a set of projective sub-measurements in $\polysub{m}{q}{d}$ with the following properties:
	\begin{enumerate}
		\item (Consistency with~$A$)\label{item:commuting-consistency}
		      On average over $(u,x) \sim \F_q^{m+1}$,
		      \[
			      A_a^{u,x} \ot I \simeq_\zeta I \ot G^x_{[g(u)=a]}.
		      \]
		\item (Strong self-consistency)\label{item:commuting-self-consistency}
		      On average over $x \sim \F_q$,
		      \[
			      G_g^x \ot I \approx_\zeta I \ot G_g^x.
		      \]
		\item (Boundedness)\label{item:commuting-boundedness}
		      There exists a positive-semidefinite matrix $Z^x$ for each $x \in \F_q$ such that
		      \[
			      \E_x \bra{\psi} (I-G^x) \ot Z^x \ket{\psi} \le \zeta
		      \]
		      and for each $x \in \F_q$ and $g \in \polyfunc{m}{q}{d}$,
		      \[
			      Z^x \ge \left(\E_u A_{g(u)}^{u,x}\right).
		      \]
	\end{enumerate}
	Let
	\[
		\nu = 30m\left(\gamma^{1/4}+\zeta^{1/4}+(d/q)^{1/4}\right).
	\]
	Then on average over independent and uniformly random $(u,x),(v,y) \sim \F_q^{m+1}$,
	\[
		G_g^x G_h^y \ot I \approx_\nu G_h^y G_g^x \ot I.
	\]
\end{theorem}
\begin{proof}
	\leanok
	\uses{lem:comm-data-processed-g, lem:normalization-condition, prop:switch-sandwich, prop:closeness-of-ip, prop:triangle-inequality-for-vectors-squared, lem:schwartz-zippel-individual}
	If at least one of $\gamma$, $\zeta$, or $d/q$ is at least~$1$, then $\nu \ge 30$ and the estimate is trivial from \ref{prop:triangle-inequality-for-vectors-squared}. We therefore assume $\gamma,\zeta,d/q \le 1$.

	Expanding the commutator square gives
	\begin{align}
		 & \E_{x,y} \sum_{g,h}
		\bra{\psi}
		\bigl(G_g^x G_h^y - G_h^y G_g^x\bigr)^\dagger
		\bigl(G_g^x G_h^y - G_h^y G_g^x\bigr)
		\ot I \ket{\psi}\nonumber  \\
		 & = 2 \E_{x,y} \sum_{g,h}
		\bra{\psi} G_g^x G_h^y G_g^x \ot I \ket{\psi}
		- 2 \E_{x,y} \sum_{g,h}
		\bra{\psi} G_g^x G_h^y G_g^x G_h^y \ot I \ket{\psi}.
		\label{eq:gcomterms}
	\end{align}
	The first term is close to $\bra{\psi} G \ot G \ket{\psi}$, where $G=\E_x G^x$, by \ref{prop:switch-sandwich} and \ref{item:commuting-self-consistency}. For the second term, \ref{prop:closeness-of-ip} and \ref{lem:normalization-condition} yield
	\begin{equation}
		\E_{x,y} \sum_{g,h}
		\bra{\psi} G_g^x G_h^y G_g^x G_h^y \ot I \ket{\psi}
		\approx_{\sqrt{\zeta}}
		\E_{x,y} \sum_{g,h}
		\bra{\psi} G_h^y G_g^x G_h^y \ot G_g^x \ket{\psi}.
		\label{eq:gcom4}
	\end{equation}
	The first Schwartz--Zippel reduction evaluates the left factor at a random point:
	\begin{align}
		\eqref{eq:gcom4}
		 & = \E_{x,y} \sum_{g,h}
		\bra{\psi} G_h^y G_g^x G_h^y \ot G_g^x \ket{\psi}\nonumber \\
		 & \approx_{\frac{dm}{q}}
		\E_{u,x,y} \sum_{a,h}
		\bra{\psi} G_h^y G^x_{[g(u)=a]} G_h^y \ot G^x_{[g(u)=a]} \ket{\psi}.
		\label{eq:evaluate-gcom-at-points}
	\end{align}
	Indeed, the difference is
	\begin{equation}\label{eq:gcom4-diff}
		\E_{u,x,y} \sum_{g \neq g',h}
		\bone[g(u)=g'(u)]
		\bra{\psi} G_h^y G_g^x G_h^y \ot G_{g'}^x \ket{\psi},
	\end{equation}
	and \ref{lem:schwartz-zippel-individual} bounds this by $dm/q$.

	Next, apply \ref{prop:closeness-of-ip} twice:
	\begin{align}
		\eqref{eq:evaluate-gcom-at-points}
		 & = \E_{u,x,y} \sum_{a,h}
		\bra{\psi} G_h^y G^x_{[g(u)=a]} G_h^y \ot G^x_{[g(u)=a]} \ket{\psi}\nonumber   \\
		 & \approx_{\sqrt{\zeta}}
		\E_{u,x,y} \sum_{a,h}
		\bra{\psi} G^x_{[g(u)=a]} G_h^y G^x_{[g(u)=a]} G_h^y \ot I \ket{\psi}\nonumber \\
		 & \approx_{\sqrt{\zeta}}
		\E_{u,x,y} \sum_{a,h}
		\bra{\psi} G^x_{[g(u)=a]} G_h^y G^x_{[g(u)=a]} \ot G_h^y \ket{\psi}.
		\label{eq:don't-understand-the-numbering-system}
	\end{align}
	A second Schwartz--Zippel reduction evaluates the right factor at an independent point:
	\begin{align}
		\eqref{eq:don't-understand-the-numbering-system}
		 & = \E_{u,x,y} \sum_{a,h}
		\bra{\psi} G^x_{[g(u)=a]} G_h^y G^x_{[g(u)=a]} \ot G_h^y \ket{\psi}\nonumber \\
		 & \approx_{\frac{dm}{q}}
		\E_{u,v,x,y} \sum_{a,b}
		\bra{\psi}
		G^x_{[g(u)=a]} G^y_{[h(v)=b]} G^x_{[g(u)=a]}
		\ot G^y_{[h(v)=b]} \ket{\psi}.
		\label{eq:evaluate-gcom-at-points-part-dos}
	\end{align}
	The difference is
	\begin{equation}
		% Paper label: eq:eq:don't-understand-the-numbering-system-diff
		\label{eq:numbering-system-diff}
		\E_{u,v,x,y} \sum_{a,h \neq h'}
		\bone[h(v)=h'(v)]
		\bra{\psi}
		G^x_{[g(u)=a]} G_h^y G^x_{[g(u)=a]} \ot G_{h'}^y \ket{\psi},
	\end{equation}
	and the same Schwartz--Zippel estimate again gives the loss $dm/q$.

	Now set
	\[
		\nu_{\mathrm{evaluation}} = 48m\bigl(\gamma^{1/2}+\zeta^{1/2}\bigr),
	\]
	the error from \ref{lem:comm-data-processed-g}. Apply \ref{lem:comm-data-processed-g} to commute the evaluated families in \ref{eq:evaluate-gcom-at-points-part-dos}; after one more use of \ref{prop:closeness-of-ip} with \ref{item:commuting-self-consistency}, the resulting quantity is $\bra{\psi} G \ot G \ket{\psi}$. Thus the second term in \ref{eq:gcomterms} is also close to $\bra{\psi} G \ot G \ket{\psi}$.

	Combining the first-term estimate, \ref{eq:gcom4}, the two Schwartz--Zippel losses, the two intermediate $\sqrt{\zeta}$ losses, and the evaluated commutation error $\nu_{\mathrm{evaluation}}^{1/2}$ gives
	\[
		30m\left(\gamma^{1/4}+\zeta^{1/4}+(d/q)^{1/4}\right),
	\]
	where the constant $30m$ accounts for the $10m$ Schwartz--Zippel evaluation losses plus $20m$ from the $\approx$-chain in the first and second terms. This proves \ref{thm:com-main}.
\end{proof}


%%FILE blueprint/src/chapter/ch09_pasting.tex
\chapter{Pasting}
\label{chap:pasting}
\label{sec:ld-pasting}

\begin{theorem}[Pasting]\label{thm:ld-pasting}
	\lean{MIPStarRE.LDT.Pasting.ldPasting}
	\leanok
	\uses{def:good-strategy, def:polymeasurements, def:simeq, def:approx_delta}
	Let $(\psi,A,B,L)$ be an $(\eps,\delta,\gamma)$-good symmetric strategy for the $(m+1,q,d)$ low individual degree test. Let $\{G^x\}_{x \in \F_q}$ be projective submeasurements in $\polysub{m}{q}{d}$ with the following properties:
	\begin{enumerate}
		\item (Completeness): \label{item:ld-pasting-completeness}
		      If $G = \E_x \sum_g G_g^x$, then
		      \[
			      \bra{\psi} G \ot I \ket{\psi} \ge 1-\kappa.
		      \]
		\item (Consistency with $A$): \label{item:ld-pasting-consistency}
		      On average over $(u,x) \sim \F_q^{m+1}$,
		      \[
			      A_a^{u,x} \ot I \simeq_\zeta I \ot G^x_{[g(u)=a]}.
		      \]
		\item (Strong self-consistency): \label{item:ld-pasting-self-consistency}
		      On average over $x \sim \F_q$,
		      \[
			      G_g^x \ot I \approx_\zeta I \ot G_g^x.
		      \]
		\item (Boundedness): \label{item:ld-pasting-boundedness}
		      There exists a positive-semidefinite matrix $Z^x$ for each $x \in \F_q$ such that
		      \[
			      \E_x \bra{\psi} (I-G^x)\ot Z^x \ket{\psi} \le \zeta
		      \]
		      and for each $x \in \F_q$ and $g \in \polyfunc{m}{q}{d}$,
		      \[
			      Z^x \ge \E_u A^{u,x}_{g(u)}.
		      \]
	\end{enumerate}
	Let $k \ge 400md$ be an integer, and set
	\[
		\nu = 100k^2m\left(\eps^{1/32}+\delta^{1/32}+\gamma^{1/32}+\zeta^{1/32}+(d/q)^{1/32}\right),
	\]
	\[
		\sigma = \kappa\left(1+\frac{1}{100m}\right) + 2\nu + e^{-k/(80000m^2)}.
	\]
	Then there exists a pasted measurement $H \in \polymeas{m+1}{q}{d}$ with the following property.
	\begin{enumerate}
		\item (Consistency with $A$): \label{item:ld-pasting-N-consistency}
		      On average over $u \sim \F_q^{m+1}$,
		      \[
			      A_a^u \ot I \simeq_\sigma I \ot H_{[h(u)=a]}.
		      \]
	\end{enumerate}
\end{theorem}
\begin{proof}
	\leanok
	\uses{lem:ld-pasting-sub-measurement}
	The argument below is the paper proof, and the unrestricted Lean theorem now
	formalizes it directly.
	Let $H$ be the submeasurement from Lemma~\ref{lem:ld-pasting-sub-measurement}, and fix an arbitrary polynomial $h^* \in \polyfunc{m+1}{q}{d}$. Define a measurement $H_{\mathrm{meas}} \in \polymeas{m+1}{q}{d}$ by
	\[
		(H_{\mathrm{meas}})_h
		=
		\begin{cases}
			H_{h^*} + (I-H) & \text{if } h=h^*, \\
			H_h             & \text{otherwise.}
		\end{cases}
	\]
	Then $\sum_h (H_{\mathrm{meas}})_h = H + (I-H) = I$. Moreover,
	\begin{align*}
		      & \E_u \sum_{a \ne b} \bra{\psi} A_a^u \ot (H_{\mathrm{meas}})_{[h(u)=b]} \ket{\psi}                            \\
		={}   & \E_u \sum_a \sum_{h:h(u)\ne a} \bra{\psi} A_a^u \ot (H_{\mathrm{meas}})_h \ket{\psi}                          \\
		={}   & \E_u \sum_a \sum_{h:h(u)\ne a} \bra{\psi} A_a^u \ot H_h \ket{\psi}
		+ \E_u \sum_{a:h^*(u)\ne a} \bra{\psi} A_a^u \ot (I-H) \ket{\psi}                                                     \\
		\le{} & \E_u \sum_a \sum_{h:h(u)\ne a} \bra{\psi} A_a^u \ot H_h \ket{\psi}
		+ \E_u \bra{\psi} I \ot (I-H) \ket{\psi}                                                                              \\
		\le{} & \nu + \left(\kappa\left(1+\frac{1}{100m}\right) + \nu + e^{-k/(80000m^2)}\right)
		\tag{by \ref{item:ld-pasting-N-consistency-sub-measurement} and \ref{item:ld-pasting-N-completeness-sub-measurement}} \\
		={}   & \sigma.
	\end{align*}
	Thus $A_a^u \ot I \simeq_\sigma I \ot (H_{\mathrm{meas}})_{[h(u)=a]}$.
\end{proof}

\begin{remark}[Restricted Lean form of the pasting theorem]\label{rem:ld-pasting-restricted-lean-form}
	The Lean declaration linked here is the restricted nontrivial-regime form of
	Theorem~\ref{thm:ld-pasting}.  In addition to the hypotheses displayed in the
	source theorem, it assumes
	\[
		\gamma \le 1,\qquad \zeta \le 1,\qquad d \le q,\qquad 0<d,\qquad 1\le k.
	\]
	The paper theorem is stated in
	\texttt{references/ldt-paper/ld-pasting.tex}, lines 12--50.  Lines 52--55
	explain that the proof may restrict to
	\(\eps,\delta,\gamma,\zeta,d/q\le 1\), since the complementary cases are
	trivial.  The unrestricted source statement is linked from
	Theorem~\ref{thm:ld-pasting}.  All complementary branches, including the
	degree-zero case, are now proved in Lean.
\end{remark}

The bound is trivial when at least one of $\eps$, $\delta$, $\gamma$, $\zeta$, or $d/q$ is at least $1$, since then $\nu \ge 1$. In the remainder of the chapter we therefore assume
\[
	\eps,\delta,\gamma,\zeta,d/q \le 1.
\]

\begin{definition}[Nontrivial-regime pasting context]\label{def:ld-pasting-context}
	\lean{MIPStarRE.LDT.Pasting.LdPastingNontrivialContext}
	\leanok
	\uses{def:good-strategy, def:polymeasurements, def:simeq,
		def:approx_delta}
	The Lean context linked here is the nontrivial-regime context used after the
	reduction explained above.  It consists of the data and hypotheses from the
	pasting theorem at the beginning of this chapter: an $(\eps,\delta,\gamma)$-good symmetric strategy
	$(\psi,A,B,L)$ for the $(m+1,q,d)$ low individual degree test, a family of
	projective submeasurements $\{G^x\}_{x \in \F_q}$ in $\polysub{m}{q}{d}$
	satisfying \ref{item:ld-pasting-completeness}--\ref{item:ld-pasting-boundedness},
	an integer $k \ge 400md$, and the quantities $\nu$ and $\sigma$ defined there.
	It also records the nontrivial-regime assumptions
	\[
		\gamma \le 1,\qquad \zeta \le 1,\qquad d \le q,\qquad 0<d,\qquad 1\le k.
	\]
	These extra fields are not assumptions of the unrestricted source theorem;
	they belong to the restricted proof stage.  The unrestricted paper-facing Lean
	theorem remains \texttt{MIPStarRE.LDT.Pasting.ldPasting}.
\end{definition}

\subsection{From Measurements to Submeasurements}

In the remainder of the chapter we only use the axis-parallel lines test in the
last coordinate direction. We therefore write $B_f^u$ for the line measurement
on the line $\{(u,x)\mid x\in\F_q\}$ and first record the consistency relations
that follow formally from the nontrivial-regime pasting context of
Definition~\ref{def:ld-pasting-context}.

Conditioning the low individual degree test on the last coordinate direction gives, on average over $(u,x) \sim \F_q^{m+1}$,
\[
	A_a^{u,x} \ot I \simeq_{(m+1)\eps} I \ot B^u_{[f(x)=a]}.
\]
Hence \ref{prop:simeq-to-approx} and $(m+1)\le 2m$ imply
\[
	A_a^{u,x} \ot I \approx_{4m\eps} I \ot B^u_{[f(x)=a]}.
\]
Using also the $(\eps,\delta,\gamma)$-goodness assumption and \ref{prop:simeq-to-approx},
\[
	I \ot A_a^{u,x} \approx_{2\delta} A_a^{u,x} \ot I \approx_{4m\eps} I \ot B^u_{[f(x)=a]}.
\]
Therefore \ref{prop:triangle-inequality-for-approx_delta} yields
\begin{equation}\label{eq:ld-abcon}
	I \ot A_a^{u,x} \approx_{8m\eps+4\delta} I \ot B^u_{[f(x)=a]}.
\end{equation}
Applying \ref{prop:triangle-sub} to \ref{item:ld-pasting-consistency} and \ref{eq:ld-abcon} gives
\begin{equation}\label{eq:ld-gbcon}
	G^x_{[g(u)=a]} \ot I \simeq_{\nu_1} I \ot B^u_{[f(x)=a]},
\end{equation}
where
\begin{equation}\label{eq:ld-nu1-def}
	\nu_1 = \zeta + \sqrt{8m\eps+4\delta}.
\end{equation}
Finally, \ref{thm:com-main} gives
\begin{equation}\label{eq:complete-slice-commutativity}
	G_g^x G_h^y \ot I \approx_{\nu_{\mathrm{com}}} G_h^y G_g^x \ot I,
\end{equation}
where
\[
	\nu_{\mathrm{com}} = 30m\left(\gamma^{1/4}+\zeta^{1/4}+(d/q)^{1/4}\right).
\]

\begin{lemma}[Slice measurements are consistent with line answers]\label{lem:ld-gbcon}
	\lean{MIPStarRE.LDT.Pasting.ldGbcon}
	\leanok
	\uses{def:ld-pasting-context, prop:simeq-to-approx,
	      prop:triangle-inequality-for-approx_delta, prop:triangle-sub}
	In the nontrivial-regime pasting context of Definition~\ref{def:ld-pasting-context}, set
	\[
		\nu_1 = \zeta + \sqrt{8m\eps+4\delta}.
	\]
	Then, on average over $(u,x) \sim \F_q^{m+1}$,
	\[
		G^x_{[g(u)=a]} \ot I \simeq_{\nu_1} I \ot B^u_{[f(x)=a]}.
	\]
\end{lemma}
\begin{proof}
	\leanok
	This is exactly \ref{eq:ld-gbcon}, with \ref{eq:ld-nu1-def}.
\end{proof}

\begin{lemma}[Lean auxiliary form of the vertical-line SDD estimate]\label{lem:lean-point-vertical-line-sdd-lifted-family}
	\lean{MIPStarRE.LDT.Pasting.pointVerticalLineSdd_liftedVerticalLine_of_axis_self}
	\leanok
	The formalization records the point-to-vertical-line state-dependent-distance
	estimate with the vertical-line side written as the submeasurement
	family \texttt{liftedVerticalLineAnswerFamily}. The error term is the same
	\(8m\eps+4\delta\) as in the complete measurement-valued construction.
\end{lemma}
\begin{proof}
	\leanok
	This is not a separate paper assertion. It is the representation of the same
	estimate obtained by replacing the complete vertical-line measurement family
	with its underlying submeasurement family, which is the form needed in
	the degree-zero branch.
\end{proof}

\begin{lemma}[Pasted submeasurement]\label{lem:ld-pasting-sub-measurement}
	\lean{MIPStarRE.LDT.Pasting.ldPastingSubMeas}
	\leanok
	\uses{def:ld-pasting-context, def:polymeasurements, def:simeq,
	      def:pasted-measurement}
	In the nontrivial-regime pasting context of Definition~\ref{def:ld-pasting-context},
	there exists a submeasurement $H \in \polysub{m+1}{q}{d}$ with the following
	properties.
	\begin{enumerate}
		\item (Consistency with $A$): \label{item:ld-pasting-N-consistency-sub-measurement}
		      On average over $u \sim \F_q^{m+1}$,
		      \[
			      A_a^u \ot I \simeq_\nu I \ot H_{[h(u)=a]}.
		      \]
		\item (Completeness): \label{item:ld-pasting-N-completeness-sub-measurement}
		      If $H = \sum_h H_h$, then
		      \[
			      \bra{\psi} H \ot I \ket{\psi}
			      \ge
			      1 - \kappa\left(1+\frac{1}{100m}\right) - \nu - e^{-k/(80000m^2)}.
		      \]
	\end{enumerate}
\end{lemma}
\begin{proof}
	\leanok
	\uses{def:pasted-measurement, lem:h-b-consistency, cor:ld-pasting-N-completeness, prop:triangle-sub}
	Let $H$ be the pasted submeasurement constructed in Definition~\ref{def:pasted-measurement}. Lemma~\ref{lem:h-b-consistency} implies
	\begin{equation}\label{eq:h-b-consistency-at-a-point}
		H_{[h(u,x)=a]} \ot I \simeq_{\nu_6} I \ot B^u_{[f(x)=a]}.
	\end{equation}
	Applying \ref{prop:triangle-sub} to \ref{eq:h-b-consistency-at-a-point} and \ref{eq:ld-abcon} gives
	\[
		H_{[h(u,x)=a]} \ot I \simeq_{\nu_6+\sqrt{8m\eps+4\delta}} I \ot A_a^{u,x}.
	\]
	Since $\sqrt{8m\eps+4\delta}\le 3m(\eps^{1/32}+\delta^{1/32})$ for $\eps,\delta\le 1$, the total error is at most
	\[
		\sqrt{8m\eps+4\delta}+\nu_6 \le 47k^2m\left(\eps^{1/32}+\delta^{1/32}+\gamma^{1/32}+\zeta^{1/32}+(d/q)^{1/32}\right)\le \nu,
	\]
	which proves \ref{item:ld-pasting-N-consistency-sub-measurement}. The completeness claim is exactly Corollary~\ref{cor:ld-pasting-N-completeness}.
\end{proof}

\subsection{The Pasted Submeasurement}

\begin{definition}[Distinct tuples]\label{def:distinct-tuples}
	\lean{MIPStarRE.LDT.Pasting.distinctTuples, MIPStarRE.LDT.Pasting.distinctTupleDistribution}
	\leanok
	For $k \ge 1$, let $\mathsf{Distinct}_k \subseteq \F_q^k$ be the set of tuples $(x_1,\dots,x_k)$ with pairwise distinct coordinates.
\end{definition}

\begin{proposition}[Distinct tuples versus independent tuples]\label{prop:ld-dnoteq}
	\lean{MIPStarRE.LDT.Pasting.ldDnoteq}
	\leanok
	\uses{def:distinct-tuples}
	If $x=(x_1,\dots,x_k)$ is uniformly random in $\F_q^k$ and $y=(y_1,\dots,y_k)$ is uniformly random in $\mathsf{Distinct}_k$, then
	\[
		d_{\mathrm{TV}}(x,y) \le \frac{k^2}{q}.
	\]
\end{proposition}
\begin{proof}
	% Lean resolves these declarations as root `PMF.*` names; the namespace is not
	% below `MIPStarRE.LDT`, so there is no longer `MIPStarRE.*` declaration name.
	\lean{PMF.sum_toReal_eq_one}
	\lean{PMF.totalVariationDistance}
	\lean{PMF.totalVariationDistance_eq_sum_max_sub}
	\lean{PMF.totalVariationDistance_uniformOfFintype_uniformOfFinset_eq}
	\lean{MIPStarRE.LDT.totalVariationDistance_uniformDistribution_uniformOnFinset_eq}
	\leanok
	The total variation distance is exactly the probability mass of the collision
	event for a uniformly random tuple, because the distinct-tuple distribution is
	the uniform distribution conditioned on landing in $\mathsf{Distinct}_k$.
	A uniformly random tuple lies outside $\mathsf{Distinct}_k$ only if two
	coordinates collide, so a union bound gives
	\[
		\Pr[x \notin \mathsf{Distinct}_k] \le \sum_{i=2}^k \frac{i-1}{q} = \frac{k(k-1)}{2q}\le \frac{k^2}{q},
	\]
	and this bounds the total variation distance.
\end{proof}

We first record the natural construction obtained by sampling $d+1$ slices and interpolating. It motivates the eventual argument but is not the construction used in the formal theorem.

\subsubsection{The First Construction}\label{sec:first-construction}

The first construction samples $d+1$ slice outcomes $g_1,\dots,g_{d+1}\in \polyfunc{m}{q}{d}$, interpolates them to a global polynomial $h \in \polyfunc{m+1}{q}{d}$, and then averages over distinct interpolation points. Its consistency is governed by the commutation of the slice measurements, but its completeness is subtler. The heuristic comparison is
\[
	\bra{\psi} H \ot I \ket{\psi} \approx \bra{\psi} G^{d+1} \ot I \ket{\psi}.
\]
Thus one would like to show
\begin{equation}\label{eq:first-construction-g-power-comparison}
	\bra{\psi} G^{d+1} \ot I \ket{\psi}
	\approx \bra{\psi} G \ot I \ket{\psi}.
\end{equation}
The naive estimate would instead suggest
\begin{equation}\label{eq:first-construction-naive-bound}
	\bra{\psi} G^{d+1} \ot I \ket{\psi} \approx 1-(d+1)\kappa,
\end{equation}
which is too weak for the induction.

The key input is that one can first show
\begin{equation}\label{eq:first-construction-power-increment}
	\bra{\psi} G^{d+2} \ot I \ket{\psi}
	\approx_\Delta \bra{\psi} G^{d+1} \ot I \ket{\psi},
\end{equation}
for a small error $\Delta = \poly(m)\cdot \poly(\eps,\delta,\gamma,\zeta,d/q)$. Writing the eigendecomposition $G=\sum_i \lambda_i \ket{v_i}\bra{v_i}$ and the induced spectral distribution $\mu(i)=\bra{\psi}(\ket{v_i}\bra{v_i}\ot I)\ket{\psi}$, this is equivalent to
\begin{equation}\label{eq:first-construction-spectral-defect}
	\Delta
	\ge \bra{\psi} G^{d+1} \ot I \ket{\psi} - \bra{\psi} G^{d+2} \ot I \ket{\psi}
	= \E_{i\sim\mu}[\lambda_i^{d+1}(1-\lambda_i)].
\end{equation}
The remaining comparison between $G$ and $G^{d+1}$ is then reduced to a scalar inequality.

\begin{lemma}[A scalar inequality for the first construction]\label{lem:looks-easy-but-took-me-a-while}
	\lean{MIPStarRE.LDT.Pasting.first_construction_scalar_inequality}
	\leanok
	For every real number $0 \le \lambda \le 1$,
	\[
		\lambda(1-\lambda^d) \le 2(\lambda^{d+1}(1-\lambda))^{1/(d+1)}.
	\]
\end{lemma}
\begin{proof}
	\leanok
	Compare $1-\lambda^d$ with $(1-\lambda)(1+\lambda+\cdots+\lambda^{d-1})$ and use that $d^{1/(d+1)} \le 2$.
\end{proof}

\begin{corollary}[First-construction completeness comparison]\label{cor:first-construction-completeness-comparison}
	\uses{def:ld-pasting-context, lem:looks-easy-but-took-me-a-while}
	Lemma~\ref{lem:looks-easy-but-took-me-a-while} converts
	\eqref{eq:first-construction-spectral-defect} into the estimate
	\[
		\E_{i\sim\mu}[\lambda_i(1-\lambda_i^d)]
		\le 2\cdot \E_{i\sim\mu}\Big[(\lambda_i^{d+1}(1-\lambda_i))^{1/(d+1)}\Big]
		\le 2\cdot \Delta^{1/(d+1)},
	\]
	where the middle step uses concavity of $a \mapsto a^{1/(d+1)}$. Equivalently,
	\[
		\bra{\psi} G \ot I \ket{\psi} - \bra{\psi} G^{d+1} \ot I \ket{\psi}
		\le 2\cdot \Delta^{1/(d+1)}.
	\]
	For constant~$d$, this is exactly the heuristic completeness comparison for the
	first construction recorded in `references/ldt-paper/ld-pasting.tex`.
\end{corollary}
\begin{proof}
	Apply Lemma~\ref{lem:looks-easy-but-took-me-a-while} pointwise to the
	spectral distribution of \(G\), and then use the concavity of
	\(a \mapsto a^{1/(d+1)}\). The displayed operator expression is the same
	quantity written in the eigenbasis of \(G\).
\end{proof}

\subsubsection{The Second Construction}\label{sec:second-construction}

\begin{definition}[The incomplete part and the completed measurement]\label{def:G-hat}
	\lean{MIPStarRE.LDT.Pasting.GHatOutcome, MIPStarRE.LDT.Pasting.gHatIdxMeas}
	\leanok
	\uses{def:polymeasurements}
	For each $x \in \F_q$, write $G^x = \sum_g G_g^x$ and define the incomplete outcome $G_\bot^x = I-G^x$. The completed measurement $\widehat G^x$ has outcome set $\polyfunc{m}{q}{d} \cup \{\bot\}$ and is given by
	\[
		\widehat G_g^x = G_g^x \quad (g \in \polyfunc{m}{q}{d}),
		\qquad
		\widehat G_\bot^x = G_\bot^x.
	\]
\end{definition}

\begin{definition}[Types]\label{def:types}
	\lean{MIPStarRE.LDT.Pasting.GHatType, MIPStarRE.LDT.Pasting.gHatTypeWeight}
	\leanok
	A type is a vector $\tau \in \{0,1\}^k$. Its weight $|\tau|$ records how many coordinates lie in $\polyfunc{m}{q}{d}$ rather than equal $\bot$.
\end{definition}

\begin{definition}[The pasted measurement]\label{def:pasted-measurement}
	\lean{MIPStarRE.LDT.Pasting.pastedInterpolationFamily, MIPStarRE.LDT.Pasting.constructedPastedSubMeas, MIPStarRE.LDT.Pasting.constructedPastedMeasurement}
	\leanok
	\uses{def:G-hat, def:types, def:distinct-tuples}
	Fix $k \ge d+1$. For $x_1,\dots,x_k \in \F_q$ and outcomes $g_1,\dots,g_k \in \polyfunc{m}{q}{d}\cup\{\bot\}$, define
	\[
		\widehat H_{g_1,\dots,g_k}^{x_1,\dots,x_k}
		=
		\widehat G_{g_1}^{x_1}\widehat G_{g_2}^{x_2}\cdots \widehat G_{g_k}^{x_k}\cdots \widehat G_{g_2}^{x_2}\widehat G_{g_1}^{x_1}.
	\]
	If $(x_1,\dots,x_k)\in \mathsf{Distinct}_k$ and $h \in \polyfunc{m+1}{q}{d}$, define
	\[
		H_h^{x_1,\dots,x_k}
		=
		\sum_{\tau:\,|\tau|\ge d+1} \widehat H_{h_\tau}^{x_1,\dots,x_k},
	\]
	where the tuple $h_\tau$ has $i$-th entry $h|_{x_i}$ when $\tau_i=1$ and $\bot$ when $\tau_i=0$. The pasted submeasurement is then
	\[
		H_h = \E_{(x_1,\dots,x_k)\sim \mathsf{Distinct}_k} H_h^{x_1,\dots,x_k}.
	\]
\end{definition}

\begin{definition}[Completed-outcome reindexing]\label{def:completed-outcome-reindexing}
	\lean{MIPStarRE.LDT.Pasting.gHatTupleOutcomeConsEquiv', MIPStarRE.LDT.Pasting.thirdSliceFrontEquiv, MIPStarRE.LDT.Pasting.splitQuestionEquivOne, MIPStarRE.LDT.Pasting.splitOutcomeEquivOne, MIPStarRE.LDT.Pasting.thirdSliceFrontOutcomeEquiv}
	\leanok
	\uses{def:G-hat, def:types, def:pasted-measurement}
	The completed outcome tuples admit the reindexings obtained by separating the
	first coordinate, moving a third slice to the front, and identifying the
	one-coordinate split with the original slice question and outcome spaces.
\end{definition}

\begin{lemma}[Vertical restriction identities]\label{lem:vertical-restriction-identities}
	\lean{MIPStarRE.LDT.Pasting.postprocess_hRestrictionToVerticalLine_eq_evaluateAt, MIPStarRE.LDT.Pasting.restrictToAxisParallelLine_apply, MIPStarRE.LDT.Pasting.interpolateCompletedSlicesFromSupport_restrictAtHeight_poly_eq_get_of_mem, MIPStarRE.LDT.Pasting.evaluateAt_averageIdxSubMeas, MIPStarRE.LDT.Pasting.hRestrictionToVerticalLine_averageIdxSubMeas, MIPStarRE.LDT.Pasting.postprocess_restrictSubMeas_outcome}
	\leanok
	\uses{def:pasted-measurement}
	Restricting an interpolated polynomial to a vertical line and then
	postprocessing at a point agrees with evaluating the corresponding slice.
	The same identities commute with averaging over indexed submeasurements.
\end{lemma}
\begin{proof}
	\lean{MIPStarRE.LDT.Pasting.postprocess_hRestrictionToVerticalLine_eq_evaluateAt}
	\lean{MIPStarRE.LDT.Pasting.restrictToAxisParallelLine_apply}
	\lean{MIPStarRE.LDT.Pasting.interpolateCompletedSlicesFromSupport_restrictAtHeight_poly_eq_get_of_mem}
	\lean{MIPStarRE.LDT.Pasting.evaluateAt_averageIdxSubMeas}
	\lean{MIPStarRE.LDT.Pasting.hRestrictionToVerticalLine_averageIdxSubMeas}
	\lean{MIPStarRE.LDT.Pasting.postprocess_restrictSubMeas_outcome}
	\leanok
	\uses{def:pasted-measurement}
	The compatibility calculation for the restriction and evaluation maps used
	in the construction of~$H$ expands the vertical
	restriction, the interpolation support condition, postprocessing, and
	averaging over indexed submeasurements.
\end{proof}

\subsection{Strong Self-Consistency and Commutation of $\widehat G$}
\label{sec:hat-consistency}

\begin{lemma}[Strong self-consistency of the complete part]\label{lem:g-complete-self-consistency}
	\lean{MIPStarRE.LDT.Pasting.gCompleteSelfConsistency}
	\leanok
	\uses{def:ld-pasting-context, prop:two-notions-of-self-consistency}
	In the nontrivial-regime pasting context of Definition~\ref{def:ld-pasting-context},
	\[
		G^x \ot I \approx_\zeta I \ot G^x.
	\]
\end{lemma}
\begin{proof}
	\leanok
	\uses{prop:two-notions-of-self-consistency}
	Because $G$ is projective, \ref{prop:two-notions-of-self-consistency} identifies the assumed strong self-consistency of the outcome-level family with
	\begin{equation}\label{eq:ld-g-self-consistency}
		\E_x \sum_g \bra{\psi} G_g^x \ot G_g^x \ket{\psi}
		\ge
		\E_x \sum_g \bra{\psi} G_g^x \ot I \ket{\psi} - \frac{\zeta}{2}.
	\end{equation}
	Hence
	\begin{align*}
		      & \E_x \|(G^x\ot I - I\ot G^x)\ket{\psi}\|^2                                                              \\
		={}   & 2\E_x \bra{\psi} G^x \ot I \ket{\psi} - 2\E_x \bra{\psi} G^x \ot G^x \ket{\psi}                         \\
		={}   & 2\E_x \sum_g \bra{\psi} G_g^x \ot I \ket{\psi} - 2\E_x \sum_{g,h} \bra{\psi} G_g^x \ot G_h^x \ket{\psi} \\
		\le{} & 2\E_x \sum_g \bra{\psi} G_g^x \ot I \ket{\psi} - 2\E_x \sum_g \bra{\psi} G_g^x \ot G_g^x \ket{\psi}
		\le \zeta
	\end{align*}
	by \ref{eq:ld-g-self-consistency}.
\end{proof}

\begin{corollary}[Strong self-consistency of the incomplete part]\label{cor:g-bot-self-consistency}
	\lean{MIPStarRE.LDT.Pasting.gBotSelfConsistency}
	\leanok
	\uses{def:ld-pasting-context, lem:g-complete-self-consistency}
	In the nontrivial-regime pasting context of Definition~\ref{def:ld-pasting-context},
	\[
		G_\bot^x \ot I \approx_\zeta I \ot G_\bot^x.
	\]
\end{corollary}
\begin{proof}
	\leanok
	\uses{lem:g-complete-self-consistency}
	Since $G_\bot^x = I-G^x$,
	\[
		G_\bot^x \ot I - I \ot G_\bot^x = I \ot G^x - G^x \ot I,
	\]
	so the claim is immediate from Lemma~\ref{lem:g-complete-self-consistency}.
\end{proof}

\subsubsection{Commutativity of $G_{\bot}$}

\begin{lemma}[Switcheroo contraction estimates]\label{lem:switcheroo-contraction-estimates}
	\lean{MIPStarRE.LDT.Pasting.switcherooAggregateFourthTerm_once_commuted_contraction_left, MIPStarRE.LDT.Pasting.switcherooAggregateFourthTerm_once_commuted_contraction_right}
	\leanok
	\uses{prop:switch-sandwich, lem:g-complete-self-consistency}
	After one commutation of a $G$-factor through an auxiliary projective
	family, the two fourth-term contractions are bounded by the same
	self-consistency and commutation errors used in the switcheroo argument.
\end{lemma}
\begin{proof}
	\lean{MIPStarRE.LDT.Pasting.switcherooAggregateFourthTerm_once_commuted_contraction_left}
	\lean{MIPStarRE.LDT.Pasting.switcherooAggregateFourthTerm_once_commuted_contraction_right}
	\leanok
	\uses{prop:switch-sandwich, lem:g-complete-self-consistency}
	After expanding the fourth switcheroo term, the inserted complete part
	forms a positive contraction.  The left and right Cauchy--Schwarz
	side conditions are therefore the corresponding sums of
	$X_gX_g^\ast$ and $X_g^\ast X_g$, where
	$X_g=\sum_o G_{\bot}^x M_o^yG_g^xM_o^y$.  Orthogonality of the
	projective slice family and the submeasurement inequalities for $M$
	and $G_{\bot}$ bound both sums by the identity.  These are precisely
	the two contraction witnesses used in the fourth-term switcheroo chain.
\end{proof}

\begin{lemma}[Commutativity with $G_g^x$ implies commutativity with $G^x$]\label{lem:commutativity-switcheroo}
	\lean{MIPStarRE.LDT.Pasting.commutativitySwitcheroo}
	\leanok
	\uses{prop:switch-sandwich, lem:g-complete-self-consistency, lem:switcheroo-contraction-estimates}
	Let $M=\{M_o^x\}$ be a projective submeasurement with outcomes in a set $\mathcal O$. Suppose that
	\begin{equation}\label{eq:M-self-consistent}
		M_o^x \ot I \approx_\omega I \ot M_o^x,
	\end{equation}
	and
	\begin{equation}\label{eq:M-commutes-with-G}
		G_g^x M_o^y \ot I \approx_\chi M_o^y G_g^x \ot I
	\end{equation}
	on average over independent uniformly random $x,y \sim \F_q$. Then
	\[
		G^x M_o^y \ot I \approx_{6\sqrt{\zeta}+6\sqrt{\omega}+4\sqrt{\chi}} M_o^y G^x \ot I.
	\]
\end{lemma}
\begin{proof}
	\leanok
	\uses{prop:switch-sandwich, lem:g-complete-self-consistency}
	Let $M = \E_y \sum_o M_o^y$. The error to bound is
	\begin{align}
		    & \E_{x,y} \sum_o \|(G^xM_o^y - M_o^yG^x)\ot I \ket{\psi}\|^2 \nonumber \\
		={} & \E_{x,y} \sum_o \bra{\psi} M_o^y (G^x)^2 M_o^y \ot I \ket{\psi}
		+ \E_{x,y} \sum_o \bra{\psi} G^x (M_o^y)^2 G^x \ot I \ket{\psi} \nonumber   \\
		    & \qquad
		- \E_{x,y} \sum_o \bra{\psi} G^x M_o^y G^x M_o^y \ot I \ket{\psi}
		- \E_{x,y} \sum_o \bra{\psi} M_o^y G^x M_o^y G^x \ot I \ket{\psi}.
		\label{eq:g-commute-with-gg-error}
	\end{align}
	We compare all four terms with $\bra{\psi} G \ot M \ket{\psi}$.

	For the first term in \ref{eq:g-commute-with-gg-error},
	\begin{align*}
		\E_{x,y} \sum_o \bra{\psi} M_o^y (G^x)^2 M_o^y \ot I \ket{\psi}
		 & =
		\E_{x,y} \sum_o \bra{\psi} M_o^y G^x M_o^y \ot I \ket{\psi}               \\
		 & =
		\E_y \sum_o \bra{\psi} M_o^y G M_o^y \ot I \ket{\psi}                     \\
		 & \approx_{2\sqrt{\omega}} \E_y \sum_o \bra{\psi} G \ot M_o^y \ket{\psi}
		\tag{by \ref{prop:switch-sandwich} and \ref{eq:M-self-consistent}}        \\
		 & =
		\bra{\psi} G \ot M \ket{\psi}.
	\end{align*}
	For the second term,
	\begin{align*}
		\E_{x,y} \sum_o \bra{\psi} G^x (M_o^y)^2 G^x \ot I \ket{\psi}
		 & =
		\E_{x,y} \sum_o \bra{\psi} G^x M_o^y G^x \ot I \ket{\psi}                     \\
		 & =
		\E_x \bra{\psi} G^x M G^x \ot I \ket{\psi}                                    \\
		 & \approx_{2\sqrt{\zeta}} \E_x \bra{\psi} M \ot G^x \ket{\psi}
		\tag{by \ref{prop:switch-sandwich} and \ref{lem:g-complete-self-consistency}} \\
		 & =
		\bra{\psi} M \ot G \ket{\psi}.
	\end{align*}

	For the third term, first
	\begin{align}
		\E_{x,y} \sum_o \bra{\psi} G^x M_o^y G^x M_o^y \ot I \ket{\psi}
		 & =
		\E_{x,y} \sum_{o,g} \bra{\psi} G^x M_o^y G_g^x G_g^x M_o^y \ot I \ket{\psi} \nonumber \\
		 & \approx_{\sqrt{\chi}}
		\E_{x,y} \sum_{o,g} \bra{\psi} G^x M_o^y G_g^x M_o^y G_g^x \ot I \ket{\psi}.
		\label{eq:split-G-and-commute}
	\end{align}
	This is the Cauchy--Schwarz estimate coming from \ref{eq:M-commutes-with-G}. Next,
	\begin{equation}\label{eq:move-G-to-right-factor}
		\eqref{eq:split-G-and-commute}
		\approx_{\sqrt{\zeta}}
		\E_{x,y} \sum_{o,g} \bra{\psi} G^x M_o^y G_g^x M_o^y \ot G_g^x \ket{\psi}.
	\end{equation}
	Here the error is bounded using \ref{item:ld-pasting-self-consistency}. Then
	\begin{equation}\label{eq:move-G-back-to-left-factor}
		\eqref{eq:move-G-to-right-factor}
		\approx_{\sqrt{\zeta}}
		\E_{x,y} \sum_{o,g} \bra{\psi} G_g^x G^x M_o^y G_g^x M_o^y \ot I \ket{\psi},
	\end{equation}
	again by strong self-consistency of $G_g^x$. Since $G$ is projective,
	\begin{align}
		\eqref{eq:move-G-back-to-left-factor}
		 & =
		\E_{x,y} \sum_{o,g} \bra{\psi} G_g^x M_o^y G_g^x M_o^y \ot I \ket{\psi} \nonumber \\
		 & \approx_{\sqrt{\chi}}
		\E_{x,y} \sum_{o,g} \bra{\psi} M_o^y G_g^x G_g^x M_o^y \ot I \ket{\psi}.
		\label{eq:commuted-G-after-transfer}
	\end{align}
	Finally,
	\begin{align*}
		\eqref{eq:commuted-G-after-transfer}
		 & =
		\E_{x,y} \sum_{o,g} \bra{\psi} M_o^y G_g^x M_o^y \ot I \ket{\psi}         \\
		 & =
		\E_y \sum_o \bra{\psi} M_o^y G M_o^y \ot I \ket{\psi}                     \\
		 & \approx_{2\sqrt{\omega}} \E_y \sum_o \bra{\psi} G \ot M_o^y \ket{\psi}
		\tag{by \ref{prop:switch-sandwich} and \ref{eq:M-self-consistent}}        \\
		 & =
		\bra{\psi} G \ot M \ket{\psi}.
	\end{align*}
	Therefore
	\begin{equation}\label{eq:third-error-term-bound}
		\E_{x,y} \sum_o \bra{\psi} G^x M_o^y G^x M_o^y \ot I \ket{\psi}
		\approx_{2\sqrt{\zeta}+2\sqrt{\omega}+2\sqrt{\chi}}
		\bra{\psi} G \ot M \ket{\psi}.
	\end{equation}
	The fourth term in \ref{eq:g-commute-with-gg-error} is the Hermitian conjugate of the third term, so it satisfies the same bound. Summing the four errors gives
	\[
		2\sqrt{\omega}+2\sqrt{\zeta}+2(2\sqrt{\zeta}+2\sqrt{\omega}+2\sqrt{\chi})
		=
		6\sqrt{\zeta}+6\sqrt{\omega}+4\sqrt{\chi}.
	\]
\end{proof}
\begin{remark}\label{rem:commutativity-switcheroo-density-swap}
	The Lean theorem corresponding to Lemma~\ref{lem:commutativity-switcheroo}
	takes one extra symmetry hypothesis,
	\texttt{hfix : swapDensity ψbi.density = ψbi.density}. This is the
	\texttt{density\_swap} component of the paper's permutation-invariance
	assumption on the symmetric two-prover strategy; the stronger
	\texttt{swap\_ev} field of \texttt{PermInvState} is not assumed by this
	theorem. The hypothesis is used to identify
	$\bra{\psi} G \ot M \ket{\psi}$ with $\bra{\psi} M \ot G \ket{\psi}$
	at the center of the triangle-inequality chain.
\end{remark}

\begin{corollary}[Commutativity of the complete part]\label{cor:commuting-with-G-complete}
	\lean{MIPStarRE.LDT.Pasting.commutingWithGComplete}
	\leanok
	\uses{def:ld-pasting-context, thm:com-main, lem:commutativity-switcheroo,
	      lem:g-complete-self-consistency}
	In the nontrivial-regime pasting context of Definition~\ref{def:ld-pasting-context},
	the following commutation relations hold:
	\[
		G_g^x G^y \ot I \approx_{\nu_2} G^y G_g^x \ot I,
		\qquad
		G^x G^y \ot I \approx_{\nu_2} G^y G^x \ot I,
	\]
	where
	\[
		\nu_2 = 36m\left(\gamma^{1/16}+\zeta^{1/16}+(d/q)^{1/16}\right).
	\]
\end{corollary}
\begin{proof}
	\leanok
	\uses{lem:commutativity-switcheroo}
	By \ref{eq:complete-slice-commutativity},
	\begin{equation}\label{eq:complete-outcome-commutativity}
		G_g^x G_h^y \ot I \approx_{\nu_{\mathrm{com}}} G_h^y G_g^x \ot I.
	\end{equation}
	Applying Lemma~\ref{lem:commutativity-switcheroo} to \ref{eq:complete-outcome-commutativity}, with $\{M_o^y\}$ equal to the outcome family $\{G_h^y\}$, gives
	\begin{equation}\label{eq:complete-total-commutes-with-outcome}
		G_g^x G^y \ot I \approx_{\theta_1} G^y G_g^x \ot I,
	\end{equation}
	where $\theta_1 = 12\sqrt{\zeta} + 4\sqrt{\nu_{\mathrm{com}}}$. Applying Lemma~\ref{lem:commutativity-switcheroo} again to \ref{eq:complete-total-commutes-with-outcome}, now with the one-outcome family $M_o^y = G^y$, gives
	\[
		G^x G^y \ot I \approx_{\theta_2} G^y G^x \ot I,
	\]
	where $\theta_2 = 12\sqrt{\zeta} + 4\sqrt{\theta_1}$.

	Since $\nu_{\mathrm{com}}=30m(\gamma^{1/4}+\zeta^{1/4}+(d/q)^{1/4})$, these errors satisfy
	\[
		\theta_1 \le 36m\left(\gamma^{1/8}+\zeta^{1/8}+(d/q)^{1/8}\right)\le \nu_2,
	\]
	\[
		\theta_2 \le 36m\left(\gamma^{1/16}+\zeta^{1/16}+(d/q)^{1/16}\right)=\nu_2.
	\]
\end{proof}

\begin{corollary}[Commutativity of the incomplete part]\label{cor:commuting-with-G-incomplete}
	\lean{MIPStarRE.LDT.Pasting.commutingWithGIncomplete}
	\leanok
	\uses{def:ld-pasting-context, cor:commuting-with-G-complete}
	In the nontrivial-regime pasting context of Definition~\ref{def:ld-pasting-context},
	\[
		G_g^x G_\bot^y \ot I \approx_{\nu_2} G_\bot^y G_g^x \ot I,
		\qquad
		G_\bot^x G_\bot^y \ot I \approx_{\nu_2} G_\bot^y G_\bot^x \ot I.
	\]
\end{corollary}
\begin{proof}
	\leanok
	\uses{cor:commuting-with-G-complete}
	The identities
	\[
		G_g^x G_\bot^y - G_\bot^y G_g^x = G^y G_g^x - G_g^x G^y
	\]
	and
	\[
		G_\bot^x G_\bot^y - G_\bot^y G_\bot^x = G^x G^y - G^y G^x
	\]
	reduce both claims to Corollary~\ref{cor:commuting-with-G-complete}.
\end{proof}

\subsubsection{Putting Everything Together}

\begin{corollary}[Strong self-consistency and commutation of $\widehat G$]\label{cor:G-hat-facts}
	\lean{MIPStarRE.LDT.Pasting.gHatFacts}
	\leanok
	\uses{lem:g-complete-self-consistency, cor:g-bot-self-consistency, cor:commuting-with-G-complete, cor:commuting-with-G-incomplete}
	The completed measurements $\widehat G^x$ satisfy
	\begin{align}
		\widehat G_g^x \ot I                & \approx_{2\zeta} I \ot \widehat G_g^x, \label{eq:g-hat-self-consistency}           \\
		\widehat G_g^x \widehat G_h^y \ot I & \approx_{\nu_3} \widehat G_h^y \widehat G_g^x \ot I, \label{eq:g-hat-commutativity}
	\end{align}
	where
	\[
		\nu_3 = 138m\left(\gamma^{1/16}+\zeta^{1/16}+(d/q)^{1/16}\right).
	\]
\end{corollary}
\begin{proof}
	\leanok
	\uses{cor:g-bot-self-consistency, cor:commuting-with-G-complete,
		cor:commuting-with-G-incomplete}
	For \ref{eq:g-hat-self-consistency}, split the outcome set into the genuine outcomes and $\bot$:
	\[
		\E_x \sum_g \|(\widehat G_g^x \ot I - I \ot \widehat G_g^x)\ket{\psi}\|^2
		\le \zeta + \zeta = 2\zeta
	\]
	by \ref{item:ld-pasting-self-consistency} and \ref{cor:g-bot-self-consistency}. For \ref{eq:g-hat-commutativity}, the same decomposition yields four cases: polynomial--polynomial, $\bot$--$\bot$, polynomial--$\bot$, and $\bot$--polynomial. These contribute $\nu_{\mathrm{com}}$ and three copies of $\nu_2$, so
	\[
		\nu_{\mathrm{com}} + 3\nu_2
		\le
		138m\left(\gamma^{1/16}+\zeta^{1/16}+(d/q)^{1/16}\right)
		= \nu_3.
	\]
\end{proof}

\begin{lemma}[Completed part bounded by a slice]\label{lem:q-sdd-complete-part-slice-bound}
	\lean{MIPStarRE.LDT.Pasting.qSDD_completePart_le_slice}
	\leanok
	\uses{def:G-hat, cor:G-hat-facts}
	For every $x \in \F_q$, the squared-distance defect of the one-outcome
	complete part is bounded by that of the original slice submeasurement:
	\[
		\bigl\|
		(G^x \ot I - I \ot G^x)\ket{\psi}
		\bigr\|^2
		\;\le\;
		\sum_g
		\bigl\|
		(G^x_g \ot I - I \ot G^x_g)\ket{\psi}
		\bigr\|^2,
	\]
	where $G^x = \sum_g G^x_g$ is the total operator of the slice
	submeasurement.
\end{lemma}
\begin{proof}
	\lean{MIPStarRE.LDT.Pasting.qSDD_completePart_le_slice}
	\leanok
	\uses{def:G-hat, cor:G-hat-facts}
	This is the monotonicity estimate for the completed part of~$\widehat G$:
	the completed contribution is one summand of the corresponding slice
	squared-distance defect, and the remaining summands are nonnegative.
\end{proof}

\subsection{Sandwiching Lemmas}
\label{sec:ld-sandwiching}

\begin{lemma}[Completed sandwich normalization]\label{lem:completed-sandwich-normalization}
	\lean{MIPStarRE.LDT.Pasting.gHatHalfSandwichLeft_split_outcome, MIPStarRE.LDT.Pasting.gHatHalfSandwichLeft_split_total, MIPStarRE.LDT.Pasting.gHatHalfSandwichRight_split_outcome, MIPStarRE.LDT.Pasting.gHatHalfSandwichRight_split_total, MIPStarRE.LDT.Pasting.gHatHalfProduct_sum_adjoint_mul_le_one, MIPStarRE.LDT.Pasting.gHatReverseHalfProduct_sum_adjoint_mul_le_one, MIPStarRE.LDT.Pasting.leftTensor_rightTensor_sum_adjoint_mul_le_one}
	\leanok
	\uses{def:G-hat, def:completed-outcome-reindexing}
	\begin{enumerate}
		\item \textbf{Head--tail split identities.} For every $k \ge 0$, point
			tuple $\bx = (x_1,\dots,x_{k+1})$, and outcome tuple
			$\bg = (g_1,\dots,g_{k+1})$,
			\[
				\widehat G^{x_1}_{g_1}\cdots
				\widehat G^{x_{k+1}}_{g_{k+1}}
				=
				\widehat G^{x_1}_{g_1}
				\bigl(
				\widehat G^{x_2}_{g_2}\cdots
				\widehat G^{x_{k+1}}_{g_{k+1}}
				\bigr),
			\]
			and the right half-sandwich analogously rotates the head to the tail.
		\item \textbf{Sum-of-adjoint-products bounds.} For every $r\ge 0$ and
			$\bx = (x_1,\dots,x_r) \in \F_q^r$,
			\[
				\sum_{g_1,\dots,g_r}
				\bigl(\widehat G^{x_1}_{g_1} \cdots
				\widehat G^{x_r}_{g_r}\bigr)^\dagger
				\bigl(\widehat G^{x_1}_{g_1} \cdots
				\widehat G^{x_r}_{g_r}\bigr)
				\;\le\; I,
			\]
			and the same bound holds for the reverse-order product family.
		\item \textbf{Generic tensor contraction.} If families~$\{A_a\}_a$
			and~$\{B_b\}_b$ satisfy $\sum_a A_a^\dagger A_a \le I$ and
			$\sum_b B_b^\dagger B_b \le I$, then
			\[
				\sum_{a,b} (A_a \ot B_b)^\dagger (A_a \ot B_b) \le I.
			\]
	\end{enumerate}
\end{lemma}
\begin{proof}
	\lean{MIPStarRE.LDT.Pasting.gHatHalfSandwichLeft_split_outcome}
	\lean{MIPStarRE.LDT.Pasting.gHatHalfSandwichLeft_split_total}
	\lean{MIPStarRE.LDT.Pasting.gHatHalfSandwichRight_split_outcome}
	\lean{MIPStarRE.LDT.Pasting.gHatHalfSandwichRight_split_total}
	\lean{MIPStarRE.LDT.Pasting.gHatHalfProduct_sum_adjoint_mul_le_one}
	\lean{MIPStarRE.LDT.Pasting.gHatReverseHalfProduct_sum_adjoint_mul_le_one}
	\lean{MIPStarRE.LDT.Pasting.leftTensor_rightTensor_sum_adjoint_mul_le_one}
	\leanok
	\uses{def:G-hat, def:completed-outcome-reindexing}
	The split identities follow from the completed-outcome equivalences.  The
	normalization inequalities are the corresponding submeasurement bounds for
	the left and right half-products after reindexing the sums.
\end{proof}

\begin{lemma}[One-step completed sandwich commutation]\label{lem:one-step-completed-sandwich-commutation}
	\lean{MIPStarRE.LDT.Pasting.gHatSelfConsistency_sddOpRel_quadThird, MIPStarRE.LDT.Pasting.commuteGHalfSandwich_step_commute, MIPStarRE.LDT.Pasting.commuteGHalfSandwich_prefixFirstSliceLeft_move, MIPStarRE.LDT.Pasting.commuteGHalfSandwich_moveSource_eq_split, MIPStarRE.LDT.Pasting.commuteGHalfSandwich_move_recursive_zero, MIPStarRE.LDT.Pasting.commuteGHalfSandwich_split_one_iff, MIPStarRE.LDT.Pasting.commuteGHalfSandwich_core_two, MIPStarRE.LDT.Pasting.commuteGHalfSandwich_moveStepMid_toTarget, MIPStarRE.LDT.Pasting.commuteGHalfSandwich_core}
	\leanok
	\uses{cor:G-hat-facts, lem:completed-sandwich-normalization,
		prop:triangle-inequality-for-approx_delta}
	Given the completed-measurement self-consistency and commutation estimates
	from Corollary~\ref{cor:G-hat-facts}, the following one-step estimates hold.
	\begin{enumerate}
		\item \textbf{Self-consistency step.} For every~$r\ge 0$,
			\[
				\E_{x,y,z}\E_{x_1,\dots,x_r}
				\sum_{g}
				\bigl\|
				\bigl(\widehat G^{z}_g \ot I -
					I \ot \widehat G^{z}_g\bigr)
				\ket{\psi}
				\bigr\|^2
				\;\le\; 2\zeta.
			\]
		\item \textbf{Commutation step.} For every~$r\ge 0$,
			\[
				\E_{x,y}\E_{x_1,\dots,x_r}
				\sum_{g,h}\sum_{g_1,\dots,g_r}
				\bigl\|
				\bigl(
					\widehat G^{x}_g \widehat G^{y}_h
					\ot
					\widehat G^{x_r}_{g_r}\cdots
					\widehat G^{x_1}_{g_1}
					-
					\widehat G^{y}_h \widehat G^{x}_g
					\ot
					\widehat G^{x_r}_{g_r}\cdots
					\widehat G^{x_1}_{g_1}
				\bigr)
				\ket{\psi}
				\bigr\|^2
				\;\le\; \nu_3.
			\]
		\item \textbf{Two-factor base case.} For $k=2$, the commutation chain
			reduces to the same tensor-placement convention with the reverse
			tail product equal to the identity, hence to the pairwise
			estimate~$\nu_3$.
		\item \textbf{Split equivalences.} The half-sandwich families are
			equivalent, under the head--tail reindexing of
			Lemma~\ref{lem:completed-sandwich-normalization}, to the
			chain-families used in the move estimates above, with ordered
			products on the first tensor factor and reverse products on the
			second tensor factor where these appear in the commutation step.
	\end{enumerate}
	These are the elementary building blocks assembled by the recursive chain
	construction in the following lemma.
\end{lemma}
\begin{proof}
	\lean{MIPStarRE.LDT.Pasting.gHatSelfConsistency_sddOpRel_quadThird}
	\lean{MIPStarRE.LDT.Pasting.commuteGHalfSandwich_step_commute}
	\lean{MIPStarRE.LDT.Pasting.commuteGHalfSandwich_prefixFirstSliceLeft_move}
	\lean{MIPStarRE.LDT.Pasting.commuteGHalfSandwich_moveSource_eq_split}
	\lean{MIPStarRE.LDT.Pasting.commuteGHalfSandwich_move_recursive_zero}
	\lean{MIPStarRE.LDT.Pasting.commuteGHalfSandwich_split_one_iff}
	\lean{MIPStarRE.LDT.Pasting.commuteGHalfSandwich_core_two}
	\lean{MIPStarRE.LDT.Pasting.commuteGHalfSandwich_moveStepMid_toTarget}
	\lean{MIPStarRE.LDT.Pasting.commuteGHalfSandwich_core}
	\leanok
	\uses{cor:G-hat-facts, lem:completed-sandwich-normalization, prop:triangle-inequality-for-approx_delta}
	For the one-step commutation move, rewrite the relevant head--tail split,
	apply the self-consistency and commutation estimates
	for~$\widehat G$, and assemble the resulting comparison by the triangle
	inequality for~$\approx_\delta$.
\end{proof}

\begin{lemma}[Commuting past multiple completed measurements]\label{lem:commute-g-half-sandwich}
	\lean{MIPStarRE.LDT.Pasting.commuteGHalfSandwich}
	\leanok
	\uses{cor:G-hat-facts, lem:one-step-completed-sandwich-commutation, prop:triangle-inequality-for-approx_delta}
	For every $k \ge 2$,
	\[
		\widehat G_{g_1}^{x_1}\widehat G_{g_2}^{x_2}\cdots \widehat G_{g_k}^{x_k} \ot I
		\approx_{\nu_4}
		\widehat G_{g_2}^{x_2}\cdots \widehat G_{g_k}^{x_k}\widehat G_{g_1}^{x_1} \ot I,
	\]
	where
	\[
		\nu_4 = 426k^2m\left(\gamma^{1/16}+\zeta^{1/16}+(d/q)^{1/16}\right).
	\]
\end{lemma}
\begin{proof}
	\leanok
	\uses{cor:G-hat-facts, lem:one-step-completed-sandwich-commutation,
		prop:triangle-inequality-for-approx_delta}
	Repeatedly move the leftmost factor across the remaining product, using \ref{eq:g-hat-self-consistency} and \ref{eq:g-hat-commutativity} at each step and summing the resulting errors with \ref{prop:triangle-inequality-for-approx_delta}. The paper's calculation gives
	\[
		3k\bigl(4k\zeta + k\nu_3\bigr)
		\le
		426k^2m\left(\gamma^{1/16}+\zeta^{1/16}+(d/q)^{1/16}\right).
	\]
\end{proof}

\begin{lemma}[Line-interpolation mismatch estimates]\label{lem:line-interpolation-mismatch-estimates}
	\lean{MIPStarRE.LDT.Pasting.tupleInterpolatedVerticalLine_ne_gives_exists_some_eval_mismatch, MIPStarRE.LDT.Pasting.interpolationEligibleSandwich_mismatch_sum_mono, MIPStarRE.LDT.Pasting.interpolationEligibleSandwich_exists_mismatch_sum_le_sum, MIPStarRE.LDT.Pasting.qBipartiteConsDefect_postprocess_eq_singleOutcome, MIPStarRE.LDT.Pasting.qBipartiteConsDefect_eq_sum_singleOutcome, MIPStarRE.LDT.Pasting.ldSandwichLineOnePointRightMeasurement_outcome_some_eq_sum, MIPStarRE.LDT.Pasting.hBConsistencyCoordMass_le_linePointDefect, MIPStarRE.LDT.Pasting.hBConsistencyBadMass_le_one, MIPStarRE.LDT.Pasting.pastedInterpolation_verticalLine_defect_le_badMass}
	\leanok
	\uses{lem:vertical-restriction-identities, lem:ld-gbcon,
		lem:ld-sandwich-line-one-point}
	If the tuple of completed slices does not interpolate to the sampled
	vertical-line value, then some active coordinate witnesses a line-value
	mismatch. Concretely, for every direction $u$ and injective tuple
	$\mathbf{x}=(x_1,\dots,x_k)$, let $\widehat{H}^{\mathbf{x}}\rvert_u$ be the
	vertical-line restriction of the pasted interpolation family. Then
	\[
		\operatorname{Defect}_\psi\bigl(
			\widehat{H}^{\mathbf{x}}\rvert_u,\,
			B^u\bigr)
		\;\le\;
		\operatorname{Bad}(\psi;u,\mathbf{x})
		\;\le\;
		\sum_{i=1}^k
		\operatorname{Defect}_\psi\bigl(
			\widehat{H}^{(i)}(u,\mathbf{x}),\,
			B^{(i)}(u,\mathbf{x})\bigr),
	\]
	where $\operatorname{Defect}_\psi(A,B)$ denotes the bipartite consistency
	defect $\max\bigl(0,\,
		\bra{\psi} A_{\mathrm{tot}}\ot B_{\mathrm{tot}}\ket{\psi}
		- \sum_a \bra{\psi} A_a\ot B_a\ket{\psi}\bigr)$,
		\[
			\operatorname{Bad}(\psi;u,\mathbf{x})
			= \sum_f \bra{\psi}
			\Bigl(\sum_{\substack{g_1,\dots,g_k\\ \exists i:\, g_i\ne\bot,\ g_i(u)\ne f(x_i)}}
				\widehat H_{g_1,\dots,g_k}^{\mathbf{x}}
			\Bigr) \ot B_f^u
			\ket{\psi},
		\]
		and $\widehat{H}^{(i)},B^{(i)}$ are the one-point line comparison
		families (Lemma~\ref{lem:ld-sandwich-line-one-point})
		restricted to coordinate $i$, and $\operatorname{Bad}(\psi;u,\mathbf{x})\le 1$.
\end{lemma}
\begin{proof}
	\lean{MIPStarRE.LDT.Pasting.tupleInterpolatedVerticalLine_ne_gives_exists_some_eval_mismatch}
	\lean{MIPStarRE.LDT.Pasting.interpolationEligibleSandwich_mismatch_sum_mono}
	\lean{MIPStarRE.LDT.Pasting.interpolationEligibleSandwich_exists_mismatch_sum_le_sum}
	\lean{MIPStarRE.LDT.Pasting.qBipartiteConsDefect_postprocess_eq_singleOutcome}
	\lean{MIPStarRE.LDT.Pasting.qBipartiteConsDefect_eq_sum_singleOutcome}
	\lean{MIPStarRE.LDT.Pasting.ldSandwichLineOnePointRightMeasurement_outcome_some_eq_sum}
	\lean{MIPStarRE.LDT.Pasting.hBConsistencyCoordMass_le_linePointDefect}
	\lean{MIPStarRE.LDT.Pasting.hBConsistencyBadMass_le_one}
	\lean{MIPStarRE.LDT.Pasting.pastedInterpolation_verticalLine_defect_le_badMass}
	\leanok
	\uses{lem:vertical-restriction-identities, lem:ld-gbcon,
		lem:ld-sandwich-line-one-point}
	The proof first identifies a failed vertical-line interpolation with a
	particular active coordinate whose slice value disagrees with the line
	answer. The operator sum over all such failures is then bounded by the sum
	of the corresponding one-coordinate line defects, using the vertical
	restriction identities above.
\end{proof}

\begin{lemma}[Averaging the line-interpolation defect]\label{lem:line-interpolation-averaging-estimates}
	\lean{MIPStarRE.LDT.Pasting.avgOver_distinct_bounded_le_avgOver_uniform_add_tv, MIPStarRE.LDT.Pasting.qBipartiteMatchMass_averageIdxSubMeas_left, MIPStarRE.LDT.Pasting.ev_opTensor_total_averageIdxSubMeas_left, MIPStarRE.LDT.Pasting.qBipartiteConsDefect_averageIdxSubMeas_left_le, MIPStarRE.LDT.Pasting.hBConsistency_fixed_u_defect_le_avgOver_distinct, MIPStarRE.LDT.Pasting.one_div_q_le_rpow_degreeRatio, MIPStarRE.LDT.Pasting.avgOver_uniform_badMass_le_k_mul_ldSandwichLineOnePointError, MIPStarRE.LDT.Pasting.avgOver_distinct_badMass_le_hBConsistencyError, MIPStarRE.LDT.Pasting.qBipartiteConsDefect_eq_false_mass_of_bool_right_true, MIPStarRE.LDT.Pasting.avgOver_distinct_pasted_defect_le_badMass}
	\leanok
	\uses{prop:ld-dnoteq, lem:line-interpolation-mismatch-estimates,
		lem:ld-sandwich-line-one-point}
	Replacing distinct tuples by independent tuples costs the total-variation
	error $k^2/q$ from Proposition~\ref{prop:ld-dnoteq}. After averaging over
	$u$ and $\mathbf{x}\sim\mathsf{Distinct}_k$, the expected bad mass and the
	expected pasted interpolation defect are bounded by the one-point sandwich
	error:
	\[
		\E_u\;
		\E_{\mathbf{x}\sim\mathsf{Distinct}_k}\;
		\operatorname{Bad}(\psi;u,\mathbf{x})
		\;\le\;
		k\cdot\nu_5 \;+\; \frac{k^2}{q},
	\]
	and consequently
	\[
		\E_u\;\E_{\mathbf{x}\sim\mathsf{Distinct}_k}\;
		\operatorname{Defect}_\psi\bigl(
			\widehat{H}^{\mathbf{x}}\rvert_u,\,
			B^u\bigr)
		\;\le\;
		k\cdot\nu_5 \;+\; \frac{k^2}{q},
	\]
	where $\nu_5 = 43km(\eps^{1/32}+\dots+(d/q)^{1/32})$
	(Lemma~\ref{lem:ld-sandwich-line-one-point}). Using
	$1/q \le (d/q)^{1/32}$ and $m\ge 1$, this is further bounded by
	$44k^2m(\eps^{1/32}+\delta^{1/32}+\gamma^{1/32}+\zeta^{1/32}+(d/q)^{1/32})$
	as used in Lemma~\ref{lem:h-b-consistency}.
\end{lemma}
\begin{proof}
	% Lean resolves these declarations as root `PMF.*` names; the namespace is not
	% below `MIPStarRE.LDT`, so there is no longer `MIPStarRE.*` declaration name.
	\lean{MIPStarRE.LDT.totalVariationDistance_eq_univ_sum}
	\lean{MIPStarRE.LDT.totalVariationDistance_eq_toPMF_sum}
	\lean{MIPStarRE.LDT.totalVariationDistance_eq_sum_max_sub}
	\lean{PMF.sum_le_sum_add_totalVariationDistance}
	\lean{MIPStarRE.LDT.avgOver_le_avgOver_add_totalVariationDistance}
	\lean{MIPStarRE.LDT.Pasting.avgOver_distinct_bounded_le_avgOver_uniform_add_tv}
	\lean{MIPStarRE.LDT.Pasting.qBipartiteMatchMass_averageIdxSubMeas_left}
	\lean{MIPStarRE.LDT.Pasting.ev_opTensor_total_averageIdxSubMeas_left}
	\lean{MIPStarRE.LDT.Pasting.qBipartiteConsDefect_averageIdxSubMeas_left_le}
	\lean{MIPStarRE.LDT.Pasting.hBConsistency_fixed_u_defect_le_avgOver_distinct}
	\lean{MIPStarRE.LDT.Pasting.one_div_q_le_rpow_degreeRatio}
	\lean{MIPStarRE.LDT.Pasting.avgOver_uniform_badMass_le_k_mul_ldSandwichLineOnePointError}
	\lean{MIPStarRE.LDT.Pasting.avgOver_distinct_badMass_le_hBConsistencyError}
	\lean{MIPStarRE.LDT.Pasting.qBipartiteConsDefect_eq_false_mass_of_bool_right_true}
	\lean{MIPStarRE.LDT.Pasting.avgOver_distinct_pasted_defect_le_badMass}
	\leanok
	\uses{prop:ld-dnoteq, lem:line-interpolation-mismatch-estimates}
	The proof compares the distinct-tuple average with the independent-tuple
	average, paying the total-variation cost from
	Proposition~\ref{prop:ld-dnoteq}. The independent average separates into
	the one-coordinate line defects from the previous lemma, and the scalar
	inequality $1/q\le (d/q)^{1/32}$ absorbs the sampling error into the stated
	pasting parameter.
\end{proof}

\begin{lemma}[Consistency of the sandwiched measurement with one line value]\label{lem:ld-sandwich-line-one-point}
	\lean{MIPStarRE.LDT.Pasting.ldSandwichLineOnePoint}
	\leanok
	\uses{lem:commute-g-half-sandwich, lem:ld-gbcon}
	For any $1 \le i \le k$,
	\[
		\E_u \E_{x_1,\dots,x_k} \sum_{g_1,\dots,g_k\,:\,g_i\ne\bot} \sum_{a\ne g_i(u)}
		\bra{\psi} \widehat H_{g_1,\dots,g_k}^{x_1,\dots,x_k} \ot B_{[f(x_i)=a]}^u \ket{\psi}
		\le \nu_5,
	\]
	where
	\[
		\nu_5 = 43km\left(\eps^{1/32}+\delta^{1/32}+\gamma^{1/32}+\zeta^{1/32}+(d/q)^{1/32}\right).
	\]
\end{lemma}
\begin{proof}
	\leanok
	\uses{lem:commute-g-half-sandwich, lem:ld-gbcon}
	Summing out the coordinates to the right of $i$ gives
	\begin{equation}\label{eq:sum-out-right-coordinates}
		\begin{aligned}
			    & \E_u \E_{x_1,\dots,x_k} \sum_{g_1,\dots,g_k\,:\,g_i\ne\bot} \sum_{a\ne g_i(u)}
			\bra{\psi} \widehat H_{g_1,\dots,g_k}^{x_1,\dots,x_k} \ot B_{[f(x_i)=a]}^u \ket{\psi} \\
			={} &
			\E_u \E_{x_1,\dots,x_i} \sum_{g_1,\dots,g_i\,:\,g_i\ne\bot} \sum_{a\ne g_i(u)}
			\bra{\psi} \widehat H_{g_1,\dots,g_i}^{x_1,\dots,x_i} \ot B_{[f(x_i)=a]}^u \ket{\psi}.
		\end{aligned}
	\end{equation}
	Write $\widehat G_{g_{<i}}^{x_{<i}} = \widehat G_{g_1}^{x_1}\cdots \widehat G_{g_{i-1}}^{x_{i-1}}$. Then
	\begin{align}
		\eqref{eq:sum-out-right-coordinates}
		 & =
		\E_u \E_{x_1,\dots,x_i} \sum_{g_1,\dots,g_i\,:\,g_i\ne\bot}
		\bra{\psi}
		\widehat G_{g_{<i}}^{x_{<i}} \widehat G_{g_i}^{x_i}\widehat G_{g_i}^{x_i} (\widehat G_{g_{<i}}^{x_{<i}})^\dagger
		\ot (I-B_{[f(x_i)=g_i(u)]}^u)\ket{\psi}\nonumber \\
		 & \approx_{\sqrt{\nu_4}}
		\E_u \E_{x_1,\dots,x_i} \sum_{g_1,\dots,g_i\,:\,g_i\ne\bot}
		\bra{\psi}
		\widehat G_{g_i}^{x_i}\widehat G_{g_{<i}}^{x_{<i}}\widehat G_{g_i}^{x_i}(\widehat G_{g_{<i}}^{x_{<i}})^\dagger
		\ot (I-B_{[f(x_i)=g_i(u)]}^u)\ket{\psi}.
		\label{eq:half-sandwich-first-commutation}
	\end{align}
	The Cauchy--Schwarz argument uses
	\begin{equation}\label{eq:completed-outcome-commutator}
		\E_u \E_{x_1,\dots,x_i} \sum_{g_1,\dots,g_i}
		\bra{\psi}
		\bigl((\widehat G_{g_{<i}}^{x_{<i}}\widehat G_{g_i}^{x_i} - \widehat G_{g_i}^{x_i}\widehat G_{g_{<i}}^{x_{<i}})
		(\widehat G_{g_i}^{x_i}(\widehat G_{g_{<i}}^{x_{<i}})^\dagger - (\widehat G_{g_{<i}}^{x_{<i}})^\dagger \widehat G_{g_i}^{x_i})\bigr)
		\ot I \ket{\psi},
	\end{equation}
	which is at most $\nu_4$ by Lemma~\ref{lem:commute-g-half-sandwich}. A second Cauchy--Schwarz step gives
	\begin{equation}\label{eq:half-sandwich-second-commutation}
		\eqref{eq:half-sandwich-first-commutation}
		\approx_{\sqrt{\nu_4}}
		\E_u \E_{x_1,\dots,x_i} \sum_{g_1,\dots,g_i\,:\,g_i\ne\bot}
		\bra{\psi}
		\widehat G_{g_i}^{x_i}\widehat G_{g_{<i}}^{x_{<i}}(\widehat G_{g_{<i}}^{x_{<i}})^\dagger \widehat G_{g_i}^{x_i}
		\ot (I-B_{[f(x_i)=g_i(u)]}^u)\ket{\psi}.
	\end{equation}
	Summing over $g_1,\dots,g_{i-1}$ collapses the middle factor to $I$, so \ref{eq:half-sandwich-second-commutation} becomes
	\[
		\E_u \E_{x_1,\dots,x_i} \sum_{g_i\,:\,g_i\ne\bot}
		\bra{\psi} \widehat G_{g_i}^{x_i} \ot (I-B_{[f(x_i)=g_i(u)]}^u)\ket{\psi}
		\le \nu_1
	\]
	by \ref{eq:ld-gbcon}. Thus the total error is
	\[
		\nu_1 + 2\sqrt{\nu_4}
		\le
		43km\left(\eps^{1/32}+\delta^{1/32}+\gamma^{1/32}+\zeta^{1/32}+(d/q)^{1/32}\right).
	\]
\end{proof}

\subsection{Consistency of $H$ with $A$}
\label{sec:consistency-of-h-with-a}

\begin{lemma}[Consistency of the pasted measurement with the line measurement]\label{lem:h-b-consistency}
	\lean{MIPStarRE.LDT.Pasting.hBConsistency}
	\leanok
	\uses{def:pasted-measurement, prop:ld-dnoteq, lem:ld-sandwich-line-one-point, lem:line-interpolation-averaging-estimates}
	The pasted submeasurement satisfies
	\[
		H_{[h|_u=f]} \ot I \simeq_{\nu_6} I \ot B_f^u,
	\]
	where
	\[
		\nu_6 = 44k^2m\left(\eps^{1/32}+\delta^{1/32}+\gamma^{1/32}+\zeta^{1/32}+(d/q)^{1/32}\right).
	\]
\end{lemma}
\begin{proof}
	\leanok
	\uses{def:pasted-measurement, prop:ld-dnoteq, lem:ld-sandwich-line-one-point,
		lem:line-interpolation-averaging-estimates}
	The inconsistency with $B^u$ expands as
	\begin{align}
		    & \E_u \sum_{f\ne f'} \bra{\psi} H_{[h|_u=f']} \ot B_f^u \ket{\psi}\nonumber                                   \\
		={} & \E_u \sum_h \sum_{f\ne h|_u} \bra{\psi} H_h \ot B_f^u \ket{\psi}\nonumber                                    \\
		={} & \E_u \E_{x_1,\dots,x_k} \sum_h \sum_{f\ne h|_u} \bra{\psi} H_h^{x_1,\dots,x_k} \ot B_f^u \ket{\psi}\nonumber \\
		={} & \E_u \E_{x_1,\dots,x_k} \sum_h \sum_{w:|w|\ge d+1} \sum_{(g_1,\dots,g_k)=h_w}
		\sum_{f\ne h|_u}
		\bra{\psi} \widehat H_{g_1,\dots,g_k}^{x_1,\dots,x_k} \ot B_f^u \ket{\psi}.
		\label{eq:h-b-consistency-expanded-sum}
	\end{align}
	If $|w|\ge d+1$ and $f \ne h|_u$, then some active coordinate $i$ must satisfy $g_i \ne \bot$ and $g_i(u)\ne f(x_i)$. Hence
	\begin{align}
		\eqref{eq:h-b-consistency-expanded-sum}
		\le{} &
		\E_u \E_{x_1,\dots,x_k} \sum_{g_1,\dots,g_k}
		\sum_{\substack{f:\exists i,\ g_i\ne\bot, \\ g_i(u)\ne f(x_i)}}
		\bra{\psi} \widehat H_{g_1,\dots,g_k}^{x_1,\dots,x_k} \ot B_f^u \ket{\psi}.
		\label{eq:h-b-consistency-bad-coordinate-union}
	\end{align}
	Passing from independent tuples to distinct tuples costs at most $k^2/q$ by Lemma~\ref{prop:ld-dnoteq}. A union bound over the offending coordinate then gives
	\[
		\eqref{eq:h-b-consistency-bad-coordinate-union}
		\le
		\sum_{i=1}^k
		\E_u \E_{y_1,\dots,y_k}
		\sum_{g_1,\dots,g_k\,:\,g_i\ne\bot}
		\sum_{a\ne g_i(u)}
		\bra{\psi} \widehat H_{g_1,\dots,g_k}^{y_1,\dots,y_k} \ot B_{[f(y_i)=a]}^u \ket{\psi}
		\le k\nu_5.
	\]
	Therefore
	\[
		\frac{k^2}{q} + k\nu_5
		\le
		44k^2m\left(\eps^{1/32}+\delta^{1/32}+\gamma^{1/32}+\zeta^{1/32}+(d/q)^{1/32}\right)
		= \nu_6.
	\]
\end{proof}

\begin{corollary}[Consistency of $H$ with $A$]\label{cor:h-a-consistency}
	\lean{MIPStarRE.LDT.Pasting.hAConsistency_submeas}
	\leanok
	\uses{lem:h-b-consistency}
	On average over $(u,x) \sim \F_q^{m+1}$,
	\[
		H_{[h(u,x)=a]} \ot I \simeq_{\nu} I \ot A_a^{u,x}.
	\]
\end{corollary}
\begin{proof}
	\leanok
	\uses{lem:h-b-consistency, prop:triangle-sub}
	Lemma~\ref{lem:h-b-consistency} implies
	\[
		H_{[h(u,x)=a]} \ot I \simeq_{\nu_6} I \ot B^u_{[f(x)=a]}.
	\]
	Applying \ref{prop:triangle-sub} to this relation and \ref{eq:ld-abcon} gives
	\[
		H_{[h(u,x)=a]} \ot I \simeq_{\nu_6+\sqrt{8m\eps+4\delta}} I \ot A_a^{u,x}.
	\]
	Since
	\[
		\sqrt{8m\eps+4\delta} + \nu_6
		\le
		47k^2m\left(\eps^{1/32}+\delta^{1/32}+\gamma^{1/32}+\zeta^{1/32}+(d/q)^{1/32}\right)
		\le \nu,
	\]
	the claim follows.
\end{proof}

\begin{lemma}[Point consistency from vertical-line consistency]
	\label{lem:h-a-consistency-from-line-consistency}
	\lean{MIPStarRE.LDT.Pasting.hAConsistency_submeas_from_lineConsistency}
	\leanok
	\uses{prop:triangle-sub}
	Let $H$ be any polynomial-valued submeasurement. Suppose that the strategy is
	good, that $\gamma,\zeta \ge 0$, that $k \ge 1$, and that the restriction of
	$H$ to every vertical line is consistent with the vertical-line measurement
	with error $\nu_6$. Then the pointwise evaluations of $H$ are consistent with
	Alice's point measurement with error $\nu$.
\end{lemma}
\begin{proof}
	\lean{MIPStarRE.LDT.Pasting.hAConsistency_submeas_from_lineConsistency}
	\leanok
	\uses{prop:triangle-sub}
	The proof transports the vertical-line consistency estimate to point
	evaluations and combines it with the good-strategy point-to-line comparison.
	The displayed numerical assumptions are used only to absorb the intermediate
	errors into the final parameter~$\nu$.
\end{proof}

\begin{corollary}[Consistency of the completed $H$ with $A$]
	\label{cor:h-a-consistency-completed}
	\lean{MIPStarRE.LDT.Pasting.hAConsistency_completed}
	\leanok
	\uses{cor:h-a-consistency, cor:ld-pasting-N-completeness,
		lem:h-a-consistency-completed-from-submeas}
	Let $H_{\mathrm{meas}}$ be the completion of $H$ at the fallback outcome.
	On average over $(u,x) \sim \F_q^{m+1}$,
	\[
		(H_{\mathrm{meas}})_{[h(u,x)=a]} \ot I \simeq_{\sigma} I \ot A_a^{u,x}.
	\]
\end{corollary}
\begin{proof}
	\leanok
	\uses{cor:h-a-consistency, cor:ld-pasting-N-completeness,
		lem:h-a-consistency-completed-from-submeas}
	Corollary~\ref{cor:h-a-consistency} gives
	\[
		H_{[h(u,x)=a]} \ot I \simeq_{\nu} I \ot A_a^{u,x}.
	\]
	Corollary~\ref{cor:ld-pasting-N-completeness} bounds the missing mass of $H$ by
	\[
		\kappa \cdot \left(1 + \frac{1}{100m}\right) + \nu + e^{-k/(80000m^2)}.
	\]
	Adding this completion mass to the $\nu$-bound gives
	\[
		\nu + \kappa \cdot \left(1 + \frac{1}{100m}\right) + \nu +
		e^{-k/(80000m^2)} = \sigma.
	\]
\end{proof}

\begin{lemma}[Completion from submeasurement consistency]
	\label{lem:h-a-consistency-completed-from-submeas}
	\lean{MIPStarRE.LDT.Pasting.hAConsistency_completed_from_submeas}
	\leanok
	\uses{cor:ld-pasting-N-completeness}
	Let $H$ be any polynomial-valued submeasurement. Suppose that the pointwise
	evaluations of $H$ are consistent with Alice's point measurement with error
	$\nu$, and that $H$ has the lower mass bound used in
	Corollary~\ref{cor:ld-pasting-N-completeness}. Then completing $H$ at the
	fallback polynomial preserves point consistency with the enlarged error
	appearing in the induction statement.
\end{lemma}
\begin{proof}
	\lean{MIPStarRE.LDT.Pasting.hAConsistency_completed_from_submeas}
	\leanok
	The completed measurement differs from the original submeasurement only at
	the fallback polynomial.  The additional term is bounded by the missing-mass
	estimate and then added to the given point-consistency error.
\end{proof}

\subsection{Completeness of $H$}
\label{sec:completeness-of-h-low-degree}

\begin{definition}[Outcomes of a fixed type]\label{def:outcomes-by-type}
	\lean{MIPStarRE.LDT.Pasting.outcomesByType, MIPStarRE.LDT.Pasting.globallyConsistentOutcomesByType, MIPStarRE.LDT.Pasting.nonglobalOutcomesByType}
	\leanok
	\uses{def:types}
	For a type $\tau \in \{0,1\}^k$, let $\mathsf{Outcomes}_\tau$ be the set of tuples $(g_1,\dots,g_k)$ such that $g_i \in \polyfunc{m}{q}{d}$ when $\tau_i=1$ and $g_i=\bot$ when $\tau_i=0$. If $x_1,\dots,x_k \in \F_q$, let $\mathsf{Global}_\tau(x)$ be the subset of $\mathsf{Outcomes}_\tau$ arising from restrictions of a single polynomial in $\polyfunc{m+1}{q}{d}$, and let $\overline{\mathsf{Global}_\tau(x)} = \mathsf{Outcomes}_\tau \setminus \mathsf{Global}_\tau(x)$.
\end{definition}

\begin{lemma}[Summing over all outcomes]\label{lem:over-all-outcomes}
	\lean{MIPStarRE.LDT.Pasting.overAllOutcomes}
	\leanok
	\uses{def:outcomes-by-type, prop:ld-dnoteq, lem:ld-sandwich-line-one-point, lem:schwartz-zippel-individual}
	If $x_1,\dots,x_k$ are sampled independently and uniformly from $\F_q$, then
	\[
		\bra{\psi} H \ot I \ket{\psi}
		\approx_{\nu_7}
		\E_{x_1,\dots,x_k} \sum_{\tau:\,|\tau|\ge d+1} \sum_{(g_1,\dots,g_k)\in \mathsf{Outcomes}_\tau}
		\bra{\psi} \widehat H_{g_1,\dots,g_k}^{x_1,\dots,x_k} \ot I \ket{\psi},
	\]
	where
	\[
		\nu_7 = 46k^2m\left(\eps^{1/32}+\delta^{1/32}+\gamma^{1/32}+\zeta^{1/32}+(d/q)^{1/32}\right).
	\]
\end{lemma}
\begin{proof}
	\leanok
	\uses{def:outcomes-by-type, prop:ld-dnoteq, lem:ld-sandwich-line-one-point,
		lem:schwartz-zippel-individual}
	Let $(y_1,\dots,y_k)\sim \mathsf{Distinct}_k$. By definition,
	\begin{align}
		\bra{\psi} H \ot I \ket{\psi}
		 & =
		\sum_h \bra{\psi} H_h \ot I \ket{\psi}\nonumber                                    \\
		 & =
		\E_{y_1,\dots,y_k} \sum_h \bra{\psi} H_h^{y_1,\dots,y_k} \ot I \ket{\psi}\nonumber \\
		 & =
		\E_{y_1,\dots,y_k} \sum_h \sum_{\tau:\,|\tau|\ge d+1}
		\sum_{(g_1,\dots,g_k)=h_\tau}
		\bra{\psi} \widehat H_{g_1,\dots,g_k}^{y_1,\dots,y_k} \ot I \ket{\psi}\nonumber    \\
		 & =
		\E_{y_1,\dots,y_k} \sum_{\tau:\,|\tau|\ge d+1}
		\sum_{(g_1,\dots,g_k)\in \mathsf{Global}_\tau(y)}
		\bra{\psi} \widehat H_{g_1,\dots,g_k}^{y_1,\dots,y_k} \ot I \ket{\psi}.
		\label{eq:sum-restricted-to-global-polynomial}
	\end{align}
	We now remove the restriction to globally consistent tuples:
	\begin{equation}\label{eq:remove-the-restriction}
		\eqref{eq:sum-restricted-to-global-polynomial}
		\approx_{\frac{k^2}{q}+k\nu_5+\frac{md}{q}}
		\E_{y_1,\dots,y_k} \sum_{\tau:\,|\tau|\ge d+1}
		\sum_{(g_1,\dots,g_k)\in \mathsf{Outcomes}_\tau}
		\bra{\psi} \widehat H_{g_1,\dots,g_k}^{y_1,\dots,y_k} \ot I \ket{\psi}.
	\end{equation}
	Since the right-hand side is at least the left-hand side, it suffices to bound the difference:
	\begin{align}
		\eqref{eq:remove-the-restriction} - \eqref{eq:sum-restricted-to-global-polynomial}
		 & =
		\E_{y_1,\dots,y_k} \sum_{\tau:\,|\tau|\ge d+1}
		\sum_{(g_1,\dots,g_k)\in \overline{\mathsf{Global}_\tau(y)}}
		\bra{\psi} \widehat H_{g_1,\dots,g_k}^{y_1,\dots,y_k} \ot I \ket{\psi}\nonumber \\
		 & =
		\E_u \E_{y_1,\dots,y_k} \sum_{\tau:\,|\tau|\ge d+1} \sum_f
		\sum_{(g_1,\dots,g_k)\in \overline{\mathsf{Global}_\tau(y)}}
		\bra{\psi} \widehat H_{g_1,\dots,g_k}^{y_1,\dots,y_k} \ot B_f^u \ket{\psi}.
		\label{eq:insert-line-measurement-resolution}
	\end{align}
	We next insert an indicator recording consistency along the sampled line:
	\begin{equation}\label{eq:insert-line-consistency-indicator}
		\eqref{eq:insert-line-measurement-resolution}
		\approx_{\frac{k^2}{q}+k\nu_5}
		\E_u \E_{y_1,\dots,y_k} \sum_{\tau:\,|\tau|\ge d+1} \sum_f
		\sum_{(g_1,\dots,g_k)\in \overline{\mathsf{Global}_\tau(y)}}
		\bra{\psi} \widehat H_{g_1,\dots,g_k}^{y_1,\dots,y_k} \ot B_f^u \ket{\psi}
		\cdot \bone[\forall i\in\tau,\ f(y_i)=g_i(u)].
	\end{equation}
	The discarded part is
	\begin{equation}\label{eq:discarded-line-inconsistency-term}
		\E_u \E_{y_1,\dots,y_k} \sum_{\tau:\,|\tau|\ge d+1} \sum_f
		\sum_{(g_1,\dots,g_k)\in \overline{\mathsf{Global}_\tau(y)}}
		\bra{\psi} \widehat H_{g_1,\dots,g_k}^{y_1,\dots,y_k} \ot B_f^u \ket{\psi}
		\cdot \bone[\exists i\in\tau,\ f(y_i)\ne g_i(u)],
	\end{equation}
	and after switching to independent tuples and applying Lemma~\ref{lem:ld-sandwich-line-one-point} coordinatewise, this contributes at most $k^2/q + k\nu_5$.

	Let $\mathsf{Consistent}_\tau(g,y,u)$ denote the event that there exists a degree-$d$ polynomial $f$ with $f(y_i)=g_i(u)$ for all $i\in\tau$. Then
	\[
		\bone[\forall i\in\tau,\ f(y_i)=g_i(u)] \le \bone[\mathsf{Consistent}_\tau(g,y,u)],
	\]
	so \ref{eq:insert-line-consistency-indicator} is bounded by
	\begin{align}
		    & \E_u \E_{y_1,\dots,y_k} \sum_{\tau:\,|\tau|\ge d+1}
		\sum_{(g_1,\dots,g_k)\in \overline{\mathsf{Global}_\tau(y)}}
		\bra{\psi} \widehat H_{g_1,\dots,g_k}^{y_1,\dots,y_k} \ot I \ket{\psi}
		\cdot \bone[\mathsf{Consistent}_\tau(g,y,u)]\nonumber     \\
		={} &
		\E_{y_1,\dots,y_k} \sum_{\tau:\,|\tau|\ge d+1}
		\sum_{(g_1,\dots,g_k)\in \overline{\mathsf{Global}_\tau(y)}}
		\bra{\psi} \widehat H_{g_1,\dots,g_k}^{y_1,\dots,y_k} \ot I \ket{\psi}
		\cdot \Pr_u[\mathsf{Consistent}_\tau(g,y,u)].
		\label{eq:consistent-indicator}
	\end{align}
	For a fixed globally inconsistent tuple, choose $d+1$ active coordinates and let $h^*$ be the unique interpolant through them. Since the tuple is not globally consistent, some active coordinate $i^*$ satisfies $g_{i^*} \ne h^*|_{y_{i^*}}$, and then Schwartz--Zippel gives
	\[
		\Pr_u[\mathsf{Consistent}_\tau(g,y,u)]
		\le
		\Pr_u[g_{i^*}(u)=h^*(u,y_{i^*})]
		\le
		\frac{md}{q}.
	\]
	Therefore \ref{eq:consistent-indicator} is at most $md/q$, which proves \ref{eq:remove-the-restriction}. A final application of Lemma~\ref{prop:ld-dnoteq} replaces distinct tuples by independent tuples, giving a total error of
	\[
		2\frac{k^2}{q} + \frac{md}{q} + k\nu_5
		\le
		46k^2m\left(\eps^{1/32}+\delta^{1/32}+\gamma^{1/32}+\zeta^{1/32}+(d/q)^{1/32}\right)
		= \nu_7.
	\]
\end{proof}

\begin{definition}[Truncated type sums]\label{def:truncated-type-sums}
	\lean{MIPStarRE.LDT.Pasting.truncatedTypeSums}
	\leanok
	\uses{def:types}
	For $1 \le \ell \le k+1$ and a tail type $\tau_{\ge \ell} \in \{0,1\}^{k-\ell+1}$, define
	\[
		S_{\tau_{\ge \ell}}
		=
		\sum_{\tau_{<\ell}\,:\,|\tau|\ge d+1}
		G^{|\tau_{<\ell}|}(I-G)^{(\ell-1)-|\tau_{<\ell}|}.
	\]
	This records the contribution of all prefixes that can still be completed to total weight at least $d+1$ once the tail $\tau_{\ge \ell}$ is fixed.
\end{definition}

\begin{lemma}[Recurrence for truncated type sums]\label{lem:truncated-type-sum-recurrence}
	\lean{MIPStarRE.LDT.Pasting.truncatedTypeSumRecurrence}
	\leanok
	\uses{def:truncated-type-sums}
	Each $S_{\tau_{\ge \ell}}$ is Hermitian, positive semidefinite, and bounded by $I$. Moreover, if $\tau_{>\ell}$ is a tail type of length $k-\ell$, then
	\[
		S_{\tau_{>\ell}}
		=
		\sum_{\tau_\ell \in \{0,1\}} S_{\tau_{\ge \ell}} G^{\tau_\ell}(I-G)^{1-\tau_\ell}.
	\]
\end{lemma}
\begin{proof}
	\leanok
	Every summand is a polynomial in the commuting positive operators $G$ and $I-G$, so $S_{\tau_{\ge \ell}}$ is Hermitian and positive semidefinite. Summing over all prefixes gives
	\[
		S_{\tau_{\ge \ell}} \le (G+(I-G))^{\ell-1}=I.
	\]
	Splitting the next type bit according to whether $\tau_\ell=1$ or $\tau_\ell=0$ gives the stated recurrence.
\end{proof}

\begin{lemma}[From the pasted sum to a polynomial in $G$]\label{lem:from-H-to-G}
	\lean{MIPStarRE.LDT.Pasting.fromHToG}
	\leanok
	\uses{cor:G-hat-facts, lem:q-sdd-complete-part-slice-bound, lem:commute-g-half-sandwich, def:pasted-measurement, def:truncated-type-sums, lem:truncated-type-sum-recurrence}
	Let $\ket{\psi_{\mathrm{bi}}}$ be a bipartite state on $\iota \otimes \iota$.  If
	$x_1,\dots,x_k$ are sampled independently and uniformly from $\F_q$, then
	\[
		\E_{x_1,\dots,x_k} \sum_{\tau:\,|\tau|\ge d+1} \sum_{(g_1,\dots,g_k)\in \mathsf{Outcomes}_\tau}
		\bra{\psi_{\mathrm{bi}}} \widehat H_{g_1,\dots,g_k}^{x_1,\dots,x_k} \ot I \ket{\psi_{\mathrm{bi}}}
		\approx_{\nu_8}
		\sum_{r=d+1}^k \binom{k}{r} \bra{\psi_{\mathrm{bi}}} G^r(I-G)^{k-r} \ot I
		\ket{\psi_{\mathrm{bi}}},
	\]
	where
	\[
		\nu_8 = 46k^2m\left(\gamma^{1/32}+\zeta^{1/32}+(d/q)^{1/32}\right).
	\]
\end{lemma}
\begin{proof}
	\leanok
	\uses{cor:G-hat-facts, lem:q-sdd-complete-part-slice-bound,
		lem:commute-g-half-sandwich, def:pasted-measurement,
		def:truncated-type-sums, lem:truncated-type-sum-recurrence}
	For a type $\tau \in \{0,1\}^k$, write $\tau_{<\ell}$, $\tau_{>\ell}$, $\tau_{\le \ell}$, and $\tau_{\ge \ell}$ for the obvious truncations, and similarly for tuples $g$. Also write
	\[
		\widehat G_{g_{\ge \ell}}^{x_{\ge \ell}}
		=
		\widehat G_{g_\ell}^{x_\ell}\cdots \widehat G_{g_k}^{x_k}.
	\]
	Then
	\begin{equation}\label{eq:G-recurrence}
		\widehat G_{g_{\ge \ell}}^{x_{\ge \ell}}
		=
		\widehat G_{g_\ell}^{x_\ell}\widehat G_{g_{>\ell}}^{x_{>\ell}},
	\end{equation}
	and
	\begin{equation}\label{eq:split-H-into-two-Gs}
		\widehat H_{g_{\ge \ell}}^{x_{\ge \ell}}
		=
		\widehat G_{g_{\ge \ell}}^{x_{\ge \ell}}
		\bigl(\widehat G_{g_{\ge \ell}}^{x_{\ge \ell}}\bigr)^\dagger.
	\end{equation}
	Write $\mathsf O_\tau = \mathsf{Outcomes}_\tau$. The main iterative step is that for each $1\le \ell \le k$,
	\begin{align}
		                                         & \E_{x_{\ge \ell}} \sum_{\tau:\,|\tau|\ge d+1} \sum_{g_{\ge \ell}\in \mathsf O_{\tau_{\ge \ell}}}
		\bra{\psi} \widehat H_{g_{\ge \ell}}^{x_{\ge \ell}}
		\ot \bigl(G^{|\tau_{<\ell}|}(I-G)^{(\ell-1)-|\tau_{<\ell}|}\bigr)\ket{\psi}\nonumber                                                        \\
		\approx_{2\sqrt{2\zeta}+2\sqrt{\nu_4}}{} &
		\E_{x_{>\ell}} \sum_{\tau:\,|\tau|\ge d+1} \sum_{g_{>\ell}\in \mathsf O_{\tau_{>\ell}}}
		\bra{\psi} \widehat H_{g_{>\ell}}^{x_{>\ell}}
		\ot \bigl(G^{|\tau_{\le \ell}|}(I-G)^{\ell-|\tau_{\le \ell}|}\bigr)\ket{\psi}.
		\label{eq:bernoulli-tail-stage-reduction}
	\end{align}
	Iterating \ref{eq:bernoulli-tail-stage-reduction} for $\ell=1,\dots,k$ produces the Bernoulli polynomial
	\[
		\sum_{r=d+1}^k \binom{k}{r} \bra{\psi} G^r(I-G)^{k-r} \ot I \ket{\psi},
	\]
	and the accumulated error is
	\[
		k(2\sqrt{2\zeta}+2\sqrt{\nu_4})
		\le
		46k^2m\left(\gamma^{1/32}+\zeta^{1/32}+(d/q)^{1/32}\right)
		= \nu_8.
	\]
	This corrects the displayed arithmetic in the paper: the iterated adjacent
	estimate contributes $k$ copies of the $2\sqrt{\nu_4}$ term, and the
	definition of $\nu_4$ already contains a factor $k^2$.

	To prove \ref{eq:bernoulli-tail-stage-reduction}, define for each $1\le \ell \le k+1$ and each tail type $\tau_{\ge \ell}$,
	\begin{equation}\label{eq:S-def}
		S_{\tau_{\ge \ell}}
		=
		\sum_{\tau_{<\ell}\,:\,|\tau|\ge d+1}
		G^{|\tau_{<\ell}|}(I-G)^{(\ell-1)-|\tau_{<\ell}|}.
	\end{equation}
	Then \ref{eq:bernoulli-tail-stage-reduction} may be rewritten as
	\begin{align}
		                                         & \E_{x_{\ge \ell}} \sum_{\tau_{\ge \ell}} \sum_{g_{\ge \ell}\in \mathsf O_{\tau_{\ge \ell}}}
		\bra{\psi} \widehat H_{g_{\ge \ell}}^{x_{\ge \ell}} \ot S_{\tau_{\ge \ell}} \ket{\psi}\nonumber                                        \\
		\approx_{2\sqrt{2\zeta}+2\sqrt{\nu_4}}{} &
		\E_{x_{>\ell}} \sum_{\tau_{>\ell}} \sum_{g_{>\ell}\in \mathsf O_{\tau_{>\ell}}}
		\bra{\psi} \widehat H_{g_{>\ell}}^{x_{>\ell}} \ot S_{\tau_{>\ell}} \ket{\psi}.
		\label{eq:bernoulli-tail-stage-reduction-expanded}
	\end{align}
	\begin{remark}[Lean proof-stage notation]
		In Lean, the type-restricted averaged submeasurement, the per-tail operator/mass, the aggregated stage quantity, and the final Bernoulli-tail scalar from this proof are recorded by the declarations above. Their indexing is shifted to the repository's 0-based convention: Lean stage $\ell$ is the paper's stage $\ell+1$.
	\end{remark}
	The operators $S_{\tau_{\ge \ell}}$ are Hermitian and positive semidefinite, and
	\begin{align}
		S_{\tau_{\ge \ell}}
		 & =
		\sum_{\tau_{<\ell}\,:\,|\tau|\ge d+1}
		G^{|\tau_{<\ell}|}(I-G)^{(\ell-1)-|\tau_{<\ell}|}\nonumber \\
		 & \le
		\sum_{\tau_{<\ell}}
		G^{|\tau_{<\ell}|}(I-G)^{(\ell-1)-|\tau_{<\ell}|}\nonumber \\
		 & =
		(G+(I-G))^{\ell-1}
		= I.
		\label{eq:S-bound}
	\end{align}
	For any $\tau_\ell \in \{0,1\}$,
	\begin{equation}\label{eq:explicit-formula-for-G-expectation}
		\E_{x_\ell} \sum_{g_\ell \in \mathsf O_{\tau_\ell}} \widehat G_{g_\ell}^{x_\ell}
		=
		G^{\tau_\ell}(I-G)^{1-\tau_\ell}.
	\end{equation}
	Therefore, for every $\tau_{>\ell}$,
	\begin{align}
		\sum_{\tau_\ell}
		S_{\tau_{\ge \ell}}
		\cdot
		\Bigl(\E_{x_\ell} \sum_{g_\ell \in \mathsf O_{\tau_\ell}} \widehat G_{g_\ell}^{x_\ell}\Bigr)
		 & =
		\sum_{\tau_\ell}
		S_{\tau_{\ge \ell}} G^{\tau_\ell}(I-G)^{1-\tau_\ell}\nonumber \\
		 & =
		S_{\tau_{>\ell}}.
		\label{eq:S-recurrence}
	\end{align}
	Moreover,
	\begin{align}
		S_{\tau_{\ge \ell}}
		\cdot
		\Bigl(\E_{x_\ell} \sum_{g_\ell \in \mathsf O_{\tau_\ell}} \widehat G_{g_\ell}^{x_\ell}\Bigr)
		\cdot
		S_{\tau_{\ge \ell}}
		\le
		\E_{x_\ell} \sum_{g_\ell \in \mathsf O_{\tau_\ell}} \widehat G_{g_\ell}^{x_\ell},
		\label{eq:S-sandwich}
	\end{align}
	because $S_{\tau_{\ge \ell}}$ commutes with $G$ and $(I-G)$ and is bounded by $I$.

	We now prove \ref{eq:bernoulli-tail-stage-reduction-expanded}. Writing $\widehat H$ as a sandwich of $\widehat G$ operators and moving the rightmost $\widehat G_{g_\ell}^{x_\ell}$ to the second tensor factor gives
	\begin{align}
		                          & \E_{x_{\ge \ell}} \sum_{\tau_{\ge \ell}} \sum_{(g_\ell,g_{>\ell})\in \mathsf O_{\tau_{\ge \ell}}}
		\bra{\psi} \widehat H_{g_{\ge \ell}}^{x_{\ge \ell}} \ot S_{\tau_{\ge \ell}}\ket{\psi}\nonumber                                \\
		={}                       &
		\E_{x_{\ge \ell}} \sum_{\tau_{\ge \ell}} \sum_{(g_\ell,g_{>\ell})\in \mathsf O_{\tau_{\ge \ell}}}
		\bra{\psi}
		\bigl(\widehat G_{g_{\ge \ell}}^{x_{\ge \ell}}
		(\widehat G_{g_{>\ell}}^{x_{>\ell}})^\dagger \widehat G_{g_\ell}^{x_\ell}\bigr)
		\ot S_{\tau_{\ge \ell}}\ket{\psi}\nonumber                                                                                    \\
		\approx_{\sqrt{2\zeta}}{} &
		\E_{x_{\ge \ell}} \sum_{\tau_{\ge \ell}} \sum_{(g_\ell,g_{>\ell})\in \mathsf O_{\tau_{\ge \ell}}}
		\bra{\psi}
		\bigl(\widehat G_{g_{\ge \ell}}^{x_{\ge \ell}}
		(\widehat G_{g_{>\ell}}^{x_{>\ell}})^\dagger\bigr)
		\ot \bigl(S_{\tau_{\ge \ell}}\widehat G_{g_\ell}^{x_\ell}\bigr)\ket{\psi}.
		\label{eq:move-rightmost-g-to-second-factor}
	\end{align}
	The corresponding Cauchy--Schwarz bound is
	\begin{multline}\label{eq:first-cauchy-schwarz-bound}
		\Bigl|
		\E_{x_{\ge \ell}} \sum_{\tau_{\ge \ell}} \sum_{(g_\ell,g_{>\ell})\in \mathsf O_{\tau_{\ge \ell}}}
		\bra{\psi}
		(\widehat G_{g_{\ge \ell}}^{x_{\ge \ell}} \ot S_{\tau_{\ge \ell}})
		\bigl((\widehat G_{g_{>\ell}}^{x_{>\ell}})^\dagger \ot I\bigr)
		(\widehat G_{g_\ell}^{x_\ell}\ot I - I \ot \widehat G_{g_\ell}^{x_\ell})
		\ket{\psi}
		\Bigr|\\
		\le
		\sqrt{
		\E_{x_{\ge \ell}} \sum_{\tau_{\ge \ell}} \sum_{(g_\ell,g_{>\ell})\in \mathsf O_{\tau_{\ge \ell}}}
		\bra{\psi} \widehat H_{g_{\ge \ell}}^{x_{\ge \ell}} \ot (S_{\tau_{\ge \ell}})^2 \ket{\psi}
		}
		\cdot
		\sqrt{
		\E_{x_{\ge \ell}} \sum_{\tau_{\ge \ell}} \sum_{(g_\ell,g_{>\ell})\in \mathsf O_{\tau_{\ge \ell}}}
		\bra{\psi}
		(\widehat G_{g_\ell}^{x_\ell}\ot I - I \ot \widehat G_{g_\ell}^{x_\ell})
		(\widehat H_{g_{>\ell}}^{x_{>\ell}} \ot I)
		(\widehat G_{g_\ell}^{x_\ell}\ot I - I \ot \widehat G_{g_\ell}^{x_\ell})
		\ket{\psi}
		}.
	\end{multline}
	The first square root is at most $1$ because $S_{\tau_{\ge \ell}}\le I$ and $\widehat H$ is a submeasurement, and the second is at most $2\zeta$ by \ref{eq:g-hat-self-consistency}.

	Next, commute the leftmost $\widehat G_{g_\ell}^{x_\ell}$ across the left half-sandwich:
	\begin{align}
		\eqref{eq:move-rightmost-g-to-second-factor}
		 & =
		\E_{x_{\ge \ell}} \sum_{\tau_{\ge \ell}} \sum_{(g_\ell,g_{>\ell})\in \mathsf O_{\tau_{\ge \ell}}}
		\bra{\psi}
		\bigl(\widehat G_{g_\ell}^{x_\ell}\widehat G_{g_{>\ell}}^{x_{>\ell}}(\widehat G_{g_{>\ell}}^{x_{>\ell}})^\dagger\bigr)
		\ot \bigl(S_{\tau_{\ge \ell}}\widehat G_{g_\ell}^{x_\ell}\bigr)\ket{\psi}\nonumber \\
		 & \approx_{\sqrt{\nu_4}}
		\E_{x_{\ge \ell}} \sum_{\tau_{\ge \ell}} \sum_{(g_\ell,g_{>\ell})\in \mathsf O_{\tau_{\ge \ell}}}
		\bra{\psi}
		\bigl(\widehat G_{g_{>\ell}}^{x_{>\ell}}\widehat G_{g_\ell}^{x_\ell}(\widehat G_{g_{>\ell}}^{x_{>\ell}})^\dagger\bigr)
		\ot \bigl(S_{\tau_{\ge \ell}}\widehat G_{g_\ell}^{x_\ell}\bigr)\ket{\psi}.
		\label{eq:commute-g-part-one}
	\end{align}
	The associated Cauchy--Schwarz estimate is
	\begin{multline}\label{eq:adjoint-commutator-cauchy-schwarz}
		\Bigl|
		\E_{x_{\ge \ell}} \sum_{\tau_{\ge \ell}} \sum_{(g_\ell,g_{>\ell})\in \mathsf O_{\tau_{\ge \ell}}}
		\bra{\psi}
		\bigl([\widehat G_{g_{>\ell}}^{x_{>\ell}},\widehat G_{g_\ell}^{x_\ell}] \ot I\bigr)
		\bigl((\widehat G_{g_{>\ell}}^{x_{>\ell}})^\dagger \ot (S_{\tau_{\ge \ell}}\widehat G_{g_\ell}^{x_\ell})\bigr)
		\ket{\psi}
		\Bigr|\\
		\le
		\sqrt{
		\E_{x_{\ge \ell}} \sum_{\tau_{\ge \ell}} \sum_{(g_\ell,g_{>\ell})\in \mathsf O_{\tau_{\ge \ell}}}
		\bra{\psi}
		[\widehat G_{g_{>\ell}}^{x_{>\ell}},\widehat G_{g_\ell}^{x_\ell}]
		[\widehat G_{g_{>\ell}}^{x_{>\ell}},\widehat G_{g_\ell}^{x_\ell}]^\dagger
		\ot I
		\ket{\psi}
		}
		\cdot
		\sqrt{
		\E_{x_{\ge \ell}} \sum_{\tau_{\ge \ell}} \sum_{(g_\ell,g_{>\ell})\in \mathsf O_{\tau_{\ge \ell}}}
		\bra{\psi}
		\widehat H_{g_{>\ell}}^{x_{>\ell}}
		\ot \bigl(\widehat G_{g_\ell}^{x_\ell}(S_{\tau_{\ge \ell}})^2\widehat G_{g_\ell}^{x_\ell}\bigr)
		\ket{\psi}
		}.
	\end{multline}
	The first square root is bounded by $\nu_4$ via Lemma~\ref{lem:commute-g-half-sandwich}, and the second by $1$ since $S_{\tau_{\ge \ell}}\le I$.

	Continue commuting:
	\begin{align}
		\eqref{eq:commute-g-part-one}
		\approx_{\sqrt{\nu_4}}
		\E_{x_{\ge \ell}} \sum_{\tau_{\ge \ell}} \sum_{(g_\ell,g_{>\ell})\in \mathsf O_{\tau_{\ge \ell}}}
		\bra{\psi}
		\bigl(\widehat G_{g_{>\ell}}^{x_{>\ell}}(\widehat G_{g_{>\ell}}^{x_{>\ell}})^\dagger \widehat G_{g_\ell}^{x_\ell}\bigr)
		\ot \bigl(S_{\tau_{\ge \ell}}\widehat G_{g_\ell}^{x_\ell}\bigr)\ket{\psi}.
		\label{eq:commute-g-part-two}
	\end{align}
	This time the Cauchy--Schwarz step is
	\begin{multline}\label{eq:second-cauchy-schwarz-bound}
		\Bigl|
		\E_{x_{\ge \ell}} \sum_{\tau_{\ge \ell}} \sum_{(g_\ell,g_{>\ell})\in \mathsf O_{\tau_{\ge \ell}}}
		\bra{\psi}
		\bigl(\widehat G_{g_{>\ell}}^{x_{>\ell}} \ot (S_{\tau_{\ge \ell}}\widehat G_{g_\ell}^{x_\ell})\bigr)
		\bigl([(\widehat G_{g_{>\ell}}^{x_{>\ell}})^\dagger,\widehat G_{g_\ell}^{x_\ell}] \ot I\bigr)
		\ket{\psi}
		\Bigr|\\
		\le
		\sqrt{
		\E_{x_{\ge \ell}} \sum_{\tau_{\ge \ell}} \sum_{(g_\ell,g_{>\ell})\in \mathsf O_{\tau_{\ge \ell}}}
		\bra{\psi}
		\widehat H_{g_{>\ell}}^{x_{>\ell}}
		\ot \bigl(S_{\tau_{\ge \ell}}\widehat G_{g_\ell}^{x_\ell}S_{\tau_{\ge \ell}}\bigr)
		\ket{\psi}
		}
		\cdot
		\sqrt{
		\E_{x_{\ge \ell}} \sum_{\tau_{\ge \ell}} \sum_{(g_\ell,g_{>\ell})\in \mathsf O_{\tau_{\ge \ell}}}
		\bra{\psi}
		[(\widehat G_{g_{>\ell}}^{x_{>\ell}})^\dagger,\widehat G_{g_\ell}^{x_\ell}]^\dagger
		[(\widehat G_{g_{>\ell}}^{x_{>\ell}})^\dagger,\widehat G_{g_\ell}^{x_\ell}]
		\ot I
		\ket{\psi}
		}.
	\end{multline}
	By \ref{eq:S-sandwich}, the first square root is at most $1$, and the second matches the first square root in \ref{eq:adjoint-commutator-cauchy-schwarz}, hence is at most $\nu_4$.

	Finally move the remaining $\widehat G_{g_\ell}^{x_\ell}$ to the second tensor factor:
	\begin{align}
		\eqref{eq:commute-g-part-two}
		\approx_{\sqrt{2\zeta}}
		\E_{x_{\ge \ell}} \sum_{\tau_{\ge \ell}} \sum_{(g_\ell,g_{>\ell})\in \mathsf O_{\tau_{\ge \ell}}}
		\bra{\psi}
		\widehat H_{g_{>\ell}}^{x_{>\ell}}
		\ot \bigl(S_{\tau_{\ge \ell}}\widehat G_{g_\ell}^{x_\ell}\widehat G_{g_\ell}^{x_\ell}\bigr)
		\ket{\psi}.
		\label{eq:remaining-g-on-second-factor}
	\end{align}
	The omitted Cauchy--Schwarz bound reuses the first square root from \ref{eq:second-cauchy-schwarz-bound} and the second square root from \ref{eq:first-cauchy-schwarz-bound}. Since $\widehat G$ is projective,
	\begin{align*}
		\eqref{eq:remaining-g-on-second-factor}
		 & =
		\E_{x_{\ge \ell}} \sum_{\tau_{\ge \ell}} \sum_{(g_\ell,g_{>\ell})\in \mathsf O_{\tau_{\ge \ell}}}
		\bra{\psi}
		\widehat H_{g_{>\ell}}^{x_{>\ell}}
		\ot \bigl(S_{\tau_{\ge \ell}}\widehat G_{g_\ell}^{x_\ell}\bigr)
		\ket{\psi} \\
		 & =
		\E_{x_{>\ell}} \sum_{\tau_{>\ell}} \sum_{g_{>\ell}\in \mathsf O_{\tau_{>\ell}}}
		\bra{\psi}
		\widehat H_{g_{>\ell}}^{x_{>\ell}}
		\ot
		\Bigl(\sum_{\tau_\ell}
		S_{\tau_{\ge \ell}}
		\Bigl(\E_{x_\ell}\sum_{g_\ell\in \mathsf O_{\tau_\ell}} \widehat G_{g_\ell}^{x_\ell}\Bigr)\Bigr)
		\ket{\psi} \\
		 & =
		\E_{x_{>\ell}} \sum_{\tau_{>\ell}} \sum_{g_{>\ell}\in \mathsf O_{\tau_{>\ell}}}
		\bra{\psi}
		\widehat H_{g_{>\ell}}^{x_{>\ell}} \ot S_{\tau_{>\ell}}
		\ket{\psi}
		\tag{by \ref{eq:S-recurrence}},
	\end{align*}
	which is exactly \ref{eq:bernoulli-tail-stage-reduction-expanded}.
\end{proof}

\begin{lemma}[Context specializations for the pasted measurement]\label{lem:pasting-context-specializations}
	\lean{MIPStarRE.LDT.Pasting.hAConsistency_submeas_of_context, MIPStarRE.LDT.Pasting.fromHToG_of_context}
	\leanok
	\uses{def:ld-pasting-context, cor:h-a-consistency, lem:from-H-to-G}
	The nontrivial-regime pasting context specializes the consistency of $H$ with $A$ and
	the conversion from the pasted sum to the polynomial in $G$ to the parameters
	and error bounds fixed in Definition~\ref{def:ld-pasting-context}.
\end{lemma}
\begin{proof}
	\lean{MIPStarRE.LDT.Pasting.hAConsistency_submeas_of_context}
	\lean{MIPStarRE.LDT.Pasting.fromHToG_of_context}
	\leanok
	\uses{def:ld-pasting-context, cor:h-a-consistency, lem:from-H-to-G}
	These are direct specializations of the two preceding construction theorems
	to the notation and numerical hypotheses packaged in the nontrivial-regime
	pasting context.
\end{proof}

\begin{lemma}[Matrix Chernoff bound for the Bernoulli polynomial]\label{lem:chernoff-bernoulli-matrix}
	\lean{MIPStarRE.LDT.Pasting.chernoffBernoulliMatrix}
	\leanok
	Let $0<\theta<1$ and integers $k,d>0$ with $k \ge 2d/\theta$. Define
	\[
		F(X) = \sum_{r=d+1}^k \binom{k}{r} X^r (I-X)^{k-r}.
	\]
	If $X$ is Hermitian, $0 \le X \le I$, and $\bra{\psi} X \ot I \ket{\psi} \ge 1-\kappa$, then
	\[
		\bra{\psi} F(X) \ot I \ket{\psi}
		\ge
		1 - \frac{\kappa}{1-\theta} - e^{-\theta^2 k/2}.
	\]
\end{lemma}
\begin{proof}
	\leanok
	Let $\rho$ be the reduced state of $\ket{\psi}$ on one prover's subsystem, and write the eigendecomposition $X = \sum_i \lambda_i \ket{v_i}\bra{v_i}$. This defines a probability distribution $\mu(i)=\bra{v_i}\rho\ket{v_i}$, with
	\[
		\E_{i\sim\mu} \lambda_i = \Tr(X\rho) = \bra{\psi} X \ot I \ket{\psi} \ge 1-\kappa.
	\]
	Equivalently, $\E_{i\sim\mu}(1-\lambda_i)\le \kappa$, so Markov's inequality
	shows that most of the spectral mass of $X$ lies on eigenvalues at least
	$\theta$:
	\begin{equation}\label{eq:spectral-mass-threshold}
		\Pr_{i\sim\mu}[\lambda_i > \theta] \ge 1-\frac{\kappa}{1-\theta}.
	\end{equation}
	For a scalar $p \in [d/k,1]$, the quantity
	\[
		F(p)=\sum_{r=d+1}^k \binom{k}{r}p^r(1-p)^{k-r}
	\]
	is the probability of observing at least $d+1$ successes in $k$
	Bernoulli$(p)$ trials. Applying the scalar Chernoff bound
	eigenvalue-by-eigenvalue therefore yields
	\begin{equation}\label{eq:bernoulli-tail-chernoff}
		F(p) \ge 1-\exp\Bigl(-2\Bigl(p-\frac{d}{k}\Bigr)^2k\Bigr).
	\end{equation}
	Hence
	\begin{align}
		\bra{\psi} F(X) \ot I \ket{\psi}
		 & = \Tr(F(X)\rho)\nonumber                                         \\
		 & = \sum_i F(\lambda_i)\bra{v_i}\rho\ket{v_i}\nonumber             \\
		 & = \E_{i\sim\mu} F(\lambda_i)\nonumber                            \\
		 & \ge \Pr_{i\sim\mu}[\lambda_i \ge \theta]\cdot F(\theta)\nonumber \\
		 & \ge \left(1-\frac{\kappa}{1-\theta}\right)\cdot F(\theta)
		\tag{by \ref{eq:spectral-mass-threshold}}\nonumber                           \\
		 & \ge \left(1-\frac{\kappa}{1-\theta}\right)\cdot
		\left(1-\exp\Bigl(-2\Bigl(\theta-\frac{d}{k}\Bigr)^2k\Bigr)\right).
		\label{eq:bernoulli-polynomial-spectral-lower-bound}
	\end{align}
	Since $(1-b)(1-c)\ge 1-b-c$ for $b,c\ge 0$, \ref{eq:bernoulli-polynomial-spectral-lower-bound} implies
	\[
		\bra{\psi} F(X) \ot I \ket{\psi}
		\ge
		1-\frac{\kappa}{1-\theta}
		- \exp\Bigl(-2\Bigl(\theta-\frac{d}{k}\Bigr)^2k\Bigr).
	\]
	Finally, $k \ge 2d/\theta$ implies $d/k \le \theta/2$, hence $(\theta-d/k)^2 \ge \theta^2/4$, which yields the claimed bound.
\end{proof}

\begin{corollary}[Completeness of the pasted submeasurement]\label{cor:ld-pasting-N-completeness}
	\lean{MIPStarRE.LDT.Pasting.ldPastingNCompleteness}
	\leanok
	\uses{lem:over-all-outcomes, lem:from-H-to-G, lem:pasting-context-specializations, lem:chernoff-bernoulli-matrix}
	If $k \ge 400md$, then
	\[
		\bra{\psi} H \ot I \ket{\psi}
		\ge
		1 - \kappa\left(1+\frac{1}{100m}\right) - \nu - e^{-k/(80000m^2)}.
	\]
\end{corollary}
\begin{proof}
	\leanok
	\uses{lem:over-all-outcomes, lem:from-H-to-G, lem:chernoff-bernoulli-matrix}
	Lemma~\ref{lem:over-all-outcomes} gives
	\[
		\bra{\psi} H \ot I \ket{\psi}
		\approx_{\nu_7}
		\E_{x_1,\dots,x_k} \sum_{\tau:\,|\tau|\ge d+1} \sum_{(g_1,\dots,g_k)\in \mathsf{Outcomes}_\tau}
		\bra{\psi} \widehat H_{g_1,\dots,g_k}^{x_1,\dots,x_k} \ot I \ket{\psi}.
	\]
	Lemma~\ref{lem:from-H-to-G} turns the right-hand side into
	\[
		\sum_{r=d+1}^k \binom{k}{r} \bra{\psi} G^r(I-G)^{k-r} \ot I \ket{\psi}
	\]
	up to error $\nu_8$. Apply Lemma~\ref{lem:chernoff-bernoulli-matrix} with $\theta = 1/(200m)$. Since $k\ge 400md$, the hypothesis $k \ge 2d/\theta$ is satisfied, and the resulting lower bound is
	\[
		1 - \frac{\kappa}{1-1/(200m)} - e^{-k/(80000m^2)}
		\ge
		1 - \kappa\left(1+\frac{1}{100m}\right) - e^{-k/(80000m^2)}.
	\]
	The paper then absorbs the approximation losses $\nu_7+\nu_8$ into $\nu$.
\end{proof}


%%FILE blueprint/src/chapter/ch10_induction.tex
\chapter{The main induction step}
\label{chap:induction}
\label{sec:induction}

\section{Restricting to one slice}

\begin{definition}[Appending a line at height $x$]\label{def:append-x}
	\lean{MIPStarRE.LDT.AxisParallelLine.appendAtHeight, MIPStarRE.LDT.AxisLinePolynomial.appendAtHeight, MIPStarRE.LDT.DiagonalLine.appendAtHeight}
	\leanok
  For $x \in \F_q$, let $\operatorname{append}_x(\ell)$ be the line in $\F_q^{m+1}$ obtained by appending the coordinate $x$ to every point of a line $\ell \subseteq \F_q^m$. If $f \colon \ell \to \F_q$, let $\operatorname{append}_x(f)$ be the polynomial on $\operatorname{append}_x(\ell)$ defined by $\operatorname{append}_x(f)(u,x)=f(u)$.
\end{definition}

\begin{definition}[$x$-restricted strategy]\label{def:restricted-strategy}
	\lean{MIPStarRE.LDT.MainInductionStep.xRestrictedStrategy, MIPStarRE.LDT.MainInductionStep.RestrictedSymStrat}
	\leanok
  \uses{def:append-x, def:symmetric-projective-strategy}
  Let $(\psi,A,B,L)$ be a symmetric strategy for the $(m+1,q,d)$ low individual degree test. For each $x \in \F_q$, the $x$-restricted strategy $(\psi,A^x,B^x,L^x)$ for the $(m,q,d)$ test is defined by
  \[
    (A^x)_a^u = A_a^{u,x},
    \qquad
    (B^x)_f^\ell = B_{\operatorname{append}_x(f)}^{\operatorname{append}_x(\ell)},
    \qquad
    (L^x)_f^\ell = L_{\operatorname{append}_x(f)}^{\operatorname{append}_x(\ell)}.
  \]
\end{definition}

\begin{lemma}[Restricted success probabilities]\label{lem:restricted-probabilities}
  \lean{MIPStarRE.LDT.MainInductionStep.restrictedProbabilities}
	\leanok
  \uses{def:restricted-strategy, def:good-strategy}
  Let $(\psi,A,B,L)$ be an $(\eps,\delta,\gamma)$-good symmetric strategy for the $(m+1,q,d)$ test. For each $x \in \F_q$, let $(\psi,A^x,B^x,L^x)$ be the restricted strategy and let $\eps_x,\delta_x,\gamma_x$ be its three failure probabilities. Then
  \[
    \E_x \eps_x \le \frac{m+1}{m}\eps,
    \qquad
    \E_x \delta_x \le \delta,
    \qquad
    \E_x \gamma_x \le \frac{m+1}{m}\gamma.
  \]
\end{lemma}
\begin{proof}
	\leanok
	The self-consistency test restricts exactly to the slice at height $x$. For the axis-parallel and diagonal-line tests, conditioning on the sampled line not being parallel to the last coordinate identifies the test with the corresponding restricted game.
\end{proof}

\begin{definition}[Average restricted errors]\label{def:average-restricted-errors}
	\lean{MIPStarRE.LDT.MainInductionStep.averageRestrictedAxisParallelError, MIPStarRE.LDT.MainInductionStep.averageRestrictedSelfConsistencyError, MIPStarRE.LDT.MainInductionStep.averageRestrictedDiagonalError}
	\leanok
	\uses{def:restricted-strategy}
	The average restricted errors are the expectations over $x \in \F_q$ of the
	axis-parallel, self-consistency, and restricted diagonal-line failure
	probabilities of the $x$-restricted strategies.
\end{definition}

\begin{lemma}[Weighted restricted-probability bounds]\label{lem:weighted-restricted-probability-bounds}
	\lean{MIPStarRE.LDT.MainInductionStep.RestrictedProbabilitiesStatement.ofWeightedBounds}
	\leanok
	\uses{lem:restricted-probabilities, def:average-restricted-errors}
	The weighted bounds from Lemma~\ref{lem:restricted-probabilities} imply the
	restricted-probability data used in the successor step.
\end{lemma}
\begin{proof}
	\leanok
	The weighted axis-parallel and diagonal-line estimates are divided by the
	transverse-direction weight, while the self-consistency estimate is exactly
	the average restricted self-consistency probability. These three estimates
	form the restricted-probability record used in the successor step.
\end{proof}

\section{The inductive construction}

\begin{theorem}[Self-improvement]\label{thm:self-improvement-in-induction-section}
  \lean{MIPStarRE.LDT.MainInductionStep.selfImprovementInInductionSection}
  \leanok
  \uses{def:good-strategy, def:polymeasurements, def:simeq, def:approx_delta}
  Let $(\psi,A,B,L)$ be an $(\eps,\delta,\gamma)$-good symmetric strategy for the $(m,q,d)$ low individual degree test. Let $G \in \polymeas{m}{q}{d}$ be a measurement with the following property:
  \begin{enumerate}
    \item (Consistency with~$A$): \label{item:si-inductive-G-consistency} On average over $u \sim \F_q^m$,
    \[
      A_a^u \ot I \simeq_\nu I \ot G_{[g(u)=a]}.
    \]
  \end{enumerate}
  Let
  \[
    \zeta = 3000m\left(\eps^{1/32}+\delta^{1/32}+(d/q)^{1/32}\right).
  \]
  Then there exists a projective submeasurement $H \in \polysub{m}{q}{d}$ with the following properties:
  \begin{enumerate}
    \item (Completeness): \label{item:si-inductive-completeness} If $H=\sum_h H_h$, then
    \[
      \bra{\psi} H \ot I \ket{\psi} \ge (1-\nu)-\zeta.
    \]
    \item (Consistency with~$A$): \label{item:si-inductive-A-consistency} On average over $u \sim \F_q^m$,
    \[
      A_a^u \ot I \simeq_\zeta I \ot H_{[h(u)=a]}.
    \]
    \item (Strong self-consistency): \label{item:si-inductive-self}
    \[
      H_h \ot I \approx_\zeta I \ot H_h.
    \]
    \item (Boundedness): \label{item:si-inductive-boundedness} There exists a positive semidefinite operator $Z$ such that
    \[
      \bra{\psi} Z \ot (I-H) \ket{\psi} \le \zeta
    \]
    and for each $h \in \polyfunc{m}{q}{d}$,
    \[
      Z \ge \E_u A_{h(u)}^u.
    \]
  \end{enumerate}
\end{theorem}
\begin{proof}
	\leanok
	\proves{thm:self-improvement-in-induction-section}
	\uses{thm:self-improvement, prop:two-notions-of-self-consistency}
	This is Theorem~\ref{thm:self-improvement}, with the conclusions reformulated in the notation used by the induction step. Since the output is projective, Lemma~\ref{prop:two-notions-of-self-consistency} converts the strong self-consistency statement to the usual bipartite formulation.
\end{proof}

\begin{remark}[Lean realization of induction self-improvement]\label{rem:self-improvement-in-induction-section-realization}
	The Lean declaration for
	Theorem~\ref{thm:self-improvement-in-induction-section} has the paper
	hypotheses, in particular that \(G\) is a complete polynomial measurement.
	The labelled theorem in
	\texttt{references/ldt-paper/inductive\_step.tex:249--286} states the input
	as a submeasurement, while the proved form in
	\texttt{references/ldt-paper/self\_improvement.tex:635--671} states it as a
	measurement.  The blueprint and Lean follow the proved measurement-valued
	form, which is also the form used at the point where the theorem is applied
	in \texttt{references/ldt-paper/inductive\_step.tex:461--485}.
	The proof applies
	\texttt{SelfImprovement.selfImprovement} to \(G\), then uses
	\texttt{selfImprovementInInductionSectionConclusion\_ofSelfImprovementConclusion}
	to transport the fields.  The Section~\ref{chap:self-improvement} SDP
	slackness input is supplied by the formalized theorem
	\texttt{sdp\_statement\_with\_slackness}; it is not an additional hypothesis of
	the induction-section theorem.
\end{remark}

\begin{theorem}[Pasting]\label{thm:ld-pasting-in-induction-section}
  \lean{MIPStarRE.LDT.MainInductionStep.ldPastingInInductionSection}
  \leanok
  \uses{def:good-strategy, def:polymeasurements, def:simeq, def:approx_delta}
  Let $(\psi,A,B,L)$ be an $(\eps,\delta,\gamma)$-good symmetric strategy for the $(m+1,q,d)$ low individual degree test. Let $\{G^x\}_{x \in \F_q}$ be projective submeasurements in $\polysub{m}{q}{d}$ with the following properties:
  \begin{enumerate}
    \item (Completeness): \label{item:ld-pasting-inductive-completeness} If $G=\E_x \sum_g G_g^x$, then
    \[
      \bra{\psi} G \ot I \ket{\psi} \ge 1-\kappa.
    \]
    \item (Consistency with~$A$): \label{item:ld-pasting-inductive-consistency} On average over $(u,x) \sim \F_q^{m+1}$,
    \[
      A_a^{u,x} \ot I \simeq_\zeta I \ot G_{[g(u)=a]}^x.
    \]
    \item (Strong self-consistency): \label{item:ld-pasting-inductive-self-consistency} On average over $x \sim \F_q$,
    \[
      G_g^x \ot I \approx_\zeta I \ot G_g^x.
    \]
    \item (Boundedness): \label{item:ld-pasting-inductive-boundedness} There are positive semidefinite operators $Z^x$ such that
    \[
      \E_x \bra{\psi} (I-G^x) \ot Z^x \ket{\psi} \le \zeta
    \]
    and for each $x \in \F_q$ and $g \in \polyfunc{m}{q}{d}$,
    \[
      Z^x \ge \E_u A_{g(u)}^{u,x}.
    \]
  \end{enumerate}
  Let $k \ge 400md$ and set
  \[
    \nu = 100k^2m\left(\eps^{1/32}+\delta^{1/32}+\gamma^{1/32}+\zeta^{1/32}+(d/q)^{1/32}\right)
  \]
  and
  \[
    \sigma = \kappa\left(1+\frac{1}{100m}\right) + 2\nu + e^{-k/(80000m^2)}.
  \]
  Then there exists a measurement $H \in \polymeas{m+1}{q}{d}$ with the following property:
  \begin{enumerate}
    \item (Consistency with~$A$): \label{item:ld-pasting-inductive-N-consistency} On average over $u \sim \F_q^{m+1}$,
    \[
      A_a^u \ot I \simeq_\sigma I \ot H_{[h(u)=a]}.
    \]
  \end{enumerate}
\end{theorem}
\begin{proof}
	\leanok
	\uses{thm:ld-pasting}
	This is exactly Theorem~\ref{thm:ld-pasting}, restated in the induction chapter
	because it is the pasting input used in the proof of
	Theorem~\ref{thm:main-induction}.  The Lean declaration linked above applies
	the unrestricted formal pasting theorem and transports its point-consistency
	field to the induction-section conclusion record.
\end{proof}

\begin{remark}[Unrestricted Lean form of induction-section pasting]\label{rem:ld-pasting-induction-unrestricted-lean-form}
	The Lean declaration linked to
	Theorem~\ref{thm:ld-pasting-in-induction-section} is now the unrestricted
	source-facing restatement.  It calls the formal Section~\ref{chap:pasting}
	theorem \texttt{Pasting.ldPasting} and therefore carries exactly the source
	hypotheses displayed above, together with the faithful formal boundary
	hypotheses needed to interpret the objects in Lean.

	The nontrivial-regime proof reduction remains represented in the Chapter~9
	declaration \texttt{Pasting.ldPastingNontrivial}.  That helper assumes
	\[
		\gamma \le 1,\qquad \zeta \le 1,\qquad d \le q,\qquad 0<d,\qquad 1\le k.
	\]
	These inequalities are not hypotheses of the paper theorem: in
	\texttt{references/ldt-paper/ld-pasting.tex}, lines 12--50 state the
	unrestricted result, while lines 52--55 use the inequalities only to reduce
	the proof to the nontrivial regime.  The formal theorem linked from the
	blueprint has already incorporated the complementary cases through
	\texttt{Pasting.ldPasting}.
\end{remark}

\begin{remark}[Lean self-improvement slice transport for the successor step]\label{rem:self-improvement-slice-transport}
	The successor step records the slice-wise self-improvement output in
	\texttt{SelfImprovementData}, and the answer-valued route records the analogous
	output in \texttt{AnswerSelfImprovementData}.  Their constructors call
	\texttt{selfImprovementInInductionSection}; they no longer carry a separate
	bundle of Section~\ref{chap:self-improvement} proof-stage data as an input.
	The \texttt{SliceStrategyTransport} structures are narrow formal interfaces: they
	assume concrete per-slice symmetric strategies and measurement-transport data,
	while averaged point-operator compatibility is derived structurally from
	point-measurement transport.  They then convert the remaining data into the
	ordinary or answer-valued self-improvement output.
\end{remark}

\begin{definition}[Pasting data for the successor step]\label{def:successor-pasting-data}
	\lean{MIPStarRE.LDT.MainInductionStep.AveragedPastingData.invokeLdPasting, MIPStarRE.LDT.MainInductionStep.assembleAveragedPastingData, MIPStarRE.LDT.MainInductionStep.assembleAveragedPastingDataOfSmallError}
	\lean{MIPStarRE.LDT.MainInductionStep.mainInductionFromAnswerStageDataOfSmallError}
	\lean{MIPStarRE.LDT.MainInductionStep.selfImprovementInInductionError_le_one_of_mainInductionError_lt_one}
	\leanok
	\uses{lem:weighted-restricted-probability-bounds, thm:ld-pasting-in-induction-section}
	The successor step forms the pasting input by combining the restricted
	probability estimates with the slice-wise self-improvement data.  The Lean
	construction applies the unrestricted induction-section pasting theorem directly,
	and the small-error constructor supplies the scalar side conditions
	\(\gamma\le 1\), \(\zeta\le 1\), and \(d\le q\) internally.  Thus the
	successor proof can call the small-error pasting constructor without carrying
	\(0<d\), \(1\le k\), \(\gamma\le 1\), \(\zeta\le 1\), or \(d\le q\) as
	separate proof inputs.  The answer-valued small-error assembly then converts
	the answer-valued slice data to this pasting interface and applies the
	checked stage-data theorem.
\end{definition}

\begin{lemma}[Base case of main induction]\label{lem:main-induction-base}
	\lean{MIPStarRE.LDT.MainInductionStep.mainInductionBaseCase}
	\lean{MIPStarRE.LDT.MainInductionStep.answerMainInductionBaseCase}
	\leanok
	\uses{def:good-strategy, def:polymeasurements, def:simeq}
	Under the hypotheses of Theorem~\ref{thm:main-induction} with $m=1$, there exists a measurement $G \in \polymeas{1}{q}{d}$ such that, on average over $u \sim \F_q$,
	\[
		A_a^u \ot I \simeq_\eps I \ot G_{[g(u)=a]},
	\]
	which is stronger than the bound required by Theorem~\ref{thm:main-induction} for $m=1$.
	Lean also proves the same one-dimensional construction for the
	answer-valued strategy interface used by the recursive-slice route.
\end{lemma}
\begin{proof}
	\leanok
	For $m=1$, there is only one axis-parallel line $\ell$ in $\F_q$, so the line measurement $B^\ell$ is already an element of $\polymeas{1}{q}{d}$. Because the strategy fails the axis-parallel lines test with probability at most $\eps$, the consistency bound follows directly from the axis-parallel consistency hypothesis.
\end{proof}

\begin{lemma}[Large-error branch of main induction]\label{lem:main-induction-large-error}
	\lean{MIPStarRE.LDT.MainInductionStep.mainInductionOfOneLeError}
	\lean{MIPStarRE.LDT.MainInductionStep.answerMainInductionOfOneLeError}
	\leanok
	\uses{def:polymeasurements, def:simeq}
	If the target error parameter in Theorem~\ref{thm:main-induction} is at least
	\(1\), then the conclusion follows from a distinguished trivial polynomial
	measurement.  Lean also proves the same large-error branch for the
	answer-valued strategy interface used by the recursive-slice route.
\end{lemma}
\begin{proof}
	\leanok
	The normalized bipartite consistency defect of two submeasurements is at most
	\(1\).  Hence any polynomial measurement, in particular the measurement
	concentrated on one polynomial, satisfies the claimed consistency relation when
	the allowed error is at least \(1\).
\end{proof}

\begin{theorem}[Main induction]\label{thm:main-induction}
	\lean{MIPStarRE.LDT.MainInductionStep.mainInduction}
	\leanok
	\uses{def:good-strategy, def:polymeasurements, def:simeq}
	This is the corrected numerical form of the theorem; the paper-gap note
	\cite{gap:issue-906-main-formal-k-bound} records the
	replacement of the printed bound \(k \ge md\) by \(k \ge 400md\).
  Let $(\psi,A,B,L)$ be an $(\eps,\delta,\gamma)$-good symmetric strategy for the $(m,q,d)$ low individual degree test. Let $k \ge 400md$ and set
  \[
    \nu = 1000k^2m^2\left(\eps^{1/1024}+\delta^{1/1024}+\gamma^{1/1024}+(d/q)^{1/1024}\right)
  \]
  and
  \[
    \sigma = m^2\left(\nu + e^{-k/(80000m^2)}\right).
  \]
  Then there exists a measurement $G \in \polymeas{m}{q}{d}$ such that, on average over $u \sim \F_q^m$,
  \[
    A_a^u \ot I \simeq_\sigma I \ot G_{[g(u)=a]}.
  \]
\end{theorem}
\begin{proof}
	\leanok
	\uses{lem:main-induction-base, lem:main-induction-large-error, prop:main-induction-successor-small-error-construction}
	The base case, the large-error branch, and the small-error successor
	construction are all formalized for the corrected large-\(k\) interface.
	The paper prints the weaker bound \(k\ge md\); the implication needed to
	apply the pasting theorem is false in the interval \(md\le k<400md\).
	Thus the factor \(400\) is recorded as a confirmed statement correction, as
	documented in \cite{gap:issue-906-main-formal-k-bound}.
\end{proof}

\begin{proposition}[Small-error successor construction for the induction theorem]\label{prop:main-induction-successor-small-error-construction}
	\lean{MIPStarRE.LDT.MainInductionStep.mainInductionSuccessorNext_ofSmallErrorConstruction}
	\leanok
	This is the nontrivial successor branch in the corrected large-\(k\)
	induction interface.  The ordinary successor theorem now reduces to the
	internal answer-valued induction theorem.  That theorem supplies the
	predecessor induction argument for the restricted slices by a genuine strong
	induction on the dimension.  The answer-valued pasting invocation is recorded
	in Proposition~\ref{prop:main-induction-successor-answer-valued-pasting},
	and the answer-valued self-improvement construction is checked by
	Proposition~\ref{prop:main-induction-successor-answer-slice-realization}.
	This is a construction internal to the proof, not an assumption of
	Theorem~\ref{thm:main-induction}.  The recursive slice route applies also
	when \(d=0\); the successor reduction no longer needs a separate degree-zero
	family construction.
\end{proposition}
\begin{proof}
	\leanok
	\uses{prop:main-induction-successor-predecessor-induction,
		prop:main-induction-successor-answer-slice-realization,
		prop:main-induction-successor-answer-valued-pasting}
	The checked assembly theorem
	\texttt{mainInductionSuccessorNext\_ofSmallErrorConstruction} reduces the
	ordinary branch to the answer-valued induction theorem through the
	answer-carrier construction.  The answer-valued induction theorem supplies
	the predecessor argument internally; the answer-valued slice self-improvement
	construction and the averaged family-field assembly are checked, and the
	answer-valued pasting construction below supplies the final pasting input.
\end{proof}

\begin{proposition}[Predecessor induction argument for successor slices]\label{prop:main-induction-successor-predecessor-induction}
	\uses{lem:restricted-probabilities}
	\lean{MIPStarRE.LDT.MainInductionStep.xRestrictedAnswerSymStratOfAnswer}
	\lean{MIPStarRE.LDT.MainInductionStep.answerSuccessorRestrictedProbabilities}
	\lean{MIPStarRE.LDT.MainInductionStep.answerMainInduction}
	\leanok
	In the small-error successor branch, the proof must apply the induction
	hypothesis in dimension \(m\) to the restricted strategies indexed by
	\(x\in\mathbb F_q\), with the same integer \(k\).  The successor large-\(k\)
	hypothesis implies the predecessor large-\(k\) side condition in the corrected
	formal interface, and the nontrivial small-error branch supplies \(k\ge1\).
	This is the genuine recursive use of the theorem; it is not an additional
	hypothesis on the successor strategy.  The answer-valued restriction map for
	an answer-valued successor strategy is now explicit in Lean; it preserves the
	full function-valued diagonal answer on the slice, so it does not use the
	ordinary low-degree realization that is unavailable for arbitrary
	answer-valued diagonal measurements.  The corresponding restricted-probability
	theorem for answer-valued successor strategies is also checked: the
	axis-parallel and diagonal slice averages are derived from the ambient
	answer-valued goodness assumptions, and the diagonal calculation uses the
	full function-valued answer rather than an ordinary polynomial-valued
	realization.  Lean now also checks the recursive application itself: after
	these restricted probabilities are constructed, the predecessor
	answer-valued induction hypothesis gives the main-induction conclusion for
	every answer-valued restricted slice.  The theorem
	\texttt{answerMainInduction} now supplies this predecessor hypothesis by
	strong induction on the dimension.  Thus the predecessor induction argument
	is a checked internal construction, not an additional theorem hypothesis.
\end{proposition}
\begin{proof}
	\leanok
	The Lean proof is contained in the simultaneous answer-valued induction
	theorem and its restricted-slice constructors.
\end{proof}

\begin{proposition}[Answer-valued slice self-improvement]\label{prop:main-induction-successor-answer-slice-realization}
	\uses{thm:self-improvement-in-induction-section}
	\lean{MIPStarRE.LDT.MainInductionStep.AnswerSelfImprovementData.slice_outputs_ofAnswerCarrier}
	\lean{MIPStarRE.LDT.MainInductionStep.AnswerSelfImprovementData.ofAnswerCarrier}
	\leanok
	The restricted slice strategy is represented in Lean by an answer-valued
	strategy.  The slice self-improvement data are constructed directly: replace
	only the diagonal part by an inert ordinary covariant measurement and apply
	the induction-section self-improvement theorem in the form whose hypotheses
	are the axis-parallel and point self-consistency estimates.  Since the
	Section~9 conclusion is independent of the diagonal-line error parameter,
	no low-degree support theorem for the answer-valued diagonal outcomes is
	needed for this step.  Such a support theorem would still be necessary for
	the stronger route which realizes the answer-valued diagonal measurement as
	an ordinary polynomial-valued covariant measurement, but that realization is
	not used here.
\end{proposition}
\begin{proof}
	\leanok
	\proves{prop:main-induction-successor-answer-slice-realization}
	\uses{thm:self-improvement-in-induction-section}
	The Lean proof is the carrier construction described in the statement.  The
	diagonal component of the carrier is a fixed covariant projective
	measurement; the hypotheses consumed by self-improvement are the
	axis-parallel and point self-consistency estimates of the answer-valued
	restricted slice.
\end{proof}

\begin{proposition}[Answer-valued pasting construction for the successor step]\label{prop:main-induction-successor-answer-valued-pasting}
	\uses{thm:ld-pasting-in-induction-section, prop:main-induction-successor-predecessor-induction, prop:main-induction-successor-answer-slice-realization}
	\lean{MIPStarRE.LDT.MainInductionStep.answerLdPastingInInductionSectionOfSmallError}
	\leanok
	After the recursive answer-valued slice measurements and the slice-wise
	self-improvement outputs have been averaged, Lean has a polynomial family
	which is complete, point-consistent with the actual answer-valued point
	measurement, strongly self-consistent, and bounded by the witnesses obtained
	from self-improvement.  The boundedness input is presently typed through the
	point-equivalent ordinary carrier, because the existing boundedness interface
	for pasting is an ordinary-strategy interface.

	The final induction-section pasting invocation is proved by first deriving
	the answer-valued Section~10 point-commutativity estimate from the
	answer-valued diagonal-line verifier relation, then invoking the Section~11
	scalar chain in the form which assumes this point-commutativity estimate.
	The proof does not replace the answer-valued strategy by the ordinary carrier
	for the diagonal-line test; the carrier is used only for the ordinary
	boundedness interface.  This construction is therefore a Lean-only internal
	construction for the simultaneous successor proof, not a new hypothesis of
	Theorem~\ref{thm:main-induction} and not a replacement for
	Theorem~\ref{thm:ld-pasting-in-induction-section}.
\end{proposition}
\begin{proof}
	\leanok
	\uses{thm:ld-pasting-in-induction-section}
	The checked reduction
	\texttt{answerMainInductionSuccessorNext\_ofRecursiveHypothesisAndAnswerPasting}
	shows that, with the predecessor answer-valued induction hypothesis in the
	successor context, the successor branch reduces to the answer-valued pasting
	theorem
	\texttt{answerLdPastingInInductionSectionOfSmallError}.  That theorem is now
	proved from the answer-valued commutativity route, the degree-zero branch, and
	the scalar absorption estimate.
\end{proof}

\paragraph{Checked status of the corrected induction interface.}
	In Lean, the base case is formalized through
	\lean{MIPStarRE.LDT.MainInductionStep.mainInductionBaseCase}.
	The answer-valued analogue needed by the simultaneous recursive route uses
	the same one-dimensional axis-parallel-line construction and does not involve
	the diagonal answer interface.
	The printed source range \(md\le k<400md\) is excluded from the corrected
	theorem statement: the missing implication to the pasting size hypothesis is
	a confirmed statement gap, and Theorem~\ref{thm:main-induction} is stated
	with the corrected large-\(k\) hypothesis.  The successor step of
	\lean{MIPStarRE.LDT.MainInductionStep.mainInduction} is isolated in its
	native \(m\to m+1\) form as the corrected large-\(k\) theorem
	\lean{MIPStarRE.LDT.MainInductionStep.mainInductionSuccessorNext}.  Its
	nontrivial branch is named as the small-error construction
	\lean{MIPStarRE.LDT.MainInductionStep.mainInductionSuccessorNext_ofSmallErrorConstruction}.
	The
	non-base branch
	\lean{MIPStarRE.LDT.MainInductionStep.mainInductionSuccessor} only decomposes
	an arbitrary parameter bundle with \(m\ne1\) into this successor form.  Neither
	declaration takes restricted-probability, recursive-slice, self-improvement, or
	pasting records as hypotheses.  The nearby formalization still names the stage data
	\lean{MIPStarRE.LDT.MainInductionStep.SliceRestrictionData, MIPStarRE.LDT.MainInductionStep.PerSliceInductionData, MIPStarRE.LDT.MainInductionStep.SelfImprovementData, MIPStarRE.LDT.MainInductionStep.AveragedPastingData}
	\lean{MIPStarRE.LDT.MainInductionStep.SelfImprovementData.SliceStrategyTransport, MIPStarRE.LDT.MainInductionStep.SelfImprovementData.ofSliceStrategyTransport}
	\lean{MIPStarRE.LDT.MainInductionStep.SelfImprovementData.SliceStrategyTransport.averagedPoint_eq_of_pointMeasurement_eq, MIPStarRE.LDT.MainInductionStep.SelfImprovementData.SliceStrategyTransport.ofPointMeasurementEq}
	that mirror the four paper stages $\text{restrict} \to \text{induct} \to
	\text{self-improve} \to \text{paste}$, but no theorem with these data as
	additional source-level hypotheses is advertised as
	Theorem~\ref{thm:main-induction}. For the answer-valued restriction route
	used in the successor branch of the main formal theorem, Lean now also exposes
	\lean{MIPStarRE.LDT.MainInductionStep.AnswerMainInductionConclusion, MIPStarRE.LDT.MainInductionStep.AnswerMainInductionHypothesis, MIPStarRE.LDT.MainInductionStep.AnswerPerSliceInductionData.ofMainInductionHypothesis},
	\lean{MIPStarRE.LDT.MainInductionStep.AnswerPerSliceInductionData.ofPerSliceInductionData, MIPStarRE.LDT.MainInductionStep.PerSliceInductionData.ofAnswer, MIPStarRE.LDT.MainInductionStep.SelfImprovementData.ofAnswerForPerSliceInductionData, MIPStarRE.LDT.MainInductionStep.SelfImprovementData.ofAnswer},
	which record the predecessor induction hypothesis for the restricted
	strategies, transport induction data in both directions between the
	answer-valued and ordinary slice interfaces, and then forget only the
	answer-valued diagonal alphabet before the existing pasting assembly. The weighted
	restricted-probability estimates are now public via
	\lean{MIPStarRE.LDT.MainInductionStep.weighted_axisParallel_bound, MIPStarRE.LDT.MainInductionStep.weighted_diagonal_bound, MIPStarRE.LDT.MainInductionStep.restrictedProbabilities},
	The Lean proof now closes the large-error branch by
	\lean{MIPStarRE.LDT.MainInductionStep.mainInductionOfOneLeError}.  The
	answer-valued large-error branch is checked by
	\lean{MIPStarRE.LDT.MainInductionStep.answerMainInductionOfOneLeError}.  The
	internal answer-valued induction theorem is
	\lean{MIPStarRE.LDT.MainInductionStep.answerMainInduction}; it is now a
	genuine strong induction on the dimension.  The surrounding development also
	contains the internal reductions
	\lean{MIPStarRE.LDT.MainInductionStep.mainInductionSuccessorNext_ofAnswerCarrier},
	\lean{MIPStarRE.LDT.MainInductionStep.mainInductionSuccessorNext_ofAnswerCarrierFromSuccessorBound},
	which derive the successor conclusion from the answer-valued
	predecessor induction hypothesis.  The answer-carrier forms construct the
	answer-valued self-improvement data directly from the restricted slices.  The
	successor-bound form also derives the predecessor
	large-\(k\) and \(k\ge1\) side conditions from the successor large-\(k\)
	hypothesis together with the small-error assumption, so these elementary
	arithmetic conditions are no
	longer additional hypotheses in the successor proof.  The answer-carrier
	route shows that the separate degree-zero
	family route is not needed for the successor reduction: the predecessor
	induction hypothesis is used for the restricted slices also when \(d=0\),
	while \(k\ge1\) comes from the nontrivial error branch.  The obsolete
	slice-transport and degree-split family routes have therefore been removed
	from the checked interface.  Thus the ordinary small-error construction theorem
	\lean{MIPStarRE.LDT.MainInductionStep.mainInductionSuccessorNext_ofSmallErrorConstruction}
	is now proved from the internal answer-valued induction theorem, and the
	answer-valued pasting theorem
	\lean{MIPStarRE.LDT.MainInductionStep.answerLdPastingInInductionSectionOfSmallError}
	is also proved by the answer-valued commutativity route, the degree-zero
	branch, and the scalar absorption estimate.  None of these constructions is a
	separate source hypothesis.  The former answer-valued self-improvement interface mismatch is
	resolved for the induction step by the checked carrier construction.  The
	older ordinary-realization constructors
	\lean{MIPStarRE.LDT.MainInductionStep.AnswerSelfImprovementData.SliceStrategyTransport.averagedPoint_eq_of_pointMeasurement_eq, MIPStarRE.LDT.MainInductionStep.AnswerSelfImprovementData.SliceStrategyTransport.good_of_restrictedGood, MIPStarRE.LDT.MainInductionStep.AnswerSelfImprovementData.SliceStrategyTransport.ofMeasurementEq}
	\lean{MIPStarRE.LDT.MainInductionStep.AnswerSelfImprovementData.ofSelfImprovementInInductionSection, MIPStarRE.LDT.MainInductionStep.AnswerSelfImprovementData.ofSliceStrategyTransport}
	remain available and show that, under the stronger ordinary-realization
	route, the Section~9 theorem can also be applied slice by slice through the
	checked transport interface.  Such a route still requires a low-degree support theorem
	for the answer-valued diagonal outcomes; otherwise the ordinary polynomial
	diagonal measurement cannot satisfy the required covariance after rebasing.
	The answer-valued ambient point-consistency averaging step is now checked by
	\lean{MIPStarRE.LDT.MainInductionStep.answer_family_pointConsistencyError_eq_avg}
	and
		\lean{MIPStarRE.LDT.MainInductionStep.answer_family_consistency_of_slice_bounds}.
		These lemmas prove the last-coordinate reindexing identity and the averaged
		consistency estimate for an ambient answer-valued strategy.  The
		answer-valued successor scalar averaging estimates are checked by
		\lean{MIPStarRE.LDT.MainInductionStep.average_answerSuccessorSliceSelfImprovementError_le},
		\lean{MIPStarRE.LDT.MainInductionStep.average_answerSuccessorSliceMainInductionNu_le}
		and
		\lean{MIPStarRE.LDT.MainInductionStep.average_answerSuccessorSliceMainInductionError_le};
		these are the \(\zeta_x\), \(\nu_x\), and \(\sigma_x\) Jensen and
		conditioning bounds for the actual
		\texttt{AnswerSymStrat} restricted successor profiles.  The theorem
		\lean{MIPStarRE.LDT.MainInductionStep.answerSuccessorRecursiveSliceMeasurements_ofMainInductionHypothesis}
		now extracts the recursive slice measurements from the predecessor
		answer-valued induction hypothesis and proves the averaged
		\(\sigma_x\)-bound.  The theorem
		\lean{MIPStarRE.LDT.MainInductionStep.answerSuccessorSelfImprovementOutputs_ofMainInductionHypothesis}
		then applies the axis-parallel/self-consistency form of self-improvement
		slice by slice, producing the projective \(\widehat G^x\)'s and
		witnesses \(Z^x\) together with the averaged \(\zeta_x\) and
		\(\sigma_x\) bounds.  These slice outputs are assembled at the
		ambient answer-valued level by
		\lean{MIPStarRE.LDT.MainInductionStep.answerSuccessorAveragedFamilyFields_ofMainInductionHypothesis},
		which proves averaged completeness, point consistency with the actual
		answer-valued point measurement, strong self-consistency, the carrier-typed
		slice boundedness input, and the \(\kappa\)- and \(\zeta\)-bounds.  This
		is not a final pasting theorem: the boundedness field uses the
		point-equivalent ordinary carrier only because the current boundedness
		interface is typed for ordinary strategies, and no claim is made that the
		carrier's dummy diagonal measurement satisfies the answer-valued diagonal
		test.  The answer-valued pasting invocation is proved as
		\lean{MIPStarRE.LDT.MainInductionStep.answerLdPastingInInductionSectionOfSmallError}.
		The checked reduction
		\lean{MIPStarRE.LDT.MainInductionStep.answerMainInductionSuccessorNext_ofRecursiveHypothesisAndAnswerPasting}
		shows that, with the predecessor answer-valued induction hypothesis
		available in the successor context, the successor branch reduces to this
		answer-valued pasting theorem.  The conditional reduction is an internal
		lemma used inside the strong-induction proof, not a proof-level
		formalization link for a source theorem.  The remaining
		ordinary family-level averaging components are also now isolated as checked
		lemmas:
	\lean{MIPStarRE.LDT.MainInductionStep.idxPolyFamily_complete_of_slice_bounds}
	\lean{MIPStarRE.LDT.MainInductionStep.idxPolyFamily_stronglySelfConsistent_of_slice_bounds}
	and
	\lean{MIPStarRE.LDT.MainInductionStep.idxPolyFamily_sliceBoundednessInput_of_slice_bounds}.
	For answer-valued self-improvement data attached to an ordinary ambient
	strategy, the corresponding pasting fields are checked directly by
	\lean{MIPStarRE.LDT.MainInductionStep.AnswerSelfImprovementData.complete_of_slice_bounds}
	\lean{MIPStarRE.LDT.MainInductionStep.AnswerSelfImprovementData.consistentWithPoints_of_slice_bounds}
	\lean{MIPStarRE.LDT.MainInductionStep.AnswerSelfImprovementData.stronglySelfConsistent_of_slice_bounds}
	and
	\lean{MIPStarRE.LDT.MainInductionStep.AnswerSelfImprovementData.sliceBoundednessInput_of_slice_bounds}.
	The direct ordinary-ambient answer-valued assembly
	\lean{MIPStarRE.LDT.MainInductionStep.mainInductionFromAnswerStageDataOfSmallErrorDirect}
	now invokes the induction-section pasting theorem from these fields and the
	answer-valued averaged scalar estimates; the exported theorem
	\lean{MIPStarRE.LDT.MainInductionStep.mainInductionFromAnswerStageDataOfSmallError}
	calls this direct proof.
	The remaining ambient answer-valued assembly is therefore the
	\texttt{AnswerSymStrat} analogue needed inside the simultaneous successor
	proof, together with the final pasting interface, not the ordinary
	answer-valued slice averaging calculation.
	The answer-valued pasting invocation and predecessor induction argument are
	now supplied internally in the proof of Theorem~\ref{thm:main-induction};
	they are not additional hypotheses about the successor strategy.
	We may assume $\eps,\delta,\gamma,d/q \le 1$, since otherwise $\nu \ge 1$ and the bound is trivial.

	The proof is by induction on $m$. The base case $m=1$ is Lemma~\ref{lem:main-induction-base}.

	For the step from $m$ to $m+1$, let $(\psi,A,B,L)$ be an $(\eps,\delta,\gamma)$-good symmetric strategy for the $(m+1,q,d)$ test, and let $k \ge 400(m+1)d$. For each $x \in \F_q$, let $(\psi,A^x,B^x,L^x)$ be the $x$-restricted strategy, with failure parameters $\eps_x,\delta_x,\gamma_x$. Apply the induction hypothesis to each restricted strategy with the same integer $k$ and define
	\[
		\nu_x = 1000k^2m^2\left(\eps_x^{1/1024}+\delta_x^{1/1024}+\gamma_x^{1/1024}+(d/q)^{1/1024}\right),
		\qquad
		\sigma_x = m^2\left(\nu_x + e^{-k/(80000m^2)}\right).
	\]
	This yields a measurement $G^x \in \polymeas{m}{q}{d}$ such that
	\[
		(A^x)_a^u \ot I \simeq_{\sigma_x} I \ot G_{[g(u)=a]}^x.
	\]

	Next define
	\[
		\zeta_x = 3000m\left(\eps_x^{1/32}+\delta_x^{1/32}+(d/q)^{1/32}\right).
	\]
	Theorem~\ref{thm:self-improvement-in-induction-section} produces a projective submeasurement $\widehat G^x \in \polysub{m}{q}{d}$ such that, for each $x$:
	\begin{enumerate}
		\item if $\widehat G^x=\sum_g \widehat G_g^x$, then
		      \[
			      \bra{\psi} \widehat G^x \ot I \ket{\psi} \ge (1-\sigma_x)-\zeta_x;
		      \]
		\item on average over $u \sim \F_q^m$,
		      \[
			      (A^x)_a^u \ot I \simeq_{\zeta_x} I \ot \widehat G_{[g(u)=a]}^x;
		      \]
		\item
		      \[
			      \widehat G_g^x \ot I \approx_{\zeta_x} I \ot \widehat G_g^x;
		      \]
		\item there exists a positive semidefinite operator $Z^x$ such that
		      \[
			      \bra{\psi} Z^x \ot (I-\widehat G^x) \ket{\psi} \le \zeta_x
		      \]
		      and for each $g \in \polyfunc{m}{q}{d}$,
		      \[
			      Z^x \ge \E_u (A^x)_{g(u)}^u.
		      \]
	\end{enumerate}

	We now average the slice-dependent parameters. By concavity of $\alpha \mapsto \alpha^c$ for $c \le 1$ and Lemma~\ref{lem:restricted-probabilities},
	\begin{align*}
		\E_x \nu_x
		 & \le 1000k^2m^2\left((\E_x \eps_x)^{1/1024}+(\E_x \delta_x)^{1/1024}+(\E_x \gamma_x)^{1/1024}+(d/q)^{1/1024}\right)                          \\
		 & \le 1000k^2m^2\left(\left(\frac{m+1}{m}\eps\right)^{1/1024}+\delta^{1/1024}+\left(\frac{m+1}{m}\gamma\right)^{1/1024}+(d/q)^{1/1024}\right) \\
		 & \le 1000k^2(m+1)^2\left(\eps^{1/1024}+\delta^{1/1024}+\gamma^{1/1024}+(d/q)^{1/1024}\right).
	\end{align*}
	Call this final quantity $\nu$. Then, defining
	\[
		\sigma = m^2\left(\nu + e^{-k/(80000m^2)}\right),
	\]
	we have
	\[
		\sigma \ge \E_x \sigma_x.
	\]
	Similarly,
	\begin{align*}
		\E_x \zeta_x
		 & \le 3000m\left((\E_x \eps_x)^{1/32}+(\E_x \delta_x)^{1/32}+(d/q)^{1/32}\right)         \\
		 & \le 3000m\left(\left(\frac{m+1}{m}\eps\right)^{1/32}+\delta^{1/32}+(d/q)^{1/32}\right) \\
		 & \le 3000(m+1)\left(\eps^{1/32}+\delta^{1/32}+(d/q)^{1/32}\right).
	\end{align*}
	Call this last quantity $\zeta$. For later use,
	\begin{align}
		\zeta
		 & = 3000(m+1)\left(\eps^{1/32}+\delta^{1/32}+(d/q)^{1/32}\right) \notag                              \\
		 & \le 1000k^2(m+1)^2\left(\eps^{1/32}+\delta^{1/32}+(d/q)^{1/32}\right) \notag                       \\
		 & \le 1000k^2(m+1)^2\left(\eps^{1/1024}+\delta^{1/1024}+\gamma^{1/1024}+(d/q)^{1/1024}\right) \notag \\
		 & = \nu. \label{eq:zeta-smaller-than-nu}
	\end{align}

	The averaged slice family $\{\widehat G^x\}_{x \in \F_q}$ therefore satisfies the four hypotheses of Theorem~\ref{thm:ld-pasting-in-induction-section} with completeness parameter $\kappa=\sigma+\zeta$. To verify the new $\nu$ bound in that theorem, observe that $(3000)^{1/32} \le 2$, so
	\[
		\zeta^{1/32} \le 2(m+1)\left(\eps^{1/1024}+\delta^{1/1024}+(d/q)^{1/1024}\right),
	\]
	and hence
	\begin{align*}
		 & 100k^2m\left(\eps^{1/32}+\delta^{1/32}+\gamma^{1/32}+\zeta^{1/32}+(d/q)^{1/32}\right)       \\
		 & \le 1000k^2(m+1)^2\left(\eps^{1/1024}+\delta^{1/1024}+\gamma^{1/1024}+(d/q)^{1/1024}\right) \\
		 & = \nu.
	\end{align*}
	Therefore Theorem~\ref{thm:ld-pasting-in-induction-section} gives a pasted measurement $H \in \polymeas{m+1}{q}{d}$ such that, on average over $u \sim \F_q^{m+1}$,
	\[
		A_a^u \ot I \simeq_{\sigma^*} I \ot H_{[h(u)=a]},
	\]
	where
	\[
		\sigma^* = (\sigma+\zeta)\left(1+\frac{1}{100m}\right) + 2\nu + e^{-k/(80000m^2)}.
	\]
	It remains to compare $\sigma^*$ with the next-stage target. The pasting
	parameter in Theorem~\ref{thm:ld-pasting-in-induction-section} is bounded by
	$\nu/5$ in this regime. Using this sharper estimate and
	\eqref{eq:zeta-smaller-than-nu}, we get
	\begin{align*}
		\sigma^*
		 & \le \left(m^2\left(\nu+e^{-k/(80000m^2)}\right)+\nu\right)
		      \left(1+\frac{1}{100m}\right) + \frac{2}{5}\nu
		      + e^{-k/(80000m^2)} \\
		 & = \left((m^2+1)\left(1+\frac{1}{100m}\right)+\frac{2}{5}\right)\nu
		     + \left(m^2\left(1+\frac{1}{100m}\right)+1\right)e^{-k/(80000m^2)}.
	\end{align*}
	For every $m\ge 1$,
	\[
		(m^2+1)\left(1+\frac{1}{100m}\right)+\frac{2}{5}\le (m+1)^2,
		\qquad
		m^2\left(1+\frac{1}{100m}\right)+1\le (m+1)^2.
	\]
	Therefore
	\[
		\sigma^* \le (m+1)^2\left(\nu+e^{-k/(80000m^2)}\right)
		\le (m+1)^2\left(\nu+e^{-k/(80000(m+1)^2)}\right),
	\]
	which is exactly the bound required for dimension $m+1$.

\begin{remark}[Lean status of the successor-step assembly]\label{rem:main-induction-successor-assembly}
	The Lean declaration for
	Theorem~\ref{thm:main-induction} has the corrected
	large-\(k\) theorem hypotheses and no additional stage-data hypotheses.  Its
	native successor step \texttt{mainInductionSuccessorNext} is now proved.  The
	weighted slice bounds feeding Lemma~\ref{lem:restricted-probabilities} remain
	exposed through the public restriction-package constructor above, but the
	recursive slice witnesses, the slice-wise self-improvement outputs, and the
	answer-valued pasting comparison are derived inside the successor proof.  They
	are not represented by a separate theorem whose assumptions are added to the
	paper statement.
\end{remark}

\section{Proof of the main formal theorem}

\begin{lemma}[Heterogeneous role-register symmetrization]\label{lem:role-register-symmetrization}
	\lean{MIPStarRE.LDT.ProjStrat.roleRegisterSymmStrategy, MIPStarRE.LDT.ProjStrat.roleRegisterSymmStrategy_is_good_three_mul, MIPStarRE.LDT.ProjStrat.roleRegisterSymmState_isNormalized}
	\leanok
	\uses{def:projective-strategy, def:symmetric-projective-strategy, def:good-strategy}
	Let $(\psi,A^{\mathrm A},B^{\mathrm A},L^{\mathrm A},A^{\mathrm B},B^{\mathrm B},L^{\mathrm B})$ be a projective strategy on two local spaces passing the $(m,q,d)$ low individual degree test with probability at least $1-\eps$.  Adding a role register and placing the two local spaces in the direct sum \(H_{\mathrm A}\oplus H_{\mathrm B}\) yields a symmetric projective strategy that is $(3\eps,3\eps,3\eps)$-good.
\end{lemma}
\begin{proof}
	\leanok
	\uses{def:projective-strategy, def:good-strategy}
	Each of the three subtests (axis-parallel, self-consistency, diagonal) occurs with probability $1/3$, so the total failure probability at most $\eps$ implies each branch is at most $3\eps$. Block-diagonalizing the measurements over the role register preserves the axis-parallel and diagonal branch averages exactly, and the cross-register self-consistency branch of the symmetrized strategy equals the original point-agreement defect, which is itself bounded by $3\eps$.
\end{proof}

\begin{remark}[Lean auxiliary data for role-register symmetrization]\label{rem:symmetrization-data}
	The formal proof of Theorem~\ref{thm:main-formal} uses the heterogeneous
	role-register strategy and its goodness proof.  This construction supports
	the final assembly and is not an additional hypothesis in
	Lemma~\ref{lem:role-register-symmetrization}.
\end{remark}

	\begin{proposition}[Lean role-register main-induction handoff for the Section 3 proof]\label{lem:main-formal-role-register-handoff}
		\lean{MIPStarRE.LDT.ProjStrat.roleRegisterSymmStrategy_sourceMainInduction}
		\leanok
		\uses{lem:role-register-symmetrization, thm:main-induction}
		For a two-space projective strategy passing the low individual degree test
		with error $\eps$, the role-register symmetrization is
		$(3\eps,3\eps,3\eps)$-good.  Under the corrected hypothesis
		$k\ge400md$, Theorem~\ref{thm:main-induction} therefore supplies the
		Section~6 polynomial measurement used in the final theorem route.
	\end{proposition}
	\begin{proof}
		\leanok
		This is a Lean-only handoff theorem: combine
		Lemma~\ref{lem:role-register-symmetrization} with the corrected
		large-\(k\) theorem \lean{MIPStarRE.LDT.MainInductionStep.mainInduction}.
	\end{proof}

	\paragraph{Removed successor restricted-recursion targets.}
		The former Section~3 Lean-only answer-valued restricted-recursion records
		have been removed along with the same-space final-theorem assembly.  The
		successor work belongs to the proof of Theorem~\ref{thm:main-induction},
		not to additional hypotheses or auxiliary inputs in
		Theorem~\ref{thm:main-formal}.

	\begin{remark}[Lean Step~6 constructions for the Section~3 proof]\label{rem:main-formal-step6-constructions}
		The old same-space base-case and completion-witness records have been
		removed.  The current route starts with the heterogeneous role-register
		application of Theorem~\ref{thm:main-induction}, extracts the two occupied
		role blocks, proves the Step~5 full-polynomial self-consistency relation,
		and then invokes the theorems for making measurements projective and for
		completion inside the two-space source route.  These constructions are
		proof steps in Theorem~\ref{thm:main-formal}, not hypotheses of its
		statement.
	\end{remark}


	\paragraph{Lean successor-dependent Step~6 targets.}
		\lean{MIPStarRE.LDT.ProjStrat.roleRegisterSymmStrategy_sourceMainInduction, MIPStarRE.LDT.ProjStrat.sourceRoleRegisterFinalPointConsistency}
		The role-register Section~6 measurement is obtained by applying the
		corrected large-\(k\) theorem \texttt{MainInductionStep.mainInduction} to
		the role-register symmetrization.  The final source-route theorem then
		carries this measurement through the two unsymmetrized point-consistency
		relations, the Step~5 self-consistency calculation, the step making
		measurements projective, completion, and the final point-evaluation
		transport.  These targets are proof-complete for the corrected large-\(k\)
		interface and are not added hypotheses of Theorem~\ref{thm:main-formal}.


	\begin{remark}[Projective-completion construction]\label{rem:main-formal-successor-completion-construction}
		The projective-completion construction is now internal to the two-space
		source route.  It is not a separate assumption in the statement of
		Theorem~\ref{thm:main-formal}.  The corrected source theorem is
		proof-complete: the general two-space role-register passage has been
		assembled, and the theorem statement records the documented \(k\ge400md\)
		and \(k>0\) boundary corrections.
\end{remark}

\begin{remark}[Lean measurement-unsymmetrization consistency]\label{rem:main-formal-unsymmetrization-consistency}
	Lean now isolates the Step~3 extraction from \texttt{inductive\_step.tex}
	lines 84--108.  In the same-space route, the principal-block map
	\texttt{extractRoleBlock} turns a role-register POVM \(G\) into the two POVMs
	\(G^{\mathrm A}\) and \(G^{\mathrm B}\) on the original prover space and
	proves that point-evaluation postprocessing commutes with this extraction.
	In the heterogeneous two-space route, the corresponding principal-block maps
	\texttt{extractRoleRegisterAlice} and \texttt{extractRoleRegisterBob}
	extract the occupied \((\mathrm A,\mathrm{inl})\) and
	\((\mathrm B,\mathrm{inr})\) blocks of a measurement on
	\(\{\mathrm A,\mathrm B\}\times(\mathcal H_{\mathrm A}\oplus\mathcal H_{\mathrm B})\).
	Lean proves the trace-level factor-two identities for arbitrary
	role-register measurements in
	\texttt{qBipartiteConsDefect\_roleRegisterProjMeas\_arbitrary\_eq\_average}
	and its two factor-two consequences.  No projectivity of arbitrary principal
	blocks is claimed, matching the paper's later step making measurements
	projective.
\end{remark}

\begin{definition}[Final error envelope]\label{def:main-formal-envelope}
	\lean{MIPStarRE.LDT.Test.mainFormalEnvelope}\leanok
	Set
	\[
		E = \eps^{1/40000} + (d/q)^{1/40000} + e^{-k/(2560000m^2)}.
	\]
\end{definition}

The remaining entries in this scalar-cascade subsection are Lean bookkeeping
lemmas for the final proof of Theorem~\ref{thm:main-formal}.  They record
explicit inequalities between the named error parameters; they are not
independent source theorems from the paper.

\begin{theorem}[Basic facts about the final error envelope]\label{thm:main-formal-envelope-basics}
	\lean{MIPStarRE.LDT.Test.mainFormalError_eq_envelope, MIPStarRE.LDT.Test.mainFormalEnvelope_nonneg}\leanok
	\uses{def:main-formal-envelope}
	If $0 \le \eps$, then
	\[
		\mainFormalError
		= 100000k^2m^4\left(\eps^{1/40000} + (d/q)^{1/40000} + e^{-k/(2560000m^2)}\right)
		= 100000k^2m^4 E,
	\]
	and $E \ge 0$.
\end{theorem}
\begin{proof}
	\leanok
	Both identities are by definition: the first unfolds $\mainFormalError$ and the second unfolds $E$. Nonnegativity of $E$ follows summand-wise from \texttt{Real.rpow\_nonneg} (using $0 \le \eps$) and \texttt{Real.exp\_nonneg}.
\end{proof}

\begin{definition}[Cascade parameters and standing hypotheses]\label{def:main-formal-error-cascade}
	\lean{MIPStarRE.LDT.Test.cascadeSigma, MIPStarRE.LDT.Test.cascadeZeta1, MIPStarRE.LDT.Test.cascadeZeta2, MIPStarRE.LDT.Test.cascadeZeta3, MIPStarRE.LDT.Test.cascadeZeta4, MIPStarRE.LDT.Test.CascadeHypotheses}\leanok
	\uses{def:main-formal-envelope}
	Assume $q > 0$, $k \ge 1$, $m \ge 1$, $0 \le \eps \le 1$, and $d \le q$. Fix an incoming error parameter $\nu$. Define
	\[
		\sigma = m^2\left(\nu + e^{-k/(80000m^2)}\right),
		\qquad
		\zeta_1 = 2\sigma + 2\sqrt{3\eps+2\sigma} + md/q,
	\]
	\[
		\zeta_2 = 200\zeta_1^{1/4} + 42\zeta_1^{1/8},
		\qquad
		\zeta_3 = 6\zeta_1 + 6\zeta_2,
		\qquad
		\zeta_4 = 2\sigma + 2\sqrt{\zeta_1+\zeta_3/2}.
	\]
	The coefficient $42$ is the Lean cascade scalar; it widens the paper's printed
	coefficient $40$ so that the completion theorem's extra $2\zeta_1$ term is
	absorbed while preserving the final error envelope.
\end{definition}

\begin{theorem}[Bound for the first cascade term]\label{thm:sigma-bound-main-formal}
	\lean{MIPStarRE.LDT.Test.sigma_bound}\leanok
	\uses{def:main-formal-envelope, def:main-formal-error-cascade}
	Under the standing hypotheses, if
	\[
		\nu \le 10000k^2m^2\left(\eps^{1/1024} + (d/q)^{1/1024}\right),
	\]
	then
	\[
		\sigma \le 10000k^2m^4 E.
	\]
\end{theorem}
\begin{proof}
	\leanok
	Multiply the assumed bound by $m^2$, replace $m^2$ by $m^4$, coarsen the exponents from $1/1024$ to $1/40000$, and enlarge the exponential denominator from $80000m^2$ to $2560000m^2$.
\end{proof}

\begin{theorem}[Cascade bounds for the later cascade terms]\label{thm:zeta-bounds-main-formal}
	\lean{MIPStarRE.LDT.Test.zeta1_bound, MIPStarRE.LDT.Test.zeta2_bound, MIPStarRE.LDT.Test.zeta3_bound, MIPStarRE.LDT.Test.zeta4_bound}\leanok
	\uses{def:main-formal-envelope, def:main-formal-error-cascade, thm:sigma-bound-main-formal}
	Under the standing hypotheses, assume
	\[
		0 \le \nu \le 10000k^2m^2\left(\eps^{1/1024} + (d/q)^{1/1024}\right).
	\]
	Define
	\[
		\sigma = m^2\left(\nu + e^{-k/(80000m^2)}\right),
	\]
	\[
		\zeta_1 = 2\sigma + 2\sqrt{3\eps + 2\sigma} + md/q,
		\qquad
		\zeta_2 = 200\zeta_1^{1/4} + 42\zeta_1^{1/8},
	\]
	\[
		\zeta_3 = 6\zeta_1 + 6\zeta_2,
		\qquad
		\zeta_4 = 2\sigma + 2\sqrt{\zeta_1 + \zeta_3/2}.
	\]
	Then
	\[
		\zeta_1, \zeta_2, \zeta_4 \le 100000k^2m^4 E,
		\qquad
		\zeta_3 \le 2 \cdot 100000k^2m^4 E.
	\]
\end{theorem}
\begin{proof}
	\leanok
	The Lean declarations first derive the native bounds for $\zeta_1,\dots,\zeta_4$
	from the cascade definitions using square-root and $rpow$ monotonicity, with
	$\zeta_2$ using the widened coefficient $42$, and then coarsen those bounds to
	the common envelope $E$. The $\zeta_2$ proof additionally uses $k \le k^2$,
	$m \le m^4$, and the nonnegativity of the earlier cascade variables.
\end{proof}

\begin{theorem}[Consolidated error cascade]\label{thm:error-cascade-main-formal}
	\lean{MIPStarRE.LDT.Test.errorCascade_le_mainFormalError}\leanok
	\uses{def:main-formal-envelope, def:main-formal-error-cascade, thm:sigma-bound-main-formal, thm:zeta-bounds-main-formal}
	Assume the standing hypotheses and
	\[
		0 \le \nu \le 10000k^2m^2\left(\eps^{1/1024} + (d/q)^{1/1024}\right).
	\]
	Let $\sigma,\zeta_1,\zeta_2,\zeta_3$ be the cascade quantities defined above. Then
	\[
		\sigma \le 100000k^2m^4 E,
		\qquad
		\zeta_1,\zeta_2 \le 100000k^2m^4 E,
		\qquad
		\zeta_3 \le 2 \cdot 100000k^2m^4 E,
	\]
	and the concrete quantity
	\[
		\zeta_4 = 2\sigma + 2\sqrt{\zeta_1 + \zeta_3/2}
	\]
	also satisfies
	\[
		\zeta_4 \le 100000k^2m^4 E.
	\]
\end{theorem}
\begin{proof}
	\leanok
	This is the direct combination of Theorem~\ref{thm:sigma-bound-main-formal} and Theorem~\ref{thm:zeta-bounds-main-formal}. The Lean declaration uses the explicit equalities naming $\sigma,\zeta_1,\zeta_2,\zeta_3$ as the corresponding cascade expressions.
\end{proof}

\begin{remark}[Lean Step 5 expansion theorem]\label{rem:main-formal-step5-sz-bound}
	For the Step~5 passage below, Lean now names the weighted polynomial-collision
	mass, proves its $md/q$ bound by specializing the existing tensor
	Schwartz--Zippel helper, and proves the algebraic expansion/reindexing
	statement comparing the evaluated consistency defect with the full-polynomial
	defect plus that collision mass. The theorem for Step~5 therefore no longer
	takes the expansion statement as an external hypothesis.
\end{remark}

\begin{remark}[Lean witness structures for the main-formal assembly]\label{rem:main-formal-lean-witness-records}
The following Lean declarations are bookkeeping infrastructure for the final
assembly witness:
They split off the vacuous branch where the theorem error is at least $1$,
derive the Step~8 scalar hypotheses in the non-vacuous branch, and isolate the
role-register Section~6 measurement behind a concrete role-witness record.
	The role-register measurement is now obtained directly from the
	\texttt{mainInduction} theorem by
	\texttt{roleRegisterSymmStrategy\_sourceMainInduction}; the successor case is
	handled by the checked Section~6 theorem, so it is not a separate boundary
	hypothesis in Section~3.  Lean also records that every non-base
parameter choice has a predecessor via
\texttt{Parameters.successorDecompositionOfNeOne}, and it retains the explicit
transport helpers for the strategy and passing proof across that predecessor
equality.  The older branch-witness records, which assembled ordinary or
answer-valued successor data from recursive slice witnesses and supplied
self-improvement inputs, have been removed.  The current completion witness
carries the concrete role witness obtained from the Section~6
theorem rather than a separate decorative branch witness.  Once the
role-measurement record is obtained,
Lean proves the two factor-two unsymmetrization links from
its symmetrized consistency estimate.  The subsequent proved lemmas use the
proved Section~5 Schwartz--Zippel theorem to convert evaluated pre-projective
consistency to full-polynomial consistency at $\zeta_1$, prove the paper
line-116 triangle step, and reconstruct the line-156 $\approx_{\zeta_3}$
relation obtained while making measurements projective from pre-projective
consistency plus left/right completion closeness at $\zeta_2$.  Lean also
derives the evaluated line-164 projective consistency.  The error-parametric
lemmas
\texttt{pointAConsistency\_ofLine169Consistency} and
\texttt{pointBConsistency\_ofLine169Consistency} show that a polynomial
line-169 estimate with error $\eta$ yields the final point error
\[
  2\sigma + 2\sqrt{\eta+\zeta_3/2}.
\]
The active formal route specializes this with the repaired value
$\eta=\zeta_1+10\zeta_1^{1/8}$.  The existing
\texttt{pointAConsistency} and \texttt{pointBConsistency} declarations are the
corresponding repaired specializations.  After the \#869 right-register
completion transport, the remaining active post-role witness asks only for
Alice's line-146 completion closeness, Bob's right-register completion
closeness, and the two repaired polynomial line-169 links.  Lean reconstructs
pre-projective full-polynomial consistency from the concrete role witness, uses
those repaired line-169 links together with the evaluated line-164 projective
consistency, and converts directly to the projective completion witness used by
the final assembly.
	It no longer asks for an arbitrary role-measurement record or an arbitrary
	Section~6 witness as an independent field.  It also no longer asks for
	stored factor-two role-block estimates, a stored Section~6 consistency input,
	or a duplicated pre-projective consistency field; those are rebuilt from the
	concrete role-register measurement and \texttt{hpass}.  Lean then applies the
	projective converse of \texttt{simeq-to-approx} inside the two-space source
	route.

As stated in \texttt{preliminaries.tex}, \texttt{prop:triangle-sub} gives a
$\delta+\sqrt{\epsilon}$ loss, so applying it directly to
$G^{\mathrm A}\simeq_{\zeta_1}G^{\mathrm B}$ and
$G^{\mathrm A}\approx_{\zeta_2}Q^{\mathrm A}$ yields
	$Q^{\mathrm A}\simeq_{\zeta_1+\sqrt{\zeta_2}}G^{\mathrm B}$, not the printed
	$\zeta_1$ estimate.  The generic corrected route would degrade the final point
	error to scale $t^{1/65536}$, outside the current \texttt{mainFormalError}
	envelope.  The sharper local pre-completion repair from
	Remark~\ref{rem:lean-line169-projectivization-match-mass} avoids that
	widening.  Its completion step uses only the formal monotonicity helper
	\texttt{ProjectivizationMatchMassMonotonicity.completeAtOutcomeProj\_left\_matchMass\_ge},
	not a separate exact match-mass witness route.
	For the corrected boundary assumptions \(k \ge 400md\) and \(k>0\), the
	source theorem now has a checked proof-level route.

\end{remark}
\begin{proof}
	\lean{MIPStarRE.LDT.Test.mainFormal}
	\leanok
	\proves{thm:main-formal}
	\uses{thm:main-induction, lem:role-register-symmetrization, lem:orthonormalization-main-lemma, prop:completing-to-measurement, prop:simeq-triangle-inequality, prop:simeq-data-processing, prop:simeq-to-approx, prop:triangle-sub, prop:triangle-inequality-for-approx_delta, lem:schwartz-zippel-individual, def:main-formal-envelope, def:main-formal-error-cascade, thm:sigma-bound-main-formal, thm:zeta-bounds-main-formal}
	Let
	\[
		(\psi,A^{\mathrm A},B^{\mathrm A},L^{\mathrm A},A^{\mathrm B},B^{\mathrm B},L^{\mathrm B})
	\]
	be a projective strategy passing the $(m,q,d)$ low individual degree test with probability at least $1-\eps$. Since the three subtests occur with probability $1/3$, this strategy is $(3\eps,3\eps,3\eps)$-good.

	Introduce a two-dimensional role register on each side and define the symmetrized state
	\[
		\ket{\psi_{\mathrm{sym}}}
		=
		\frac{1}{\sqrt{2}}
		\left(
		\ket{0}_{\mathrm{A}'}\ket{1}_{\mathrm{B}'}\ket{\psi}_{\mathrm{AB}}
		+
		\ket{1}_{\mathrm{A}'}\ket{0}_{\mathrm{B}'}\ket{\psi_{\mathrm{swap}}}_{\mathrm{AB}}
		\right).
	\]
	Define the symmetrized point measurement by
	\[
		(A_{\mathrm{sym}})_a^u
		=
		\ket{0}\bra{0} \ot A_a^{\mathrm A,u}
		+
		\ket{1}\bra{1} \ot A_a^{\mathrm B,u},
	\]
	and similarly for $B_{\mathrm{sym}}$ and $L_{\mathrm{sym}}$. This yields a symmetric $(3\eps,3\eps,3\eps)$-good strategy. Set
	\[
		\nu = 10000k^2m^2\left(\eps^{1/1024}+(d/q)^{1/1024}\right),
		\qquad
		\sigma = m^2\left(\nu+e^{-k/(80000m^2)}\right).
	\]
	Theorem~\ref{thm:main-induction} applied to the symmetrized strategy produces a measurement $G=\{G_g\} \in \polymeas{m}{q}{d}$ such that
	\begin{equation}\label{eq:symmetrized-induction-point-consistency}
		(A_{\mathrm{sym}})_a^u \ot I \simeq_\sigma I \ot G_{[g(u)=a]},
		\qquad
		G_{[g(u)=a]} \ot I \simeq_\sigma I \ot (A_{\mathrm{sym}})_a^u.
	\end{equation}

	Now unsymmetrize by defining
	\[
		G_g^{\mathrm A} = (\bra{0} \ot I)\, G_g\, (\ket{0} \ot I),
		\qquad
		G_g^{\mathrm B} = (\bra{1} \ot I)\, G_g\, (\ket{1} \ot I).
	\]
	These are POVMs on the original prover space. Since $A_{\mathrm{sym}}$ is a convex combination of $(\ket{0}\bra{0} \ot I)$ and $(\ket{1}\bra{1} \ot I)$ weighted by the test probabilities, the cross-terms $\bra{0} \cdot A_{\mathrm{sym}} \cdot \ket{1}$ vanish. Expanding the consistency of $G$ gives
	\begin{gather}
		G_{[g(u)=a]}^{\mathrm A} \ot I \simeq_{2\sigma} I \ot A_a^{\mathrm B,u}, \label{eq:cons-a}\\
		I \ot G_{[g(u)=a]}^{\mathrm B} \simeq_{2\sigma} A_a^{\mathrm A,u} \ot I. \label{eq:cons-b}
	\end{gather}
	Since the original strategy is $(3\eps,3\eps,3\eps)$-good,
	\[
		A_a^{\mathrm A,u} \ot I \simeq_{3\eps} I \ot A_a^{\mathrm B,u}.
	\]
	Therefore Lemma~\ref{prop:simeq-triangle-inequality} yields
	\[
		G_{[g(u)=a]}^{\mathrm A} \ot I \simeq_{2\sigma + 2\sqrt{3\eps+2\sigma}} I \ot G_{[g(u)=a]}^{\mathrm B}.
	\]
	Expanding this consistency relation and applying Lemma~\ref{lem:schwartz-zippel-individual} to replace point evaluations by polynomial evaluations (at cost $md/q$) gives
	\begin{equation}\label{eq:G-self-consistency}
		G_g^{\mathrm A} \ot I \simeq_{\zeta_1} I \ot G_g^{\mathrm B},
	\end{equation}
	where
	\[
		\zeta_1 = 2\sigma + 2\sqrt{3\eps+2\sigma} + \frac{md}{q}.
	\]
	The Lean Step~5 theorem
	\lean{MIPStarRE.LDT.Test.MainFormalStep5ExpansionBound}
	\lean{MIPStarRE.LDT.Test.mainFormalStep5_expansionBound}
	\lean{MIPStarRE.LDT.Test.mainFormalStep5_selfConsistency_ofExpansionBound}
	formalizes this expansion/reindexing step and combines it with the proved
	Schwartz--Zippel tensor bound
	\lean{MIPStarRE.LDT.Preliminaries.polynomialCollisionMass_le_mdq}.  The
	Lean statement is now formulated for a bipartite state on
	\(H_A\otimes H_B\), so this particular Schwartz--Zippel loss is no longer a
	same-space restriction in the source-theorem route.

	Apply Lemma~\ref{lem:orthonormalization-main-lemma} to \ref{eq:G-self-consistency}. This gives projective submeasurements $P^{\mathrm A},P^{\mathrm B} \in \polysub{m}{q}{d}$ with
	\[
		G_g^{\mathrm A} \ot I \approx_{100\zeta_1^{1/4}} P_g^{\mathrm A} \ot I,
		\qquad
		I \ot G_g^{\mathrm B} \approx_{100\zeta_1^{1/4}} I \ot P_g^{\mathrm B}.
	\]
	Completing them by Lemma~\ref{prop:completing-to-measurement} yields projective measurements $Q^{\mathrm A},Q^{\mathrm B} \in \polymeas{m}{q}{d}$ such that
	\begin{align}
		G_g^{\mathrm A} \ot I & \approx_{\zeta_2} Q_g^{\mathrm A} \ot I, \label{eq:G-with-Q-A} \\
		I \ot G_g^{\mathrm B} & \approx_{\zeta_2} I \ot Q_g^{\mathrm B}, \notag
	\end{align}
	where
	\[
		\zeta_2 = 200\zeta_1^{1/4} + 42\zeta_1^{1/8}.
	\]
	The paper prints the coefficient $40$ in this display.  The formal Lean
	cascade uses the widened coefficient $42$ to absorb the literal completion
	error $2\cdot(100\zeta_1^{1/4})+4\sqrt{100\zeta_1^{1/4}}+2\zeta_1$:
	under the non-vacuous scalar regime $0 \le \zeta_1 \le 1$, the extra term
	$2\zeta_1$ is bounded by $2\zeta_1^{1/8}$.

	By \ref{eq:G-self-consistency} and Lemma~\ref{prop:simeq-to-approx},
	\[
		G_g^{\mathrm A} \ot I \approx_{2\zeta_1} I \ot G_g^{\mathrm B}.
	\]
	Hence Lemma~\ref{prop:triangle-inequality-for-approx_delta} gives
	\[
		Q_g^{\mathrm A} \ot I \approx_{\zeta_3} I \ot Q_g^{\mathrm B},
	\]
	where
	\[
		\zeta_3 = 6\zeta_1 + 6\zeta_2.
	\]
	Because both measurements are projective, Lemma~\ref{prop:simeq-to-approx} converts this to
	\begin{equation}\label{eq:completed-polynomial-consistency}
		Q_g^{\mathrm A} \ot I \simeq_{\zeta_3/2} I \ot Q_g^{\mathrm B}.
	\end{equation}
	By data processing,
	\begin{equation}\label{eq:point-evaluation-data-processing}
		Q_{[g(u)=a]}^{\mathrm A} \ot I \simeq_{\zeta_3/2} I \ot Q_{[g(u)=a]}^{\mathrm B}.
	\end{equation}

	Next, Lemma~\ref{prop:triangle-sub}, applied to \ref{eq:G-self-consistency} and \ref{eq:G-with-Q-A}, implies
	\[
		Q_g^{\mathrm A} \ot I \simeq_{\zeta_1} I \ot G_g^{\mathrm B}.
	\]
	Data processing then gives
	\begin{equation}\label{eq:left-completed-point-with-original-polynomial}
		Q_{[g(u)=a]}^{\mathrm A} \ot I \simeq_{\zeta_1} I \ot G_{[g(u)=a]}^{\mathrm B}.
	\end{equation}
	Applying Lemma~\ref{prop:simeq-triangle-inequality} to \ref{eq:cons-b}, \ref{eq:left-completed-point-with-original-polynomial}, and \ref{eq:point-evaluation-data-processing} yields
	\begin{equation}\label{eq:final-left-point-consistency}
		A_a^{\mathrm A,u} \ot I \simeq_{\zeta_4} I \ot Q_{[g(u)=a]}^{\mathrm B},
	\end{equation}
	where
	\[
		\zeta_4 = 2\sigma + 2\sqrt{\zeta_1+\zeta_3/2}.
	\]
	The same argument with the roles reversed gives
	\begin{equation}\label{eq:final-right-point-consistency}
		I \ot A_a^{\mathrm B,u} \simeq_{\zeta_4} Q_{[g(u)=a]}^{\mathrm A} \ot I.
	\end{equation}

	Remark~\ref{rem:main-formal-lean-witness-records} records the Lean witness
	structures that separate the scalar cascade and the Step~5 Schwartz--Zippel
	handoff from the remaining projective-stage transport work.  The scalar side
	discharges the paper's line~71--73 coarsening and the Step~8 hypotheses in the
	non-vacuous branch; the Step~5 theorem now turns the evaluated estimate at
	line~116 into full-polynomial consistency at $\zeta_1$.  The remaining Lean
	witness is the completion-and-line-169 construction: supply the actual
	role-register measurement, Alice's line-146 completion estimate, Bob's
	left-register completion estimate, and the two exact polynomial-level
	$\zeta_1$ links.  Lean transports the Bob-side completion estimate to the
	line-147 right-register form, reconstructs the duplicated Step~6 pre-projective
	consistency field from these inputs and \texttt{hpass}, then data-processes the
	line-169 links to line~172 and derives the two $\zeta_4$ point goals and the
	$\zeta_3/2$ self-consistency target.

	With $E$ as in Definition~\ref{def:main-formal-envelope}, the quantities $\sigma,\zeta_1,\zeta_2,\zeta_3,\zeta_4$ form the cascade of Definition~\ref{def:main-formal-error-cascade}. Theorem~\ref{thm:sigma-bound-main-formal} and Theorem~\ref{thm:zeta-bounds-main-formal} give the final inflation to the scale $100000k^2m^4 E$. The underlying estimates, with the formal widened $\zeta_2$, are:
	\[
		\sigma \le 10000k^2m^4\left(\eps^{1/1024}+(d/q)^{1/1024}+e^{-k/(80000m^2)}\right).
	\]
	Then
	\[
		\zeta_1 \le 20204k^2m^4\left(\eps^{1/2048}+(d/q)^{1/2048}+e^{-k/(160000m^2)}\right),
	\]
	and
	\[
		\zeta_2 \le 2568km\left(\eps^{1/16384}+(d/q)^{1/16384}+e^{-k/(1280000m^2)}\right).
	\]
	Consequently,
	\[
		\zeta_3 \le 150000k^2m^4\left(\eps^{1/16384}+(d/q)^{1/16384}+e^{-k/(1280000m^2)}\right),
	\]
	and
	\[
		\zeta_4 \le 40000k^2m^4\left(\eps^{1/32768}+(d/q)^{1/32768}+e^{-k/(2560000m^2)}\right).
	\]
	Both $\zeta_3/2$ and $\zeta_4$ are therefore bounded by
	\[
		100000k^2m^4\left(\eps^{1/40000}+(d/q)^{1/40000}+e^{-k/(2560000m^2)}\right).
	\]
	Hence \ref{eq:final-left-point-consistency}, \ref{eq:final-right-point-consistency}, and \ref{eq:completed-polynomial-consistency} give the three claimed conclusions, with $Q^{\mathrm A}$ and $Q^{\mathrm B}$ renamed as the theorem's $G^{\mathrm A}$ and $G^{\mathrm B}$.
\end{proof}


%%FILE blueprint/src/content.tex
% Main blueprint content.
% Imported by both print.tex and web.tex.

\input{chapter/ch01_overview}
\input{chapter/ch02_test}
\input{chapter/ch03_preliminaries}
\input{chapter/ch04_projective}
\input{chapter/ch05_expansion}
\input{chapter/ch06_variance}
\input{chapter/ch07_self_improvement}
\input{chapter/ch08_commutativity}
\input{chapter/ch09_pasting}
\input{chapter/ch10_induction}
